% ═══════════════════════════════════════════════════════════════════════
\section{La dimension de similitude généralisée}
% ═══════════════════════════════════════════════════════════════════════

\begin{frame}{Des copies de rapports différents}
  \small Si $E = w_1(E)\cup\cdots\cup w_N(E)$, où $w_i$ est une
  similitude de rapport $r_i$, sans chevauchement, l'argument
  heuristique de la section 5 donne
  \[
    M_d(E) = r_1^d\,M_d(E)+r_2^d\,M_d(E)+\cdots+r_N^d\,M_d(E),
  \]
  possible avec une taille finie non nulle seulement si
  $r_1^d+\cdots+r_N^d = 1$.
\end{frame}

\begin{frame}{Proposition 3.5 --- L'équation de Moran}
  \begin{Enonce}{Proposition 3.5 --- L'équation de Moran}
    Soient $N\geq2$ et $r_1,\dots,r_N\in\intoo{0}{1}$. L'équation
    \[
      r_1^d+r_2^d+\cdots+r_N^d = 1 \tag{3.1}
    \]
    possède une unique solution $d>0$.
  \end{Enonce}
  \begin{AuTableau}
    \footnotesize $\varphi(d)=\sum r_i^d=\sum\eu^{d\ln r_i}$ est continue
    et strictement décroissante ; $\varphi(0)=N>1$, $\varphi(d)\to0$ ;
    valeurs intermédiaires.
  \end{AuTableau}
\end{frame}

\begin{frame}{Outil 3.2 --- Résoudre l'équation de Moran}
  \begin{Outil}{Résoudre l'équation de Moran}
    \small
    \begin{itemize}
      \item (3.1) ne se résout pas algébriquement en général, mais très
        bien graphiquement. Dans Desmos : tracer
        $y = 2(1/2)^x+3(1/3)^x$ et $y=1$, puis cliquer sur l'intersection.
      \item Le graphe illustre la démonstration : courbe décroissante,
        qui part de $N$ en $x=0$ et tend vers $0$.
      \item WolframAlpha : \texttt{solve 2(1/2)\^{}d + 3(1/3)\^{}d = 1}.
      \item En $5^{\text{e}}$ secondaire : comparer avec la résolution
        par logarithmes du cas $Nr^d=1$ --- ce que les logarithmes
        permettent, et ce qu'ils ne permettent pas.
    \end{itemize}
  \end{Outil}
\end{frame}

\begin{frame}{Définitions 3.4 et 3.5 --- Ne pas se chevaucher}
  \begin{Enonce}{Définition 3.4 --- Condition de l'ouvert (Moran)}
    \small Le SIF $\{w_1,\dots,w_N\}$ satisfait la \defi{condition de
    l'ouvert} s'il existe un ouvert borné non vide $V$ tel que
    (a) $w_i(V)\subseteq V$ pour tout $i$, et (b) les $w_i(V)$ sont deux
    à deux disjoints.
  \end{Enonce}
  \small Sierpiński : $V$ = l'intérieur du triangle de départ. Ses trois
  images sont disjointes, car les triangles ne se touchent qu'en des
  sommets, sur leurs frontières. Ouvert et frontière (chapitre 2)
  interviennent directement.
  \begin{Enonce}{Définition 3.5 --- Dimension de similitude généralisée}
    \small Soit $E$ l'attracteur d'un SIF de similitudes de rapports
    $r_1,\dots,r_N$ qui satisfait la condition de l'ouvert. Sa
    \defi{dimension de similitude généralisée} est l'unique solution
    $d$ de (3.1).
  \end{Enonce}
\end{frame}

\begin{frame}{Figure 3.12 --- Trois SIF de similitudes}
  \FigBox[0.5\textheight]{%
    \begin{tikzpicture}[scale=3.3]
      \begin{scope}
        \draw[sectioncolor!60, dashed, thin] (0,0) rectangle (1,1);
        % Chevauchement : disjoint
% Fichier généré par generer.py : ne pas modifier à la main.
\fill[nuage]
  (0.256,0.129)rectangle++(0.009,0.009) (0.121,0.059)rectangle++(0.009,0.009) (0.576,0.026)rectangle++(0.009,0.009) (0.274,0.53)rectangle++(0.009,0.009) (0.129,0.252)rectangle++(0.009,0.009)
  (0.06,0.119)rectangle++(0.009,0.009) (0.546,0.055)rectangle++(0.009,0.009) (0.78,0.024)rectangle++(0.009,0.009) (0.372,0.529)rectangle++(0.009,0.009) (0.176,0.772)rectangle++(0.009,0.009)
  (0.602,0.368)rectangle++(0.009,0.009) (0.287,0.174)rectangle++(0.009,0.009) (0.135,0.601)rectangle++(0.009,0.009) (0.063,0.806)rectangle++(0.009,0.009) (0.028,0.385)rectangle++(0.009,0.009)
  (0.011,0.702)rectangle++(0.009,0.009) (0.003,0.855)rectangle++(0.009,0.009) (0.519,0.408)rectangle++(0.009,0.009) (0.247,0.193)rectangle++(0.009,0.009) (0.116,0.611)rectangle++(0.009,0.009)
  (0.053,0.811)rectangle++(0.009,0.009) (0.023,0.387)rectangle++(0.009,0.009) (0.529,0.183)rectangle++(0.009,0.009) (0.252,0.086)rectangle++(0.009,0.009) (0.118,0.039)rectangle++(0.009,0.009)
  (0.574,0.016)rectangle++(0.009,0.009) (0.793,0.005)rectangle++(0.009,0.009) (0.898,0)rectangle++(0.009,0.009) (0.949,-0.002)rectangle++(0.009,0.009) (0.453,0.517)rectangle++(0.009,0.009)
  (0.215,0.766)rectangle++(0.009,0.009) (0.101,0.885)rectangle++(0.009,0.009) (0.046,0.943)rectangle++(0.009,0.009) (0.02,0.97)rectangle++(0.009,0.009) (0.007,0.463)rectangle++(0.009,0.009)
  (0.001,0.22)rectangle++(0.009,0.009) (0.518,0.103)rectangle++(0.009,0.009) (0.766,0.047)rectangle++(0.009,0.009) (0.366,0.54)rectangle++(0.009,0.009) (0.173,0.257)rectangle++(0.009,0.009)
  (0.081,0.121)rectangle++(0.009,0.009) (0.036,0.576)rectangle++(0.009,0.009) (0.535,0.274)rectangle++(0.009,0.009) (0.255,0.649)rectangle++(0.009,0.009) (0.12,0.309)rectangle++(0.009,0.009)
  (0.575,0.146)rectangle++(0.009,0.009) (0.274,0.588)rectangle++(0.009,0.009) (0.649,0.28)rectangle++(0.009,0.009) (0.309,0.652)rectangle++(0.009,0.009) (0.146,0.311)rectangle++(0.009,0.009)
  (0.588,0.147)rectangle++(0.009,0.009) (0.28,0.588)rectangle++(0.009,0.009) (0.652,0.28)rectangle++(0.009,0.009) (0.311,0.652)rectangle++(0.009,0.009) (0.667,0.311)rectangle++(0.009,0.009)
  (0.318,0.147)rectangle++(0.009,0.009) (0.15,0.068)rectangle++(0.009,0.009) (0.07,0.55)rectangle++(0.009,0.009) (0.031,0.262)rectangle++(0.009,0.009) (0.013,0.123)rectangle++(0.009,0.009)
  (0.524,0.057)rectangle++(0.009,0.009) (0.249,0.545)rectangle++(0.009,0.009) (0.117,0.779)rectangle++(0.009,0.009) (0.574,0.372)rectangle++(0.009,0.009) (0.793,0.176)rectangle++(0.009,0.009)
  (0.378,0.602)rectangle++(0.009,0.009) (0.179,0.287)rectangle++(0.009,0.009) (0.604,0.135)rectangle++(0.009,0.009) (0.807,0.063)rectangle++(0.009,0.009) (0.905,0.028)rectangle++(0.009,0.009)
  (0.432,0.011)rectangle++(0.009,0.009) (0.205,0.003)rectangle++(0.009,0.009) (0.096,-0.001)rectangle++(0.009,0.009) (0.044,0.517)rectangle++(0.009,0.009) (0.019,0.766)rectangle++(0.009,0.009)
  (0.007,0.885)rectangle++(0.009,0.009) (0.001,0.943)rectangle++(0.009,0.009) (-0.002,0.45)rectangle++(0.009,0.009) (-0.003,0.734)rectangle++(0.009,0.009) (-0.004,0.87)rectangle++(0.009,0.009)
  (0.516,0.415)rectangle++(0.009,0.009) (0.245,0.197)rectangle++(0.009,0.009) (0.635,0.092)rectangle++(0.009,0.009) (0.303,0.042)rectangle++(0.009,0.009) (0.143,0.018)rectangle++(0.009,0.009)
  (0.066,0.006)rectangle++(0.009,0.009) (0.029,0.001)rectangle++(0.009,0.009) (0.012,0.518)rectangle++(0.009,0.009) (0.003,0.246)rectangle++(0.009,0.009) (0.519,0.116)rectangle++(0.009,0.009)
  (0.247,0.053)rectangle++(0.009,0.009) (0.116,0.023)rectangle++(0.009,0.009) (0.053,0.529)rectangle++(0.009,0.009) (0.023,0.251)rectangle++(0.009,0.009) (0.009,0.118)rectangle++(0.009,0.009)
  (0.522,0.054)rectangle++(0.009,0.009) (0.248,0.544)rectangle++(0.009,0.009) (0.637,0.259)rectangle++(0.009,0.009) (0.303,0.642)rectangle++(0.009,0.009) (0.663,0.306)rectangle++(0.009,0.009)
  (0.836,0.144)rectangle++(0.009,0.009) (0.399,0.587)rectangle++(0.009,0.009) (0.709,0.279)rectangle++(0.009,0.009) (0.338,0.652)rectangle++(0.009,0.009) (0.68,0.311)rectangle++(0.009,0.009)
  (0.324,0.147)rectangle++(0.009,0.009) (0.153,0.068)rectangle++(0.009,0.009) (0.071,0.55)rectangle++(0.009,0.009) (0.032,0.782)rectangle++(0.009,0.009) (0.013,0.893)rectangle++(0.009,0.009)
  (0.004,0.946)rectangle++(0.009,0.009) (0.52,0.452)rectangle++(0.009,0.009) (0.247,0.215)rectangle++(0.009,0.009) (0.116,0.621)rectangle++(0.009,0.009) (0.053,0.296)rectangle++(0.009,0.009)
  (0.543,0.14)rectangle++(0.009,0.009) (0.258,0.065)rectangle++(0.009,0.009) (0.642,0.029)rectangle++(0.009,0.009) (0.306,0.011)rectangle++(0.009,0.009) (0.664,0.003)rectangle++(0.009,0.009)
  (0.317,0.519)rectangle++(0.009,0.009) (0.67,0.247)rectangle++(0.009,0.009) (0.319,0.116)rectangle++(0.009,0.009) (0.151,0.573)rectangle++(0.009,0.009) (0.07,0.793)rectangle++(0.009,0.009)
  (0.031,0.898)rectangle++(0.009,0.009) (0.013,0.429)rectangle++(0.009,0.009) (0.524,0.203)rectangle++(0.009,0.009) (0.249,0.615)rectangle++(0.009,0.009) (0.117,0.293)rectangle++(0.009,0.009)
  (0.054,0.138)rectangle++(0.009,0.009) (0.024,0.584)rectangle++(0.009,0.009) (0.529,0.278)rectangle++(0.009,0.009) (0.772,0.131)rectangle++(0.009,0.009) (0.368,0.581)rectangle++(0.009,0.009)
  (0.694,0.276)rectangle++(0.009,0.009) (0.331,0.13)rectangle++(0.009,0.009) (0.157,0.06)rectangle++(0.009,0.009) (0.593,0.027)rectangle++(0.009,0.009) (0.282,0.53)rectangle++(0.009,0.009)
  (0.653,0.252)rectangle++(0.009,0.009) (0.311,0.639)rectangle++(0.009,0.009) (0.147,0.304)rectangle++(0.009,0.009) (0.068,0.144)rectangle++(0.009,0.009) (0.55,0.067)rectangle++(0.009,0.009)
  (0.782,0.03)rectangle++(0.009,0.009) (0.893,0.012)rectangle++(0.009,0.009) (0.426,0.523)rectangle++(0.009,0.009) (0.202,0.249)rectangle++(0.009,0.009) (0.095,0.117)rectangle++(0.009,0.009)
  (0.043,0.054)rectangle++(0.009,0.009) (0.538,0.024)rectangle++(0.009,0.009) (0.256,0.009)rectangle++(0.009,0.009) (0.641,0.002)rectangle++(0.009,0.009) (0.305,-0.001)rectangle++(0.009,0.009)
  (0.144,-0.003)rectangle++(0.009,0.009) (0.587,-0.004)rectangle++(0.009,0.009) (0.279,0.516)rectangle++(0.009,0.009) (0.132,0.245)rectangle++(0.009,0.009) (0.061,0.635)rectangle++(0.009,0.009)
  (0.547,0.303)rectangle++(0.009,0.009) (0.78,0.143)rectangle++(0.009,0.009) (0.372,0.586)rectangle++(0.009,0.009) (0.176,0.279)rectangle++(0.009,0.009) (0.602,0.132)rectangle++(0.009,0.009)
  (0.287,0.581)rectangle++(0.009,0.009) (0.655,0.276)rectangle++(0.009,0.009) (0.312,0.13)rectangle++(0.009,0.009) (0.148,0.06)rectangle++(0.009,0.009) (0.588,0.027)rectangle++(0.009,0.009)
  (0.28,0.01)rectangle++(0.009,0.009) (0.652,0.003)rectangle++(0.009,0.009) (0.831,-0.001)rectangle++(0.009,0.009) (0.396,0.517)rectangle++(0.009,0.009) (0.188,0.766)rectangle++(0.009,0.009)
  (0.608,0.365)rectangle++(0.009,0.009) (0.289,0.693)rectangle++(0.009,0.009) (0.137,0.33)rectangle++(0.009,0.009) (0.583,0.156)rectangle++(0.009,0.009) (0.278,0.593)rectangle++(0.009,0.009)
  (0.651,0.282)rectangle++(0.009,0.009) (0.31,0.133)rectangle++(0.009,0.009) (0.667,0.062)rectangle++(0.009,0.009) (0.318,0.547)rectangle++(0.009,0.009) (0.15,0.26)rectangle++(0.009,0.009)
  (0.07,0.643)rectangle++(0.009,0.009) (0.031,0.306)rectangle++(0.009,0.009) (0.013,0.665)rectangle++(0.009,0.009) (0.004,0.837)rectangle++(0.009,0.009) (-0.001,0.399)rectangle++(0.009,0.009)
  (-0.003,0.709)rectangle++(0.009,0.009) (-0.004,0.858)rectangle++(0.009,0.009) (-0.004,0.41)rectangle++(0.009,0.009) (-0.004,0.194)rectangle++(0.009,0.009) (0.516,0.091)rectangle++(0.009,0.009)
  (0.245,0.561)rectangle++(0.009,0.009) (0.115,0.267)rectangle++(0.009,0.009) (0.053,0.126)rectangle++(0.009,0.009) (0.023,0.058)rectangle++(0.009,0.009) (0.009,0.026)rectangle++(0.009,0.009)
  (0.522,0.01)rectangle++(0.009,0.009) (0.768,0.002)rectangle++(0.009,0.009) (0.886,-0.001)rectangle++(0.009,0.009) (0.423,-0.003)rectangle++(0.009,0.009) (0.201,-0.004)rectangle++(0.009,0.009)
  (0.094,0.516)rectangle++(0.009,0.009) (0.563,0.245)rectangle++(0.009,0.009) (0.788,0.115)rectangle++(0.009,0.009) (0.896,0.053)rectangle++(0.009,0.009) (0.428,0.543)rectangle++(0.009,0.009)
  (0.203,0.778)rectangle++(0.009,0.009) (0.095,0.891)rectangle++(0.009,0.009) (0.043,0.945)rectangle++(0.009,0.009) (0.538,0.451)rectangle++(0.009,0.009) (0.256,0.214)rectangle++(0.009,0.009)
  (0.121,0.101)rectangle++(0.009,0.009) (0.056,0.566)rectangle++(0.009,0.009) (0.024,0.269)rectangle++(0.009,0.009) (0.529,0.127)rectangle++(0.009,0.009) (0.252,0.579)rectangle++(0.009,0.009)
  (0.118,0.275)rectangle++(0.009,0.009) (0.575,0.13)rectangle++(0.009,0.009) (0.793,0.06)rectangle++(0.009,0.009) (0.379,0.546)rectangle++(0.009,0.009) (0.179,0.78)rectangle++(0.009,0.009)
  (0.084,0.892)rectangle++(0.009,0.009) (0.038,0.426)rectangle++(0.009,0.009) (0.016,0.722)rectangle++(0.009,0.009) (0.525,0.344)rectangle++(0.009,0.009) (0.77,0.163)rectangle++(0.009,0.009)
  (0.367,0.596)rectangle++(0.009,0.009) (0.174,0.804)rectangle++(0.009,0.009) (0.081,0.903)rectangle++(0.009,0.009) (0.037,0.951)rectangle++(0.009,0.009) (0.015,0.974)rectangle++(0.009,0.009)
  (0.525,0.465)rectangle++(0.009,0.009) (0.77,0.221)rectangle++(0.009,0.009) (0.367,0.624)rectangle++(0.009,0.009) (0.694,0.297)rectangle++(0.009,0.009) (0.331,0.14)rectangle++(0.009,0.009)
  (0.676,0.065)rectangle++(0.009,0.009) (0.322,0.549)rectangle++(0.009,0.009) (0.152,0.781)rectangle++(0.009,0.009) (0.071,0.373)rectangle++(0.009,0.009) (0.552,0.177)rectangle++(0.009,0.009)
  (0.262,0.082)rectangle++(0.009,0.009) (0.644,0.037)rectangle++(0.009,0.009) (0.307,0.536)rectangle++(0.009,0.009) (0.665,0.255)rectangle++(0.009,0.009) (0.317,0.12)rectangle++(0.009,0.009)
  (0.15,0.055)rectangle++(0.009,0.009) (0.59,0.024)rectangle++(0.009,0.009) (0.281,0.529)rectangle++(0.009,0.009) (0.652,0.252)rectangle++(0.009,0.009) (0.311,0.118)rectangle++(0.009,0.009)
  (0.147,0.055)rectangle++(0.009,0.009) (0.068,0.024)rectangle++(0.009,0.009) (0.03,0.529)rectangle++(0.009,0.009) (0.012,0.252)rectangle++(0.009,0.009) (0.004,0.638)rectangle++(0.009,0.009)
  (-0.001,0.304)rectangle++(0.009,0.009) (-0.003,0.664)rectangle++(0.009,0.009) (-0.004,0.316)rectangle++(0.009,0.009) (-0.004,0.669)rectangle++(0.009,0.009) (0.516,0.319)rectangle++(0.009,0.009)
  (0.245,0.151)rectangle++(0.009,0.009) (0.115,0.07)rectangle++(0.009,0.009) (0.053,0.551)rectangle++(0.009,0.009) (0.543,0.262)rectangle++(0.009,0.009) (0.778,0.124)rectangle++(0.009,0.009)
  (0.891,0.057)rectangle++(0.009,0.009) (0.425,0.545)rectangle++(0.009,0.009) (0.722,0.259)rectangle++(0.009,0.009) (0.344,0.122)rectangle++(0.009,0.009) (0.683,0.056)rectangle++(0.009,0.009)
  (0.325,0.025)rectangle++(0.009,0.009) (0.154,0.53)rectangle++(0.009,0.009) (0.592,0.252)rectangle++(0.009,0.009) (0.282,0.639)rectangle++(0.009,0.009) (0.133,0.304)rectangle++(0.009,0.009)
  (0.061,0.664)rectangle++(0.009,0.009) (0.547,0.316)rectangle++(0.009,0.009) (0.26,0.669)rectangle++(0.009,0.009) (0.123,0.839)rectangle++(0.009,0.009) (0.577,0.4)rectangle++(0.009,0.009)
  (0.794,0.19)rectangle++(0.009,0.009) (0.379,0.609)rectangle++(0.009,0.009) (0.18,0.29)rectangle++(0.009,0.009) (0.604,0.137)rectangle++(0.009,0.009) (0.288,0.583)rectangle++(0.009,0.009)
  (0.656,0.278)rectangle++(0.009,0.009) (0.832,0.131)rectangle++(0.009,0.009) (0.397,0.581)rectangle++(0.009,0.009) (0.188,0.796)rectangle++(0.009,0.009) (0.088,0.38)rectangle++(0.009,0.009)
  (0.04,0.18)rectangle++(0.009,0.009) (0.537,0.084)rectangle++(0.009,0.009) (0.255,0.038)rectangle++(0.009,0.009) (0.12,0.536)rectangle++(0.009,0.009) (0.575,0.255)rectangle++(0.009,0.009)
  (0.794,0.12)rectangle++(0.009,0.009) (0.379,0.575)rectangle++(0.009,0.009) (0.179,0.794)rectangle++(0.009,0.009) (0.084,0.379)rectangle++(0.009,0.009) (0.558,0.179)rectangle++(0.009,0.009)
  (0.265,0.084)rectangle++(0.009,0.009) (0.645,0.038)rectangle++(0.009,0.009) (0.307,0.016)rectangle++(0.009,0.009) (0.145,0.525)rectangle++(0.009,0.009) (0.587,0.25)rectangle++(0.009,0.009)
  (0.8,0.118)rectangle++(0.009,0.009) (0.381,0.574)rectangle++(0.009,0.009) (0.181,0.793)rectangle++(0.009,0.009) (0.604,0.378)rectangle++(0.009,0.009) (0.808,0.179)rectangle++(0.009,0.009)
  (0.905,0.084)rectangle++(0.009,0.009) (0.432,0.038)rectangle++(0.009,0.009) (0.205,0.016)rectangle++(0.009,0.009) (0.096,0.525)rectangle++(0.009,0.009) (0.044,0.25)rectangle++(0.009,0.009)
  (0.019,0.638)rectangle++(0.009,0.009) (0.007,0.824)rectangle++(0.009,0.009) (0.521,0.393)rectangle++(0.009,0.009) (0.768,0.186)rectangle++(0.009,0.009) (0.366,0.607)rectangle++(0.009,0.009)
  (0.173,0.809)rectangle++(0.009,0.009) (0.081,0.386)rectangle++(0.009,0.009) (0.036,0.703)rectangle++(0.009,0.009) (0.535,0.335)rectangle++(0.009,0.009) (0.255,0.158)rectangle++(0.009,0.009)
  (0.12,0.074)rectangle++(0.009,0.009) (0.055,0.033)rectangle++(0.009,0.009) (0.024,0.534)rectangle++(0.009,0.009) (0.009,0.774)rectangle++(0.009,0.009) (0.522,0.369)rectangle++(0.009,0.009)
  (0.248,0.175)rectangle++(0.009,0.009) (0.117,0.082)rectangle++(0.009,0.009) (0.574,0.037)rectangle++(0.009,0.009) (0.273,0.015)rectangle++(0.009,0.009) (0.129,0.005)rectangle++(0.009,0.009)
  (0.059,0)rectangle++(0.009,0.009) (0.026,0.518)rectangle++(0.009,0.009) (0.01,0.246)rectangle++(0.009,0.009) (0.523,0.116)rectangle++(0.009,0.009) (0.768,0.053)rectangle++(0.009,0.009)
  (0.367,0.023)rectangle++(0.009,0.009) (0.174,0.009)rectangle++(0.009,0.009) (0.081,0.002)rectangle++(0.009,0.009) (0.037,-0.001)rectangle++(0.009,0.009) (0.015,0.517)rectangle++(0.009,0.009)
  (0.005,0.766)rectangle++(0.009,0.009) (0,0.885)rectangle++(0.009,0.009) (-0.002,0.943)rectangle++(0.009,0.009) (-0.003,0.45)rectangle++(0.009,0.009) (-0.004,0.214)rectangle++(0.009,0.009)
  (-0.004,0.62)rectangle++(0.009,0.009) (0.516,0.295)rectangle++(0.009,0.009) (0.245,0.659)rectangle++(0.009,0.009) (0.635,0.314)rectangle++(0.009,0.009) (0.823,0.148)rectangle++(0.009,0.009)
  (0.393,0.589)rectangle++(0.009,0.009) (0.186,0.8)rectangle++(0.009,0.009) (0.087,0.382)rectangle++(0.009,0.009) (0.039,0.181)rectangle++(0.009,0.009) (0.017,0.085)rectangle++(0.009,0.009)
  (0.526,0.038)rectangle++(0.009,0.009) (0.25,0.536)rectangle++(0.009,0.009) (0.638,0.255)rectangle++(0.009,0.009) (0.304,0.12)rectangle++(0.009,0.009) (0.143,0.575)rectangle++(0.009,0.009)
  (0.587,0.274)rectangle++(0.009,0.009) (0.279,0.129)rectangle++(0.009,0.009) (0.132,0.58)rectangle++(0.009,0.009) (0.581,0.276)rectangle++(0.009,0.009) (0.796,0.13)rectangle++(0.009,0.009)
  (0.9,0.06)rectangle++(0.009,0.009) (0.43,0.027)rectangle++(0.009,0.009) (0.724,0.01)rectangle++(0.009,0.009) (0.865,0.003)rectangle++(0.009,0.009) (0.413,-0.001)rectangle++(0.009,0.009)
  (0.716,-0.003)rectangle++(0.009,0.009) (0.861,-0.004)rectangle++(0.009,0.009) (0.411,0.516)rectangle++(0.009,0.009) (0.195,0.765)rectangle++(0.009,0.009) (0.611,0.365)rectangle++(0.009,0.009)
  (0.291,0.693)rectangle++(0.009,0.009) (0.657,0.33)rectangle++(0.009,0.009) (0.833,0.156)rectangle++(0.009,0.009) (0.918,0.073)rectangle++(0.009,0.009) (0.438,0.553)rectangle++(0.009,0.009)
  (0.208,0.263)rectangle++(0.009,0.009) (0.617,0.124)rectangle++(0.009,0.009) (0.814,0.057)rectangle++(0.009,0.009) (0.388,0.545)rectangle++(0.009,0.009) (0.184,0.779)rectangle++(0.009,0.009)
  (0.606,0.372)rectangle++(0.009,0.009) (0.289,0.176)rectangle++(0.009,0.009) (0.656,0.082)rectangle++(0.009,0.009) (0.833,0.037)rectangle++(0.009,0.009) (0.397,0.535)rectangle++(0.009,0.009)
  (0.188,0.775)rectangle++(0.009,0.009) (0.608,0.37)rectangle++(0.009,0.009) (0.81,0.175)rectangle++(0.009,0.009) (0.386,0.602)rectangle++(0.009,0.009) (0.183,0.286)rectangle++(0.009,0.009)
  (0.086,0.135)rectangle++(0.009,0.009) (0.039,0.583)rectangle++(0.009,0.009) (0.016,0.797)rectangle++(0.009,0.009) (0.005,0.9)rectangle++(0.009,0.009) (0,0.43)rectangle++(0.009,0.009)
  (-0.002,0.724)rectangle++(0.009,0.009) (-0.003,0.345)rectangle++(0.009,0.009) (0.516,0.163)rectangle++(0.009,0.009) (0.245,0.076)rectangle++(0.009,0.009) (0.635,0.034)rectangle++(0.009,0.009)
  (0.823,0.014)rectangle++(0.009,0.009) (0.393,0.524)rectangle++(0.009,0.009) (0.186,0.769)rectangle++(0.009,0.009) (0.087,0.367)rectangle++(0.009,0.009) (0.559,0.174)rectangle++(0.009,0.009)
  (0.786,0.081)rectangle++(0.009,0.009) (0.375,0.037)rectangle++(0.009,0.009) (0.178,0.015)rectangle++(0.009,0.009) (0.083,0.525)rectangle++(0.009,0.009) (0.037,0.77)rectangle++(0.009,0.009)
  (0.536,0.367)rectangle++(0.009,0.009) (0.255,0.174)rectangle++(0.009,0.009) (0.12,0.601)rectangle++(0.009,0.009) (0.575,0.286)rectangle++(0.009,0.009) (0.794,0.135)rectangle++(0.009,0.009)
  (0.899,0.062)rectangle++(0.009,0.009) (0.949,0.028)rectangle++(0.009,0.009) (0.453,0.531)rectangle++(0.009,0.009) (0.735,0.253)rectangle++(0.009,0.009) (0.351,0.639)rectangle++(0.009,0.009)
  (0.166,0.824)rectangle++(0.009,0.009) (0.077,0.913)rectangle++(0.009,0.009) (0.035,0.436)rectangle++(0.009,0.009) (0.534,0.207)rectangle++(0.009,0.009) (0.254,0.617)rectangle++(0.009,0.009)
  (0.64,0.294)rectangle++(0.009,0.009) (0.825,0.139)rectangle++(0.009,0.009) (0.394,0.584)rectangle++(0.009,0.009) (0.187,0.278)rectangle++(0.009,0.009) (0.087,0.131)rectangle++(0.009,0.009)
  (0.04,0.061)rectangle++(0.009,0.009) (0.017,0.547)rectangle++(0.009,0.009) (0.526,0.26)rectangle++(0.009,0.009) (0.25,0.123)rectangle++(0.009,0.009) (0.638,0.056)rectangle++(0.009,0.009)
  (0.304,0.025)rectangle++(0.009,0.009) (0.143,0.53)rectangle++(0.009,0.009) (0.067,0.772)rectangle++(0.009,0.009) (0.55,0.368)rectangle++(0.009,0.009) (0.261,0.694)rectangle++(0.009,0.009)
  (0.123,0.331)rectangle++(0.009,0.009) (0.577,0.157)rectangle++(0.009,0.009) (0.275,0.073)rectangle++(0.009,0.009) (0.129,0.033)rectangle++(0.009,0.009) (0.58,0.013)rectangle++(0.009,0.009)
  (0.276,0.004)rectangle++(0.009,0.009) (0.13,0)rectangle++(0.009,0.009) (0.06,-0.003)rectangle++(0.009,0.009) (0.547,-0.004)rectangle++(0.009,0.009) (0.26,0.516)rectangle++(0.009,0.009)
  (0.122,0.245)rectangle++(0.009,0.009) (0.056,0.115)rectangle++(0.009,0.009) (0.025,0.573)rectangle++(0.009,0.009) (0.01,0.793)rectangle++(0.009,0.009) (0.002,0.898)rectangle++(0.009,0.009)
  (-0.001,0.429)rectangle++(0.009,0.009) (-0.003,0.203)rectangle++(0.009,0.009) (0.516,0.095)rectangle++(0.009,0.009) (0.245,0.563)rectangle++(0.009,0.009) (0.635,0.268)rectangle++(0.009,0.009)
  (0.303,0.646)rectangle++(0.009,0.009) (0.143,0.308)rectangle++(0.009,0.009) (0.586,0.145)rectangle++(0.009,0.009) (0.279,0.587)rectangle++(0.009,0.009) (0.132,0.28)rectangle++(0.009,0.009)
  (0.581,0.132)rectangle++(0.009,0.009) (0.276,0.581)rectangle++(0.009,0.009) (0.13,0.797)rectangle++(0.009,0.009) (0.06,0.9)rectangle++(0.009,0.009) (0.027,0.95)rectangle++(0.009,0.009)
  (0.01,0.973)rectangle++(0.009,0.009) (0.523,0.465)rectangle++(0.009,0.009) (0.249,0.221)rectangle++(0.009,0.009) (0.117,0.104)rectangle++(0.009,0.009) (0.054,0.047)rectangle++(0.009,0.009)
  (0.023,0.54)rectangle++(0.009,0.009) (0.529,0.257)rectangle++(0.009,0.009) (0.772,0.121)rectangle++(0.009,0.009) (0.368,0.576)rectangle++(0.009,0.009) (0.174,0.794)rectangle++(0.009,0.009)
  (0.601,0.379)rectangle++(0.009,0.009) (0.806,0.179)rectangle++(0.009,0.009) (0.385,0.084)rectangle++(0.009,0.009) (0.182,0.038)rectangle++(0.009,0.009) (0.605,0.016)rectangle++(0.009,0.009)
  (0.288,0.525)rectangle++(0.009,0.009) (0.656,0.25)rectangle++(0.009,0.009) (0.833,0.118)rectangle++(0.009,0.009) (0.917,0.054)rectangle++(0.009,0.009) (0.958,0.024)rectangle++(0.009,0.009)
  (0.977,0.009)rectangle++(0.009,0.009) (0.467,0.002)rectangle++(0.009,0.009) (0.742,-0.001)rectangle++(0.009,0.009) (0.354,0.517)rectangle++(0.009,0.009) (0.167,0.766)rectangle++(0.009,0.009)
  (0.078,0.365)rectangle++(0.009,0.009) (0.555,0.173)rectangle++(0.009,0.009) (0.264,0.601)rectangle++(0.009,0.009) (0.124,0.286)rectangle++(0.009,0.009) (0.577,0.135)rectangle++(0.009,0.009)
  (0.795,0.062)rectangle++(0.009,0.009) (0.379,0.028)rectangle++(0.009,0.009) (0.18,0.531)rectangle++(0.009,0.009) (0.084,0.253)rectangle++(0.009,0.009) (0.038,0.639)rectangle++(0.009,0.009)
  (0.016,0.824)rectangle++(0.009,0.009) (0.525,0.393)rectangle++(0.009,0.009) (0.25,0.186)rectangle++(0.009,0.009) (0.118,0.607)rectangle++(0.009,0.009) (0.054,0.809)rectangle++(0.009,0.009)
  (0.024,0.386)rectangle++(0.009,0.009) (0.009,0.183)rectangle++(0.009,0.009) (0.002,0.085)rectangle++(0.009,0.009) (0.519,0.039)rectangle++(0.009,0.009) (0.767,0.016)rectangle++(0.009,0.009)
  (0.366,0.005)rectangle++(0.009,0.009) (0.173,0.52)rectangle++(0.009,0.009) (0.081,0.247)rectangle++(0.009,0.009) (0.036,0.116)rectangle++(0.009,0.009) (0.015,0.574)rectangle++(0.009,0.009)
  (0.005,0.793)rectangle++(0.009,0.009) (0,0.898)rectangle++(0.009,0.009) (0.518,0.429)rectangle++(0.009,0.009) (0.246,0.724)rectangle++(0.009,0.009) (0.116,0.865)rectangle++(0.009,0.009)
  (0.053,0.413)rectangle++(0.009,0.009) (0.023,0.196)rectangle++(0.009,0.009) (0.529,0.092)rectangle++(0.009,0.009) (0.251,0.562)rectangle++(0.009,0.009) (0.638,0.267)rectangle++(0.009,0.009)
  (0.304,0.646)rectangle++(0.009,0.009) (0.144,0.308)rectangle++(0.009,0.009) (0.067,0.665)rectangle++(0.009,0.009) (0.55,0.317)rectangle++(0.009,0.009) (0.781,0.15)rectangle++(0.009,0.009)
  (0.893,0.07)rectangle++(0.009,0.009) (0.426,0.551)rectangle++(0.009,0.009) (0.722,0.262)rectangle++(0.009,0.009) (0.344,0.124)rectangle++(0.009,0.009) (0.683,0.057)rectangle++(0.009,0.009)
  (0.845,0.025)rectangle++(0.009,0.009) (0.403,0.01)rectangle++(0.009,0.009) (0.191,0.522)rectangle++(0.009,0.009) (0.089,0.248)rectangle++(0.009,0.009) (0.041,0.637)rectangle++(0.009,0.009)
  (0.537,0.303)rectangle++(0.009,0.009) (0.255,0.663)rectangle++(0.009,0.009) (0.12,0.836)rectangle++(0.009,0.009) (0.575,0.399)rectangle++(0.009,0.009) (0.274,0.709)rectangle++(0.009,0.009)
  (0.129,0.338)rectangle++(0.009,0.009) (0.06,0.16)rectangle++(0.009,0.009) (0.546,0.074)rectangle++(0.009,0.009) (0.26,0.033)rectangle++(0.009,0.009) (0.642,0.014)rectangle++(0.009,0.009)
  (0.306,0.524)rectangle++(0.009,0.009) (0.665,0.249)rectangle++(0.009,0.009) (0.837,0.117)rectangle++(0.009,0.009) (0.399,0.574)rectangle++(0.009,0.009) (0.189,0.793)rectangle++(0.009,0.009)
  (0.089,0.378)rectangle++(0.009,0.009) (0.04,0.699)rectangle++(0.009,0.009) (0.017,0.853)rectangle++(0.009,0.009) (0.006,0.407)rectangle++(0.009,0.009) (0,0.193)rectangle++(0.009,0.009)
  (-0.002,0.09)rectangle++(0.009,0.009) (-0.003,0.041)rectangle++(0.009,0.009) (0.516,0.017)rectangle++(0.009,0.009) (0.245,0.526)rectangle++(0.009,0.009) (0.635,0.25)rectangle++(0.009,0.009)
  (0.303,0.118)rectangle++(0.009,0.009) (0.143,0.054)rectangle++(0.009,0.009) (0.066,0.544)rectangle++(0.009,0.009) (0.549,0.259)rectangle++(0.009,0.009) (0.261,0.122)rectangle++(0.009,0.009)
  (0.123,0.576)rectangle++(0.009,0.009) (0.577,0.274)rectangle++(0.009,0.009) (0.275,0.649)rectangle++(0.009,0.009) (0.649,0.309)rectangle++(0.009,0.009) (0.309,0.146)rectangle++(0.009,0.009)
  (0.146,0.068)rectangle++(0.009,0.009) (0.068,0.03)rectangle++(0.009,0.009) (0.03,0.532)rectangle++(0.009,0.009) (0.012,0.253)rectangle++(0.009,0.009) (0.003,0.639)rectangle++(0.009,0.009)
  (-0.001,0.824)rectangle++(0.009,0.009) (0.517,0.393)rectangle++(0.009,0.009) (0.766,0.186)rectangle++(0.009,0.009) (0.365,0.607)rectangle++(0.009,0.009) (0.693,0.289)rectangle++(0.009,0.009)
  (0.85,0.136)rectangle++(0.009,0.009) (0.406,0.063)rectangle++(0.009,0.009) (0.192,0.028)rectangle++(0.009,0.009) (0.09,0.011)rectangle++(0.009,0.009) (0.561,0.003)rectangle++(0.009,0.009)
  (0.787,-0.001)rectangle++(0.009,0.009) (0.375,0.517)rectangle++(0.009,0.009) (0.698,0.246)rectangle++(0.009,0.009) (0.333,0.636)rectangle++(0.009,0.009) (0.677,0.303)rectangle++(0.009,0.009)
  (0.323,0.663)rectangle++(0.009,0.009) (0.673,0.316)rectangle++(0.009,0.009) (0.321,0.669)rectangle++(0.009,0.009) (0.672,0.319)rectangle++(0.009,0.009) (0.32,0.151)rectangle++(0.009,0.009)
  (0.151,0.07)rectangle++(0.009,0.009) (0.07,0.551)rectangle++(0.009,0.009) (0.031,0.782)rectangle++(0.009,0.009) (0.533,0.373)rectangle++(0.009,0.009) (0.253,0.697)rectangle++(0.009,0.009)
  (0.119,0.332)rectangle++(0.009,0.009) (0.055,0.677)rectangle++(0.009,0.009) (0.544,0.323)rectangle++(0.009,0.009) (0.259,0.153)rectangle++(0.009,0.009) (0.122,0.071)rectangle++(0.009,0.009)
  (0.056,0.032)rectangle++(0.009,0.009) (0.025,0.013)rectangle++(0.009,0.009) (0.009,0.524)rectangle++(0.009,0.009) (0.522,0.249)rectangle++(0.009,0.009) (0.248,0.117)rectangle++(0.009,0.009)
  (0.637,0.054)rectangle++(0.009,0.009) (0.303,0.544)rectangle++(0.009,0.009) (0.143,0.779)rectangle++(0.009,0.009) (0.066,0.371)rectangle++(0.009,0.009) (0.55,0.176)rectangle++(0.009,0.009)
  (0.781,0.082)rectangle++(0.009,0.009) (0.373,0.037)rectangle++(0.009,0.009) (0.697,0.015)rectangle++(0.009,0.009) (0.332,0.005)rectangle++(0.009,0.009) (0.677,0)rectangle++(0.009,0.009)
  (0.323,0.518)rectangle++(0.009,0.009) (0.153,0.766)rectangle++(0.009,0.009) (0.591,0.365)rectangle++(0.009,0.009) (0.281,0.693)rectangle++(0.009,0.009) (0.133,0.85)rectangle++(0.009,0.009)
  (0.061,0.406)rectangle++(0.009,0.009) (0.027,0.712)rectangle++(0.009,0.009) (0.011,0.86)rectangle++(0.009,0.009) (0.523,0.41)rectangle++(0.009,0.009) (0.249,0.195)rectangle++(0.009,0.009)
  (0.637,0.091)rectangle++(0.009,0.009) (0.823,0.041)rectangle++(0.009,0.009) (0.393,0.538)rectangle++(0.009,0.009) (0.186,0.256)rectangle++(0.009,0.009) (0.087,0.64)rectangle++(0.009,0.009)
  (0.039,0.825)rectangle++(0.009,0.009) (0.537,0.394)rectangle++(0.009,0.009) (0.255,0.187)rectangle++(0.009,0.009) (0.12,0.607)rectangle++(0.009,0.009) (0.055,0.289)rectangle++(0.009,0.009)
  (0.024,0.656)rectangle++(0.009,0.009) (0.009,0.833)rectangle++(0.009,0.009) (0.522,0.397)rectangle++(0.009,0.009) (0.248,0.188)rectangle++(0.009,0.009) (0.117,0.088)rectangle++(0.009,0.009)
  (0.054,0.04)rectangle++(0.009,0.009) (0.023,0.017)rectangle++(0.009,0.009) (0.529,0.006)rectangle++(0.009,0.009) (0.772,0)rectangle++(0.009,0.009) (0.368,0.518)rectangle++(0.009,0.009)
  (0.174,0.246)rectangle++(0.009,0.009) (0.081,0.636)rectangle++(0.009,0.009) (0.037,0.303)rectangle++(0.009,0.009) (0.015,0.663)rectangle++(0.009,0.009) (0.005,0.316)rectangle++(0.009,0.009)
  (0.52,0.149)rectangle++(0.009,0.009) (0.247,0.589)rectangle++(0.009,0.009) (0.116,0.801)rectangle++(0.009,0.009) (0.054,0.382)rectangle++(0.009,0.009) (0.543,0.181)rectangle++(0.009,0.009)
  (0.258,0.605)rectangle++(0.009,0.009) (0.122,0.288)rectangle++(0.009,0.009) (0.576,0.136)rectangle++(0.009,0.009) (0.274,0.063)rectangle++(0.009,0.009) (0.129,0.028)rectangle++(0.009,0.009)
  (0.58,0.011)rectangle++(0.009,0.009) (0.276,0.003)rectangle++(0.009,0.009) (0.13,-0.001)rectangle++(0.009,0.009) (0.06,0.517)rectangle++(0.009,0.009) (0.027,0.766)rectangle++(0.009,0.009)
  (0.53,0.365)rectangle++(0.009,0.009) (0.252,0.173)rectangle++(0.009,0.009) (0.119,0.081)rectangle++(0.009,0.009) (0.575,0.036)rectangle++(0.009,0.009) (0.793,0.015)rectangle++(0.009,0.009)
  (0.379,0.005)rectangle++(0.009,0.009) (0.179,0.52)rectangle++(0.009,0.009) (0.604,0.247)rectangle++(0.009,0.009) (0.287,0.116)rectangle++(0.009,0.009) (0.656,0.054)rectangle++(0.009,0.009)
  (0.312,0.543)rectangle++(0.009,0.009) (0.668,0.258)rectangle++(0.009,0.009) (0.838,0.122)rectangle++(0.009,0.009) (0.4,0.056)rectangle++(0.009,0.009) (0.71,0.025)rectangle++(0.009,0.009)
  (0.338,0.529)rectangle++(0.009,0.009) (0.16,0.772)rectangle++(0.009,0.009) (0.594,0.368)rectangle++(0.009,0.009) (0.283,0.694)rectangle++(0.009,0.009) (0.134,0.851)rectangle++(0.009,0.009)
  (0.062,0.406)rectangle++(0.009,0.009) (0.027,0.713)rectangle++(0.009,0.009) (0.011,0.34)rectangle++(0.009,0.009) (0.003,0.681)rectangle++(0.009,0.009) (-0.001,0.844)rectangle++(0.009,0.009)
  (-0.003,0.923)rectangle++(0.009,0.009) (0.516,0.441)rectangle++(0.009,0.009) (0.245,0.729)rectangle++(0.009,0.009) (0.115,0.868)rectangle++(0.009,0.009) (0.573,0.414)rectangle++(0.009,0.009)
  (0.273,0.196)rectangle++(0.009,0.009) (0.649,0.092)rectangle++(0.009,0.009) (0.309,0.562)rectangle++(0.009,0.009) (0.666,0.267)rectangle++(0.009,0.009) (0.317,0.646)rectangle++(0.009,0.009)
  (0.15,0.828)rectangle++(0.009,0.009) (0.59,0.395)rectangle++(0.009,0.009) (0.801,0.187)rectangle++(0.009,0.009) (0.902,0.088)rectangle++(0.009,0.009) (0.431,0.04)rectangle++(0.009,0.009)
  (0.204,0.017)rectangle++(0.009,0.009) (0.096,0.526)rectangle++(0.009,0.009) (0.564,0.25)rectangle++(0.009,0.009) (0.788,0.118)rectangle++(0.009,0.009) (0.376,0.574)rectangle++(0.009,0.009)
  (0.178,0.793)rectangle++(0.009,0.009) (0.083,0.378)rectangle++(0.009,0.009) (0.558,0.179)rectangle++(0.009,0.009) (0.785,0.084)rectangle++(0.009,0.009) (0.375,0.038)rectangle++(0.009,0.009)
  (0.177,0.536)rectangle++(0.009,0.009) (0.603,0.255)rectangle++(0.009,0.009) (0.287,0.64)rectangle++(0.009,0.009) (0.135,0.305)rectangle++(0.009,0.009) (0.063,0.144)rectangle++(0.009,0.009)
  (0.028,0.067)rectangle++(0.009,0.009) (0.531,0.03)rectangle++(0.009,0.009) (0.253,0.012)rectangle++(0.009,0.009) (0.639,0.003)rectangle++(0.009,0.009) (0.824,-0.001)rectangle++(0.009,0.009)
  (0.393,-0.003)rectangle++(0.009,0.009) (0.186,0.516)rectangle++(0.009,0.009) (0.607,0.246)rectangle++(0.009,0.009) (0.289,0.636)rectangle++(0.009,0.009) (0.656,0.303)rectangle++(0.009,0.009)
  (0.313,0.663)rectangle++(0.009,0.009) (0.668,0.316)rectangle++(0.009,0.009) (0.838,0.149)rectangle++(0.009,0.009) (0.4,0.589)rectangle++(0.009,0.009) (0.71,0.281)rectangle++(0.009,0.009)
  (0.338,0.652)rectangle++(0.009,0.009) (0.68,0.311)rectangle++(0.009,0.009) (0.324,0.667)rectangle++(0.009,0.009) (0.153,0.318)rectangle++(0.009,0.009) (0.591,0.15)rectangle++(0.009,0.009)
  (0.281,0.07)rectangle++(0.009,0.009) (0.133,0.551)rectangle++(0.009,0.009) (0.061,0.782)rectangle++(0.009,0.009) (0.027,0.893)rectangle++(0.009,0.009) (0.531,0.426)rectangle++(0.009,0.009)
  (0.252,0.202)rectangle++(0.009,0.009) (0.119,0.615)rectangle++(0.009,0.009) (0.055,0.293)rectangle++(0.009,0.009) (0.024,0.658)rectangle++(0.009,0.009) (0.529,0.314)rectangle++(0.009,0.009)
  (0.252,0.668)rectangle++(0.009,0.009) (0.118,0.838)rectangle++(0.009,0.009) (0.055,0.4)rectangle++(0.009,0.009) (0.544,0.19)rectangle++(0.009,0.009) (0.779,0.089)rectangle++(0.009,0.009)
  (0.371,0.56)rectangle++(0.009,0.009) (0.176,0.787)rectangle++(0.009,0.009) (0.082,0.895)rectangle++(0.009,0.009) (0.557,0.427)rectangle++(0.009,0.009) (0.785,0.203)rectangle++(0.009,0.009)
  (0.374,0.095)rectangle++(0.009,0.009) (0.697,0.043)rectangle++(0.009,0.009) (0.852,0.018)rectangle++(0.009,0.009) (0.407,0.527)rectangle++(0.009,0.009) (0.193,0.25)rectangle++(0.009,0.009)
  (0.09,0.118)rectangle++(0.009,0.009) (0.561,0.054)rectangle++(0.009,0.009) (0.267,0.544)rectangle++(0.009,0.009) (0.126,0.779)rectangle++(0.009,0.009) (0.058,0.371)rectangle++(0.009,0.009)
  (0.026,0.696)rectangle++(0.009,0.009) (0.53,0.332)rectangle++(0.009,0.009) (0.772,0.157)rectangle++(0.009,0.009) (0.368,0.593)rectangle++(0.009,0.009) (0.694,0.282)rectangle++(0.009,0.009)
  (0.851,0.133)rectangle++(0.009,0.009) (0.406,0.582)rectangle++(0.009,0.009) (0.193,0.277)rectangle++(0.009,0.009) (0.61,0.131)rectangle++(0.009,0.009) (0.291,0.58)rectangle++(0.009,0.009)
  (0.137,0.796)rectangle++(0.009,0.009) (0.063,0.38)rectangle++(0.009,0.009) (0.028,0.7)rectangle++(0.009,0.009) (0.011,0.854)rectangle++(0.009,0.009) (0.523,0.407)rectangle++(0.009,0.009)
  (0.769,0.193)rectangle++(0.009,0.009) (0.367,0.09)rectangle++(0.009,0.009) (0.174,0.561)rectangle++(0.009,0.009) (0.081,0.267)rectangle++(0.009,0.009) (0.037,0.646)rectangle++(0.009,0.009)
  (0.015,0.828)rectangle++(0.009,0.009) (0.005,0.915)rectangle++(0.009,0.009) (0,0.957)rectangle++(0.009,0.009) (0.518,0.457)rectangle++(0.009,0.009) (0.246,0.737)rectangle++(0.009,0.009)
  (0.116,0.871)rectangle++(0.009,0.009) (0.573,0.416)rectangle++(0.009,0.009) (0.793,0.197)rectangle++(0.009,0.009) (0.378,0.092)rectangle++(0.009,0.009) (0.179,0.042)rectangle++(0.009,0.009)
  (0.604,0.018)rectangle++(0.009,0.009) (0.807,0.006)rectangle++(0.009,0.009) (0.385,0.521)rectangle++(0.009,0.009) (0.183,0.248)rectangle++(0.009,0.009) (0.605,0.116)rectangle++(0.009,0.009)
  (0.808,0.054)rectangle++(0.009,0.009) (0.386,0.543)rectangle++(0.009,0.009) (0.183,0.778)rectangle++(0.009,0.009) (0.085,0.891)rectangle++(0.009,0.009) (0.039,0.945)rectangle++(0.009,0.009)
  (0.536,0.451)rectangle++(0.009,0.009) (0.255,0.734)rectangle++(0.009,0.009) (0.12,0.35)rectangle++(0.009,0.009) (0.055,0.686)rectangle++(0.009,0.009) (0.544,0.327)rectangle++(0.009,0.009)
  (0.259,0.155)rectangle++(0.009,0.009) (0.122,0.072)rectangle++(0.009,0.009) (0.576,0.032)rectangle++(0.009,0.009) (0.794,0.013)rectangle++(0.009,0.009) (0.899,0.004)rectangle++(0.009,0.009)
  (0.949,0)rectangle++(0.009,0.009) (0.453,-0.003)rectangle++(0.009,0.009) (0.215,-0.004)rectangle++(0.009,0.009) (0.621,-0.004)rectangle++(0.009,0.009) (0.296,0.516)rectangle++(0.009,0.009)
  (0.14,0.765)rectangle++(0.009,0.009) (0.585,0.365)rectangle++(0.009,0.009) (0.278,0.693)rectangle++(0.009,0.009) (0.131,0.33)rectangle++(0.009,0.009) (0.061,0.156)rectangle++(0.009,0.009)
  (0.027,0.073)rectangle++(0.009,0.009) (0.011,0.033)rectangle++(0.009,0.009) (0.003,0.013)rectangle++(0.009,0.009) (0.519,0.004)rectangle++(0.009,0.009) (0.247,0.52)rectangle++(0.009,0.009)
  (0.636,0.247)rectangle++(0.009,0.009) (0.303,0.116)rectangle++(0.009,0.009) (0.663,0.053)rectangle++(0.009,0.009) (0.316,0.543)rectangle++(0.009,0.009) (0.669,0.258)rectangle++(0.009,0.009)
  (0.319,0.642)rectangle++(0.009,0.009) (0.151,0.306)rectangle++(0.009,0.009) (0.59,0.144)rectangle++(0.009,0.009) (0.801,0.067)rectangle++(0.009,0.009) (0.382,0.55)rectangle++(0.009,0.009)
  (0.181,0.782)rectangle++(0.009,0.009) (0.085,0.373)rectangle++(0.009,0.009) (0.038,0.177)rectangle++(0.009,0.009) (0.016,0.082)rectangle++(0.009,0.009) (0.005,0.557)rectangle++(0.009,0.009)
  (0.52,0.265)rectangle++(0.009,0.009) (0.767,0.125)rectangle++(0.009,0.009) (0.366,0.058)rectangle++(0.009,0.009) (0.173,0.545)rectangle++(0.009,0.009) (0.081,0.779)rectangle++(0.009,0.009)
  (0.036,0.372)rectangle++(0.009,0.009) (0.015,0.696)rectangle++(0.009,0.009) (0.525,0.332)rectangle++(0.009,0.009) (0.25,0.157)rectangle++(0.009,0.009) (0.117,0.073)rectangle++(0.009,0.009)
  (0.054,0.553)rectangle++(0.009,0.009) (0.024,0.783)rectangle++(0.009,0.009) (0.529,0.373)rectangle++(0.009,0.009) (0.252,0.697)rectangle++(0.009,0.009) (0.118,0.332)rectangle++(0.009,0.009)
  (0.054,0.677)rectangle++(0.009,0.009) (0.024,0.323)rectangle++(0.009,0.009) (0.009,0.153)rectangle++(0.009,0.009) (0.522,0.071)rectangle++(0.009,0.009) (0.248,0.032)rectangle++(0.009,0.009)
  (0.117,0.533)rectangle++(0.009,0.009) (0.574,0.253)rectangle++(0.009,0.009) (0.793,0.119)rectangle++(0.009,0.009) (0.378,0.575)rectangle++(0.009,0.009) (0.699,0.274)rectangle++(0.009,0.009)
  (0.333,0.129)rectangle++(0.009,0.009) (0.678,0.06)rectangle++(0.009,0.009) (0.843,0.026)rectangle++(0.009,0.009) (0.402,0.53)rectangle++(0.009,0.009) (0.191,0.772)rectangle++(0.009,0.009)
  (0.089,0.368)rectangle++(0.009,0.009) (0.56,0.174)rectangle++(0.009,0.009) (0.267,0.601)rectangle++(0.009,0.009) (0.646,0.286)rectangle++(0.009,0.009) (0.308,0.655)rectangle++(0.009,0.009)
  (0.145,0.832)rectangle++(0.009,0.009) (0.067,0.397)rectangle++(0.009,0.009) (0.03,0.708)rectangle++(0.009,0.009) (0.012,0.338)rectangle++(0.009,0.009) (0.523,0.16)rectangle++(0.009,0.009)
  (0.249,0.074)rectangle++(0.009,0.009) (0.117,0.033)rectangle++(0.009,0.009) (0.054,0.534)rectangle++(0.009,0.009) (0.544,0.254)rectangle++(0.009,0.009) (0.779,0.119)rectangle++(0.009,0.009)
  (0.891,0.055)rectangle++(0.009,0.009) (0.426,0.024)rectangle++(0.009,0.009) (0.202,0.009)rectangle++(0.009,0.009) (0.095,0.522)rectangle++(0.009,0.009) (0.043,0.768)rectangle++(0.009,0.009)
  (0.018,0.886)rectangle++(0.009,0.009) (0.006,0.423)rectangle++(0.009,0.009) (0.001,0.201)rectangle++(0.009,0.009) (0.518,0.094)rectangle++(0.009,0.009) (0.246,0.043)rectangle++(0.009,0.009)
  (0.636,0.018)rectangle++(0.009,0.009) (0.303,0.006)rectangle++(0.009,0.009) (0.143,0.521)rectangle++(0.009,0.009) (0.066,0.248)rectangle++(0.009,0.009) (0.549,0.117)rectangle++(0.009,0.009)
  (0.781,0.054)rectangle++(0.009,0.009) (0.373,0.023)rectangle++(0.009,0.009) (0.177,0.009)rectangle++(0.009,0.009) (0.082,0.002)rectangle++(0.009,0.009) (0.037,-0.001)rectangle++(0.009,0.009)
  (0.536,-0.003)rectangle++(0.009,0.009) (0.775,-0.004)rectangle++(0.009,0.009) (0.89,-0.004)rectangle++(0.009,0.009) (0.425,0.516)rectangle++(0.009,0.009) (0.721,0.245)rectangle++(0.009,0.009)
  (0.344,0.115)rectangle++(0.009,0.009) (0.683,0.053)rectangle++(0.009,0.009) (0.325,0.023)rectangle++(0.009,0.009) (0.154,0.529)rectangle++(0.009,0.009) (0.072,0.251)rectangle++(0.009,0.009)
  (0.032,0.118)rectangle++(0.009,0.009) (0.013,0.574)rectangle++(0.009,0.009) (0.524,0.273)rectangle++(0.009,0.009) (0.769,0.129)rectangle++(0.009,0.009) (0.367,0.06)rectangle++(0.009,0.009)
  (0.174,0.026)rectangle++(0.009,0.009) (0.081,0.53)rectangle++(0.009,0.009) (0.037,0.252)rectangle++(0.009,0.009) (0.015,0.639)rectangle++(0.009,0.009) (0.005,0.304)rectangle++(0.009,0.009)
  (0,0.664)rectangle++(0.009,0.009) (-0.002,0.316)rectangle++(0.009,0.009) (-0.003,0.149)rectangle++(0.009,0.009) (-0.004,0.069)rectangle++(0.009,0.009) (-0.004,0.551)rectangle++(0.009,0.009)
  (-0.004,0.262)rectangle++(0.009,0.009) (0.516,0.123)rectangle++(0.009,0.009) (0.245,0.577)rectangle++(0.009,0.009) (0.115,0.795)rectangle++(0.009,0.009) (0.053,0.379)rectangle++(0.009,0.009)
  (0.543,0.18)rectangle++(0.009,0.009) (0.258,0.084)rectangle++(0.009,0.009) (0.122,0.038)rectangle++(0.009,0.009) (0.056,0.016)rectangle++(0.009,0.009) (0.025,0.005)rectangle++(0.009,0.009)
  (0.529,0)rectangle++(0.009,0.009) (0.252,0.518)rectangle++(0.009,0.009) (0.119,0.766)rectangle++(0.009,0.009) (0.575,0.365)rectangle++(0.009,0.009) (0.793,0.173)rectangle++(0.009,0.009)
  (0.379,0.601)rectangle++(0.009,0.009) (0.179,0.286)rectangle++(0.009,0.009) (0.084,0.655)rectangle++(0.009,0.009) (0.038,0.832)rectangle++(0.009,0.009) (0.536,0.397)rectangle++(0.009,0.009)
  (0.255,0.708)rectangle++(0.009,0.009) (0.64,0.338)rectangle++(0.009,0.009) (0.305,0.68)rectangle++(0.009,0.009) (0.144,0.844)rectangle++(0.009,0.009) (0.067,0.923)rectangle++(0.009,0.009)
  (0.03,0.961)rectangle++(0.009,0.009) (0.532,0.459)rectangle++(0.009,0.009) (0.253,0.218)rectangle++(0.009,0.009) (0.639,0.102)rectangle++(0.009,0.009) (0.304,0.047)rectangle++(0.009,0.009)
  (0.144,0.54)rectangle++(0.009,0.009) (0.067,0.257)rectangle++(0.009,0.009) (0.03,0.641)rectangle++(0.009,0.009) (0.532,0.305)rectangle++(0.009,0.009) (0.253,0.664)rectangle++(0.009,0.009)
  (0.639,0.316)rectangle++(0.009,0.009) (0.304,0.67)rectangle++(0.009,0.009) (0.144,0.319)rectangle++(0.009,0.009) (0.067,0.151)rectangle++(0.009,0.009) (0.03,0.59)rectangle++(0.009,0.009)
  (0.012,0.801)rectangle++(0.009,0.009) (0.003,0.382)rectangle++(0.009,0.009) (0.519,0.181)rectangle++(0.009,0.009) (0.247,0.605)rectangle++(0.009,0.009) (0.116,0.288)rectangle++(0.009,0.009)
  (0.053,0.656)rectangle++(0.009,0.009) (0.023,0.832)rectangle++(0.009,0.009) (0.009,0.917)rectangle++(0.009,0.009) (0.002,0.958)rectangle++(0.009,0.009) (-0.001,0.977)rectangle++(0.009,0.009)
  (-0.003,0.467)rectangle++(0.009,0.009) (-0.004,0.742)rectangle++(0.009,0.009) (-0.004,0.354)rectangle++(0.009,0.009) (-0.004,0.167)rectangle++(0.009,0.009) (-0.004,0.598)rectangle++(0.009,0.009)
  (0.516,0.285)rectangle++(0.009,0.009) (0.765,0.134)rectangle++(0.009,0.009) (0.365,0.582)rectangle++(0.009,0.009) (0.173,0.797)rectangle++(0.009,0.009) (0.081,0.38)rectangle++(0.009,0.009)
  (0.036,0.7)rectangle++(0.009,0.009) (0.015,0.854)rectangle++(0.009,0.009) (0.005,0.927)rectangle++(0.009,0.009) (0,0.443)rectangle++(0.009,0.009) (0.518,0.21)rectangle++(0.009,0.009)
  (0.766,0.099)rectangle++(0.009,0.009) (0.365,0.045)rectangle++(0.009,0.009) (0.173,0.539)rectangle++(0.009,0.009) (0.081,0.256)rectangle++(0.009,0.009) (0.036,0.121)rectangle++(0.009,0.009)
  (0.015,0.576)rectangle++(0.009,0.009) (0.525,0.274)rectangle++(0.009,0.009) (0.25,0.129)rectangle++(0.009,0.009) (0.117,0.58)rectangle++(0.009,0.009) (0.054,0.796)rectangle++(0.009,0.009)
  (0.544,0.38)rectangle++(0.009,0.009) (0.779,0.18)rectangle++(0.009,0.009) (0.371,0.084)rectangle++(0.009,0.009) (0.176,0.038)rectangle++(0.009,0.009) (0.082,0.536)rectangle++(0.009,0.009)
  (0.037,0.775)rectangle++(0.009,0.009) (0.015,0.89)rectangle++(0.009,0.009) (0.005,0.945)rectangle++(0.009,0.009) (0.52,0.451)rectangle++(0.009,0.009) (0.247,0.214)rectangle++(0.009,0.009)
  (0.116,0.62)rectangle++(0.009,0.009) (0.054,0.815)rectangle++(0.009,0.009) (0.023,0.389)rectangle++(0.009,0.009) (0.009,0.184)rectangle++(0.009,0.009) (0.002,0.086)rectangle++(0.009,0.009)
  (-0.001,0.559)rectangle++(0.009,0.009) (0.517,0.266)rectangle++(0.009,0.009) (0.246,0.125)rectangle++(0.009,0.009) (0.636,0.058)rectangle++(0.009,0.009) (0.303,0.545)rectangle++(0.009,0.009)
  (0.143,0.779)rectangle++(0.009,0.009) (0.066,0.892)rectangle++(0.009,0.009) (0.549,0.426)rectangle++(0.009,0.009) (0.261,0.202)rectangle++(0.009,0.009) (0.643,0.095)rectangle++(0.009,0.009)
  (0.826,0.043)rectangle++(0.009,0.009) (0.394,0.538)rectangle++(0.009,0.009) (0.187,0.256)rectangle++(0.009,0.009) (0.087,0.121)rectangle++(0.009,0.009) (0.04,0.056)rectangle++(0.009,0.009)
  (0.017,0.024)rectangle++(0.009,0.009) (0.006,0.009)rectangle++(0.009,0.009) (0.52,0.002)rectangle++(0.009,0.009) (0.247,0.519)rectangle++(0.009,0.009) (0.116,0.767)rectangle++(0.009,0.009)
  (0.054,0.366)rectangle++(0.009,0.009) (0.543,0.173)rectangle++(0.009,0.009) (0.778,0.081)rectangle++(0.009,0.009) (0.891,0.036)rectangle++(0.009,0.009) (0.425,0.015)rectangle++(0.009,0.009)
  (0.202,0.005)rectangle++(0.009,0.009) (0.615,0)rectangle++(0.009,0.009) (0.293,-0.002)rectangle++(0.009,0.009) (0.138,0.517)rectangle++(0.009,0.009) (0.064,0.766)rectangle++(0.009,0.009)
  (0.028,0.885)rectangle++(0.009,0.009) (0.011,0.423)rectangle++(0.009,0.009) (0.003,0.2)rectangle++(0.009,0.009) (0.519,0.094)rectangle++(0.009,0.009) (0.247,0.563)rectangle++(0.009,0.009)
  (0.636,0.268)rectangle++(0.009,0.009) (0.303,0.126)rectangle++(0.009,0.009) (0.663,0.058)rectangle++(0.009,0.009) (0.316,0.026)rectangle++(0.009,0.009) (0.149,0.01)rectangle++(0.009,0.009)
  (0.069,0.002)rectangle++(0.009,0.009) (0.031,-0.001)rectangle++(0.009,0.009) (0.013,0.517)rectangle++(0.009,0.009) (0.524,0.246)rectangle++(0.009,0.009) (0.249,0.636)rectangle++(0.009,0.009)
  (0.637,0.303)rectangle++(0.009,0.009) (0.824,0.143)rectangle++(0.009,0.009) (0.913,0.066)rectangle++(0.009,0.009) (0.956,0.029)rectangle++(0.009,0.009) (0.456,0.532)rectangle++(0.009,0.009)
  (0.217,0.773)rectangle++(0.009,0.009) (0.102,0.369)rectangle++(0.009,0.009) (0.046,0.175)rectangle++(0.009,0.009) (0.54,0.081)rectangle++(0.009,0.009) (0.777,0.037)rectangle++(0.009,0.009)
  (0.891,0.015)rectangle++(0.009,0.009) (0.425,0.005)rectangle++(0.009,0.009) (0.202,0)rectangle++(0.009,0.009) (0.614,-0.002)rectangle++(0.009,0.009) (0.293,0.517)rectangle++(0.009,0.009)
  (0.138,0.766)rectangle++(0.009,0.009) (0.064,0.365)rectangle++(0.009,0.009) (0.548,0.173)rectangle++(0.009,0.009) (0.261,0.081)rectangle++(0.009,0.009) (0.643,0.036)rectangle++(0.009,0.009)
  (0.826,0.015)rectangle++(0.009,0.009) (0.394,0.525)rectangle++(0.009,0.009) (0.707,0.25)rectangle++(0.009,0.009) (0.857,0.117)rectangle++(0.009,0.009) (0.929,0.054)rectangle++(0.009,0.009)
  (0.444,0.024)rectangle++(0.009,0.009) (0.211,0.009)rectangle++(0.009,0.009) (0.619,0.002)rectangle++(0.009,0.009) (0.295,-0.001)rectangle++(0.009,0.009) (0.659,-0.003)rectangle++(0.009,0.009)
  (0.834,-0.004)rectangle++(0.009,0.009) (0.398,0.516)rectangle++(0.009,0.009) (0.189,0.765)rectangle++(0.009,0.009) (0.608,0.365)rectangle++(0.009,0.009) (0.29,0.173)rectangle++(0.009,0.009)
  (0.657,0.081)rectangle++(0.009,0.009) (0.313,0.556)rectangle++(0.009,0.009) (0.148,0.265)rectangle++(0.009,0.009) (0.069,0.645)rectangle++(0.009,0.009) (0.031,0.307)rectangle++(0.009,0.009)
  (0.532,0.145)rectangle++(0.009,0.009) (0.253,0.587)rectangle++(0.009,0.009) (0.119,0.8)rectangle++(0.009,0.009) (0.055,0.901)rectangle++(0.009,0.009) (0.024,0.43)rectangle++(0.009,0.009)
  (0.529,0.204)rectangle++(0.009,0.009) (0.252,0.616)rectangle++(0.009,0.009) (0.638,0.293)rectangle++(0.009,0.009) (0.304,0.138)rectangle++(0.009,0.009) (0.664,0.064)rectangle++(0.009,0.009)
  (0.316,0.028)rectangle++(0.009,0.009) (0.149,0.531)rectangle++(0.009,0.009) (0.589,0.253)rectangle++(0.009,0.009) (0.801,0.119)rectangle++(0.009,0.009) (0.382,0.055)rectangle++(0.009,0.009)
  (0.701,0.024)rectangle++(0.009,0.009) (0.334,0.009)rectangle++(0.009,0.009) (0.158,0.002)rectangle++(0.009,0.009) (0.074,0.519)rectangle++(0.009,0.009) (0.553,0.247)rectangle++(0.009,0.009)
  (0.783,0.116)rectangle++(0.009,0.009) (0.374,0.053)rectangle++(0.009,0.009) (0.177,0.023)rectangle++(0.009,0.009) (0.083,0.009)rectangle++(0.009,0.009) (0.557,0.002)rectangle++(0.009,0.009)
  (0.265,0.519)rectangle++(0.009,0.009) (0.125,0.767)rectangle++(0.009,0.009) (0.578,0.366)rectangle++(0.009,0.009) (0.795,0.173)rectangle++(0.009,0.009) (0.379,0.601)rectangle++(0.009,0.009)
  (0.18,0.806)rectangle++(0.009,0.009) (0.084,0.905)rectangle++(0.009,0.009) (0.038,0.432)rectangle++(0.009,0.009) (0.536,0.205)rectangle++(0.009,0.009) (0.255,0.096)rectangle++(0.009,0.009)
  (0.64,0.044)rectangle++(0.009,0.009) (0.305,0.539)rectangle++(0.009,0.009) (0.664,0.256)rectangle++(0.009,0.009) (0.316,0.121)rectangle++(0.009,0.009) (0.15,0.576)rectangle++(0.009,0.009)
  (0.589,0.274)rectangle++(0.009,0.009) (0.281,0.649)rectangle++(0.009,0.009) (0.132,0.829)rectangle++(0.009,0.009) (0.061,0.916)rectangle++(0.009,0.009) (0.547,0.437)rectangle++(0.009,0.009)
  (0.26,0.208)rectangle++(0.009,0.009) (0.123,0.617)rectangle++(0.009,0.009) (0.056,0.294)rectangle++(0.009,0.009) (0.545,0.139)rectangle++(0.009,0.009) (0.259,0.064)rectangle++(0.009,0.009)
  (0.122,0.029)rectangle++(0.009,0.009) (0.056,0.531)rectangle++(0.009,0.009) (0.025,0.773)rectangle++(0.009,0.009) (0.529,0.369)rectangle++(0.009,0.009) (0.252,0.175)rectangle++(0.009,0.009)
  (0.119,0.081)rectangle++(0.009,0.009) (0.055,0.037)rectangle++(0.009,0.009) (0.024,0.535)rectangle++(0.009,0.009) (0.009,0.255)rectangle++(0.009,0.009) (0.522,0.12)rectangle++(0.009,0.009)
  (0.248,0.575)rectangle++(0.009,0.009) (0.117,0.274)rectangle++(0.009,0.009) (0.574,0.129)rectangle++(0.009,0.009) (0.273,0.06)rectangle++(0.009,0.009) (0.649,0.026)rectangle++(0.009,0.009)
  (0.309,0.01)rectangle++(0.009,0.009) (0.146,0.523)rectangle++(0.009,0.009) (0.068,0.249)rectangle++(0.009,0.009) (0.03,0.637)rectangle++(0.009,0.009) (0.532,0.303)rectangle++(0.009,0.009)
  (0.773,0.143)rectangle++(0.009,0.009) (0.369,0.586)rectangle++(0.009,0.009) (0.175,0.799)rectangle++(0.009,0.009) (0.601,0.381)rectangle++(0.009,0.009) (0.286,0.181)rectangle++(0.009,0.009)
  (0.655,0.084)rectangle++(0.009,0.009) (0.312,0.038)rectangle++(0.009,0.009) (0.147,0.016)rectangle++(0.009,0.009) (0.588,0.005)rectangle++(0.009,0.009) (0.28,0.52)rectangle++(0.009,0.009)
  (0.132,0.767)rectangle++(0.009,0.009) (0.061,0.366)rectangle++(0.009,0.009) (0.027,0.693)rectangle++(0.009,0.009) (0.011,0.85)rectangle++(0.009,0.009) (0.003,0.926)rectangle++(0.009,0.009)
  (-0.001,0.962)rectangle++(0.009,0.009) (0.517,0.459)rectangle++(0.009,0.009) (0.246,0.738)rectangle++(0.009,0.009) (0.636,0.352)rectangle++(0.009,0.009) (0.823,0.167)rectangle++(0.009,0.009)
  (0.913,0.078)rectangle++(0.009,0.009) (0.956,0.035)rectangle++(0.009,0.009) (0.976,0.014)rectangle++(0.009,0.009) (0.986,0.005)rectangle++(0.009,0.009) (0.471,0)rectangle++(0.009,0.009)
  (0.224,0.518)rectangle++(0.009,0.009) (0.105,0.246)rectangle++(0.009,0.009) (0.048,0.116)rectangle++(0.009,0.009) (0.021,0.573)rectangle++(0.009,0.009) (0.008,0.793)rectangle++(0.009,0.009)
  (0.001,0.898)rectangle++(0.009,0.009) (0.518,0.429)rectangle++(0.009,0.009) (0.246,0.203)rectangle++(0.009,0.009) (0.636,0.095)rectangle++(0.009,0.009) (0.303,0.043)rectangle++(0.009,0.009)
  (0.663,0.019)rectangle++(0.009,0.009) (0.316,0.007)rectangle++(0.009,0.009) (0.669,0.001)rectangle++(0.009,0.009) (0.319,0.518)rectangle++(0.009,0.009) (0.151,0.766)rectangle++(0.009,0.009)
  (0.07,0.365)rectangle++(0.009,0.009) (0.551,0.173)rectangle++(0.009,0.009) (0.262,0.601)rectangle++(0.009,0.009) (0.124,0.806)rectangle++(0.009,0.009) (0.577,0.385)rectangle++(0.009,0.009)
  (0.275,0.702)rectangle++(0.009,0.009) (0.129,0.335)rectangle++(0.009,0.009) (0.06,0.158)rectangle++(0.009,0.009) (0.026,0.594)rectangle++(0.009,0.009) (0.01,0.803)rectangle++(0.009,0.009)
  (0.003,0.383)rectangle++(0.009,0.009) (-0.001,0.701)rectangle++(0.009,0.009) (-0.003,0.854)rectangle++(0.009,0.009) (-0.004,0.408)rectangle++(0.009,0.009) (0.516,0.193)rectangle++(0.009,0.009)
  (0.245,0.09)rectangle++(0.009,0.009) (0.115,0.561)rectangle++(0.009,0.009) (0.573,0.267)rectangle++(0.009,0.009) (0.273,0.646)rectangle++(0.009,0.009) (0.129,0.828)rectangle++(0.009,0.009)
  (0.059,0.395)rectangle++(0.009,0.009) (0.546,0.187)rectangle++(0.009,0.009) (0.78,0.088)rectangle++(0.009,0.009) (0.892,0.04)rectangle++(0.009,0.009) (0.426,0.537)rectangle++(0.009,0.009)
  (0.202,0.255)rectangle++(0.009,0.009) (0.095,0.64)rectangle++(0.009,0.009) (0.043,0.825)rectangle++(0.009,0.009) (0.018,0.394)rectangle++(0.009,0.009) (0.006,0.707)rectangle++(0.009,0.009)
  (0.521,0.337)rectangle++(0.009,0.009) (0.248,0.679)rectangle++(0.009,0.009) (0.117,0.324)rectangle++(0.009,0.009) (0.054,0.673)rectangle++(0.009,0.009) (0.023,0.841)rectangle++(0.009,0.009)
  (0.009,0.921)rectangle++(0.009,0.009) (0.522,0.44)rectangle++(0.009,0.009) (0.768,0.209)rectangle++(0.009,0.009) (0.366,0.098)rectangle++(0.009,0.009) (0.174,0.045)rectangle++(0.009,0.009)
  (0.081,0.539)rectangle++(0.009,0.009) (0.037,0.776)rectangle++(0.009,0.009) (0.015,0.89)rectangle++(0.009,0.009) (0.005,0.425)rectangle++(0.009,0.009) (0,0.202)rectangle++(0.009,0.009)
  (0.518,0.094)rectangle++(0.009,0.009) (0.766,0.043)rectangle++(0.009,0.009) (0.365,0.538)rectangle++(0.009,0.009) (0.173,0.776)rectangle++(0.009,0.009) (0.081,0.89)rectangle++(0.009,0.009)
  (0.036,0.945)rectangle++(0.009,0.009) (0.015,0.971)rectangle++(0.009,0.009) (0.005,0.984)rectangle++(0.009,0.009) (0,0.99)rectangle++(0.009,0.009) (-0.002,0.473)rectangle++(0.009,0.009)
  (0.517,0.225)rectangle++(0.009,0.009) (0.246,0.625)rectangle++(0.009,0.009) (0.116,0.298)rectangle++(0.009,0.009) (0.573,0.141)rectangle++(0.009,0.009) (0.273,0.065)rectangle++(0.009,0.009)
  (0.129,0.549)rectangle++(0.009,0.009) (0.059,0.781)rectangle++(0.009,0.009) (0.026,0.373)rectangle++(0.009,0.009) (0.01,0.177)rectangle++(0.009,0.009) (0.003,0.082)rectangle++(0.009,0.009)
  (-0.001,0.557)rectangle++(0.009,0.009) (-0.003,0.265)rectangle++(0.009,0.009) (-0.004,0.645)rectangle++(0.009,0.009) (-0.004,0.827)rectangle++(0.009,0.009) (0.516,0.395)rectangle++(0.009,0.009)
  (0.245,0.707)rectangle++(0.009,0.009) (0.115,0.857)rectangle++(0.009,0.009) (0.573,0.409)rectangle++(0.009,0.009) (0.793,0.194)rectangle++(0.009,0.009) (0.378,0.091)rectangle++(0.009,0.009)
  (0.179,0.561)rectangle++(0.009,0.009) (0.084,0.787)rectangle++(0.009,0.009) (0.038,0.375)rectangle++(0.009,0.009) (0.016,0.178)rectangle++(0.009,0.009) (0.005,0.083)rectangle++(0.009,0.009)
  (0,0.558)rectangle++(0.009,0.009) (-0.002,0.265)rectangle++(0.009,0.009) (-0.003,0.125)rectangle++(0.009,0.009) (-0.004,0.578)rectangle++(0.009,0.009) (-0.004,0.275)rectangle++(0.009,0.009)
  (0.516,0.13)rectangle++(0.009,0.009) (0.245,0.58)rectangle++(0.009,0.009) (0.115,0.796)rectangle++(0.009,0.009) (0.053,0.38)rectangle++(0.009,0.009) (0.023,0.18)rectangle++(0.009,0.009)
  (0.529,0.084)rectangle++(0.009,0.009) (0.771,0.038)rectangle++(0.009,0.009) (0.368,0.536)rectangle++(0.009,0.009) (0.174,0.775)rectangle++(0.009,0.009) (0.081,0.89)rectangle++(0.009,0.009)
  (0.037,0.945)rectangle++(0.009,0.009) (0.015,0.451)rectangle++(0.009,0.009) (0.005,0.734)rectangle++(0.009,0.009) (0,0.87)rectangle++(0.009,0.009) (0.518,0.415)rectangle++(0.009,0.009)
  (0.246,0.197)rectangle++(0.009,0.009) (0.116,0.092)rectangle++(0.009,0.009) (0.053,0.562)rectangle++(0.009,0.009) (0.543,0.267)rectangle++(0.009,0.009) (0.778,0.126)rectangle++(0.009,0.009)
  (0.371,0.058)rectangle++(0.009,0.009) (0.176,0.026)rectangle++(0.009,0.009) (0.082,0.01)rectangle++(0.009,0.009) (0.037,0.522)rectangle++(0.009,0.009) (0.535,0.248)rectangle++(0.009,0.009)
  (0.775,0.117)rectangle++(0.009,0.009) (0.37,0.054)rectangle++(0.009,0.009) (0.695,0.023)rectangle++(0.009,0.009) (0.851,0.009)rectangle++(0.009,0.009) (0.406,0.522)rectangle++(0.009,0.009)
  (0.193,0.768)rectangle++(0.009,0.009) (0.61,0.366)rectangle++(0.009,0.009) (0.291,0.694)rectangle++(0.009,0.009) (0.657,0.331)rectangle++(0.009,0.009) (0.313,0.676)rectangle++(0.009,0.009)
  (0.148,0.322)rectangle++(0.009,0.009) (0.069,0.672)rectangle++(0.009,0.009) (0.551,0.32)rectangle++(0.009,0.009) (0.262,0.671)rectangle++(0.009,0.009) (0.123,0.32)rectangle++(0.009,0.009)
  (0.057,0.151)rectangle++(0.009,0.009) (0.545,0.07)rectangle++(0.009,0.009) (0.259,0.551)rectangle++(0.009,0.009) (0.122,0.782)rectangle++(0.009,0.009) (0.056,0.373)rectangle++(0.009,0.009)
  (0.545,0.177)rectangle++(0.009,0.009) (0.259,0.603)rectangle++(0.009,0.009) (0.122,0.807)rectangle++(0.009,0.009) (0.576,0.385)rectangle++(0.009,0.009) (0.794,0.182)rectangle++(0.009,0.009)
  (0.379,0.085)rectangle++(0.009,0.009) (0.7,0.039)rectangle++(0.009,0.009) (0.853,0.016)rectangle++(0.009,0.009) (0.407,0.525)rectangle++(0.009,0.009) (0.193,0.77)rectangle++(0.009,0.009)
  (0.09,0.887)rectangle++(0.009,0.009) (0.041,0.424)rectangle++(0.009,0.009) (0.017,0.721)rectangle++(0.009,0.009) (0.006,0.344)rectangle++(0.009,0.009) (0.001,0.683)rectangle++(0.009,0.009)
  (-0.002,0.325)rectangle++(0.009,0.009) (-0.003,0.154)rectangle++(0.009,0.009) (-0.004,0.071)rectangle++(0.009,0.009) (0.516,0.032)rectangle++(0.009,0.009) (0.765,0.013)rectangle++(0.009,0.009)
  (0.365,0.004)rectangle++(0.009,0.009) (0.173,0)rectangle++(0.009,0.009) (0.081,-0.003)rectangle++(0.009,0.009) (0.556,-0.004)rectangle++(0.009,0.009) (0.265,-0.004)rectangle++(0.009,0.009)
  (0.125,-0.004)rectangle++(0.009,0.009) (0.058,-0.004)rectangle++(0.009,0.009) (0.025,-0.004)rectangle++(0.009,0.009) (0.01,0.516)rectangle++(0.009,0.009) (0.522,0.245)rectangle++(0.009,0.009)
  (0.248,0.115)rectangle++(0.009,0.009) (0.117,0.573)rectangle++(0.009,0.009) (0.574,0.273)rectangle++(0.009,0.009) (0.793,0.129)rectangle++(0.009,0.009) (0.898,0.059)rectangle++(0.009,0.009)
  (0.429,0.026)rectangle++(0.009,0.009) (0.724,0.01)rectangle++(0.009,0.009) (0.865,0.003)rectangle++(0.009,0.009) (0.413,0.519)rectangle++(0.009,0.009) (0.716,0.247)rectangle++(0.009,0.009)
  (0.341,0.116)rectangle++(0.009,0.009) (0.681,0.053)rectangle++(0.009,0.009) (0.325,0.023)rectangle++(0.009,0.009) (0.154,0.009)rectangle++(0.009,0.009) (0.591,0.002)rectangle++(0.009,0.009)
  (0.282,0.519)rectangle++(0.009,0.009) (0.133,0.767)rectangle++(0.009,0.009) (0.581,0.366)rectangle++(0.009,0.009) (0.277,0.173)rectangle++(0.009,0.009) (0.13,0.081)rectangle++(0.009,0.009)
  (0.06,0.556)rectangle++(0.009,0.009) (0.547,0.265)rectangle++(0.009,0.009) (0.78,0.125)rectangle++(0.009,0.009) (0.892,0.058)rectangle++(0.009,0.009) (0.946,0.025)rectangle++(0.009,0.009)
  (0.972,0.01)rectangle++(0.009,0.009) (0.464,0.002)rectangle++(0.009,0.009) (0.22,-0.001)rectangle++(0.009,0.009) (0.623,-0.003)rectangle++(0.009,0.009) (0.297,0.516)rectangle++(0.009,0.009)
  (0.14,0.245)rectangle++(0.009,0.009) (0.585,0.115)rectangle++(0.009,0.009) (0.798,0.053)rectangle++(0.009,0.009) (0.381,0.543)rectangle++(0.009,0.009) (0.7,0.258)rectangle++(0.009,0.009)
  (0.854,0.122)rectangle++(0.009,0.009) (0.408,0.056)rectangle++(0.009,0.009) (0.193,0.545)rectangle++(0.009,0.009) (0.09,0.259)rectangle++(0.009,0.009) (0.561,0.122)rectangle++(0.009,0.009)
  (0.267,0.576)rectangle++(0.009,0.009) (0.126,0.794)rectangle++(0.009,0.009) (0.058,0.899)rectangle++(0.009,0.009) (0.546,0.429)rectangle++(0.009,0.009) (0.78,0.204)rectangle++(0.009,0.009)
  (0.892,0.095)rectangle++(0.009,0.009) (0.426,0.043)rectangle++(0.009,0.009) (0.202,0.539)rectangle++(0.009,0.009) (0.095,0.256)rectangle++(0.009,0.009) (0.043,0.641)rectangle++(0.009,0.009)
  (0.018,0.825)rectangle++(0.009,0.009) (0.006,0.394)rectangle++(0.009,0.009) (0.001,0.187)rectangle++(0.009,0.009) (-0.002,0.607)rectangle++(0.009,0.009) (-0.003,0.289)rectangle++(0.009,0.009)
  (-0.004,0.656)rectangle++(0.009,0.009) (0.516,0.313)rectangle++(0.009,0.009) (0.245,0.668)rectangle++(0.009,0.009) (0.635,0.318)rectangle++(0.009,0.009) (0.303,0.67)rectangle++(0.009,0.009)
  (0.143,0.319)rectangle++(0.009,0.009) (0.586,0.151)rectangle++(0.009,0.009) (0.279,0.07)rectangle++(0.009,0.009) (0.652,0.031)rectangle++(0.009,0.009) (0.31,0.013)rectangle++(0.009,0.009)
  (0.147,0.524)rectangle++(0.009,0.009) (0.068,0.769)rectangle++(0.009,0.009) (0.03,0.367)rectangle++(0.009,0.009) (0.012,0.174)rectangle++(0.009,0.009) (0.004,0.601)rectangle++(0.009,0.009)
  (-0.001,0.286)rectangle++(0.009,0.009) (-0.003,0.655)rectangle++(0.009,0.009) (-0.004,0.312)rectangle++(0.009,0.009) (-0.004,0.667)rectangle++(0.009,0.009) (-0.004,0.838)rectangle++(0.009,0.009)
  (-0.004,0.92)rectangle++(0.009,0.009) (-0.004,0.959)rectangle++(0.009,0.009) (0.516,0.458)rectangle++(0.009,0.009) (0.765,0.218)rectangle++(0.009,0.009) (0.365,0.102)rectangle++(0.009,0.009)
  (0.173,0.047)rectangle++(0.009,0.009) (0.601,0.02)rectangle++(0.009,0.009) (0.806,0.007)rectangle++(0.009,0.009) (0.385,0.521)rectangle++(0.009,0.009) (0.182,0.248)rectangle++(0.009,0.009);
\fill[nuage]
  (0.085,0.637)rectangle++(0.009,0.009) (0.039,0.823)rectangle++(0.009,0.009) (0.016,0.913)rectangle++(0.009,0.009) (0.005,0.956)rectangle++(0.009,0.009) (0,0.976)rectangle++(0.009,0.009)
  (-0.002,0.986)rectangle++(0.009,0.009) (-0.003,0.471)rectangle++(0.009,0.009) (-0.004,0.224)rectangle++(0.009,0.009) (-0.004,0.105)rectangle++(0.009,0.009) (0.516,0.048)rectangle++(0.009,0.009)
  (0.245,0.541)rectangle++(0.009,0.009) (0.635,0.257)rectangle++(0.009,0.009) (0.823,0.121)rectangle++(0.009,0.009) (0.393,0.576)rectangle++(0.009,0.009) (0.186,0.274)rectangle++(0.009,0.009)
  (0.607,0.129)rectangle++(0.009,0.009) (0.289,0.06)rectangle++(0.009,0.009) (0.136,0.026)rectangle++(0.009,0.009) (0.583,0.01)rectangle++(0.009,0.009) (0.278,0.003)rectangle++(0.009,0.009)
  (0.131,-0.001)rectangle++(0.009,0.009) (0.06,0.517)rectangle++(0.009,0.009) (0.547,0.246)rectangle++(0.009,0.009) (0.26,0.116)rectangle++(0.009,0.009) (0.642,0.053)rectangle++(0.009,0.009)
  (0.306,0.543)rectangle++(0.009,0.009) (0.665,0.258)rectangle++(0.009,0.009) (0.837,0.122)rectangle++(0.009,0.009) (0.399,0.576)rectangle++(0.009,0.009) (0.189,0.274)rectangle++(0.009,0.009)
  (0.609,0.129)rectangle++(0.009,0.009) (0.29,0.06)rectangle++(0.009,0.009) (0.137,0.026)rectangle++(0.009,0.009) (0.063,0.01)rectangle++(0.009,0.009) (0.028,0.523)rectangle++(0.009,0.009)
  (0.531,0.249)rectangle++(0.009,0.009) (0.253,0.637)rectangle++(0.009,0.009) (0.119,0.303)rectangle++(0.009,0.009) (0.055,0.143)rectangle++(0.009,0.009) (0.024,0.586)rectangle++(0.009,0.009)
  (0.009,0.279)rectangle++(0.009,0.009) (0.002,0.652)rectangle++(0.009,0.009) (-0.001,0.83)rectangle++(0.009,0.009) (-0.003,0.396)rectangle++(0.009,0.009) (0.516,0.188)rectangle++(0.009,0.009)
  (0.245,0.608)rectangle++(0.009,0.009) (0.115,0.809)rectangle++(0.009,0.009) (0.053,0.386)rectangle++(0.009,0.009) (0.543,0.183)rectangle++(0.009,0.009) (0.258,0.606)rectangle++(0.009,0.009)
  (0.122,0.808)rectangle++(0.009,0.009) (0.056,0.906)rectangle++(0.009,0.009) (0.545,0.432)rectangle++(0.009,0.009) (0.259,0.725)rectangle++(0.009,0.009) (0.122,0.346)rectangle++(0.009,0.009)
  (0.576,0.164)rectangle++(0.009,0.009) (0.274,0.076)rectangle++(0.009,0.009) (0.649,0.034)rectangle++(0.009,0.009) (0.829,0.014)rectangle++(0.009,0.009) (0.396,0.524)rectangle++(0.009,0.009)
  (0.708,0.249)rectangle++(0.009,0.009) (0.337,0.117)rectangle++(0.009,0.009) (0.16,0.054)rectangle++(0.009,0.009) (0.074,0.544)rectangle++(0.009,0.009) (0.553,0.259)rectangle++(0.009,0.009)
  (0.783,0.122)rectangle++(0.009,0.009) (0.374,0.576)rectangle++(0.009,0.009) (0.697,0.274)rectangle++(0.009,0.009) (0.332,0.649)rectangle++(0.009,0.009) (0.157,0.309)rectangle++(0.009,0.009)
  (0.593,0.146)rectangle++(0.009,0.009) (0.802,0.068)rectangle++(0.009,0.009) (0.903,0.03)rectangle++(0.009,0.009) (0.431,0.532)rectangle++(0.009,0.009) (0.205,0.773)rectangle++(0.009,0.009)
  (0.096,0.369)rectangle++(0.009,0.009) (0.044,0.175)rectangle++(0.009,0.009) (0.019,0.081)rectangle++(0.009,0.009) (0.007,0.037)rectangle++(0.009,0.009) (0.001,0.535)rectangle++(0.009,0.009)
  (-0.002,0.255)rectangle++(0.009,0.009) (-0.003,0.12)rectangle++(0.009,0.009) (0.516,0.055)rectangle++(0.009,0.009) (0.245,0.024)rectangle++(0.009,0.009) (0.115,0.529)rectangle++(0.009,0.009)
  (0.053,0.772)rectangle++(0.009,0.009) (0.023,0.368)rectangle++(0.009,0.009) (0.009,0.174)rectangle++(0.009,0.009) (0.002,0.081)rectangle++(0.009,0.009) (0.519,0.037)rectangle++(0.009,0.009)
  (0.767,0.015)rectangle++(0.009,0.009) (0.886,0.005)rectangle++(0.009,0.009) (0.423,0.52)rectangle++(0.009,0.009) (0.721,0.247)rectangle++(0.009,0.009) (0.344,0.636)rectangle++(0.009,0.009)
  (0.683,0.303)rectangle++(0.009,0.009) (0.325,0.143)rectangle++(0.009,0.009) (0.674,0.066)rectangle++(0.009,0.009) (0.841,0.03)rectangle++(0.009,0.009) (0.401,0.532)rectangle++(0.009,0.009)
  (0.71,0.253)rectangle++(0.009,0.009) (0.339,0.119)rectangle++(0.009,0.009) (0.68,0.055)rectangle++(0.009,0.009) (0.844,0.024)rectangle++(0.009,0.009) (0.923,0.009)rectangle++(0.009,0.009)
  (0.441,0.522)rectangle++(0.009,0.009) (0.209,0.248)rectangle++(0.009,0.009) (0.098,0.117)rectangle++(0.009,0.009) (0.045,0.054)rectangle++(0.009,0.009) (0.019,0.023)rectangle++(0.009,0.009)
  (0.007,0.529)rectangle++(0.009,0.009) (0.001,0.772)rectangle++(0.009,0.009) (-0.002,0.368)rectangle++(0.009,0.009) (-0.003,0.694)rectangle++(0.009,0.009) (-0.004,0.851)rectangle++(0.009,0.009)
  (0.516,0.406)rectangle++(0.009,0.009) (0.245,0.193)rectangle++(0.009,0.009) (0.115,0.61)rectangle++(0.009,0.009) (0.053,0.811)rectangle++(0.009,0.009) (0.543,0.387)rectangle++(0.009,0.009)
  (0.258,0.703)rectangle++(0.009,0.009) (0.122,0.855)rectangle++(0.009,0.009) (0.576,0.408)rectangle++(0.009,0.009) (0.274,0.714)rectangle++(0.009,0.009) (0.129,0.86)rectangle++(0.009,0.009)
  (0.58,0.411)rectangle++(0.009,0.009) (0.276,0.715)rectangle++(0.009,0.009) (0.13,0.341)rectangle++(0.009,0.009) (0.06,0.681)rectangle++(0.009,0.009) (0.027,0.845)rectangle++(0.009,0.009)
  (0.01,0.923)rectangle++(0.009,0.009) (0.523,0.441)rectangle++(0.009,0.009) (0.249,0.209)rectangle++(0.009,0.009) (0.117,0.618)rectangle++(0.009,0.009) (0.054,0.294)rectangle++(0.009,0.009)
  (0.023,0.139)rectangle++(0.009,0.009) (0.009,0.584)rectangle++(0.009,0.009) (0.002,0.278)rectangle++(0.009,0.009) (-0.001,0.651)rectangle++(0.009,0.009) (0.517,0.31)rectangle++(0.009,0.009)
  (0.246,0.147)rectangle++(0.009,0.009) (0.116,0.588)rectangle++(0.009,0.009) (0.573,0.28)rectangle++(0.009,0.009) (0.273,0.652)rectangle++(0.009,0.009) (0.129,0.831)rectangle++(0.009,0.009)
  (0.579,0.396)rectangle++(0.009,0.009) (0.796,0.188)rectangle++(0.009,0.009) (0.38,0.088)rectangle++(0.009,0.009) (0.18,0.56)rectangle++(0.009,0.009) (0.084,0.266)rectangle++(0.009,0.009)
  (0.038,0.646)rectangle++(0.009,0.009) (0.536,0.308)rectangle++(0.009,0.009) (0.255,0.145)rectangle++(0.009,0.009) (0.12,0.587)rectangle++(0.009,0.009) (0.055,0.8)rectangle++(0.009,0.009)
  (0.024,0.901)rectangle++(0.009,0.009) (0.009,0.43)rectangle++(0.009,0.009) (0.002,0.204)rectangle++(0.009,0.009) (0.519,0.096)rectangle++(0.009,0.009) (0.767,0.044)rectangle++(0.009,0.009)
  (0.366,0.539)rectangle++(0.009,0.009) (0.173,0.256)rectangle++(0.009,0.009) (0.081,0.641)rectangle++(0.009,0.009) (0.036,0.825)rectangle++(0.009,0.009) (0.535,0.394)rectangle++(0.009,0.009)
  (0.775,0.187)rectangle++(0.009,0.009) (0.889,0.087)rectangle++(0.009,0.009) (0.425,0.04)rectangle++(0.009,0.009) (0.201,0.017)rectangle++(0.009,0.009) (0.614,0.006)rectangle++(0.009,0.009)
  (0.293,0.52)rectangle++(0.009,0.009) (0.658,0.247)rectangle++(0.009,0.009) (0.834,0.116)rectangle++(0.009,0.009) (0.918,0.054)rectangle++(0.009,0.009) (0.958,0.023)rectangle++(0.009,0.009)
  (0.978,0.009)rectangle++(0.009,0.009) (0.987,0.002)rectangle++(0.009,0.009) (0.991,-0.001)rectangle++(0.009,0.009) (0.994,-0.003)rectangle++(0.009,0.009) (0.475,-0.004)rectangle++(0.009,0.009)
  (0.745,-0.004)rectangle++(0.009,0.009) (0.875,-0.004)rectangle++(0.009,0.009) (0.938,-0.004)rectangle++(0.009,0.009) (0.448,0.516)rectangle++(0.009,0.009) (0.733,0.245)rectangle++(0.009,0.009)
  (0.869,0.115)rectangle++(0.009,0.009) (0.415,0.053)rectangle++(0.009,0.009) (0.197,0.023)rectangle++(0.009,0.009) (0.612,0.009)rectangle++(0.009,0.009) (0.291,0.522)rectangle++(0.009,0.009)
  (0.138,0.768)rectangle++(0.009,0.009) (0.064,0.886)rectangle++(0.009,0.009) (0.028,0.943)rectangle++(0.009,0.009) (0.011,0.97)rectangle++(0.009,0.009) (0.523,0.463)rectangle++(0.009,0.009)
  (0.249,0.22)rectangle++(0.009,0.009) (0.117,0.623)rectangle++(0.009,0.009) (0.574,0.297)rectangle++(0.009,0.009) (0.273,0.66)rectangle++(0.009,0.009) (0.649,0.315)rectangle++(0.009,0.009)
  (0.829,0.149)rectangle++(0.009,0.009) (0.396,0.589)rectangle++(0.009,0.009) (0.188,0.28)rectangle++(0.009,0.009) (0.608,0.132)rectangle++(0.009,0.009) (0.289,0.581)rectangle++(0.009,0.009)
  (0.137,0.797)rectangle++(0.009,0.009) (0.063,0.38)rectangle++(0.009,0.009) (0.028,0.7)rectangle++(0.009,0.009) (0.011,0.854)rectangle++(0.009,0.009) (0.523,0.407)rectangle++(0.009,0.009)
  (0.249,0.193)rectangle++(0.009,0.009) (0.117,0.09)rectangle++(0.009,0.009) (0.054,0.041)rectangle++(0.009,0.009) (0.024,0.537)rectangle++(0.009,0.009) (0.009,0.256)rectangle++(0.009,0.009)
  (0.002,0.64)rectangle++(0.009,0.009) (-0.001,0.305)rectangle++(0.009,0.009) (0.517,0.144)rectangle++(0.009,0.009) (0.766,0.067)rectangle++(0.009,0.009) (0.365,0.03)rectangle++(0.009,0.009)
  (0.693,0.012)rectangle++(0.009,0.009) (0.33,0.523)rectangle++(0.009,0.009) (0.156,0.249)rectangle++(0.009,0.009) (0.073,0.117)rectangle++(0.009,0.009) (0.553,0.054)rectangle++(0.009,0.009)
  (0.783,0.024)rectangle++(0.009,0.009) (0.373,0.529)rectangle++(0.009,0.009) (0.177,0.252)rectangle++(0.009,0.009) (0.083,0.118)rectangle++(0.009,0.009) (0.557,0.054)rectangle++(0.009,0.009)
  (0.265,0.024)rectangle++(0.009,0.009) (0.125,0.529)rectangle++(0.009,0.009) (0.058,0.772)rectangle++(0.009,0.009) (0.025,0.888)rectangle++(0.009,0.009) (0.01,0.424)rectangle++(0.009,0.009)
  (0.522,0.201)rectangle++(0.009,0.009) (0.248,0.614)rectangle++(0.009,0.009) (0.117,0.812)rectangle++(0.009,0.009) (0.574,0.388)rectangle++(0.009,0.009) (0.273,0.184)rectangle++(0.009,0.009)
  (0.129,0.606)rectangle++(0.009,0.009) (0.059,0.288)rectangle++(0.009,0.009) (0.026,0.136)rectangle++(0.009,0.009) (0.01,0.063)rectangle++(0.009,0.009) (0.523,0.028)rectangle++(0.009,0.009)
  (0.248,0.011)rectangle++(0.009,0.009) (0.117,0.003)rectangle++(0.009,0.009) (0.054,0.519)rectangle++(0.009,0.009) (0.023,0.767)rectangle++(0.009,0.009) (0.009,0.886)rectangle++(0.009,0.009)
  (0.002,0.943)rectangle++(0.009,0.009) (-0.001,0.45)rectangle++(0.009,0.009) (0.517,0.214)rectangle++(0.009,0.009) (0.766,0.1)rectangle++(0.009,0.009) (0.885,0.046)rectangle++(0.009,0.009)
  (0.423,0.02)rectangle++(0.009,0.009) (0.72,0.007)rectangle++(0.009,0.009) (0.863,0.001)rectangle++(0.009,0.009) (0.412,0.518)rectangle++(0.009,0.009) (0.715,0.246)rectangle++(0.009,0.009)
  (0.341,0.116)rectangle++(0.009,0.009) (0.161,0.573)rectangle++(0.009,0.009) (0.075,0.793)rectangle++(0.009,0.009) (0.034,0.898)rectangle++(0.009,0.009) (0.014,0.429)rectangle++(0.009,0.009)
  (0.524,0.203)rectangle++(0.009,0.009) (0.249,0.095)rectangle++(0.009,0.009) (0.117,0.563)rectangle++(0.009,0.009) (0.574,0.268)rectangle++(0.009,0.009) (0.273,0.126)rectangle++(0.009,0.009)
  (0.649,0.058)rectangle++(0.009,0.009) (0.829,0.026)rectangle++(0.009,0.009) (0.916,0.01)rectangle++(0.009,0.009) (0.957,0.002)rectangle++(0.009,0.009) (0.977,-0.001)rectangle++(0.009,0.009)
  (0.467,0.517)rectangle++(0.009,0.009) (0.742,0.246)rectangle++(0.009,0.009) (0.874,0.116)rectangle++(0.009,0.009) (0.417,0.573)rectangle++(0.009,0.009) (0.198,0.273)rectangle++(0.009,0.009)
  (0.093,0.129)rectangle++(0.009,0.009) (0.042,0.059)rectangle++(0.009,0.009) (0.018,0.546)rectangle++(0.009,0.009) (0.526,0.26)rectangle++(0.009,0.009) (0.77,0.122)rectangle++(0.009,0.009)
  (0.367,0.576)rectangle++(0.009,0.009) (0.694,0.274)rectangle++(0.009,0.009) (0.331,0.129)rectangle++(0.009,0.009) (0.676,0.06)rectangle++(0.009,0.009) (0.322,0.026)rectangle++(0.009,0.009)
  (0.152,0.01)rectangle++(0.009,0.009) (0.591,0.003)rectangle++(0.009,0.009) (0.801,-0.001)rectangle++(0.009,0.009) (0.382,0.517)rectangle++(0.009,0.009) (0.701,0.246)rectangle++(0.009,0.009)
  (0.854,0.116)rectangle++(0.009,0.009) (0.408,0.053)rectangle++(0.009,0.009) (0.193,0.543)rectangle++(0.009,0.009) (0.61,0.258)rectangle++(0.009,0.009) (0.291,0.122)rectangle++(0.009,0.009)
  (0.137,0.056)rectangle++(0.009,0.009) (0.064,0.025)rectangle++(0.009,0.009) (0.548,0.009)rectangle++(0.009,0.009) (0.261,0.002)rectangle++(0.009,0.009) (0.643,-0.001)rectangle++(0.009,0.009)
  (0.306,0.517)rectangle++(0.009,0.009) (0.145,0.246)rectangle++(0.009,0.009) (0.067,0.116)rectangle++(0.009,0.009) (0.03,0.053)rectangle++(0.009,0.009) (0.012,0.543)rectangle++(0.009,0.009)
  (0.523,0.258)rectangle++(0.009,0.009) (0.249,0.642)rectangle++(0.009,0.009) (0.117,0.306)rectangle++(0.009,0.009) (0.574,0.144)rectangle++(0.009,0.009) (0.793,0.067)rectangle++(0.009,0.009)
  (0.378,0.03)rectangle++(0.009,0.009) (0.179,0.012)rectangle++(0.009,0.009) (0.604,0.003)rectangle++(0.009,0.009) (0.287,-0.001)rectangle++(0.009,0.009) (0.656,-0.003)rectangle++(0.009,0.009)
  (0.312,-0.004)rectangle++(0.009,0.009) (0.148,-0.004)rectangle++(0.009,0.009) (0.069,-0.004)rectangle++(0.009,0.009) (0.031,0.516)rectangle++(0.009,0.009) (0.012,0.245)rectangle++(0.009,0.009)
  (0.524,0.115)rectangle++(0.009,0.009) (0.249,0.573)rectangle++(0.009,0.009) (0.117,0.273)rectangle++(0.009,0.009) (0.574,0.129)rectangle++(0.009,0.009) (0.273,0.579)rectangle++(0.009,0.009)
  (0.649,0.276)rectangle++(0.009,0.009) (0.309,0.65)rectangle++(0.009,0.009) (0.666,0.31)rectangle++(0.009,0.009) (0.317,0.666)rectangle++(0.009,0.009) (0.15,0.317)rectangle++(0.009,0.009)
  (0.07,0.67)rectangle++(0.009,0.009) (0.031,0.319)rectangle++(0.009,0.009) (0.013,0.151)rectangle++(0.009,0.009) (0.004,0.07)rectangle++(0.009,0.009) (-0.001,0.551)rectangle++(0.009,0.009)
  (0.517,0.262)rectangle++(0.009,0.009) (0.766,0.124)rectangle++(0.009,0.009) (0.885,0.057)rectangle++(0.009,0.009) (0.423,0.545)rectangle++(0.009,0.009) (0.201,0.259)rectangle++(0.009,0.009)
  (0.614,0.122)rectangle++(0.009,0.009) (0.292,0.056)rectangle++(0.009,0.009) (0.138,0.025)rectangle++(0.009,0.009) (0.584,0.01)rectangle++(0.009,0.009) (0.278,0.522)rectangle++(0.009,0.009)
  (0.131,0.768)rectangle++(0.009,0.009) (0.581,0.366)rectangle++(0.009,0.009) (0.276,0.174)rectangle++(0.009,0.009) (0.13,0.081)rectangle++(0.009,0.009) (0.58,0.037)rectangle++(0.009,0.009)
  (0.276,0.535)rectangle++(0.009,0.009) (0.13,0.255)rectangle++(0.009,0.009) (0.06,0.12)rectangle++(0.009,0.009) (0.547,0.055)rectangle++(0.009,0.009) (0.26,0.544)rectangle++(0.009,0.009)
  (0.122,0.259)rectangle++(0.009,0.009) (0.576,0.122)rectangle++(0.009,0.009) (0.274,0.576)rectangle++(0.009,0.009) (0.649,0.274)rectangle++(0.009,0.009) (0.309,0.649)rectangle++(0.009,0.009)
  (0.666,0.309)rectangle++(0.009,0.009) (0.837,0.146)rectangle++(0.009,0.009) (0.92,0.068)rectangle++(0.009,0.009) (0.959,0.03)rectangle++(0.009,0.009) (0.458,0.532)rectangle++(0.009,0.009)
  (0.738,0.253)rectangle++(0.009,0.009) (0.872,0.119)rectangle++(0.009,0.009) (0.936,0.055)rectangle++(0.009,0.009) (0.447,0.544)rectangle++(0.009,0.009) (0.732,0.259)rectangle++(0.009,0.009)
  (0.869,0.122)rectangle++(0.009,0.009) (0.415,0.576)rectangle++(0.009,0.009) (0.197,0.794)rectangle++(0.009,0.009) (0.092,0.379)rectangle++(0.009,0.009) (0.042,0.18)rectangle++(0.009,0.009)
  (0.018,0.604)rectangle++(0.009,0.009) (0.006,0.287)rectangle++(0.009,0.009) (0.521,0.136)rectangle++(0.009,0.009) (0.248,0.583)rectangle++(0.009,0.009) (0.116,0.797)rectangle++(0.009,0.009)
  (0.054,0.9)rectangle++(0.009,0.009) (0.023,0.95)rectangle++(0.009,0.009) (0.009,0.974)rectangle++(0.009,0.009) (0.522,0.465)rectangle++(0.009,0.009) (0.248,0.741)rectangle++(0.009,0.009)
  (0.117,0.353)rectangle++(0.009,0.009) (0.574,0.167)rectangle++(0.009,0.009) (0.273,0.598)rectangle++(0.009,0.009) (0.649,0.285)rectangle++(0.009,0.009) (0.309,0.134)rectangle++(0.009,0.009)
  (0.146,0.582)rectangle++(0.009,0.009) (0.068,0.277)rectangle++(0.009,0.009) (0.55,0.131)rectangle++(0.009,0.009) (0.262,0.58)rectangle++(0.009,0.009) (0.123,0.276)rectangle++(0.009,0.009)
  (0.057,0.65)rectangle++(0.009,0.009) (0.025,0.83)rectangle++(0.009,0.009) (0.53,0.396)rectangle++(0.009,0.009) (0.252,0.708)rectangle++(0.009,0.009) (0.639,0.337)rectangle++(0.009,0.009)
  (0.304,0.16)rectangle++(0.009,0.009) (0.144,0.074)rectangle++(0.009,0.009) (0.587,0.033)rectangle++(0.009,0.009) (0.799,0.014)rectangle++(0.009,0.009) (0.381,0.524)rectangle++(0.009,0.009)
  (0.181,0.769)rectangle++(0.009,0.009) (0.604,0.367)rectangle++(0.009,0.009) (0.808,0.174)rectangle++(0.009,0.009) (0.905,0.081)rectangle++(0.009,0.009) (0.432,0.037)rectangle++(0.009,0.009)
  (0.205,0.015)rectangle++(0.009,0.009) (0.616,0.005)rectangle++(0.009,0.009) (0.293,0)rectangle++(0.009,0.009) (0.658,-0.002)rectangle++(0.009,0.009) (0.314,0.517)rectangle++(0.009,0.009)
  (0.148,0.246)rectangle++(0.009,0.009) (0.069,0.636)rectangle++(0.009,0.009) (0.551,0.303)rectangle++(0.009,0.009) (0.782,0.143)rectangle++(0.009,0.009) (0.373,0.586)rectangle++(0.009,0.009)
  (0.177,0.279)rectangle++(0.009,0.009) (0.602,0.132)rectangle++(0.009,0.009) (0.287,0.061)rectangle++(0.009,0.009) (0.135,0.027)rectangle++(0.009,0.009) (0.063,0.531)rectangle++(0.009,0.009)
  (0.028,0.252)rectangle++(0.009,0.009) (0.531,0.119)rectangle++(0.009,0.009) (0.253,0.575)rectangle++(0.009,0.009) (0.119,0.274)rectangle++(0.009,0.009) (0.055,0.649)rectangle++(0.009,0.009)
  (0.024,0.829)rectangle++(0.009,0.009) (0.009,0.916)rectangle++(0.009,0.009) (0.002,0.437)rectangle++(0.009,0.009) (-0.001,0.208)rectangle++(0.009,0.009) (0.517,0.097)rectangle++(0.009,0.009)
  (0.246,0.044)rectangle++(0.009,0.009) (0.116,0.539)rectangle++(0.009,0.009) (0.573,0.256)rectangle++(0.009,0.009) (0.273,0.641)rectangle++(0.009,0.009) (0.649,0.305)rectangle++(0.009,0.009)
  (0.309,0.664)rectangle++(0.009,0.009) (0.666,0.316)rectangle++(0.009,0.009) (0.317,0.15)rectangle++(0.009,0.009) (0.15,0.589)rectangle++(0.009,0.009) (0.07,0.801)rectangle++(0.009,0.009)
  (0.031,0.382)rectangle++(0.009,0.009) (0.013,0.701)rectangle++(0.009,0.009) (0.004,0.334)rectangle++(0.009,0.009) (0.519,0.158)rectangle++(0.009,0.009) (0.247,0.594)rectangle++(0.009,0.009)
  (0.636,0.283)rectangle++(0.009,0.009) (0.303,0.133)rectangle++(0.009,0.009) (0.663,0.062)rectangle++(0.009,0.009) (0.316,0.547)rectangle++(0.009,0.009) (0.149,0.78)rectangle++(0.009,0.009)
  (0.589,0.372)rectangle++(0.009,0.009) (0.801,0.176)rectangle++(0.009,0.009) (0.382,0.602)rectangle++(0.009,0.009) (0.181,0.287)rectangle++(0.009,0.009) (0.085,0.655)rectangle++(0.009,0.009)
  (0.558,0.312)rectangle++(0.009,0.009) (0.266,0.668)rectangle++(0.009,0.009) (0.125,0.318)rectangle++(0.009,0.009) (0.578,0.15)rectangle++(0.009,0.009) (0.795,0.07)rectangle++(0.009,0.009)
  (0.379,0.551)rectangle++(0.009,0.009) (0.18,0.262)rectangle++(0.009,0.009) (0.604,0.124)rectangle++(0.009,0.009) (0.808,0.057)rectangle++(0.009,0.009) (0.385,0.545)rectangle++(0.009,0.009)
  (0.183,0.259)rectangle++(0.009,0.009) (0.085,0.642)rectangle++(0.009,0.009) (0.039,0.306)rectangle++(0.009,0.009) (0.536,0.144)rectangle++(0.009,0.009) (0.775,0.067)rectangle++(0.009,0.009)
  (0.37,0.55)rectangle++(0.009,0.009) (0.175,0.782)rectangle++(0.009,0.009) (0.082,0.373)rectangle++(0.009,0.009) (0.037,0.177)rectangle++(0.009,0.009) (0.015,0.602)rectangle++(0.009,0.009)
  (0.005,0.807)rectangle++(0.009,0.009) (0,0.385)rectangle++(0.009,0.009) (-0.002,0.182)rectangle++(0.009,0.009) (0.517,0.085)rectangle++(0.009,0.009) (0.766,0.039)rectangle++(0.009,0.009)
  (0.365,0.536)rectangle++(0.009,0.009) (0.173,0.775)rectangle++(0.009,0.009) (0.081,0.89)rectangle++(0.009,0.009) (0.036,0.945)rectangle++(0.009,0.009) (0.535,0.451)rectangle++(0.009,0.009)
  (0.255,0.734)rectangle++(0.009,0.009) (0.12,0.87)rectangle++(0.009,0.009) (0.055,0.935)rectangle++(0.009,0.009) (0.024,0.967)rectangle++(0.009,0.009) (0.529,0.462)rectangle++(0.009,0.009)
  (0.772,0.219)rectangle++(0.009,0.009) (0.368,0.623)rectangle++(0.009,0.009) (0.174,0.297)rectangle++(0.009,0.009) (0.081,0.14)rectangle++(0.009,0.009) (0.557,0.065)rectangle++(0.009,0.009)
  (0.265,0.029)rectangle++(0.009,0.009) (0.645,0.011)rectangle++(0.009,0.009) (0.307,0.523)rectangle++(0.009,0.009) (0.145,0.249)rectangle++(0.009,0.009) (0.067,0.117)rectangle++(0.009,0.009)
  (0.55,0.054)rectangle++(0.009,0.009) (0.782,0.024)rectangle++(0.009,0.009) (0.893,0.009)rectangle++(0.009,0.009) (0.426,0.522)rectangle++(0.009,0.009) (0.202,0.248)rectangle++(0.009,0.009)
  (0.095,0.117)rectangle++(0.009,0.009) (0.043,0.054)rectangle++(0.009,0.009) (0.018,0.023)rectangle++(0.009,0.009) (0.006,0.009)rectangle++(0.009,0.009) (0.521,0.002)rectangle++(0.009,0.009)
  (0.248,0.519)rectangle++(0.009,0.009) (0.117,0.767)rectangle++(0.009,0.009) (0.574,0.366)rectangle++(0.009,0.009) (0.793,0.173)rectangle++(0.009,0.009) (0.378,0.081)rectangle++(0.009,0.009)
  (0.179,0.556)rectangle++(0.009,0.009) (0.604,0.265)rectangle++(0.009,0.009) (0.807,0.125)rectangle++(0.009,0.009) (0.905,0.058)rectangle++(0.009,0.009) (0.432,0.545)rectangle++(0.009,0.009)
  (0.725,0.259)rectangle++(0.009,0.009) (0.346,0.642)rectangle++(0.009,0.009) (0.684,0.306)rectangle++(0.009,0.009) (0.326,0.144)rectangle++(0.009,0.009) (0.674,0.067)rectangle++(0.009,0.009)
  (0.841,0.03)rectangle++(0.009,0.009) (0.921,0.012)rectangle++(0.009,0.009) (0.44,0.003)rectangle++(0.009,0.009) (0.729,-0.001)rectangle++(0.009,0.009) (0.348,0.517)rectangle++(0.009,0.009)
  (0.164,0.766)rectangle++(0.009,0.009) (0.597,0.365)rectangle++(0.009,0.009) (0.804,0.173)rectangle++(0.009,0.009) (0.384,0.601)rectangle++(0.009,0.009) (0.702,0.286)rectangle++(0.009,0.009)
  (0.335,0.135)rectangle++(0.009,0.009) (0.678,0.062)rectangle++(0.009,0.009) (0.323,0.548)rectangle++(0.009,0.009) (0.153,0.781)rectangle++(0.009,0.009) (0.071,0.892)rectangle++(0.009,0.009)
  (0.032,0.946)rectangle++(0.009,0.009) (0.013,0.972)rectangle++(0.009,0.009) (0.004,0.984)rectangle++(0.009,0.009) (0.52,0.47)rectangle++(0.009,0.009) (0.247,0.743)rectangle++(0.009,0.009)
  (0.636,0.354)rectangle++(0.009,0.009) (0.303,0.688)rectangle++(0.009,0.009) (0.143,0.328)rectangle++(0.009,0.009) (0.066,0.155)rectangle++(0.009,0.009) (0.03,0.072)rectangle++(0.009,0.009)
  (0.532,0.032)rectangle++(0.009,0.009) (0.253,0.013)rectangle++(0.009,0.009) (0.119,0.524)rectangle++(0.009,0.009) (0.055,0.249)rectangle++(0.009,0.009) (0.024,0.637)rectangle++(0.009,0.009)
  (0.529,0.304)rectangle++(0.009,0.009) (0.772,0.143)rectangle++(0.009,0.009) (0.888,0.066)rectangle++(0.009,0.009) (0.424,0.55)rectangle++(0.009,0.009) (0.721,0.261)rectangle++(0.009,0.009)
  (0.344,0.643)rectangle++(0.009,0.009) (0.163,0.306)rectangle++(0.009,0.009) (0.596,0.145)rectangle++(0.009,0.009) (0.284,0.587)rectangle++(0.009,0.009) (0.134,0.279)rectangle++(0.009,0.009)
  (0.062,0.652)rectangle++(0.009,0.009) (0.027,0.311)rectangle++(0.009,0.009) (0.011,0.667)rectangle++(0.009,0.009) (0.003,0.838)rectangle++(0.009,0.009) (-0.001,0.92)rectangle++(0.009,0.009)
  (-0.003,0.439)rectangle++(0.009,0.009) (-0.004,0.208)rectangle++(0.009,0.009) (-0.004,0.098)rectangle++(0.009,0.009) (-0.004,0.565)rectangle++(0.009,0.009) (0.516,0.269)rectangle++(0.009,0.009)
  (0.765,0.127)rectangle++(0.009,0.009) (0.885,0.058)rectangle++(0.009,0.009) (0.422,0.546)rectangle++(0.009,0.009) (0.72,0.26)rectangle++(0.009,0.009) (0.343,0.642)rectangle++(0.009,0.009)
  (0.163,0.306)rectangle++(0.009,0.009) (0.076,0.665)rectangle++(0.009,0.009) (0.034,0.317)rectangle++(0.009,0.009) (0.534,0.15)rectangle++(0.009,0.009) (0.254,0.069)rectangle++(0.009,0.009)
  (0.12,0.551)rectangle++(0.009,0.009) (0.055,0.262)rectangle++(0.009,0.009) (0.544,0.123)rectangle++(0.009,0.009) (0.259,0.057)rectangle++(0.009,0.009) (0.642,0.025)rectangle++(0.009,0.009)
  (0.306,0.01)rectangle++(0.009,0.009) (0.144,0.002)rectangle++(0.009,0.009) (0.067,0.519)rectangle++(0.009,0.009) (0.55,0.247)rectangle++(0.009,0.009) (0.262,0.636)rectangle++(0.009,0.009)
  (0.123,0.823)rectangle++(0.009,0.009) (0.057,0.913)rectangle++(0.009,0.009) (0.545,0.436)rectangle++(0.009,0.009) (0.259,0.207)rectangle++(0.009,0.009) (0.122,0.617)rectangle++(0.009,0.009)
  (0.576,0.294)rectangle++(0.009,0.009) (0.794,0.139)rectangle++(0.009,0.009) (0.379,0.064)rectangle++(0.009,0.009) (0.7,0.028)rectangle++(0.009,0.009) (0.333,0.531)rectangle++(0.009,0.009)
  (0.158,0.773)rectangle++(0.009,0.009) (0.073,0.889)rectangle++(0.009,0.009) (0.033,0.424)rectangle++(0.009,0.009) (0.013,0.721)rectangle++(0.009,0.009) (0.004,0.344)rectangle++(0.009,0.009)
  (0,0.163)rectangle++(0.009,0.009) (-0.003,0.596)rectangle++(0.009,0.009) (0.516,0.284)rectangle++(0.009,0.009) (0.246,0.654)rectangle++(0.009,0.009) (0.116,0.311)rectangle++(0.009,0.009)
  (0.053,0.147)rectangle++(0.009,0.009) (0.023,0.068)rectangle++(0.009,0.009) (0.009,0.55)rectangle++(0.009,0.009) (0.002,0.262)rectangle++(0.009,0.009) (-0.001,0.643)rectangle++(0.009,0.009)
  (-0.003,0.826)rectangle++(0.009,0.009) (-0.004,0.394)rectangle++(0.009,0.009) (0.516,0.187)rectangle++(0.009,0.009) (0.245,0.607)rectangle++(0.009,0.009) (0.115,0.809)rectangle++(0.009,0.009)
  (0.573,0.386)rectangle++(0.009,0.009) (0.273,0.183)rectangle++(0.009,0.009) (0.649,0.085)rectangle++(0.009,0.009) (0.309,0.559)rectangle++(0.009,0.009) (0.146,0.786)rectangle++(0.009,0.009)
  (0.068,0.895)rectangle++(0.009,0.009) (0.03,0.427)rectangle++(0.009,0.009) (0.532,0.203)rectangle++(0.009,0.009) (0.253,0.095)rectangle++(0.009,0.009) (0.119,0.563)rectangle++(0.009,0.009)
  (0.055,0.788)rectangle++(0.009,0.009) (0.544,0.376)rectangle++(0.009,0.009) (0.259,0.178)rectangle++(0.009,0.009) (0.122,0.083)rectangle++(0.009,0.009) (0.576,0.038)rectangle++(0.009,0.009)
  (0.794,0.016)rectangle++(0.009,0.009) (0.899,0.005)rectangle++(0.009,0.009) (0.429,0)rectangle++(0.009,0.009) (0.724,-0.002)rectangle++(0.009,0.009) (0.345,0.517)rectangle++(0.009,0.009)
  (0.683,0.246)rectangle++(0.009,0.009) (0.326,0.116)rectangle++(0.009,0.009) (0.154,0.573)rectangle++(0.009,0.009) (0.072,0.273)rectangle++(0.009,0.009) (0.552,0.129)rectangle++(0.009,0.009)
  (0.263,0.579)rectangle++(0.009,0.009) (0.124,0.796)rectangle++(0.009,0.009) (0.577,0.38)rectangle++(0.009,0.009) (0.795,0.18)rectangle++(0.009,0.009) (0.899,0.084)rectangle++(0.009,0.009)
  (0.949,0.038)rectangle++(0.009,0.009) (0.453,0.016)rectangle++(0.009,0.009) (0.215,0.005)rectangle++(0.009,0.009) (0.621,0)rectangle++(0.009,0.009) (0.816,-0.002)rectangle++(0.009,0.009)
  (0.389,0.517)rectangle++(0.009,0.009) (0.184,0.246)rectangle++(0.009,0.009) (0.086,0.116)rectangle++(0.009,0.009) (0.559,0.053)rectangle++(0.009,0.009) (0.786,0.023)rectangle++(0.009,0.009)
  (0.375,0.009)rectangle++(0.009,0.009) (0.698,0.002)rectangle++(0.009,0.009) (0.333,-0.001)rectangle++(0.009,0.009) (0.157,-0.003)rectangle++(0.009,0.009) (0.593,-0.004)rectangle++(0.009,0.009)
  (0.802,-0.004)rectangle++(0.009,0.009) (0.383,-0.004)rectangle++(0.009,0.009) (0.701,-0.004)rectangle++(0.009,0.009) (0.854,-0.004)rectangle++(0.009,0.009) (0.408,-0.004)rectangle++(0.009,0.009)
  (0.193,-0.004)rectangle++(0.009,0.009) (0.61,-0.004)rectangle++(0.009,0.009) (0.811,-0.004)rectangle++(0.009,0.009) (0.387,0.516)rectangle++(0.009,0.009) (0.183,0.245)rectangle++(0.009,0.009)
  (0.086,0.115)rectangle++(0.009,0.009) (0.039,0.573)rectangle++(0.009,0.009) (0.016,0.273)rectangle++(0.009,0.009) (0.005,0.129)rectangle++(0.009,0.009) (0,0.579)rectangle++(0.009,0.009)
  (-0.002,0.796)rectangle++(0.009,0.009) (-0.003,0.38)rectangle++(0.009,0.009) (0.516,0.18)rectangle++(0.009,0.009) (0.245,0.604)rectangle++(0.009,0.009) (0.635,0.288)rectangle++(0.009,0.009)
  (0.303,0.656)rectangle++(0.009,0.009) (0.143,0.312)rectangle++(0.009,0.009) (0.586,0.148)rectangle++(0.009,0.009) (0.799,0.069)rectangle++(0.009,0.009) (0.381,0.031)rectangle++(0.009,0.009)
  (0.701,0.012)rectangle++(0.009,0.009) (0.854,0.004)rectangle++(0.009,0.009) (0.928,-0.001)rectangle++(0.009,0.009) (0.963,-0.003)rectangle++(0.009,0.009) (0.46,-0.004)rectangle++(0.009,0.009)
  (0.738,-0.004)rectangle++(0.009,0.009) (0.352,-0.004)rectangle++(0.009,0.009) (0.167,0.516)rectangle++(0.009,0.009) (0.078,0.245)rectangle++(0.009,0.009) (0.035,0.635)rectangle++(0.009,0.009)
  (0.014,0.823)rectangle++(0.009,0.009) (0.005,0.393)rectangle++(0.009,0.009) (0,0.706)rectangle++(0.009,0.009) (-0.002,0.857)rectangle++(0.009,0.009) (0.517,0.409)rectangle++(0.009,0.009)
  (0.246,0.714)rectangle++(0.009,0.009) (0.636,0.34)rectangle++(0.009,0.009) (0.303,0.161)rectangle++(0.009,0.009) (0.143,0.075)rectangle++(0.009,0.009) (0.066,0.034)rectangle++(0.009,0.009)
  (0.029,0.534)rectangle++(0.009,0.009) (0.012,0.774)rectangle++(0.009,0.009) (0.523,0.369)rectangle++(0.009,0.009) (0.249,0.695)rectangle++(0.009,0.009) (0.637,0.331)rectangle++(0.009,0.009)
  (0.823,0.157)rectangle++(0.009,0.009) (0.393,0.073)rectangle++(0.009,0.009) (0.186,0.033)rectangle++(0.009,0.009) (0.087,0.533)rectangle++(0.009,0.009) (0.039,0.774)rectangle++(0.009,0.009)
  (0.537,0.369)rectangle++(0.009,0.009) (0.775,0.175)rectangle++(0.009,0.009) (0.89,0.082)rectangle++(0.009,0.009) (0.945,0.037)rectangle++(0.009,0.009) (0.451,0.535)rectangle++(0.009,0.009)
  (0.214,0.775)rectangle++(0.009,0.009) (0.1,0.889)rectangle++(0.009,0.009) (0.566,0.425)rectangle++(0.009,0.009) (0.269,0.201)rectangle++(0.009,0.009) (0.647,0.094)rectangle++(0.009,0.009)
  (0.828,0.043)rectangle++(0.009,0.009) (0.395,0.538)rectangle++(0.009,0.009) (0.707,0.256)rectangle++(0.009,0.009) (0.857,0.121)rectangle++(0.009,0.009) (0.409,0.576)rectangle++(0.009,0.009)
  (0.194,0.274)rectangle++(0.009,0.009) (0.091,0.129)rectangle++(0.009,0.009) (0.561,0.06)rectangle++(0.009,0.009) (0.267,0.026)rectangle++(0.009,0.009) (0.126,0.01)rectangle++(0.009,0.009)
  (0.578,0.003)rectangle++(0.009,0.009) (0.275,-0.001)rectangle++(0.009,0.009) (0.13,-0.003)rectangle++(0.009,0.009) (0.06,-0.004)rectangle++(0.009,0.009) (0.026,0.516)rectangle++(0.009,0.009)
  (0.01,0.765)rectangle++(0.009,0.009) (0.003,0.365)rectangle++(0.009,0.009) (-0.001,0.693)rectangle++(0.009,0.009) (0.517,0.33)rectangle++(0.009,0.009) (0.246,0.156)rectangle++(0.009,0.009)
  (0.636,0.073)rectangle++(0.009,0.009) (0.823,0.033)rectangle++(0.009,0.009) (0.393,0.533)rectangle++(0.009,0.009) (0.706,0.254)rectangle++(0.009,0.009) (0.337,0.639)rectangle++(0.009,0.009)
  (0.679,0.305)rectangle++(0.009,0.009) (0.324,0.664)rectangle++(0.009,0.009) (0.153,0.316)rectangle++(0.009,0.009) (0.071,0.669)rectangle++(0.009,0.009) (0.032,0.319)rectangle++(0.009,0.009)
  (0.013,0.151)rectangle++(0.009,0.009) (0.524,0.07)rectangle++(0.009,0.009) (0.249,0.031)rectangle++(0.009,0.009) (0.637,0.013)rectangle++(0.009,0.009) (0.824,0.004)rectangle++(0.009,0.009)
  (0.393,0.519)rectangle++(0.009,0.009) (0.186,0.767)rectangle++(0.009,0.009) (0.087,0.366)rectangle++(0.009,0.009) (0.039,0.693)rectangle++(0.009,0.009) (0.017,0.85)rectangle++(0.009,0.009)
  (0.006,0.926)rectangle++(0.009,0.009) (0.52,0.442)rectangle++(0.009,0.009) (0.247,0.73)rectangle++(0.009,0.009) (0.636,0.348)rectangle++(0.009,0.009) (0.303,0.165)rectangle++(0.009,0.009)
  (0.143,0.077)rectangle++(0.009,0.009) (0.066,0.034)rectangle++(0.009,0.009) (0.55,0.014)rectangle++(0.009,0.009) (0.261,0.004)rectangle++(0.009,0.009) (0.123,0)rectangle++(0.009,0.009)
  (0.577,-0.002)rectangle++(0.009,0.009) (0.795,-0.004)rectangle++(0.009,0.009) (0.899,-0.004)rectangle++(0.009,0.009) (0.429,0.516)rectangle++(0.009,0.009) (0.204,0.765)rectangle++(0.009,0.009)
  (0.095,0.885)rectangle++(0.009,0.009) (0.563,0.422)rectangle++(0.009,0.009) (0.268,0.2)rectangle++(0.009,0.009) (0.126,0.094)rectangle++(0.009,0.009) (0.578,0.043)rectangle++(0.009,0.009)
  (0.275,0.538)rectangle++(0.009,0.009) (0.13,0.776)rectangle++(0.009,0.009) (0.58,0.37)rectangle++(0.009,0.009) (0.276,0.175)rectangle++(0.009,0.009) (0.65,0.082)rectangle++(0.009,0.009)
  (0.83,0.037)rectangle++(0.009,0.009) (0.396,0.535)rectangle++(0.009,0.009) (0.188,0.775)rectangle++(0.009,0.009) (0.608,0.369)rectangle++(0.009,0.009) (0.809,0.175)rectangle++(0.009,0.009)
  (0.386,0.602)rectangle++(0.009,0.009) (0.183,0.286)rectangle++(0.009,0.009) (0.086,0.655)rectangle++(0.009,0.009) (0.039,0.832)rectangle++(0.009,0.009) (0.536,0.397)rectangle++(0.009,0.009)
  (0.775,0.188)rectangle++(0.009,0.009) (0.37,0.088)rectangle++(0.009,0.009) (0.695,0.04)rectangle++(0.009,0.009) (0.331,0.537)rectangle++(0.009,0.009) (0.677,0.255)rectangle++(0.009,0.009)
  (0.322,0.64)rectangle++(0.009,0.009) (0.152,0.825)rectangle++(0.009,0.009) (0.071,0.394)rectangle++(0.009,0.009) (0.032,0.187)rectangle++(0.009,0.009) (0.533,0.087)rectangle++(0.009,0.009)
  (0.253,0.56)rectangle++(0.009,0.009) (0.639,0.266)rectangle++(0.009,0.009) (0.305,0.645)rectangle++(0.009,0.009) (0.664,0.307)rectangle++(0.009,0.009) (0.316,0.665)rectangle++(0.009,0.009)
  (0.669,0.317)rectangle++(0.009,0.009) (0.319,0.67)rectangle++(0.009,0.009) (0.671,0.319)rectangle++(0.009,0.009) (0.32,0.671)rectangle++(0.009,0.009) (0.151,0.32)rectangle++(0.009,0.009)
  (0.07,0.151)rectangle++(0.009,0.009) (0.031,0.07)rectangle++(0.009,0.009) (0.013,0.551)rectangle++(0.009,0.009) (0.004,0.782)rectangle++(0.009,0.009) (-0.001,0.893)rectangle++(0.009,0.009)
  (0.517,0.426)rectangle++(0.009,0.009) (0.246,0.722)rectangle++(0.009,0.009) (0.636,0.344)rectangle++(0.009,0.009) (0.303,0.683)rectangle++(0.009,0.009) (0.663,0.325)rectangle++(0.009,0.009)
  (0.316,0.674)rectangle++(0.009,0.009) (0.669,0.321)rectangle++(0.009,0.009) (0.839,0.152)rectangle++(0.009,0.009) (0.4,0.591)rectangle++(0.009,0.009) (0.19,0.281)rectangle++(0.009,0.009)
  (0.089,0.653)rectangle++(0.009,0.009) (0.04,0.831)rectangle++(0.009,0.009) (0.017,0.916)rectangle++(0.009,0.009) (0.006,0.438)rectangle++(0.009,0.009) (0.52,0.208)rectangle++(0.009,0.009)
  (0.767,0.097)rectangle++(0.009,0.009) (0.366,0.564)rectangle++(0.009,0.009) (0.173,0.269)rectangle++(0.009,0.009) (0.081,0.647)rectangle++(0.009,0.009) (0.556,0.308)rectangle++(0.009,0.009)
  (0.265,0.666)rectangle++(0.009,0.009) (0.125,0.317)rectangle++(0.009,0.009) (0.058,0.67)rectangle++(0.009,0.009) (0.025,0.839)rectangle++(0.009,0.009) (0.53,0.4)rectangle++(0.009,0.009)
  (0.772,0.19)rectangle++(0.009,0.009) (0.368,0.089)rectangle++(0.009,0.009) (0.694,0.04)rectangle++(0.009,0.009) (0.331,0.537)rectangle++(0.009,0.009) (0.157,0.775)rectangle++(0.009,0.009)
  (0.593,0.37)rectangle++(0.009,0.009) (0.802,0.175)rectangle++(0.009,0.009) (0.383,0.082)rectangle++(0.009,0.009) (0.701,0.037)rectangle++(0.009,0.009) (0.334,0.015)rectangle++(0.009,0.009)
  (0.158,0.005)rectangle++(0.009,0.009) (0.074,0.52)rectangle++(0.009,0.009) (0.553,0.247)rectangle++(0.009,0.009) (0.263,0.116)rectangle++(0.009,0.009) (0.124,0.054)rectangle++(0.009,0.009)
  (0.057,0.543)rectangle++(0.009,0.009) (0.025,0.258)rectangle++(0.009,0.009) (0.01,0.122)rectangle++(0.009,0.009) (0.002,0.056)rectangle++(0.009,0.009) (0.519,0.025)rectangle++(0.009,0.009)
  (0.247,0.009)rectangle++(0.009,0.009) (0.636,0.002)rectangle++(0.009,0.009) (0.303,0.519)rectangle++(0.009,0.009) (0.143,0.767)rectangle++(0.009,0.009) (0.066,0.366)rectangle++(0.009,0.009)
  (0.03,0.693)rectangle++(0.009,0.009) (0.012,0.33)rectangle++(0.009,0.009) (0.003,0.676)rectangle++(0.009,0.009) (0.519,0.322)rectangle++(0.009,0.009) (0.767,0.152)rectangle++(0.009,0.009)
  (0.886,0.071)rectangle++(0.009,0.009) (0.943,0.032)rectangle++(0.009,0.009) (0.45,0.533)rectangle++(0.009,0.009) (0.734,0.253)rectangle++(0.009,0.009) (0.35,0.119)rectangle++(0.009,0.009)
  (0.166,0.055)rectangle++(0.009,0.009) (0.077,0.024)rectangle++(0.009,0.009) (0.555,0.009)rectangle++(0.009,0.009) (0.264,0.002)rectangle++(0.009,0.009) (0.124,0.519)rectangle++(0.009,0.009)
  (0.057,0.767)rectangle++(0.009,0.009) (0.025,0.366)rectangle++(0.009,0.009) (0.01,0.173)rectangle++(0.009,0.009) (0.522,0.081)rectangle++(0.009,0.009) (0.248,0.556)rectangle++(0.009,0.009)
  (0.117,0.785)rectangle++(0.009,0.009) (0.054,0.894)rectangle++(0.009,0.009) (0.543,0.427)rectangle++(0.009,0.009) (0.259,0.723)rectangle++(0.009,0.009) (0.122,0.345)rectangle++(0.009,0.009)
  (0.056,0.163)rectangle++(0.009,0.009) (0.025,0.596)rectangle++(0.009,0.009) (0.009,0.804)rectangle++(0.009,0.009) (0.522,0.383)rectangle++(0.009,0.009) (0.248,0.182)rectangle++(0.009,0.009)
  (0.637,0.085)rectangle++(0.009,0.009) (0.303,0.558)rectangle++(0.009,0.009) (0.143,0.266)rectangle++(0.009,0.009) (0.066,0.645)rectangle++(0.009,0.009) (0.03,0.827)rectangle++(0.009,0.009)
  (0.012,0.915)rectangle++(0.009,0.009) (0.523,0.437)rectangle++(0.009,0.009) (0.249,0.727)rectangle++(0.009,0.009) (0.117,0.867)rectangle++(0.009,0.009) (0.054,0.934)rectangle++(0.009,0.009)
  (0.544,0.446)rectangle++(0.009,0.009) (0.259,0.732)rectangle++(0.009,0.009) (0.122,0.869)rectangle++(0.009,0.009) (0.056,0.935)rectangle++(0.009,0.009) (0.025,0.446)rectangle++(0.009,0.009)
  (0.009,0.732)rectangle++(0.009,0.009) (0.002,0.869)rectangle++(0.009,0.009) (-0.001,0.935)rectangle++(0.009,0.009) (0.517,0.446)rectangle++(0.009,0.009) (0.246,0.732)rectangle++(0.009,0.009)
  (0.636,0.349)rectangle++(0.009,0.009) (0.823,0.165)rectangle++(0.009,0.009) (0.913,0.077)rectangle++(0.009,0.009) (0.436,0.555)rectangle++(0.009,0.009) (0.207,0.264)rectangle++(0.009,0.009)
  (0.097,0.124)rectangle++(0.009,0.009) (0.044,0.057)rectangle++(0.009,0.009) (0.019,0.545)rectangle++(0.009,0.009) (0.007,0.259)rectangle++(0.009,0.009) (0.001,0.642)rectangle++(0.009,0.009)
  (-0.002,0.306)rectangle++(0.009,0.009) (-0.003,0.664)rectangle++(0.009,0.009) (-0.004,0.317)rectangle++(0.009,0.009) (-0.004,0.67)rectangle++(0.009,0.009) (-0.004,0.839)rectangle++(0.009,0.009)
  (0.516,0.4)rectangle++(0.009,0.009) (0.245,0.71)rectangle++(0.009,0.009) (0.635,0.338)rectangle++(0.009,0.009) (0.303,0.16)rectangle++(0.009,0.009) (0.143,0.595)rectangle++(0.009,0.009)
  (0.066,0.283)rectangle++(0.009,0.009) (0.549,0.134)rectangle++(0.009,0.009) (0.261,0.582)rectangle++(0.009,0.009) (0.123,0.277)rectangle++(0.009,0.009) (0.057,0.651)rectangle++(0.009,0.009)
  (0.025,0.31)rectangle++(0.009,0.009) (0.01,0.666)rectangle++(0.009,0.009) (0.522,0.318)rectangle++(0.009,0.009) (0.248,0.15)rectangle++(0.009,0.009) (0.637,0.07)rectangle++(0.009,0.009)
  (0.303,0.551)rectangle++(0.009,0.009) (0.663,0.262)rectangle++(0.009,0.009) (0.836,0.124)rectangle++(0.009,0.009) (0.919,0.057)rectangle++(0.009,0.009) (0.959,0.025)rectangle++(0.009,0.009)
  (0.458,0.01)rectangle++(0.009,0.009) (0.737,0.002)rectangle++(0.009,0.009) (0.872,-0.001)rectangle++(0.009,0.009) (0.416,0.517)rectangle++(0.009,0.009) (0.197,0.246)rectangle++(0.009,0.009)
  (0.612,0.116)rectangle++(0.009,0.009) (0.812,0.053)rectangle++(0.009,0.009) (0.907,0.023)rectangle++(0.009,0.009) (0.433,0.009)rectangle++(0.009,0.009) (0.206,0.522)rectangle++(0.009,0.009)
  (0.096,0.248)rectangle++(0.009,0.009) (0.044,0.117)rectangle++(0.009,0.009) (0.019,0.574)rectangle++(0.009,0.009) (0.527,0.273)rectangle++(0.009,0.009) (0.25,0.129)rectangle++(0.009,0.009)
  (0.118,0.059)rectangle++(0.009,0.009) (0.054,0.546)rectangle++(0.009,0.009) (0.544,0.26)rectangle++(0.009,0.009) (0.259,0.642)rectangle++(0.009,0.009) (0.642,0.306)rectangle++(0.009,0.009)
  (0.306,0.145)rectangle++(0.009,0.009) (0.664,0.067)rectangle++(0.009,0.009) (0.317,0.03)rectangle++(0.009,0.009) (0.67,0.012)rectangle++(0.009,0.009) (0.319,0.523)rectangle++(0.009,0.009)
  (0.151,0.249)rectangle++(0.009,0.009) (0.59,0.117)rectangle++(0.009,0.009) (0.801,0.054)rectangle++(0.009,0.009) (0.902,0.024)rectangle++(0.009,0.009) (0.431,0.009)rectangle++(0.009,0.009)
  (0.724,0.002)rectangle++(0.009,0.009) (0.865,-0.001)rectangle++(0.009,0.009) (0.413,-0.003)rectangle++(0.009,0.009) (0.196,0.516)rectangle++(0.009,0.009) (0.612,0.245)rectangle++(0.009,0.009)
  (0.811,0.115)rectangle++(0.009,0.009) (0.907,0.053)rectangle++(0.009,0.009) (0.433,0.023)rectangle++(0.009,0.009) (0.726,0.009)rectangle++(0.009,0.009) (0.346,0.522)rectangle++(0.009,0.009)
  (0.164,0.768)rectangle++(0.009,0.009) (0.076,0.366)rectangle++(0.009,0.009) (0.034,0.694)rectangle++(0.009,0.009) (0.534,0.331)rectangle++(0.009,0.009) (0.254,0.156)rectangle++(0.009,0.009)
  (0.12,0.073)rectangle++(0.009,0.009) (0.575,0.033)rectangle++(0.009,0.009) (0.274,0.013)rectangle++(0.009,0.009) (0.649,0.004)rectangle++(0.009,0.009) (0.309,0)rectangle++(0.009,0.009)
  (0.146,-0.003)rectangle++(0.009,0.009) (0.068,-0.004)rectangle++(0.009,0.009) (0.03,-0.004)rectangle++(0.009,0.009) (0.532,-0.004)rectangle++(0.009,0.009) (0.253,-0.004)rectangle++(0.009,0.009)
  (0.119,-0.004)rectangle++(0.009,0.009) (0.055,-0.004)rectangle++(0.009,0.009) (0.024,0.516)rectangle++(0.009,0.009) (0.529,0.245)rectangle++(0.009,0.009) (0.772,0.115)rectangle++(0.009,0.009)
  (0.368,0.573)rectangle++(0.009,0.009) (0.174,0.273)rectangle++(0.009,0.009) (0.081,0.129)rectangle++(0.009,0.009) (0.557,0.059)rectangle++(0.009,0.009) (0.785,0.026)rectangle++(0.009,0.009)
  (0.374,0.53)rectangle++(0.009,0.009) (0.697,0.252)rectangle++(0.009,0.009) (0.852,0.119)rectangle++(0.009,0.009) (0.407,0.055)rectangle++(0.009,0.009) (0.713,0.024)rectangle++(0.009,0.009)
  (0.86,0.009)rectangle++(0.009,0.009) (0.41,0.522)rectangle++(0.009,0.009) (0.195,0.248)rectangle++(0.009,0.009) (0.611,0.117)rectangle++(0.009,0.009) (0.291,0.054)rectangle++(0.009,0.009)
  (0.657,0.023)rectangle++(0.009,0.009) (0.833,0.009)rectangle++(0.009,0.009) (0.398,0.522)rectangle++(0.009,0.009) (0.189,0.248)rectangle++(0.009,0.009) (0.608,0.117)rectangle++(0.009,0.009)
  (0.81,0.054)rectangle++(0.009,0.009) (0.386,0.543)rectangle++(0.009,0.009) (0.183,0.259)rectangle++(0.009,0.009) (0.086,0.122)rectangle++(0.009,0.009) (0.039,0.056)rectangle++(0.009,0.009)
  (0.536,0.025)rectangle++(0.009,0.009) (0.255,0.529)rectangle++(0.009,0.009) (0.12,0.252)rectangle++(0.009,0.009) (0.055,0.639)rectangle++(0.009,0.009) (0.544,0.304)rectangle++(0.009,0.009)
  (0.259,0.664)rectangle++(0.009,0.009) (0.642,0.316)rectangle++(0.009,0.009) (0.306,0.669)rectangle++(0.009,0.009) (0.144,0.839)rectangle++(0.009,0.009) (0.067,0.4)rectangle++(0.009,0.009)
  (0.55,0.19)rectangle++(0.009,0.009) (0.782,0.089)rectangle++(0.009,0.009) (0.893,0.04)rectangle++(0.009,0.009) (0.946,0.017)rectangle++(0.009,0.009) (0.452,0.526)rectangle++(0.009,0.009)
  (0.215,0.77)rectangle++(0.009,0.009) (0.101,0.887)rectangle++(0.009,0.009) (0.046,0.424)rectangle++(0.009,0.009) (0.54,0.201)rectangle++(0.009,0.009) (0.257,0.614)rectangle++(0.009,0.009)
  (0.641,0.292)rectangle++(0.009,0.009) (0.305,0.658)rectangle++(0.009,0.009) (0.144,0.314)rectangle++(0.009,0.009) (0.587,0.148)rectangle++(0.009,0.009) (0.279,0.069)rectangle++(0.009,0.009)
  (0.652,0.031)rectangle++(0.009,0.009) (0.311,0.532)rectangle++(0.009,0.009) (0.147,0.773)rectangle++(0.009,0.009) (0.588,0.369)rectangle++(0.009,0.009) (0.28,0.175)rectangle++(0.009,0.009)
  (0.132,0.602)rectangle++(0.009,0.009) (0.061,0.806)rectangle++(0.009,0.009) (0.547,0.385)rectangle++(0.009,0.009) (0.26,0.702)rectangle++(0.009,0.009) (0.643,0.335)rectangle++(0.009,0.009)
  (0.306,0.158)rectangle++(0.009,0.009) (0.145,0.074)rectangle++(0.009,0.009) (0.067,0.553)rectangle++(0.009,0.009) (0.55,0.263)rectangle++(0.009,0.009) (0.262,0.124)rectangle++(0.009,0.009)
  (0.123,0.577)rectangle++(0.009,0.009) (0.057,0.275)rectangle++(0.009,0.009) (0.025,0.13)rectangle++(0.009,0.009) (0.53,0.06)rectangle++(0.009,0.009) (0.252,0.026)rectangle++(0.009,0.009)
  (0.119,0.01)rectangle++(0.009,0.009) (0.055,0.003)rectangle++(0.009,0.009) (0.024,0.519)rectangle++(0.009,0.009) (0.009,0.767)rectangle++(0.009,0.009) (0.002,0.886)rectangle++(0.009,0.009)
  (-0.001,0.423)rectangle++(0.009,0.009) (0.517,0.201)rectangle++(0.009,0.009) (0.766,0.094)rectangle++(0.009,0.009) (0.885,0.043)rectangle++(0.009,0.009) (0.423,0.538)rectangle++(0.009,0.009)
  (0.72,0.256)rectangle++(0.009,0.009) (0.343,0.641)rectangle++(0.009,0.009) (0.163,0.305)rectangle++(0.009,0.009) (0.076,0.664)rectangle++(0.009,0.009) (0.554,0.316)rectangle++(0.009,0.009)
  (0.784,0.15)rectangle++(0.009,0.009) (0.374,0.069)rectangle++(0.009,0.009) (0.177,0.551)rectangle++(0.009,0.009) (0.603,0.262)rectangle++(0.009,0.009) (0.807,0.123)rectangle++(0.009,0.009)
  (0.385,0.057)rectangle++(0.009,0.009) (0.182,0.025)rectangle++(0.009,0.009) (0.085,0.01)rectangle++(0.009,0.009) (0.039,0.002)rectangle++(0.009,0.009) (0.536,-0.001)rectangle++(0.009,0.009)
  (0.255,0.517)rectangle++(0.009,0.009) (0.64,0.246)rectangle++(0.009,0.009) (0.305,0.116)rectangle++(0.009,0.009) (0.144,0.573)rectangle++(0.009,0.009) (0.587,0.273)rectangle++(0.009,0.009)
  (0.279,0.129)rectangle++(0.009,0.009) (0.132,0.059)rectangle++(0.009,0.009) (0.581,0.026)rectangle++(0.009,0.009) (0.796,0.01)rectangle++(0.009,0.009) (0.9,0.003)rectangle++(0.009,0.009)
  (0.43,0.519)rectangle++(0.009,0.009) (0.724,0.247)rectangle++(0.009,0.009) (0.345,0.116)rectangle++(0.009,0.009) (0.683,0.053)rectangle++(0.009,0.009) (0.326,0.023)rectangle++(0.009,0.009)
  (0.154,0.529)rectangle++(0.009,0.009) (0.072,0.772)rectangle++(0.009,0.009) (0.032,0.368)rectangle++(0.009,0.009) (0.013,0.174)rectangle++(0.009,0.009) (0.004,0.081)rectangle++(0.009,0.009)
  (0,0.037)rectangle++(0.009,0.009) (0.517,0.015)rectangle++(0.009,0.009) (0.246,0.525)rectangle++(0.009,0.009) (0.116,0.77)rectangle++(0.009,0.009) (0.053,0.887)rectangle++(0.009,0.009)
  (0.023,0.943)rectangle++(0.009,0.009) (0.009,0.451)rectangle++(0.009,0.009) (0.522,0.214)rectangle++(0.009,0.009) (0.248,0.62)rectangle++(0.009,0.009) (0.117,0.295)rectangle++(0.009,0.009)
  (0.054,0.659)rectangle++(0.009,0.009) (0.543,0.314)rectangle++(0.009,0.009) (0.779,0.148)rectangle++(0.009,0.009) (0.371,0.069)rectangle++(0.009,0.009) (0.696,0.031)rectangle++(0.009,0.009)
  (0.852,0.012)rectangle++(0.009,0.009) (0.926,0.004)rectangle++(0.009,0.009) (0.962,-0.001)rectangle++(0.009,0.009) (0.46,0.517)rectangle++(0.009,0.009) (0.218,0.766)rectangle++(0.009,0.009)
  (0.622,0.365)rectangle++(0.009,0.009) (0.296,0.693)rectangle++(0.009,0.009) (0.14,0.85)rectangle++(0.009,0.009) (0.065,0.926)rectangle++(0.009,0.009) (0.029,0.962)rectangle++(0.009,0.009)
  (0.011,0.979)rectangle++(0.009,0.009) (0.003,0.988)rectangle++(0.009,0.009) (-0.001,0.992)rectangle++(0.009,0.009) (-0.003,0.474)rectangle++(0.009,0.009) (-0.004,0.745)rectangle++(0.009,0.009)
  (-0.004,0.875)rectangle++(0.009,0.009) (-0.004,0.938)rectangle++(0.009,0.009) (-0.004,0.448)rectangle++(0.009,0.009) (-0.004,0.733)rectangle++(0.009,0.009) (-0.004,0.869)rectangle++(0.009,0.009)
  (-0.004,0.935)rectangle++(0.009,0.009) (0.516,0.446)rectangle++(0.009,0.009) (0.245,0.732)rectangle++(0.009,0.009) (0.635,0.349)rectangle++(0.009,0.009) (0.303,0.165)rectangle++(0.009,0.009)
  (0.663,0.077)rectangle++(0.009,0.009) (0.316,0.035)rectangle++(0.009,0.009) (0.669,0.014)rectangle++(0.009,0.009) (0.319,0.525)rectangle++(0.009,0.009) (0.671,0.249)rectangle++(0.009,0.009)
  (0.32,0.637)rectangle++(0.009,0.009) (0.671,0.304)rectangle++(0.009,0.009) (0.32,0.663)rectangle++(0.009,0.009) (0.151,0.316)rectangle++(0.009,0.009) (0.07,0.669)rectangle++(0.009,0.009)
  (0.031,0.839)rectangle++(0.009,0.009) (0.533,0.4)rectangle++(0.009,0.009) (0.253,0.71)rectangle++(0.009,0.009) (0.639,0.338)rectangle++(0.009,0.009) (0.825,0.16)rectangle++(0.009,0.009)
  (0.913,0.074)rectangle++(0.009,0.009) (0.956,0.033)rectangle++(0.009,0.009) (0.977,0.014)rectangle++(0.009,0.009) (0.466,0.524)rectangle++(0.009,0.009) (0.222,0.769)rectangle++(0.009,0.009)
  (0.104,0.367)rectangle++(0.009,0.009) (0.048,0.694)rectangle++(0.009,0.009) (0.02,0.331)rectangle++(0.009,0.009) (0.007,0.676)rectangle++(0.009,0.009) (0.001,0.322)rectangle++(0.009,0.009)
  (0.518,0.152)rectangle++(0.009,0.009) (0.246,0.071)rectangle++(0.009,0.009) (0.116,0.552)rectangle++(0.009,0.009) (0.573,0.262)rectangle++(0.009,0.009) (0.273,0.644)rectangle++(0.009,0.009)
  (0.129,0.827)rectangle++(0.009,0.009) (0.059,0.914)rectangle++(0.009,0.009) (0.026,0.957)rectangle++(0.009,0.009) (0.01,0.457)rectangle++(0.009,0.009) (0.003,0.217)rectangle++(0.009,0.009)
  (-0.001,0.102)rectangle++(0.009,0.009) (-0.003,0.567)rectangle++(0.009,0.009) (0.516,0.27)rectangle++(0.009,0.009) (0.765,0.127)rectangle++(0.009,0.009) (0.365,0.059)rectangle++(0.009,0.009)
  (0.693,0.026)rectangle++(0.009,0.009) (0.33,0.53)rectangle++(0.009,0.009) (0.676,0.252)rectangle++(0.009,0.009) (0.322,0.639)rectangle++(0.009,0.009) (0.672,0.304)rectangle++(0.009,0.009)
  (0.32,0.144)rectangle++(0.009,0.009) (0.671,0.067)rectangle++(0.009,0.009) (0.32,0.03)rectangle++(0.009,0.009) (0.151,0.532)rectangle++(0.009,0.009) (0.07,0.253)rectangle++(0.009,0.009)
  (0.031,0.639)rectangle++(0.009,0.009) (0.013,0.304)rectangle++(0.009,0.009) (0.004,0.664)rectangle++(0.009,0.009) (-0.001,0.316)rectangle++(0.009,0.009) (0.517,0.149)rectangle++(0.009,0.009)
  (0.246,0.589)rectangle++(0.009,0.009) (0.636,0.281)rectangle++(0.009,0.009) (0.303,0.132)rectangle++(0.009,0.009) (0.143,0.581)rectangle++(0.009,0.009) (0.066,0.797)rectangle++(0.009,0.009)
  (0.029,0.38)rectangle++(0.009,0.009) (0.012,0.7)rectangle++(0.009,0.009) (0.003,0.854)rectangle++(0.009,0.009) (0.519,0.407)rectangle++(0.009,0.009) (0.247,0.193)rectangle++(0.009,0.009)
  (0.636,0.09)rectangle++(0.009,0.009) (0.303,0.561)rectangle++(0.009,0.009) (0.143,0.787)rectangle++(0.009,0.009) (0.066,0.375)rectangle++(0.009,0.009) (0.55,0.178)rectangle++(0.009,0.009)
  (0.781,0.083)rectangle++(0.009,0.009) (0.373,0.558)rectangle++(0.009,0.009) (0.177,0.785)rectangle++(0.009,0.009) (0.602,0.375)rectangle++(0.009,0.009) (0.287,0.697)rectangle++(0.009,0.009)
  (0.655,0.332)rectangle++(0.009,0.009) (0.312,0.677)rectangle++(0.009,0.009) (0.148,0.843)rectangle++(0.009,0.009) (0.588,0.402)rectangle++(0.009,0.009) (0.28,0.711)rectangle++(0.009,0.009)
  (0.132,0.339)rectangle++(0.009,0.009) (0.061,0.16)rectangle++(0.009,0.009) (0.027,0.595)rectangle++(0.009,0.009) (0.011,0.283)rectangle++(0.009,0.009) (0.003,0.134)rectangle++(0.009,0.009)
  (-0.001,0.582)rectangle++(0.009,0.009) (-0.003,0.797)rectangle++(0.009,0.009) (0.516,0.38)rectangle++(0.009,0.009) (0.245,0.7)rectangle++(0.009,0.009) (0.115,0.334)rectangle++(0.009,0.009)
  (0.053,0.678)rectangle++(0.009,0.009) (0.023,0.323)rectangle++(0.009,0.009) (0.009,0.153)rectangle++(0.009,0.009) (0.522,0.071)rectangle++(0.009,0.009) (0.248,0.032)rectangle++(0.009,0.009)
  (0.117,0.013)rectangle++(0.009,0.009) (0.054,0.524)rectangle++(0.009,0.009) (0.023,0.769)rectangle++(0.009,0.009) (0.009,0.887)rectangle++(0.009,0.009) (0.002,0.423)rectangle++(0.009,0.009)
  (0.519,0.201)rectangle++(0.009,0.009) (0.247,0.094)rectangle++(0.009,0.009) (0.116,0.043)rectangle++(0.009,0.009) (0.053,0.018)rectangle++(0.009,0.009) (0.023,0.526)rectangle++(0.009,0.009)
  (0.009,0.77)rectangle++(0.009,0.009) (0.002,0.887)rectangle++(0.009,0.009) (-0.001,0.424)rectangle++(0.009,0.009) (-0.003,0.721)rectangle++(0.009,0.009) (-0.004,0.344)rectangle++(0.009,0.009)
  (0.516,0.163)rectangle++(0.009,0.009) (0.765,0.076)rectangle++(0.009,0.009) (0.365,0.554)rectangle++(0.009,0.009) (0.173,0.264)rectangle++(0.009,0.009) (0.081,0.644)rectangle++(0.009,0.009)
  (0.556,0.307)rectangle++(0.009,0.009) (0.265,0.145)rectangle++(0.009,0.009) (0.125,0.067)rectangle++(0.009,0.009) (0.058,0.55)rectangle++(0.009,0.009) (0.025,0.262)rectangle++(0.009,0.009)
  (0.01,0.123)rectangle++(0.009,0.009) (0.002,0.577)rectangle++(0.009,0.009) (-0.001,0.795)rectangle++(0.009,0.009) (-0.003,0.899)rectangle++(0.009,0.009) (-0.004,0.429)rectangle++(0.009,0.009)
  (0.516,0.204)rectangle++(0.009,0.009) (0.245,0.095)rectangle++(0.009,0.009) (0.115,0.563)rectangle++(0.009,0.009) (0.053,0.788)rectangle++(0.009,0.009) (0.023,0.376)rectangle++(0.009,0.009)
  (0.009,0.698)rectangle++(0.009,0.009) (0.522,0.333)rectangle++(0.009,0.009) (0.248,0.677)rectangle++(0.009,0.009) (0.117,0.843)rectangle++(0.009,0.009) (0.574,0.402)rectangle++(0.009,0.009)
  (0.793,0.191)rectangle++(0.009,0.009) (0.378,0.089)rectangle++(0.009,0.009) (0.699,0.04)rectangle++(0.009,0.009) (0.333,0.537)rectangle++(0.009,0.009) (0.158,0.775)rectangle++(0.009,0.009)
  (0.073,0.37)rectangle++(0.009,0.009) (0.033,0.175)rectangle++(0.009,0.009) (0.533,0.082)rectangle++(0.009,0.009) (0.254,0.557)rectangle++(0.009,0.009) (0.119,0.265)rectangle++(0.009,0.009)
  (0.575,0.125)rectangle++(0.009,0.009) (0.274,0.058)rectangle++(0.009,0.009) (0.129,0.545)rectangle++(0.009,0.009) (0.06,0.259)rectangle++(0.009,0.009) (0.026,0.122)rectangle++(0.009,0.009)
  (0.01,0.576)rectangle++(0.009,0.009) (0.003,0.794)rectangle++(0.009,0.009) (-0.001,0.379)rectangle++(0.009,0.009) (-0.003,0.7)rectangle++(0.009,0.009) (-0.004,0.853)rectangle++(0.009,0.009)
  (0.516,0.407)rectangle++(0.009,0.009) (0.245,0.193)rectangle++(0.009,0.009) (0.115,0.09)rectangle++(0.009,0.009) (0.053,0.041)rectangle++(0.009,0.009) (0.543,0.017)rectangle++(0.009,0.009)
  (0.778,0.006)rectangle++(0.009,0.009) (0.371,0.001)rectangle++(0.009,0.009) (0.176,0.518)rectangle++(0.009,0.009) (0.082,0.766)rectangle++(0.009,0.009) (0.037,0.365)rectangle++(0.009,0.009)
  (0.015,0.173)rectangle++(0.009,0.009) (0.005,0.601)rectangle++(0.009,0.009) (0,0.806)rectangle++(0.009,0.009) (-0.002,0.385)rectangle++(0.009,0.009) (-0.003,0.182)rectangle++(0.009,0.009)
  (-0.004,0.605)rectangle++(0.009,0.009) (-0.004,0.808)rectangle++(0.009,0.009) (0.516,0.386)rectangle++(0.009,0.009) (0.245,0.703)rectangle++(0.009,0.009) (0.635,0.335)rectangle++(0.009,0.009)
  (0.823,0.158)rectangle++(0.009,0.009) (0.913,0.074)rectangle++(0.009,0.009) (0.436,0.033)rectangle++(0.009,0.009) (0.207,0.014)rectangle++(0.009,0.009) (0.097,0.004)rectangle++(0.009,0.009)
  (0.044,0.52)rectangle++(0.009,0.009) (0.539,0.247)rectangle++(0.009,0.009) (0.256,0.636)rectangle++(0.009,0.009) (0.121,0.823)rectangle++(0.009,0.009) (0.576,0.393)rectangle++(0.009,0.009)
  (0.274,0.706)rectangle++(0.009,0.009) (0.129,0.337)rectangle++(0.009,0.009) (0.06,0.679)rectangle++(0.009,0.009) (0.026,0.324)rectangle++(0.009,0.009) (0.53,0.153)rectangle++(0.009,0.009)
  (0.252,0.591)rectangle++(0.009,0.009) (0.639,0.281)rectangle++(0.009,0.009) (0.304,0.653)rectangle++(0.009,0.009) (0.144,0.831)rectangle++(0.009,0.009) (0.067,0.917)rectangle++(0.009,0.009)
  (0.03,0.958)rectangle++(0.009,0.009) (0.012,0.457)rectangle++(0.009,0.009) (0.523,0.217)rectangle++(0.009,0.009) (0.249,0.102)rectangle++(0.009,0.009) (0.117,0.567)rectangle++(0.009,0.009)
  (0.574,0.27)rectangle++(0.009,0.009) (0.793,0.127)rectangle++(0.009,0.009) (0.378,0.059)rectangle++(0.009,0.009) (0.699,0.026)rectangle++(0.009,0.009) (0.853,0.01)rectangle++(0.009,0.009)
  (0.407,0.002)rectangle++(0.009,0.009) (0.193,-0.001)rectangle++(0.009,0.009) (0.61,-0.003)rectangle++(0.009,0.009) (0.811,-0.004)rectangle++(0.009,0.009) (0.387,-0.004)rectangle++(0.009,0.009)
  (0.183,0.516)rectangle++(0.009,0.009) (0.086,0.245)rectangle++(0.009,0.009) (0.039,0.635)rectangle++(0.009,0.009) (0.536,0.303)rectangle++(0.009,0.009) (0.255,0.663)rectangle++(0.009,0.009)
  (0.12,0.836)rectangle++(0.009,0.009) (0.575,0.399)rectangle++(0.009,0.009) (0.794,0.189)rectangle++(0.009,0.009) (0.379,0.608)rectangle++(0.009,0.009) (0.699,0.29)rectangle++(0.009,0.009)
  (0.853,0.137)rectangle++(0.009,0.009) (0.407,0.583)rectangle++(0.009,0.009) (0.193,0.798)rectangle++(0.009,0.009) (0.61,0.381)rectangle++(0.009,0.009) (0.291,0.7)rectangle++(0.009,0.009)
  (0.137,0.334)rectangle++(0.009,0.009) (0.064,0.158)rectangle++(0.009,0.009) (0.548,0.073)rectangle++(0.009,0.009) (0.261,0.033)rectangle++(0.009,0.009) (0.123,0.013)rectangle++(0.009,0.009)
  (0.057,0.524)rectangle++(0.009,0.009) (0.025,0.249)rectangle++(0.009,0.009) (0.01,0.117)rectangle++(0.009,0.009) (0.522,0.054)rectangle++(0.009,0.009) (0.768,0.024)rectangle++(0.009,0.009)
  (0.366,0.529)rectangle++(0.009,0.009) (0.694,0.252)rectangle++(0.009,0.009) (0.331,0.638)rectangle++(0.009,0.009) (0.156,0.304)rectangle++(0.009,0.009) (0.073,0.664)rectangle++(0.009,0.009)
  (0.033,0.836)rectangle++(0.009,0.009) (0.013,0.919)rectangle++(0.009,0.009) (0.004,0.959)rectangle++(0.009,0.009) (0.52,0.458)rectangle++(0.009,0.009) (0.767,0.217)rectangle++(0.009,0.009)
  (0.886,0.102)rectangle++(0.009,0.009) (0.423,0.047)rectangle++(0.009,0.009) (0.201,0.54)rectangle++(0.009,0.009) (0.614,0.257)rectangle++(0.009,0.009) (0.292,0.641)rectangle++(0.009,0.009)
  (0.138,0.305)rectangle++(0.009,0.009) (0.064,0.144)rectangle++(0.009,0.009) (0.028,0.587)rectangle++(0.009,0.009) (0.011,0.799)rectangle++(0.009,0.009) (0.003,0.381)rectangle++(0.009,0.009)
  (0.519,0.181)rectangle++(0.009,0.009) (0.767,0.084)rectangle++(0.009,0.009) (0.366,0.038)rectangle++(0.009,0.009) (0.693,0.016)rectangle++(0.009,0.009) (0.85,0.005)rectangle++(0.009,0.009)
  (0.926,0)rectangle++(0.009,0.009) (0.962,-0.002)rectangle++(0.009,0.009) (0.459,0.517)rectangle++(0.009,0.009) (0.218,0.246)rectangle++(0.009,0.009) (0.102,0.636)rectangle++(0.009,0.009)
  (0.567,0.303)rectangle++(0.009,0.009) (0.79,0.143)rectangle++(0.009,0.009) (0.377,0.066)rectangle++(0.009,0.009) (0.178,0.029)rectangle++(0.009,0.009) (0.083,0.532)rectangle++(0.009,0.009)
  (0.038,0.253)rectangle++(0.009,0.009) (0.536,0.119)rectangle++(0.009,0.009) (0.775,0.055)rectangle++(0.009,0.009) (0.37,0.544)rectangle++(0.009,0.009) (0.175,0.779)rectangle++(0.009,0.009)
  (0.602,0.371)rectangle++(0.009,0.009) (0.806,0.176)rectangle++(0.009,0.009) (0.385,0.602)rectangle++(0.009,0.009) (0.702,0.287)rectangle++(0.009,0.009) (0.335,0.135)rectangle++(0.009,0.009)
  (0.158,0.063)rectangle++(0.009,0.009) (0.074,0.548)rectangle++(0.009,0.009) (0.033,0.781)rectangle++(0.009,0.009) (0.534,0.372)rectangle++(0.009,0.009) (0.254,0.176)rectangle++(0.009,0.009)
  (0.639,0.082)rectangle++(0.009,0.009) (0.825,0.037)rectangle++(0.009,0.009) (0.393,0.016)rectangle++(0.009,0.009) (0.187,0.525)rectangle++(0.009,0.009) (0.087,0.25)rectangle++(0.009,0.009)
  (0.56,0.118)rectangle++(0.009,0.009) (0.266,0.574)rectangle++(0.009,0.009) (0.125,0.273)rectangle++(0.009,0.009) (0.578,0.129)rectangle++(0.009,0.009) (0.275,0.579)rectangle++(0.009,0.009)
  (0.13,0.276)rectangle++(0.009,0.009) (0.06,0.65)rectangle++(0.009,0.009) (0.546,0.31)rectangle++(0.009,0.009) (0.26,0.666)rectangle++(0.009,0.009) (0.642,0.317)rectangle++(0.009,0.009)
  (0.306,0.15)rectangle++(0.009,0.009) (0.145,0.07)rectangle++(0.009,0.009) (0.067,0.031)rectangle++(0.009,0.009) (0.55,0.013)rectangle++(0.009,0.009) (0.262,0.524)rectangle++(0.009,0.009);
\fill[nuage]
  (0.643,0.249)rectangle++(0.009,0.009) (0.306,0.117)rectangle++(0.009,0.009) (0.145,0.574)rectangle++(0.009,0.009) (0.067,0.793)rectangle++(0.009,0.009) (0.03,0.898)rectangle++(0.009,0.009)
  (0.532,0.429)rectangle++(0.009,0.009) (0.253,0.724)rectangle++(0.009,0.009) (0.119,0.345)rectangle++(0.009,0.009) (0.055,0.163)rectangle++(0.009,0.009) (0.024,0.596)rectangle++(0.009,0.009)
  (0.529,0.284)rectangle++(0.009,0.009) (0.252,0.654)rectangle++(0.009,0.009) (0.118,0.312)rectangle++(0.009,0.009) (0.575,0.147)rectangle++(0.009,0.009) (0.273,0.588)rectangle++(0.009,0.009)
  (0.129,0.28)rectangle++(0.009,0.009) (0.06,0.132)rectangle++(0.009,0.009) (0.026,0.581)rectangle++(0.009,0.009) (0.53,0.277)rectangle++(0.009,0.009) (0.252,0.65)rectangle++(0.009,0.009)
  (0.119,0.31)rectangle++(0.009,0.009) (0.055,0.146)rectangle++(0.009,0.009) (0.024,0.068)rectangle++(0.009,0.009) (0.009,0.03)rectangle++(0.009,0.009) (0.002,0.012)rectangle++(0.009,0.009)
  (0.519,0.004)rectangle++(0.009,0.009) (0.247,0.519)rectangle++(0.009,0.009) (0.116,0.247)rectangle++(0.009,0.009) (0.053,0.636)rectangle++(0.009,0.009) (0.023,0.303)rectangle++(0.009,0.009)
  (0.529,0.143)rectangle++(0.009,0.009) (0.771,0.066)rectangle++(0.009,0.009) (0.368,0.55)rectangle++(0.009,0.009) (0.174,0.261)rectangle++(0.009,0.009) (0.081,0.123)rectangle++(0.009,0.009)
  (0.037,0.057)rectangle++(0.009,0.009) (0.015,0.545)rectangle++(0.009,0.009) (0.525,0.259)rectangle++(0.009,0.009) (0.25,0.642)rectangle++(0.009,0.009) (0.117,0.826)rectangle++(0.009,0.009)
  (0.574,0.394)rectangle++(0.009,0.009) (0.793,0.187)rectangle++(0.009,0.009) (0.378,0.607)rectangle++(0.009,0.009) (0.179,0.809)rectangle++(0.009,0.009) (0.084,0.906)rectangle++(0.009,0.009)
  (0.038,0.433)rectangle++(0.009,0.009) (0.536,0.205)rectangle++(0.009,0.009) (0.775,0.096)rectangle++(0.009,0.009) (0.37,0.564)rectangle++(0.009,0.009) (0.695,0.268)rectangle++(0.009,0.009)
  (0.851,0.126)rectangle++(0.009,0.009) (0.406,0.058)rectangle++(0.009,0.009) (0.193,0.026)rectangle++(0.009,0.009) (0.09,0.01)rectangle++(0.009,0.009) (0.561,0.002)rectangle++(0.009,0.009)
  (0.267,-0.001)rectangle++(0.009,0.009) (0.126,0.517)rectangle++(0.009,0.009) (0.058,0.766)rectangle++(0.009,0.009) (0.026,0.885)rectangle++(0.009,0.009) (0.01,0.423)rectangle++(0.009,0.009)
  (0.002,0.201)rectangle++(0.009,0.009) (0.519,0.094)rectangle++(0.009,0.009) (0.247,0.043)rectangle++(0.009,0.009) (0.636,0.018)rectangle++(0.009,0.009) (0.303,0.006)rectangle++(0.009,0.009)
  (0.143,0.001)rectangle++(0.009,0.009) (0.066,-0.002)rectangle++(0.009,0.009) (0.03,-0.003)rectangle++(0.009,0.009) (0.012,0.516)rectangle++(0.009,0.009) (0.003,0.245)rectangle++(0.009,0.009)
  (-0.001,0.115)rectangle++(0.009,0.009) (-0.003,0.053)rectangle++(0.009,0.009) (-0.004,0.543)rectangle++(0.009,0.009) (-0.004,0.778)rectangle++(0.009,0.009) (0.516,0.371)rectangle++(0.009,0.009)
  (0.245,0.696)rectangle++(0.009,0.009) (0.635,0.332)rectangle++(0.009,0.009) (0.823,0.157)rectangle++(0.009,0.009) (0.393,0.593)rectangle++(0.009,0.009) (0.186,0.282)rectangle++(0.009,0.009)
  (0.087,0.653)rectangle++(0.009,0.009) (0.039,0.831)rectangle++(0.009,0.009) (0.017,0.917)rectangle++(0.009,0.009) (0.526,0.438)rectangle++(0.009,0.009) (0.77,0.208)rectangle++(0.009,0.009)
  (0.367,0.617)rectangle++(0.009,0.009) (0.174,0.814)rectangle++(0.009,0.009) (0.081,0.388)rectangle++(0.009,0.009) (0.557,0.184)rectangle++(0.009,0.009) (0.785,0.086)rectangle++(0.009,0.009)
  (0.374,0.039)rectangle++(0.009,0.009) (0.177,0.536)rectangle++(0.009,0.009) (0.603,0.255)rectangle++(0.009,0.009) (0.287,0.64)rectangle++(0.009,0.009) (0.135,0.825)rectangle++(0.009,0.009)
  (0.063,0.914)rectangle++(0.009,0.009) (0.028,0.436)rectangle++(0.009,0.009) (0.011,0.727)rectangle++(0.009,0.009) (0.003,0.867)rectangle++(0.009,0.009) (0.519,0.414)rectangle++(0.009,0.009)
  (0.767,0.196)rectangle++(0.009,0.009) (0.366,0.092)rectangle++(0.009,0.009) (0.173,0.042)rectangle++(0.009,0.009) (0.081,0.018)rectangle++(0.009,0.009) (0.036,0.006)rectangle++(0.009,0.009)
  (0.535,0.001)rectangle++(0.009,0.009) (0.255,-0.002)rectangle++(0.009,0.009) (0.12,-0.003)rectangle++(0.009,0.009) (0.055,0.516)rectangle++(0.009,0.009) (0.024,0.245)rectangle++(0.009,0.009)
  (0.009,0.115)rectangle++(0.009,0.009) (0.002,0.573)rectangle++(0.009,0.009) (-0.001,0.273)rectangle++(0.009,0.009) (-0.003,0.649)rectangle++(0.009,0.009) (0.516,0.309)rectangle++(0.009,0.009)
  (0.245,0.666)rectangle++(0.009,0.009) (0.635,0.317)rectangle++(0.009,0.009) (0.303,0.15)rectangle++(0.009,0.009) (0.143,0.07)rectangle++(0.009,0.009) (0.586,0.031)rectangle++(0.009,0.009)
  (0.279,0.533)rectangle++(0.009,0.009) (0.132,0.773)rectangle++(0.009,0.009) (0.581,0.369)rectangle++(0.009,0.009) (0.276,0.695)rectangle++(0.009,0.009) (0.13,0.331)rectangle++(0.009,0.009)
  (0.06,0.157)rectangle++(0.009,0.009) (0.027,0.073)rectangle++(0.009,0.009) (0.01,0.553)rectangle++(0.009,0.009) (0.003,0.783)rectangle++(0.009,0.009) (0.519,0.373)rectangle++(0.009,0.009)
  (0.247,0.177)rectangle++(0.009,0.009) (0.636,0.083)rectangle++(0.009,0.009) (0.303,0.557)rectangle++(0.009,0.009) (0.143,0.785)rectangle++(0.009,0.009) (0.586,0.375)rectangle++(0.009,0.009)
  (0.799,0.177)rectangle++(0.009,0.009) (0.381,0.083)rectangle++(0.009,0.009) (0.181,0.557)rectangle++(0.009,0.009) (0.604,0.265)rectangle++(0.009,0.009) (0.288,0.645)rectangle++(0.009,0.009)
  (0.136,0.307)rectangle++(0.009,0.009) (0.583,0.145)rectangle++(0.009,0.009) (0.277,0.587)rectangle++(0.009,0.009) (0.131,0.28)rectangle++(0.009,0.009) (0.06,0.652)rectangle++(0.009,0.009)
  (0.027,0.311)rectangle++(0.009,0.009) (0.01,0.147)rectangle++(0.009,0.009) (0.003,0.068)rectangle++(0.009,0.009) (-0.001,0.03)rectangle++(0.009,0.009) (-0.003,0.012)rectangle++(0.009,0.009)
  (-0.004,0.524)rectangle++(0.009,0.009) (-0.004,0.769)rectangle++(0.009,0.009) (-0.004,0.887)rectangle++(0.009,0.009) (-0.004,0.943)rectangle++(0.009,0.009) (0.516,0.45)rectangle++(0.009,0.009)
  (0.245,0.734)rectangle++(0.009,0.009) (0.115,0.35)rectangle++(0.009,0.009) (0.053,0.686)rectangle++(0.009,0.009) (0.023,0.847)rectangle++(0.009,0.009) (0.009,0.404)rectangle++(0.009,0.009)
  (0.002,0.712)rectangle++(0.009,0.009) (0.519,0.339)rectangle++(0.009,0.009) (0.247,0.68)rectangle++(0.009,0.009) (0.116,0.324)rectangle++(0.009,0.009) (0.573,0.153)rectangle++(0.009,0.009)
  (0.273,0.591)rectangle++(0.009,0.009) (0.649,0.281)rectangle++(0.009,0.009) (0.309,0.133)rectangle++(0.009,0.009) (0.666,0.061)rectangle++(0.009,0.009) (0.837,0.027)rectangle++(0.009,0.009)
  (0.4,0.531)rectangle++(0.009,0.009) (0.189,0.252)rectangle++(0.009,0.009) (0.089,0.119)rectangle++(0.009,0.009) (0.56,0.055)rectangle++(0.009,0.009) (0.267,0.544)rectangle++(0.009,0.009)
  (0.126,0.259)rectangle++(0.009,0.009) (0.058,0.122)rectangle++(0.009,0.009) (0.025,0.056)rectangle++(0.009,0.009) (0.53,0.025)rectangle++(0.009,0.009) (0.772,0.009)rectangle++(0.009,0.009)
  (0.888,0.002)rectangle++(0.009,0.009) (0.424,0.519)rectangle++(0.009,0.009) (0.201,0.247)rectangle++(0.009,0.009) (0.614,0.116)rectangle++(0.009,0.009) (0.292,0.573)rectangle++(0.009,0.009)
  (0.138,0.793)rectangle++(0.009,0.009) (0.584,0.378)rectangle++(0.009,0.009) (0.278,0.699)rectangle++(0.009,0.009) (0.131,0.333)rectangle++(0.009,0.009) (0.061,0.678)rectangle++(0.009,0.009)
  (0.027,0.323)rectangle++(0.009,0.009) (0.01,0.673)rectangle++(0.009,0.009) (0.003,0.321)rectangle++(0.009,0.009) (-0.001,0.152)rectangle++(0.009,0.009) (0.517,0.07)rectangle++(0.009,0.009)
  (0.246,0.031)rectangle++(0.009,0.009) (0.116,0.013)rectangle++(0.009,0.009) (0.573,0.004)rectangle++(0.009,0.009) (0.793,-0.001)rectangle++(0.009,0.009) (0.378,0.517)rectangle++(0.009,0.009)
  (0.699,0.246)rectangle++(0.009,0.009) (0.333,0.636)rectangle++(0.009,0.009) (0.158,0.303)rectangle++(0.009,0.009) (0.073,0.143)rectangle++(0.009,0.009) (0.553,0.066)rectangle++(0.009,0.009)
  (0.263,0.029)rectangle++(0.009,0.009) (0.644,0.012)rectangle++(0.009,0.009) (0.307,0.523)rectangle++(0.009,0.009) (0.145,0.249)rectangle++(0.009,0.009) (0.067,0.117)rectangle++(0.009,0.009)
  (0.55,0.054)rectangle++(0.009,0.009) (0.262,0.024)rectangle++(0.009,0.009) (0.643,0.009)rectangle++(0.009,0.009) (0.826,0.002)rectangle++(0.009,0.009) (0.914,-0.001)rectangle++(0.009,0.009)
  (0.957,-0.003)rectangle++(0.009,0.009) (0.457,-0.004)rectangle++(0.009,0.009) (0.217,0.516)rectangle++(0.009,0.009) (0.102,0.245)rectangle++(0.009,0.009) (0.047,0.635)rectangle++(0.009,0.009)
  (0.02,0.823)rectangle++(0.009,0.009) (0.527,0.393)rectangle++(0.009,0.009) (0.251,0.706)rectangle++(0.009,0.009) (0.118,0.337)rectangle++(0.009,0.009) (0.054,0.159)rectangle++(0.009,0.009)
  (0.024,0.594)rectangle++(0.009,0.009) (0.529,0.283)rectangle++(0.009,0.009) (0.252,0.133)rectangle++(0.009,0.009) (0.118,0.062)rectangle++(0.009,0.009) (0.055,0.027)rectangle++(0.009,0.009)
  (0.024,0.531)rectangle++(0.009,0.009) (0.529,0.252)rectangle++(0.009,0.009) (0.252,0.639)rectangle++(0.009,0.009) (0.118,0.304)rectangle++(0.009,0.009) (0.575,0.144)rectangle++(0.009,0.009)
  (0.273,0.067)rectangle++(0.009,0.009) (0.129,0.03)rectangle++(0.009,0.009) (0.06,0.532)rectangle++(0.009,0.009) (0.546,0.253)rectangle++(0.009,0.009) (0.26,0.119)rectangle++(0.009,0.009)
  (0.122,0.055)rectangle++(0.009,0.009) (0.576,0.024)rectangle++(0.009,0.009) (0.794,0.009)rectangle++(0.009,0.009) (0.379,0.002)rectangle++(0.009,0.009) (0.18,-0.001)rectangle++(0.009,0.009)
  (0.084,0.517)rectangle++(0.009,0.009) (0.558,0.246)rectangle++(0.009,0.009) (0.265,0.116)rectangle++(0.009,0.009) (0.125,0.053)rectangle++(0.009,0.009) (0.058,0.543)rectangle++(0.009,0.009)
  (0.025,0.778)rectangle++(0.009,0.009) (0.01,0.891)rectangle++(0.009,0.009) (0.522,0.425)rectangle++(0.009,0.009) (0.248,0.722)rectangle++(0.009,0.009) (0.117,0.344)rectangle++(0.009,0.009)
  (0.054,0.683)rectangle++(0.009,0.009) (0.023,0.845)rectangle++(0.009,0.009) (0.529,0.403)rectangle++(0.009,0.009) (0.772,0.191)rectangle++(0.009,0.009) (0.368,0.089)rectangle++(0.009,0.009)
  (0.174,0.561)rectangle++(0.009,0.009) (0.081,0.787)rectangle++(0.009,0.009) (0.037,0.895)rectangle++(0.009,0.009) (0.015,0.947)rectangle++(0.009,0.009) (0.005,0.972)rectangle++(0.009,0.009)
  (0.52,0.464)rectangle++(0.009,0.009) (0.247,0.221)rectangle++(0.009,0.009) (0.116,0.104)rectangle++(0.009,0.009) (0.054,0.567)rectangle++(0.009,0.009) (0.023,0.27)rectangle++(0.009,0.009)
  (0.529,0.127)rectangle++(0.009,0.009) (0.252,0.059)rectangle++(0.009,0.009) (0.118,0.026)rectangle++(0.009,0.009) (0.054,0.53)rectangle++(0.009,0.009) (0.544,0.252)rectangle++(0.009,0.009)
  (0.779,0.119)rectangle++(0.009,0.009) (0.371,0.055)rectangle++(0.009,0.009) (0.176,0.544)rectangle++(0.009,0.009) (0.082,0.779)rectangle++(0.009,0.009) (0.557,0.371)rectangle++(0.009,0.009)
  (0.265,0.176)rectangle++(0.009,0.009) (0.125,0.602)rectangle++(0.009,0.009) (0.058,0.807)rectangle++(0.009,0.009) (0.545,0.385)rectangle++(0.009,0.009) (0.259,0.702)rectangle++(0.009,0.009)
  (0.642,0.335)rectangle++(0.009,0.009) (0.826,0.158)rectangle++(0.009,0.009) (0.394,0.074)rectangle++(0.009,0.009) (0.187,0.553)rectangle++(0.009,0.009) (0.087,0.783)rectangle++(0.009,0.009)
  (0.04,0.374)rectangle++(0.009,0.009) (0.017,0.177)rectangle++(0.009,0.009) (0.006,0.603)rectangle++(0.009,0.009) (0,0.807)rectangle++(0.009,0.009) (-0.002,0.905)rectangle++(0.009,0.009)
  (-0.003,0.952)rectangle++(0.009,0.009) (0.516,0.455)rectangle++(0.009,0.009) (0.765,0.216)rectangle++(0.009,0.009) (0.885,0.101)rectangle++(0.009,0.009) (0.942,0.046)rectangle++(0.009,0.009)
  (0.97,0.02)rectangle++(0.009,0.009) (0.463,0.007)rectangle++(0.009,0.009) (0.22,0.521)rectangle++(0.009,0.009) (0.623,0.248)rectangle++(0.009,0.009) (0.817,0.117)rectangle++(0.009,0.009)
  (0.39,0.574)rectangle++(0.009,0.009) (0.185,0.793)rectangle++(0.009,0.009) (0.606,0.378)rectangle++(0.009,0.009) (0.289,0.699)rectangle++(0.009,0.009) (0.136,0.333)rectangle++(0.009,0.009)
  (0.063,0.158)rectangle++(0.009,0.009) (0.028,0.073)rectangle++(0.009,0.009) (0.531,0.033)rectangle++(0.009,0.009) (0.253,0.533)rectangle++(0.009,0.009) (0.639,0.254)rectangle++(0.009,0.009)
  (0.304,0.119)rectangle++(0.009,0.009) (0.664,0.055)rectangle++(0.009,0.009) (0.836,0.024)rectangle++(0.009,0.009) (0.399,0.009)rectangle++(0.009,0.009) (0.709,0.002)rectangle++(0.009,0.009)
  (0.858,-0.001)rectangle++(0.009,0.009) (0.93,-0.003)rectangle++(0.009,0.009) (0.964,-0.004)rectangle++(0.009,0.009) (0.98,-0.004)rectangle++(0.009,0.009) (0.468,0.516)rectangle++(0.009,0.009)
  (0.742,0.245)rectangle++(0.009,0.009) (0.354,0.635)rectangle++(0.009,0.009) (0.168,0.823)rectangle++(0.009,0.009) (0.078,0.913)rectangle++(0.009,0.009) (0.555,0.436)rectangle++(0.009,0.009)
  (0.264,0.727)rectangle++(0.009,0.009) (0.124,0.347)rectangle++(0.009,0.009) (0.057,0.684)rectangle++(0.009,0.009) (0.545,0.326)rectangle++(0.009,0.009) (0.779,0.154)rectangle++(0.009,0.009)
  (0.372,0.072)rectangle++(0.009,0.009) (0.176,0.552)rectangle++(0.009,0.009) (0.602,0.263)rectangle++(0.009,0.009) (0.807,0.124)rectangle++(0.009,0.009) (0.385,0.057)rectangle++(0.009,0.009)
  (0.702,0.025)rectangle++(0.009,0.009) (0.335,0.01)rectangle++(0.009,0.009) (0.678,0.002)rectangle++(0.009,0.009) (0.843,-0.001)rectangle++(0.009,0.009) (0.402,0.517)rectangle++(0.009,0.009)
  (0.711,0.246)rectangle++(0.009,0.009) (0.859,0.116)rectangle++(0.009,0.009) (0.41,0.053)rectangle++(0.009,0.009) (0.194,0.543)rectangle++(0.009,0.009) (0.091,0.778)rectangle++(0.009,0.009)
  (0.041,0.891)rectangle++(0.009,0.009) (0.537,0.425)rectangle++(0.009,0.009) (0.256,0.202)rectangle++(0.009,0.009) (0.64,0.095)rectangle++(0.009,0.009) (0.305,0.563)rectangle++(0.009,0.009)
  (0.144,0.788)rectangle++(0.009,0.009) (0.587,0.376)rectangle++(0.009,0.009) (0.279,0.178)rectangle++(0.009,0.009) (0.132,0.603)rectangle++(0.009,0.009) (0.061,0.807)rectangle++(0.009,0.009)
  (0.027,0.385)rectangle++(0.009,0.009) (0.011,0.183)rectangle++(0.009,0.009) (0.003,0.085)rectangle++(0.009,0.009) (-0.001,0.039)rectangle++(0.009,0.009) (-0.003,0.536)rectangle++(0.009,0.009)
  (-0.004,0.255)rectangle++(0.009,0.009) (-0.004,0.12)rectangle++(0.009,0.009) (0.516,0.055)rectangle++(0.009,0.009) (0.765,0.024)rectangle++(0.009,0.009) (0.885,0.009)rectangle++(0.009,0.009)
  (0.942,0.002)rectangle++(0.009,0.009) (0.45,0.519)rectangle++(0.009,0.009) (0.214,0.247)rectangle++(0.009,0.009) (0.1,0.636)rectangle++(0.009,0.009) (0.566,0.303)rectangle++(0.009,0.009)
  (0.269,0.143)rectangle++(0.009,0.009) (0.127,0.066)rectangle++(0.009,0.009) (0.059,0.55)rectangle++(0.009,0.009) (0.026,0.781)rectangle++(0.009,0.009) (0.01,0.373)rectangle++(0.009,0.009)
  (0.002,0.177)rectangle++(0.009,0.009) (0.519,0.082)rectangle++(0.009,0.009) (0.767,0.037)rectangle++(0.009,0.009) (0.366,0.536)rectangle++(0.009,0.009) (0.173,0.255)rectangle++(0.009,0.009)
  (0.601,0.12)rectangle++(0.009,0.009) (0.286,0.055)rectangle++(0.009,0.009) (0.135,0.024)rectangle++(0.009,0.009) (0.582,0.009)rectangle++(0.009,0.009) (0.277,0.002)rectangle++(0.009,0.009)
  (0.131,0.519)rectangle++(0.009,0.009) (0.58,0.247)rectangle++(0.009,0.009) (0.796,0.116)rectangle++(0.009,0.009) (0.38,0.573)rectangle++(0.009,0.009) (0.18,0.273)rectangle++(0.009,0.009)
  (0.604,0.129)rectangle++(0.009,0.009) (0.808,0.059)rectangle++(0.009,0.009) (0.385,0.026)rectangle++(0.009,0.009) (0.703,0.01)rectangle++(0.009,0.009) (0.855,0.003)rectangle++(0.009,0.009)
  (0.408,0.519)rectangle++(0.009,0.009) (0.714,0.247)rectangle++(0.009,0.009) (0.34,0.636)rectangle++(0.009,0.009) (0.161,0.823)rectangle++(0.009,0.009) (0.075,0.913)rectangle++(0.009,0.009)
  (0.554,0.436)rectangle++(0.009,0.009) (0.783,0.207)rectangle++(0.009,0.009) (0.374,0.617)rectangle++(0.009,0.009) (0.177,0.294)rectangle++(0.009,0.009) (0.083,0.659)rectangle++(0.009,0.009)
  (0.557,0.314)rectangle++(0.009,0.009) (0.785,0.148)rectangle++(0.009,0.009) (0.375,0.589)rectangle++(0.009,0.009) (0.697,0.28)rectangle++(0.009,0.009) (0.852,0.132)rectangle++(0.009,0.009)
  (0.407,0.581)rectangle++(0.009,0.009) (0.193,0.277)rectangle++(0.009,0.009) (0.09,0.65)rectangle++(0.009,0.009) (0.041,0.31)rectangle++(0.009,0.009) (0.017,0.146)rectangle++(0.009,0.009)
  (0.006,0.068)rectangle++(0.009,0.009) (0.001,0.03)rectangle++(0.009,0.009) (0.518,0.012)rectangle++(0.009,0.009) (0.246,0.524)rectangle++(0.009,0.009) (0.116,0.249)rectangle++(0.009,0.009)
  (0.053,0.637)rectangle++(0.009,0.009) (0.543,0.303)rectangle++(0.009,0.009) (0.778,0.143)rectangle++(0.009,0.009) (0.371,0.066)rectangle++(0.009,0.009) (0.176,0.55)rectangle++(0.009,0.009)
  (0.082,0.781)rectangle++(0.009,0.009) (0.557,0.373)rectangle++(0.009,0.009) (0.265,0.177)rectangle++(0.009,0.009) (0.125,0.602)rectangle++(0.009,0.009) (0.058,0.287)rectangle++(0.009,0.009)
  (0.025,0.655)rectangle++(0.009,0.009) (0.01,0.832)rectangle++(0.009,0.009) (0.002,0.397)rectangle++(0.009,0.009) (-0.001,0.708)rectangle++(0.009,0.009) (0.517,0.338)rectangle++(0.009,0.009)
  (0.766,0.16)rectangle++(0.009,0.009) (0.365,0.594)rectangle++(0.009,0.009) (0.173,0.803)rectangle++(0.009,0.009) (0.601,0.383)rectangle++(0.009,0.009) (0.286,0.182)rectangle++(0.009,0.009)
  (0.135,0.085)rectangle++(0.009,0.009) (0.062,0.038)rectangle++(0.009,0.009) (0.028,0.536)rectangle++(0.009,0.009) (0.011,0.255)rectangle++(0.009,0.009) (0.523,0.12)rectangle++(0.009,0.009)
  (0.769,0.055)rectangle++(0.009,0.009) (0.367,0.024)rectangle++(0.009,0.009) (0.174,0.009)rectangle++(0.009,0.009) (0.081,0.002)rectangle++(0.009,0.009) (0.037,0.519)rectangle++(0.009,0.009)
  (0.015,0.767)rectangle++(0.009,0.009) (0.005,0.366)rectangle++(0.009,0.009) (0.52,0.173)rectangle++(0.009,0.009) (0.247,0.601)rectangle++(0.009,0.009) (0.116,0.286)rectangle++(0.009,0.009)
  (0.054,0.655)rectangle++(0.009,0.009) (0.023,0.312)rectangle++(0.009,0.009) (0.009,0.147)rectangle++(0.009,0.009) (0.002,0.068)rectangle++(0.009,0.009) (-0.001,0.551)rectangle++(0.009,0.009)
  (-0.003,0.782)rectangle++(0.009,0.009) (-0.004,0.373)rectangle++(0.009,0.009) (-0.004,0.177)rectangle++(0.009,0.009) (-0.004,0.602)rectangle++(0.009,0.009) (0.516,0.287)rectangle++(0.009,0.009)
  (0.765,0.135)rectangle++(0.009,0.009) (0.885,0.063)rectangle++(0.009,0.009) (0.942,0.028)rectangle++(0.009,0.009) (0.45,0.011)rectangle++(0.009,0.009) (0.214,0.003)rectangle++(0.009,0.009)
  (0.1,-0.001)rectangle++(0.009,0.009) (0.046,0.517)rectangle++(0.009,0.009) (0.02,0.246)rectangle++(0.009,0.009) (0.007,0.116)rectangle++(0.009,0.009) (0.001,0.573)rectangle++(0.009,0.009)
  (-0.002,0.793)rectangle++(0.009,0.009) (-0.003,0.898)rectangle++(0.009,0.009) (-0.004,0.429)rectangle++(0.009,0.009) (0.516,0.203)rectangle++(0.009,0.009) (0.245,0.615)rectangle++(0.009,0.009)
  (0.115,0.813)rectangle++(0.009,0.009) (0.053,0.908)rectangle++(0.009,0.009) (0.543,0.433)rectangle++(0.009,0.009) (0.778,0.206)rectangle++(0.009,0.009) (0.891,0.096)rectangle++(0.009,0.009)
  (0.425,0.044)rectangle++(0.009,0.009) (0.202,0.539)rectangle++(0.009,0.009) (0.615,0.256)rectangle++(0.009,0.009) (0.293,0.641)rectangle++(0.009,0.009) (0.138,0.825)rectangle++(0.009,0.009)
  (0.584,0.394)rectangle++(0.009,0.009) (0.798,0.187)rectangle++(0.009,0.009) (0.381,0.607)rectangle++(0.009,0.009) (0.18,0.289)rectangle++(0.009,0.009) (0.604,0.136)rectangle++(0.009,0.009)
  (0.808,0.063)rectangle++(0.009,0.009) (0.905,0.028)rectangle++(0.009,0.009) (0.432,0.531)rectangle++(0.009,0.009) (0.205,0.253)rectangle++(0.009,0.009) (0.096,0.119)rectangle++(0.009,0.009)
  (0.044,0.575)rectangle++(0.009,0.009) (0.019,0.794)rectangle++(0.009,0.009) (0.007,0.379)rectangle++(0.009,0.009) (0.521,0.179)rectangle++(0.009,0.009) (0.248,0.604)rectangle++(0.009,0.009)
  (0.637,0.287)rectangle++(0.009,0.009) (0.303,0.656)rectangle++(0.009,0.009) (0.663,0.312)rectangle++(0.009,0.009) (0.316,0.148)rectangle++(0.009,0.009) (0.149,0.589)rectangle++(0.009,0.009)
  (0.069,0.28)rectangle++(0.009,0.009) (0.551,0.132)rectangle++(0.009,0.009) (0.262,0.581)rectangle++(0.009,0.009) (0.643,0.277)rectangle++(0.009,0.009) (0.307,0.65)rectangle++(0.009,0.009)
  (0.145,0.83)rectangle++(0.009,0.009) (0.067,0.396)rectangle++(0.009,0.009) (0.03,0.188)rectangle++(0.009,0.009) (0.012,0.088)rectangle++(0.009,0.009) (0.523,0.04)rectangle++(0.009,0.009)
  (0.249,0.537)rectangle++(0.009,0.009) (0.117,0.775)rectangle++(0.009,0.009) (0.054,0.89)rectangle++(0.009,0.009) (0.544,0.425)rectangle++(0.009,0.009) (0.259,0.202)rectangle++(0.009,0.009)
  (0.642,0.094)rectangle++(0.009,0.009) (0.826,0.043)rectangle++(0.009,0.009) (0.394,0.018)rectangle++(0.009,0.009) (0.707,0.006)rectangle++(0.009,0.009) (0.857,0.001)rectangle++(0.009,0.009)
  (0.409,0.518)rectangle++(0.009,0.009) (0.714,0.246)rectangle++(0.009,0.009) (0.86,0.116)rectangle++(0.009,0.009) (0.411,0.053)rectangle++(0.009,0.009) (0.195,0.023)rectangle++(0.009,0.009)
  (0.091,0.009)rectangle++(0.009,0.009) (0.041,0.002)rectangle++(0.009,0.009) (0.018,-0.001)rectangle++(0.009,0.009) (0.006,-0.003)rectangle++(0.009,0.009) (0.521,-0.004)rectangle++(0.009,0.009)
  (0.248,0.516)rectangle++(0.009,0.009) (0.636,0.245)rectangle++(0.009,0.009) (0.303,0.635)rectangle++(0.009,0.009) (0.143,0.303)rectangle++(0.009,0.009) (0.066,0.663)rectangle++(0.009,0.009)
  (0.03,0.316)rectangle++(0.009,0.009) (0.012,0.669)rectangle++(0.009,0.009) (0.523,0.319)rectangle++(0.009,0.009) (0.249,0.151)rectangle++(0.009,0.009) (0.117,0.07)rectangle++(0.009,0.009)
  (0.574,0.031)rectangle++(0.009,0.009) (0.793,0.013)rectangle++(0.009,0.009) (0.378,0.004)rectangle++(0.009,0.009) (0.179,0.519)rectangle++(0.009,0.009) (0.604,0.247)rectangle++(0.009,0.009)
  (0.287,0.116)rectangle++(0.009,0.009) (0.136,0.053)rectangle++(0.009,0.009) (0.063,0.023)rectangle++(0.009,0.009) (0.548,0.009)rectangle++(0.009,0.009) (0.261,0.002)rectangle++(0.009,0.009)
  (0.123,-0.001)rectangle++(0.009,0.009) (0.057,0.517)rectangle++(0.009,0.009) (0.025,0.246)rectangle++(0.009,0.009) (0.01,0.116)rectangle++(0.009,0.009) (0.002,0.053)rectangle++(0.009,0.009)
  (-0.001,0.543)rectangle++(0.009,0.009) (0.517,0.258)rectangle++(0.009,0.009) (0.246,0.122)rectangle++(0.009,0.009) (0.116,0.056)rectangle++(0.009,0.009) (0.573,0.025)rectangle++(0.009,0.009)
  (0.273,0.009)rectangle++(0.009,0.009) (0.129,0.002)rectangle++(0.009,0.009) (0.579,-0.001)rectangle++(0.009,0.009) (0.276,-0.003)rectangle++(0.009,0.009) (0.65,-0.004)rectangle++(0.009,0.009)
  (0.31,-0.004)rectangle++(0.009,0.009) (0.146,0.516)rectangle++(0.009,0.009) (0.068,0.765)rectangle++(0.009,0.009) (0.55,0.365)rectangle++(0.009,0.009) (0.262,0.693)rectangle++(0.009,0.009)
  (0.123,0.33)rectangle++(0.009,0.009) (0.057,0.156)rectangle++(0.009,0.009) (0.025,0.593)rectangle++(0.009,0.009) (0.01,0.802)rectangle++(0.009,0.009) (0.002,0.383)rectangle++(0.009,0.009)
  (-0.001,0.181)rectangle++(0.009,0.009) (-0.003,0.085)rectangle++(0.009,0.009) (-0.004,0.558)rectangle++(0.009,0.009) (-0.004,0.266)rectangle++(0.009,0.009) (0.516,0.125)rectangle++(0.009,0.009)
  (0.245,0.058)rectangle++(0.009,0.009) (0.635,0.025)rectangle++(0.009,0.009) (0.823,0.01)rectangle++(0.009,0.009) (0.393,0.002)rectangle++(0.009,0.009) (0.186,0.519)rectangle++(0.009,0.009)
  (0.087,0.767)rectangle++(0.009,0.009) (0.039,0.886)rectangle++(0.009,0.009) (0.537,0.423)rectangle++(0.009,0.009) (0.775,0.201)rectangle++(0.009,0.009) (0.89,0.094)rectangle++(0.009,0.009)
  (0.425,0.043)rectangle++(0.009,0.009) (0.202,0.018)rectangle++(0.009,0.009) (0.094,0.526)rectangle++(0.009,0.009) (0.043,0.77)rectangle++(0.009,0.009) (0.018,0.367)rectangle++(0.009,0.009)
  (0.526,0.174)rectangle++(0.009,0.009) (0.25,0.601)rectangle++(0.009,0.009) (0.638,0.286)rectangle++(0.009,0.009) (0.304,0.135)rectangle++(0.009,0.009) (0.143,0.062)rectangle++(0.009,0.009)
  (0.067,0.548)rectangle++(0.009,0.009) (0.03,0.781)rectangle++(0.009,0.009) (0.532,0.372)rectangle++(0.009,0.009) (0.253,0.176)rectangle++(0.009,0.009) (0.639,0.082)rectangle++(0.009,0.009)
  (0.304,0.557)rectangle++(0.009,0.009) (0.664,0.265)rectangle++(0.009,0.009) (0.316,0.125)rectangle++(0.009,0.009) (0.669,0.058)rectangle++(0.009,0.009) (0.839,0.025)rectangle++(0.009,0.009)
  (0.4,0.01)rectangle++(0.009,0.009) (0.19,0.002)rectangle++(0.009,0.009) (0.089,0.519)rectangle++(0.009,0.009) (0.56,0.247)rectangle++(0.009,0.009) (0.267,0.636)rectangle++(0.009,0.009)
  (0.646,0.303)rectangle++(0.009,0.009) (0.828,0.143)rectangle++(0.009,0.009) (0.915,0.066)rectangle++(0.009,0.009) (0.437,0.03)rectangle++(0.009,0.009) (0.727,0.012)rectangle++(0.009,0.009)
  (0.347,0.523)rectangle++(0.009,0.009) (0.164,0.769)rectangle++(0.009,0.009) (0.596,0.367)rectangle++(0.009,0.009) (0.284,0.174)rectangle++(0.009,0.009) (0.654,0.081)rectangle++(0.009,0.009)
  (0.312,0.557)rectangle++(0.009,0.009) (0.147,0.785)rectangle++(0.009,0.009) (0.588,0.374)rectangle++(0.009,0.009) (0.28,0.177)rectangle++(0.009,0.009) (0.652,0.083)rectangle++(0.009,0.009)
  (0.311,0.037)rectangle++(0.009,0.009) (0.147,0.016)rectangle++(0.009,0.009) (0.068,0.005)rectangle++(0.009,0.009) (0.03,0.52)rectangle++(0.009,0.009) (0.012,0.247)rectangle++(0.009,0.009)
  (0.004,0.116)rectangle++(0.009,0.009) (0.519,0.054)rectangle++(0.009,0.009) (0.767,0.023)rectangle++(0.009,0.009) (0.366,0.529)rectangle++(0.009,0.009) (0.693,0.252)rectangle++(0.009,0.009)
  (0.85,0.118)rectangle++(0.009,0.009) (0.406,0.574)rectangle++(0.009,0.009) (0.192,0.793)rectangle++(0.009,0.009) (0.61,0.378)rectangle++(0.009,0.009) (0.29,0.179)rectangle++(0.009,0.009)
  (0.657,0.084)rectangle++(0.009,0.009) (0.313,0.558)rectangle++(0.009,0.009) (0.148,0.265)rectangle++(0.009,0.009) (0.589,0.125)rectangle++(0.009,0.009) (0.28,0.578)rectangle++(0.009,0.009)
  (0.132,0.275)rectangle++(0.009,0.009) (0.581,0.13)rectangle++(0.009,0.009) (0.277,0.58)rectangle++(0.009,0.009) (0.65,0.276)rectangle++(0.009,0.009) (0.31,0.65)rectangle++(0.009,0.009)
  (0.146,0.83)rectangle++(0.009,0.009) (0.588,0.396)rectangle++(0.009,0.009) (0.28,0.188)rectangle++(0.009,0.009) (0.652,0.088)rectangle++(0.009,0.009) (0.311,0.04)rectangle++(0.009,0.009)
  (0.147,0.537)rectangle++(0.009,0.009) (0.588,0.255)rectangle++(0.009,0.009) (0.8,0.12)rectangle++(0.009,0.009) (0.902,0.055)rectangle++(0.009,0.009) (0.43,0.544)rectangle++(0.009,0.009)
  (0.204,0.259)rectangle++(0.009,0.009) (0.616,0.122)rectangle++(0.009,0.009) (0.813,0.056)rectangle++(0.009,0.009) (0.908,0.025)rectangle++(0.009,0.009) (0.433,0.009)rectangle++(0.009,0.009)
  (0.206,0.002)rectangle++(0.009,0.009) (0.096,-0.001)rectangle++(0.009,0.009) (0.564,-0.003)rectangle++(0.009,0.009) (0.268,0.516)rectangle++(0.009,0.009) (0.126,0.765)rectangle++(0.009,0.009)
  (0.578,0.365)rectangle++(0.009,0.009) (0.795,0.173)rectangle++(0.009,0.009) (0.899,0.081)rectangle++(0.009,0.009) (0.429,0.036)rectangle++(0.009,0.009) (0.724,0.015)rectangle++(0.009,0.009)
  (0.345,0.005)rectangle++(0.009,0.009) (0.683,0)rectangle++(0.009,0.009) (0.846,-0.002)rectangle++(0.009,0.009) (0.404,-0.003)rectangle++(0.009,0.009) (0.191,-0.004)rectangle++(0.009,0.009)
  (0.09,0.516)rectangle++(0.009,0.009) (0.041,0.765)rectangle++(0.009,0.009) (0.017,0.365)rectangle++(0.009,0.009) (0.526,0.173)rectangle++(0.009,0.009) (0.25,0.081)rectangle++(0.009,0.009)
  (0.638,0.036)rectangle++(0.009,0.009) (0.824,0.015)rectangle++(0.009,0.009) (0.393,0.005)rectangle++(0.009,0.009) (0.186,0.52)rectangle++(0.009,0.009) (0.087,0.767)rectangle++(0.009,0.009)
  (0.559,0.366)rectangle++(0.009,0.009) (0.266,0.173)rectangle++(0.009,0.009) (0.125,0.081)rectangle++(0.009,0.009) (0.578,0.036)rectangle++(0.009,0.009) (0.275,0.535)rectangle++(0.009,0.009)
  (0.13,0.255)rectangle++(0.009,0.009) (0.06,0.12)rectangle++(0.009,0.009) (0.546,0.055)rectangle++(0.009,0.009) (0.26,0.024)rectangle++(0.009,0.009) (0.122,0.529)rectangle++(0.009,0.009)
  (0.056,0.772)rectangle++(0.009,0.009) (0.025,0.888)rectangle++(0.009,0.009) (0.01,0.424)rectangle++(0.009,0.009) (0.002,0.201)rectangle++(0.009,0.009) (0.519,0.094)rectangle++(0.009,0.009)
  (0.247,0.043)rectangle++(0.009,0.009) (0.116,0.538)rectangle++(0.009,0.009) (0.053,0.256)rectangle++(0.009,0.009) (0.023,0.121)rectangle++(0.009,0.009) (0.009,0.576)rectangle++(0.009,0.009)
  (0.002,0.274)rectangle++(0.009,0.009) (-0.001,0.649)rectangle++(0.009,0.009) (0.517,0.309)rectangle++(0.009,0.009) (0.766,0.146)rectangle++(0.009,0.009) (0.885,0.068)rectangle++(0.009,0.009)
  (0.423,0.03)rectangle++(0.009,0.009) (0.2,0.012)rectangle++(0.009,0.009) (0.094,0.523)rectangle++(0.009,0.009) (0.563,0.249)rectangle++(0.009,0.009) (0.268,0.117)rectangle++(0.009,0.009)
  (0.646,0.054)rectangle++(0.009,0.009) (0.828,0.024)rectangle++(0.009,0.009) (0.395,0.009)rectangle++(0.009,0.009) (0.707,0.002)rectangle++(0.009,0.009) (0.857,-0.001)rectangle++(0.009,0.009)
  (0.409,-0.003)rectangle++(0.009,0.009) (0.194,0.516)rectangle++(0.009,0.009) (0.611,0.245)rectangle++(0.009,0.009) (0.811,0.115)rectangle++(0.009,0.009) (0.907,0.053)rectangle++(0.009,0.009)
  (0.433,0.023)rectangle++(0.009,0.009) (0.205,0.009)rectangle++(0.009,0.009) (0.096,0.002)rectangle++(0.009,0.009) (0.044,0.519)rectangle++(0.009,0.009) (0.019,0.767)rectangle++(0.009,0.009)
  (0.527,0.366)rectangle++(0.009,0.009) (0.77,0.173)rectangle++(0.009,0.009) (0.367,0.601)rectangle++(0.009,0.009) (0.694,0.286)rectangle++(0.009,0.009) (0.851,0.135)rectangle++(0.009,0.009)
  (0.406,0.062)rectangle++(0.009,0.009) (0.713,0.028)rectangle++(0.009,0.009) (0.34,0.531)rectangle++(0.009,0.009) (0.681,0.253)rectangle++(0.009,0.009) (0.844,0.119)rectangle++(0.009,0.009)
  (0.403,0.575)rectangle++(0.009,0.009) (0.191,0.274)rectangle++(0.009,0.009) (0.089,0.129)rectangle++(0.009,0.009) (0.561,0.06)rectangle++(0.009,0.009) (0.267,0.546)rectangle++(0.009,0.009)
  (0.646,0.26)rectangle++(0.009,0.009) (0.828,0.122)rectangle++(0.009,0.009) (0.915,0.056)rectangle++(0.009,0.009) (0.957,0.025)rectangle++(0.009,0.009) (0.457,0.01)rectangle++(0.009,0.009)
  (0.737,0.002)rectangle++(0.009,0.009) (0.351,-0.001)rectangle++(0.009,0.009) (0.166,0.517)rectangle++(0.009,0.009) (0.598,0.246)rectangle++(0.009,0.009) (0.284,0.636)rectangle++(0.009,0.009)
  (0.134,0.303)rectangle++(0.009,0.009) (0.062,0.663)rectangle++(0.009,0.009) (0.027,0.836)rectangle++(0.009,0.009) (0.531,0.399)rectangle++(0.009,0.009) (0.252,0.709)rectangle++(0.009,0.009)
  (0.119,0.338)rectangle++(0.009,0.009) (0.055,0.16)rectangle++(0.009,0.009) (0.544,0.074)rectangle++(0.009,0.009) (0.259,0.033)rectangle++(0.009,0.009) (0.122,0.534)rectangle++(0.009,0.009)
  (0.576,0.254)rectangle++(0.009,0.009) (0.274,0.119)rectangle++(0.009,0.009) (0.129,0.575)rectangle++(0.009,0.009) (0.06,0.794)rectangle++(0.009,0.009) (0.026,0.899)rectangle++(0.009,0.009)
  (0.01,0.949)rectangle++(0.009,0.009) (0.003,0.453)rectangle++(0.009,0.009) (-0.001,0.735)rectangle++(0.009,0.009) (0.517,0.351)rectangle++(0.009,0.009) (0.246,0.686)rectangle++(0.009,0.009)
  (0.116,0.327)rectangle++(0.009,0.009) (0.573,0.155)rectangle++(0.009,0.009) (0.273,0.592)rectangle++(0.009,0.009) (0.129,0.282)rectangle++(0.009,0.009) (0.059,0.133)rectangle++(0.009,0.009)
  (0.026,0.061)rectangle++(0.009,0.009) (0.01,0.547)rectangle++(0.009,0.009) (0.523,0.26)rectangle++(0.009,0.009) (0.248,0.123)rectangle++(0.009,0.009) (0.637,0.057)rectangle++(0.009,0.009)
  (0.303,0.545)rectangle++(0.009,0.009) (0.143,0.259)rectangle++(0.009,0.009) (0.066,0.122)rectangle++(0.009,0.009) (0.03,0.576)rectangle++(0.009,0.009) (0.532,0.274)rectangle++(0.009,0.009)
  (0.253,0.649)rectangle++(0.009,0.009) (0.639,0.309)rectangle++(0.009,0.009) (0.824,0.146)rectangle++(0.009,0.009) (0.913,0.068)rectangle++(0.009,0.009) (0.436,0.03)rectangle++(0.009,0.009)
  (0.727,0.012)rectangle++(0.009,0.009) (0.347,0.523)rectangle++(0.009,0.009) (0.684,0.249)rectangle++(0.009,0.009) (0.846,0.117)rectangle++(0.009,0.009) (0.404,0.054)rectangle++(0.009,0.009)
  (0.191,0.024)rectangle++(0.009,0.009) (0.09,0.529)rectangle++(0.009,0.009) (0.041,0.252)rectangle++(0.009,0.009) (0.017,0.118)rectangle++(0.009,0.009) (0.006,0.054)rectangle++(0.009,0.009)
  (0,0.024)rectangle++(0.009,0.009) (0.518,0.009)rectangle++(0.009,0.009) (0.246,0.002)rectangle++(0.009,0.009) (0.116,-0.001)rectangle++(0.009,0.009) (0.053,0.517)rectangle++(0.009,0.009)
  (0.023,0.246)rectangle++(0.009,0.009) (0.529,0.116)rectangle++(0.009,0.009) (0.251,0.573)rectangle++(0.009,0.009) (0.118,0.273)rectangle++(0.009,0.009) (0.054,0.649)rectangle++(0.009,0.009)
  (0.544,0.309)rectangle++(0.009,0.009) (0.779,0.146)rectangle++(0.009,0.009) (0.891,0.068)rectangle++(0.009,0.009) (0.426,0.03)rectangle++(0.009,0.009) (0.202,0.012)rectangle++(0.009,0.009)
  (0.095,0.003)rectangle++(0.009,0.009) (0.563,-0.001)rectangle++(0.009,0.009) (0.268,-0.003)rectangle++(0.009,0.009) (0.126,0.516)rectangle++(0.009,0.009) (0.058,0.246)rectangle++(0.009,0.009)
  (0.026,0.636)rectangle++(0.009,0.009) (0.53,0.303)rectangle++(0.009,0.009) (0.772,0.143)rectangle++(0.009,0.009) (0.888,0.066)rectangle++(0.009,0.009) (0.944,0.029)rectangle++(0.009,0.009)
  (0.971,0.012)rectangle++(0.009,0.009) (0.464,0.523)rectangle++(0.009,0.009) (0.74,0.249)rectangle++(0.009,0.009) (0.873,0.117)rectangle++(0.009,0.009) (0.417,0.574)rectangle++(0.009,0.009)
  (0.198,0.273)rectangle++(0.009,0.009) (0.093,0.129)rectangle++(0.009,0.009) (0.042,0.059)rectangle++(0.009,0.009) (0.018,0.026)rectangle++(0.009,0.009) (0.526,0.01)rectangle++(0.009,0.009)
  (0.77,0.003)rectangle++(0.009,0.009) (0.887,-0.001)rectangle++(0.009,0.009) (0.424,-0.003)rectangle++(0.009,0.009) (0.201,0.516)rectangle++(0.009,0.009) (0.614,0.245)rectangle++(0.009,0.009)
  (0.292,0.635)rectangle++(0.009,0.009) (0.138,0.823)rectangle++(0.009,0.009) (0.064,0.393)rectangle++(0.009,0.009) (0.028,0.706)rectangle++(0.009,0.009) (0.011,0.857)rectangle++(0.009,0.009)
  (0.003,0.929)rectangle++(0.009,0.009) (-0.001,0.443)rectangle++(0.009,0.009) (-0.003,0.731)rectangle++(0.009,0.009) (-0.004,0.348)rectangle++(0.009,0.009) (0.516,0.165)rectangle++(0.009,0.009)
  (0.245,0.597)rectangle++(0.009,0.009) (0.115,0.804)rectangle++(0.009,0.009) (0.053,0.904)rectangle++(0.009,0.009) (0.023,0.951)rectangle++(0.009,0.009) (0.009,0.974)rectangle++(0.009,0.009)
  (0.522,0.465)rectangle++(0.009,0.009) (0.248,0.221)rectangle++(0.009,0.009) (0.117,0.104)rectangle++(0.009,0.009) (0.574,0.047)rectangle++(0.009,0.009) (0.793,0.02)rectangle++(0.009,0.009)
  (0.898,0.007)rectangle++(0.009,0.009) (0.429,0.521)rectangle++(0.009,0.009) (0.724,0.248)rectangle++(0.009,0.009) (0.345,0.637)rectangle++(0.009,0.009) (0.683,0.303)rectangle++(0.009,0.009)
  (0.846,0.143)rectangle++(0.009,0.009) (0.924,0.066)rectangle++(0.009,0.009) (0.441,0.55)rectangle++(0.009,0.009) (0.209,0.261)rectangle++(0.009,0.009) (0.098,0.123)rectangle++(0.009,0.009)
  (0.045,0.057)rectangle++(0.009,0.009) (0.019,0.545)rectangle++(0.009,0.009) (0.527,0.259)rectangle++(0.009,0.009) (0.251,0.122)rectangle++(0.009,0.009) (0.638,0.056)rectangle++(0.009,0.009)
  (0.304,0.025)rectangle++(0.009,0.009) (0.664,0.009)rectangle++(0.009,0.009) (0.316,0.522)rectangle++(0.009,0.009) (0.669,0.248)rectangle++(0.009,0.009) (0.839,0.117)rectangle++(0.009,0.009)
  (0.4,0.054)rectangle++(0.009,0.009) (0.19,0.543)rectangle++(0.009,0.009) (0.089,0.779)rectangle++(0.009,0.009) (0.56,0.371)rectangle++(0.009,0.009) (0.267,0.176)rectangle++(0.009,0.009)
  (0.126,0.082)rectangle++(0.009,0.009) (0.058,0.557)rectangle++(0.009,0.009) (0.025,0.265)rectangle++(0.009,0.009) (0.53,0.125)rectangle++(0.009,0.009) (0.772,0.058)rectangle++(0.009,0.009)
  (0.368,0.025)rectangle++(0.009,0.009) (0.694,0.01)rectangle++(0.009,0.009) (0.331,0.522)rectangle++(0.009,0.009) (0.157,0.248)rectangle++(0.009,0.009) (0.593,0.117)rectangle++(0.009,0.009)
  (0.282,0.574)rectangle++(0.009,0.009) (0.653,0.273)rectangle++(0.009,0.009) (0.311,0.649)rectangle++(0.009,0.009) (0.147,0.309)rectangle++(0.009,0.009) (0.068,0.146)rectangle++(0.009,0.009)
  (0.55,0.068)rectangle++(0.009,0.009) (0.262,0.03)rectangle++(0.009,0.009) (0.643,0.012)rectangle++(0.009,0.009) (0.306,0.003)rectangle++(0.009,0.009) (0.145,-0.001)rectangle++(0.009,0.009)
  (0.067,-0.003)rectangle++(0.009,0.009) (0.03,-0.004)rectangle++(0.009,0.009) (0.532,-0.004)rectangle++(0.009,0.009) (0.773,-0.004)rectangle++(0.009,0.009) (0.369,-0.004)rectangle++(0.009,0.009)
  (0.175,-0.004)rectangle++(0.009,0.009) (0.081,0.516)rectangle++(0.009,0.009) (0.037,0.765)rectangle++(0.009,0.009) (0.535,0.365)rectangle++(0.009,0.009) (0.255,0.693)rectangle++(0.009,0.009)
  (0.64,0.33)rectangle++(0.009,0.009) (0.825,0.156)rectangle++(0.009,0.009) (0.914,0.073)rectangle++(0.009,0.009) (0.436,0.033)rectangle++(0.009,0.009) (0.207,0.013)rectangle++(0.009,0.009)
  (0.097,0.004)rectangle++(0.009,0.009) (0.564,0)rectangle++(0.009,0.009) (0.788,-0.003)rectangle++(0.009,0.009) (0.376,-0.004)rectangle++(0.009,0.009) (0.178,0.516)rectangle++(0.009,0.009)
  (0.603,0.245)rectangle++(0.009,0.009) (0.807,0.115)rectangle++(0.009,0.009) (0.385,0.573)rectangle++(0.009,0.009) (0.183,0.273)rectangle++(0.009,0.009) (0.085,0.649)rectangle++(0.009,0.009)
  (0.039,0.829)rectangle++(0.009,0.009) (0.536,0.396)rectangle++(0.009,0.009) (0.255,0.188)rectangle++(0.009,0.009) (0.12,0.608)rectangle++(0.009,0.009) (0.055,0.289)rectangle++(0.009,0.009)
  (0.024,0.137)rectangle++(0.009,0.009) (0.009,0.063)rectangle++(0.009,0.009) (0.002,0.028)rectangle++(0.009,0.009) (-0.001,0.011)rectangle++(0.009,0.009) (-0.003,0.523)rectangle++(0.009,0.009)
  (-0.004,0.249)rectangle++(0.009,0.009) (-0.004,0.637)rectangle++(0.009,0.009) (0.516,0.303)rectangle++(0.009,0.009) (0.245,0.143)rectangle++(0.009,0.009) (0.115,0.066)rectangle++(0.009,0.009)
  (0.053,0.55)rectangle++(0.009,0.009) (0.023,0.781)rectangle++(0.009,0.009) (0.009,0.373)rectangle++(0.009,0.009) (0.522,0.177)rectangle++(0.009,0.009) (0.768,0.082)rectangle++(0.009,0.009)
  (0.366,0.557)rectangle++(0.009,0.009) (0.174,0.785)rectangle++(0.009,0.009) (0.601,0.375)rectangle++(0.009,0.009) (0.806,0.177)rectangle++(0.009,0.009) (0.385,0.083)rectangle++(0.009,0.009)
  (0.182,0.037)rectangle++(0.009,0.009) (0.085,0.536)rectangle++(0.009,0.009) (0.039,0.775)rectangle++(0.009,0.009) (0.536,0.37)rectangle++(0.009,0.009) (0.255,0.175)rectangle++(0.009,0.009)
  (0.64,0.082)rectangle++(0.009,0.009) (0.305,0.557)rectangle++(0.009,0.009) (0.664,0.265)rectangle++(0.009,0.009) (0.836,0.125)rectangle++(0.009,0.009) (0.399,0.058)rectangle++(0.009,0.009)
  (0.189,0.545)rectangle++(0.009,0.009) (0.088,0.259)rectangle++(0.009,0.009) (0.04,0.642)rectangle++(0.009,0.009) (0.537,0.306)rectangle++(0.009,0.009) (0.255,0.144)rectangle++(0.009,0.009)
  (0.12,0.587)rectangle++(0.009,0.009) (0.055,0.799)rectangle++(0.009,0.009) (0.544,0.381)rectangle++(0.009,0.009) (0.779,0.181)rectangle++(0.009,0.009) (0.372,0.604)rectangle++(0.009,0.009)
  (0.176,0.808)rectangle++(0.009,0.009) (0.082,0.385)rectangle++(0.009,0.009) (0.037,0.183)rectangle++(0.009,0.009) (0.015,0.085)rectangle++(0.009,0.009) (0.005,0.559)rectangle++(0.009,0.009)
  (0.52,0.266)rectangle++(0.009,0.009) (0.767,0.125)rectangle++(0.009,0.009) (0.366,0.578)rectangle++(0.009,0.009) (0.693,0.275)rectangle++(0.009,0.009) (0.33,0.13)rectangle++(0.009,0.009)
  (0.156,0.58)rectangle++(0.009,0.009) (0.073,0.796)rectangle++(0.009,0.009) (0.553,0.38)rectangle++(0.009,0.009) (0.263,0.18)rectangle++(0.009,0.009) (0.644,0.084)rectangle++(0.009,0.009)
  (0.307,0.558)rectangle++(0.009,0.009) (0.665,0.265)rectangle++(0.009,0.009) (0.837,0.125)rectangle++(0.009,0.009) (0.399,0.058)rectangle++(0.009,0.009) (0.709,0.025)rectangle++(0.009,0.009)
  (0.338,0.01)rectangle++(0.009,0.009) (0.68,0.002)rectangle++(0.009,0.009) (0.844,-0.001)rectangle++(0.009,0.009) (0.923,-0.003)rectangle++(0.009,0.009) (0.961,-0.004)rectangle++(0.009,0.009)
  (0.459,-0.004)rectangle++(0.009,0.009) (0.218,0.516)rectangle++(0.009,0.009) (0.622,0.245)rectangle++(0.009,0.009) (0.296,0.115)rectangle++(0.009,0.009) (0.14,0.053)rectangle++(0.009,0.009)
  (0.585,0.023)rectangle++(0.009,0.009) (0.278,0.009)rectangle++(0.009,0.009) (0.131,0.002)rectangle++(0.009,0.009) (0.061,-0.001)rectangle++(0.009,0.009) (0.027,0.517)rectangle++(0.009,0.009)
  (0.531,0.246)rectangle++(0.009,0.009) (0.252,0.636)rectangle++(0.009,0.009) (0.119,0.823)rectangle++(0.009,0.009) (0.575,0.393)rectangle++(0.009,0.009) (0.794,0.186)rectangle++(0.009,0.009)
  (0.379,0.607)rectangle++(0.009,0.009) (0.179,0.289)rectangle++(0.009,0.009) (0.084,0.656)rectangle++(0.009,0.009) (0.038,0.313)rectangle++(0.009,0.009) (0.016,0.148)rectangle++(0.009,0.009)
  (0.525,0.069)rectangle++(0.009,0.009) (0.25,0.551)rectangle++(0.009,0.009) (0.118,0.782)rectangle++(0.009,0.009) (0.054,0.893)rectangle++(0.009,0.009) (0.544,0.426)rectangle++(0.009,0.009)
  (0.259,0.722)rectangle++(0.009,0.009) (0.642,0.344)rectangle++(0.009,0.009) (0.306,0.163)rectangle++(0.009,0.009) (0.144,0.076)rectangle++(0.009,0.009) (0.587,0.034)rectangle++(0.009,0.009)
  (0.279,0.534)rectangle++(0.009,0.009) (0.132,0.254)rectangle++(0.009,0.009) (0.061,0.12)rectangle++(0.009,0.009) (0.547,0.055)rectangle++(0.009,0.009) (0.26,0.024)rectangle++(0.009,0.009)
  (0.123,0.009)rectangle++(0.009,0.009) (0.056,0.002)rectangle++(0.009,0.009) (0.545,-0.001)rectangle++(0.009,0.009) (0.779,-0.003)rectangle++(0.009,0.009) (0.372,0.516)rectangle++(0.009,0.009)
  (0.696,0.245)rectangle++(0.009,0.009) (0.332,0.635)rectangle++(0.009,0.009) (0.157,0.303)rectangle++(0.009,0.009) (0.593,0.143)rectangle++(0.009,0.009) (0.282,0.066)rectangle++(0.009,0.009)
  (0.653,0.029)rectangle++(0.009,0.009) (0.311,0.532)rectangle++(0.009,0.009) (0.147,0.773)rectangle++(0.009,0.009) (0.068,0.889)rectangle++(0.009,0.009) (0.03,0.944)rectangle++(0.009,0.009)
  (0.012,0.971)rectangle++(0.009,0.009) (0.004,0.464)rectangle++(0.009,0.009) (-0.001,0.74)rectangle++(0.009,0.009) (-0.003,0.353)rectangle++(0.009,0.009) (-0.004,0.687)rectangle++(0.009,0.009)
  (0.516,0.327)rectangle++(0.009,0.009) (0.765,0.155)rectangle++(0.009,0.009) (0.365,0.592)rectangle++(0.009,0.009) (0.693,0.282)rectangle++(0.009,0.009) (0.33,0.653)rectangle++(0.009,0.009)
  (0.676,0.311)rectangle++(0.009,0.009) (0.322,0.147)rectangle++(0.009,0.009) (0.152,0.588)rectangle++(0.009,0.009) (0.591,0.28)rectangle++(0.009,0.009) (0.801,0.132)rectangle++(0.009,0.009)
  (0.382,0.581)rectangle++(0.009,0.009) (0.181,0.797)rectangle++(0.009,0.009) (0.085,0.9)rectangle++(0.009,0.009) (0.038,0.95)rectangle++(0.009,0.009) (0.016,0.973)rectangle++(0.009,0.009)
  (0.525,0.465)rectangle++(0.009,0.009) (0.25,0.741)rectangle++(0.009,0.009) (0.118,0.353)rectangle++(0.009,0.009) (0.054,0.687)rectangle++(0.009,0.009) (0.024,0.848)rectangle++(0.009,0.009)
  (0.009,0.924)rectangle++(0.009,0.009) (0.522,0.441)rectangle++(0.009,0.009) (0.248,0.21)rectangle++(0.009,0.009) (0.117,0.098)rectangle++(0.009,0.009) (0.574,0.045)rectangle++(0.009,0.009)
  (0.273,0.539)rectangle++(0.009,0.009) (0.129,0.256)rectangle++(0.009,0.009) (0.059,0.641)rectangle++(0.009,0.009) (0.546,0.305)rectangle++(0.009,0.009) (0.78,0.144)rectangle++(0.009,0.009)
  (0.372,0.587)rectangle++(0.009,0.009) (0.176,0.279)rectangle++(0.009,0.009) (0.082,0.132)rectangle++(0.009,0.009) (0.557,0.061)rectangle++(0.009,0.009) (0.265,0.027)rectangle++(0.009,0.009)
  (0.645,0.011)rectangle++(0.009,0.009) (0.307,0.003)rectangle++(0.009,0.009) (0.145,-0.001)rectangle++(0.009,0.009) (0.067,0.517)rectangle++(0.009,0.009) (0.55,0.246)rectangle++(0.009,0.009)
  (0.262,0.636)rectangle++(0.009,0.009) (0.123,0.303)rectangle++(0.009,0.009) (0.577,0.143)rectangle++(0.009,0.009) (0.275,0.586)rectangle++(0.009,0.009) (0.129,0.279)rectangle++(0.009,0.009)
  (0.58,0.132)rectangle++(0.009,0.009) (0.276,0.581)rectangle++(0.009,0.009) (0.13,0.796)rectangle++(0.009,0.009) (0.06,0.38)rectangle++(0.009,0.009) (0.027,0.18)rectangle++(0.009,0.009)
  (0.01,0.084)rectangle++(0.009,0.009) (0.003,0.038)rectangle++(0.009,0.009) (-0.001,0.016)rectangle++(0.009,0.009) (-0.003,0.525)rectangle++(0.009,0.009) (0.516,0.25)rectangle++(0.009,0.009)
  (0.245,0.638)rectangle++(0.009,0.009) (0.635,0.304)rectangle++(0.009,0.009) (0.303,0.663)rectangle++(0.009,0.009) (0.143,0.836)rectangle++(0.009,0.009) (0.586,0.399)rectangle++(0.009,0.009)
  (0.799,0.189)rectangle++(0.009,0.009) (0.381,0.608)rectangle++(0.009,0.009) (0.181,0.81)rectangle++(0.009,0.009) (0.084,0.386)rectangle++(0.009,0.009) (0.038,0.703)rectangle++(0.009,0.009)
  (0.016,0.855)rectangle++(0.009,0.009) (0.525,0.408)rectangle++(0.009,0.009) (0.25,0.714)rectangle++(0.009,0.009) (0.118,0.34)rectangle++(0.009,0.009) (0.574,0.161)rectangle++(0.009,0.009)
  (0.273,0.075)rectangle++(0.009,0.009) (0.649,0.034)rectangle++(0.009,0.009) (0.309,0.014)rectangle++(0.009,0.009) (0.146,0.004)rectangle++(0.009,0.009) (0.068,0.52)rectangle++(0.009,0.009)
  (0.55,0.247)rectangle++(0.009,0.009) (0.262,0.116)rectangle++(0.009,0.009) (0.123,0.573)rectangle++(0.009,0.009) (0.057,0.793)rectangle++(0.009,0.009) (0.545,0.378)rectangle++(0.009,0.009)
  (0.779,0.179)rectangle++(0.009,0.009) (0.372,0.604)rectangle++(0.009,0.009) (0.696,0.287)rectangle++(0.009,0.009) (0.332,0.656)rectangle++(0.009,0.009) (0.157,0.832)rectangle++(0.009,0.009)
  (0.073,0.397)rectangle++(0.009,0.009) (0.033,0.188)rectangle++(0.009,0.009) (0.013,0.088)rectangle++(0.009,0.009) (0.004,0.56)rectangle++(0.009,0.009) (0,0.786)rectangle++(0.009,0.009)
  (0.517,0.375)rectangle++(0.009,0.009) (0.246,0.178)rectangle++(0.009,0.009) (0.116,0.603)rectangle++(0.009,0.009) (0.053,0.807)rectangle++(0.009,0.009) (0.543,0.385)rectangle++(0.009,0.009)
  (0.258,0.182)rectangle++(0.009,0.009) (0.642,0.085)rectangle++(0.009,0.009) (0.826,0.039)rectangle++(0.009,0.009) (0.394,0.536)rectangle++(0.009,0.009) (0.707,0.255)rectangle++(0.009,0.009)
  (0.337,0.64)rectangle++(0.009,0.009) (0.679,0.305)rectangle++(0.009,0.009) (0.324,0.144)rectangle++(0.009,0.009) (0.153,0.587)rectangle++(0.009,0.009) (0.071,0.799)rectangle++(0.009,0.009)
  (0.032,0.901)rectangle++(0.009,0.009) (0.013,0.95)rectangle++(0.009,0.009) (0.524,0.454)rectangle++(0.009,0.009) (0.249,0.215)rectangle++(0.009,0.009) (0.637,0.101)rectangle++(0.009,0.009)
  (0.304,0.046)rectangle++(0.009,0.009) (0.143,0.54)rectangle++(0.009,0.009) (0.066,0.257)rectangle++(0.009,0.009) (0.55,0.121)rectangle++(0.009,0.009) (0.781,0.056)rectangle++(0.009,0.009)
  (0.893,0.024)rectangle++(0.009,0.009) (0.426,0.009)rectangle++(0.009,0.009) (0.202,0.002)rectangle++(0.009,0.009) (0.615,-0.001)rectangle++(0.009,0.009) (0.293,-0.003)rectangle++(0.009,0.009)
  (0.658,-0.004)rectangle++(0.009,0.009) (0.834,-0.004)rectangle++(0.009,0.009) (0.918,-0.004)rectangle++(0.009,0.009) (0.438,-0.004)rectangle++(0.009,0.009) (0.728,-0.004)rectangle++(0.009,0.009)
  (0.867,-0.004)rectangle++(0.009,0.009) (0.414,0.516)rectangle++(0.009,0.009) (0.716,0.245)rectangle++(0.009,0.009) (0.341,0.635)rectangle++(0.009,0.009) (0.682,0.303)rectangle++(0.009,0.009)
  (0.325,0.663)rectangle++(0.009,0.009) (0.674,0.316)rectangle++(0.009,0.009) (0.841,0.149)rectangle++(0.009,0.009) (0.401,0.069)rectangle++(0.009,0.009) (0.19,0.551)rectangle++(0.009,0.009)
  (0.089,0.262)rectangle++(0.009,0.009) (0.04,0.643)rectangle++(0.009,0.009) (0.537,0.307)rectangle++(0.009,0.009) (0.775,0.145)rectangle++(0.009,0.009) (0.37,0.587)rectangle++(0.009,0.009)
  (0.695,0.279)rectangle++(0.009,0.009) (0.331,0.132)rectangle++(0.009,0.009) (0.677,0.061)rectangle++(0.009,0.009) (0.322,0.547)rectangle++(0.009,0.009) (0.152,0.26)rectangle++(0.009,0.009)
  (0.071,0.123)rectangle++(0.009,0.009) (0.032,0.056)rectangle++(0.009,0.009) (0.013,0.545)rectangle++(0.009,0.009) (0.004,0.259)rectangle++(0.009,0.009) (-0.001,0.122)rectangle++(0.009,0.009)
  (-0.003,0.056)rectangle++(0.009,0.009) (-0.004,0.545)rectangle++(0.009,0.009) (-0.004,0.259)rectangle++(0.009,0.009) (0.516,0.122)rectangle++(0.009,0.009) (0.245,0.576)rectangle++(0.009,0.009)
  (0.115,0.274)rectangle++(0.009,0.009) (0.053,0.649)rectangle++(0.009,0.009) (0.023,0.829)rectangle++(0.009,0.009) (0.009,0.916)rectangle++(0.009,0.009) (0.002,0.437)rectangle++(0.009,0.009)
  (0.519,0.208)rectangle++(0.009,0.009) (0.247,0.097)rectangle++(0.009,0.009) (0.116,0.564)rectangle++(0.009,0.009) (0.053,0.789)rectangle++(0.009,0.009) (0.023,0.376)rectangle++(0.009,0.009)
  (0.529,0.178)rectangle++(0.009,0.009) (0.251,0.603)rectangle++(0.009,0.009) (0.118,0.807)rectangle++(0.009,0.009) (0.054,0.385)rectangle++(0.009,0.009) (0.024,0.703)rectangle++(0.009,0.009)
  (0.009,0.335)rectangle++(0.009,0.009) (0.002,0.158)rectangle++(0.009,0.009) (-0.001,0.074)rectangle++(0.009,0.009) (-0.003,0.553)rectangle++(0.009,0.009) (-0.004,0.263)rectangle++(0.009,0.009)
  (-0.004,0.644)rectangle++(0.009,0.009) (-0.004,0.827)rectangle++(0.009,0.009) (-0.004,0.915)rectangle++(0.009,0.009) (0.516,0.437)rectangle++(0.009,0.009) (0.765,0.207)rectangle++(0.009,0.009)
  (0.365,0.097)rectangle++(0.009,0.009) (0.173,0.044)rectangle++(0.009,0.009) (0.081,0.019)rectangle++(0.009,0.009) (0.036,0.527)rectangle++(0.009,0.009) (0.535,0.25)rectangle++(0.009,0.009)
  (0.255,0.118)rectangle++(0.009,0.009) (0.12,0.054)rectangle++(0.009,0.009) (0.055,0.544)rectangle++(0.009,0.009) (0.024,0.779)rectangle++(0.009,0.009) (0.009,0.371)rectangle++(0.009,0.009)
  (0.002,0.696)rectangle++(0.009,0.009) (0.519,0.332)rectangle++(0.009,0.009) (0.767,0.157)rectangle++(0.009,0.009) (0.366,0.593)rectangle++(0.009,0.009) (0.173,0.282)rectangle++(0.009,0.009)
  (0.081,0.133)rectangle++(0.009,0.009) (0.036,0.062)rectangle++(0.009,0.009) (0.015,0.547)rectangle++(0.009,0.009) (0.005,0.78)rectangle++(0.009,0.009) (0,0.372)rectangle++(0.009,0.009)
  (-0.002,0.176)rectangle++(0.009,0.009) (-0.003,0.082)rectangle++(0.009,0.009) (-0.004,0.557)rectangle++(0.009,0.009) (0.516,0.265)rectangle++(0.009,0.009) (0.765,0.125)rectangle++(0.009,0.009)
  (0.885,0.058)rectangle++(0.009,0.009) (0.422,0.545)rectangle++(0.009,0.009) (0.2,0.779)rectangle++(0.009,0.009) (0.094,0.372)rectangle++(0.009,0.009) (0.563,0.176)rectangle++(0.009,0.009)
  (0.268,0.082)rectangle++(0.009,0.009) (0.646,0.037)rectangle++(0.009,0.009) (0.308,0.535)rectangle++(0.009,0.009) (0.665,0.255)rectangle++(0.009,0.009) (0.317,0.64)rectangle++(0.009,0.009)
  (0.15,0.305)rectangle++(0.009,0.009) (0.59,0.144)rectangle++(0.009,0.009) (0.281,0.067)rectangle++(0.009,0.009) (0.132,0.55)rectangle++(0.009,0.009) (0.061,0.782)rectangle++(0.009,0.009)
  (0.027,0.373)rectangle++(0.009,0.009) (0.011,0.177)rectangle++(0.009,0.009) (0.003,0.602)rectangle++(0.009,0.009) (-0.001,0.807)rectangle++(0.009,0.009) (-0.003,0.905)rectangle++(0.009,0.009)
  (-0.004,0.432)rectangle++(0.009,0.009) (-0.004,0.205)rectangle++(0.009,0.009) (-0.004,0.616)rectangle++(0.009,0.009) (0.516,0.293)rectangle++(0.009,0.009) (0.245,0.658)rectangle++(0.009,0.009)
  (0.635,0.314)rectangle++(0.009,0.009) (0.303,0.148)rectangle++(0.009,0.009) (0.663,0.069)rectangle++(0.009,0.009) (0.316,0.551)rectangle++(0.009,0.009) (0.149,0.782)rectangle++(0.009,0.009)
  (0.589,0.373)rectangle++(0.009,0.009) (0.801,0.177)rectangle++(0.009,0.009) (0.902,0.082)rectangle++(0.009,0.009) (0.431,0.557)rectangle++(0.009,0.009) (0.204,0.265)rectangle++(0.009,0.009)
  (0.096,0.645)rectangle++(0.009,0.009) (0.044,0.827)rectangle++(0.009,0.009) (0.019,0.915)rectangle++(0.009,0.009) (0.007,0.957)rectangle++(0.009,0.009) (0.521,0.457)rectangle++(0.009,0.009)
  (0.768,0.217)rectangle++(0.009,0.009) (0.366,0.102)rectangle++(0.009,0.009) (0.173,0.047)rectangle++(0.009,0.009) (0.081,0.02)rectangle++(0.009,0.009) (0.036,0.527)rectangle++(0.009,0.009)
  (0.015,0.771)rectangle++(0.009,0.009) (0.005,0.888)rectangle++(0.009,0.009) (0,0.944)rectangle++(0.009,0.009) (0.518,0.451)rectangle++(0.009,0.009) (0.766,0.214)rectangle++(0.009,0.009)
  (0.885,0.1)rectangle++(0.009,0.009) (0.943,0.046)rectangle++(0.009,0.009) (0.45,0.54)rectangle++(0.009,0.009) (0.214,0.777)rectangle++(0.009,0.009) (0.1,0.89)rectangle++(0.009,0.009)
  (0.046,0.425)rectangle++(0.009,0.009) (0.02,0.722)rectangle++(0.009,0.009) (0.527,0.344)rectangle++(0.009,0.009) (0.251,0.163)rectangle++(0.009,0.009) (0.638,0.076)rectangle++(0.009,0.009)
  (0.304,0.554)rectangle++(0.009,0.009) (0.144,0.784)rectangle++(0.009,0.009) (0.587,0.374)rectangle++(0.009,0.009) (0.799,0.177)rectangle++(0.009,0.009) (0.381,0.083)rectangle++(0.009,0.009)
  (0.701,0.037)rectangle++(0.009,0.009) (0.854,0.016)rectangle++(0.009,0.009) (0.928,0.005)rectangle++(0.009,0.009) (0.963,0)rectangle++(0.009,0.009) (0.46,0.518)rectangle++(0.009,0.009)
  (0.218,0.246)rectangle++(0.009,0.009) (0.102,0.116)rectangle++(0.009,0.009) (0.047,0.053)rectangle++(0.009,0.009) (0.54,0.023)rectangle++(0.009,0.009) (0.257,0.529)rectangle++(0.009,0.009)
  (0.121,0.771)rectangle++(0.009,0.009) (0.056,0.368)rectangle++(0.009,0.009) (0.024,0.174)rectangle++(0.009,0.009) (0.009,0.081)rectangle++(0.009,0.009) (0.002,0.037)rectangle++(0.009,0.009)
  (-0.001,0.015)rectangle++(0.009,0.009) (-0.003,0.005)rectangle++(0.009,0.009) (0.516,0)rectangle++(0.009,0.009) (0.245,0.518)rectangle++(0.009,0.009) (0.115,0.766)rectangle++(0.009,0.009)
  (0.573,0.365)rectangle++(0.009,0.009) (0.273,0.693)rectangle++(0.009,0.009) (0.129,0.85)rectangle++(0.009,0.009) (0.059,0.406)rectangle++(0.009,0.009) (0.026,0.712)rectangle++(0.009,0.009)
  (0.53,0.34)rectangle++(0.009,0.009) (0.252,0.681)rectangle++(0.009,0.009) (0.119,0.324)rectangle++(0.009,0.009) (0.575,0.153)rectangle++(0.009,0.009) (0.793,0.071)rectangle++(0.009,0.009)
  (0.379,0.032)rectangle++(0.009,0.009) (0.179,0.013)rectangle++(0.009,0.009) (0.604,0.004)rectangle++(0.009,0.009) (0.807,0)rectangle++(0.009,0.009) (0.385,-0.003)rectangle++(0.009,0.009)
  (0.183,-0.004)rectangle++(0.009,0.009) (0.085,0.516)rectangle++(0.009,0.009) (0.559,0.245)rectangle++(0.009,0.009) (0.266,0.115)rectangle++(0.009,0.009) (0.125,0.573)rectangle++(0.009,0.009)
  (0.058,0.793)rectangle++(0.009,0.009) (0.025,0.378)rectangle++(0.009,0.009) (0.53,0.179)rectangle++(0.009,0.009) (0.252,0.084)rectangle++(0.009,0.009) (0.639,0.038)rectangle++(0.009,0.009)
  (0.824,0.016)rectangle++(0.009,0.009) (0.913,0.005)rectangle++(0.009,0.009) (0.436,0)rectangle++(0.009,0.009) (0.727,-0.002)rectangle++(0.009,0.009) (0.867,-0.003)rectangle++(0.009,0.009)
  (0.934,-0.004)rectangle++(0.009,0.009) (0.446,0.516)rectangle++(0.009,0.009) (0.732,0.245)rectangle++(0.009,0.009) (0.869,0.115)rectangle++(0.009,0.009) (0.415,0.053)rectangle++(0.009,0.009)
  (0.717,0.023)rectangle++(0.009,0.009) (0.342,0.009)rectangle++(0.009,0.009) (0.162,0.522)rectangle++(0.009,0.009) (0.595,0.248)rectangle++(0.009,0.009) (0.283,0.117)rectangle++(0.009,0.009)
  (0.134,0.574)rectangle++(0.009,0.009) (0.062,0.793)rectangle++(0.009,0.009) (0.547,0.378)rectangle++(0.009,0.009) (0.26,0.699)rectangle++(0.009,0.009) (0.123,0.853)rectangle++(0.009,0.009)
  (0.057,0.927)rectangle++(0.009,0.009) (0.025,0.963)rectangle++(0.009,0.009) (0.01,0.98)rectangle++(0.009,0.009) (0.522,0.468)rectangle++(0.009,0.009) (0.248,0.222)rectangle++(0.009,0.009)
  (0.117,0.104)rectangle++(0.009,0.009) (0.054,0.048)rectangle++(0.009,0.009) (0.543,0.021)rectangle++(0.009,0.009) (0.259,0.008)rectangle++(0.009,0.009) (0.642,0.001)rectangle++(0.009,0.009)
  (0.826,-0.002)rectangle++(0.009,0.009) (0.394,0.517)rectangle++(0.009,0.009) (0.187,0.246)rectangle++(0.009,0.009) (0.607,0.116)rectangle++(0.009,0.009) (0.809,0.053)rectangle++(0.009,0.009)
  (0.386,0.543)rectangle++(0.009,0.009) (0.183,0.778)rectangle++(0.009,0.009) (0.605,0.371)rectangle++(0.009,0.009) (0.288,0.176)rectangle++(0.009,0.009) (0.656,0.082)rectangle++(0.009,0.009)
  (0.833,0.037)rectangle++(0.009,0.009) (0.397,0.015)rectangle++(0.009,0.009) (0.708,0.005)rectangle++(0.009,0.009) (0.858,0)rectangle++(0.009,0.009) (0.409,-0.002)rectangle++(0.009,0.009)
  (0.194,0.517)rectangle++(0.009,0.009) (0.091,0.246)rectangle++(0.009,0.009) (0.561,0.116)rectangle++(0.009,0.009) (0.787,0.053)rectangle++(0.009,0.009) (0.895,0.023)rectangle++(0.009,0.009)
  (0.947,0.009)rectangle++(0.009,0.009) (0.972,0.002)rectangle++(0.009,0.009) (0.464,-0.001)rectangle++(0.009,0.009) (0.741,-0.003)rectangle++(0.009,0.009) (0.353,-0.004)rectangle++(0.009,0.009)
  (0.167,-0.004)rectangle++(0.009,0.009) (0.078,-0.004)rectangle++(0.009,0.009) (0.555,-0.004)rectangle++(0.009,0.009) (0.264,-0.004)rectangle++(0.009,0.009) (0.644,-0.004)rectangle++(0.009,0.009)
  (0.827,-0.004)rectangle++(0.009,0.009) (0.915,-0.004)rectangle++(0.009,0.009) (0.957,-0.004)rectangle++(0.009,0.009) (0.977,-0.004)rectangle++(0.009,0.009) (0.467,-0.004)rectangle++(0.009,0.009)
  (0.222,-0.004)rectangle++(0.009,0.009) (0.624,-0.004)rectangle++(0.009,0.009) (0.297,-0.004)rectangle++(0.009,0.009) (0.14,0.516)rectangle++(0.009,0.009) (0.585,0.245)rectangle++(0.009,0.009)
  (0.278,0.115)rectangle++(0.009,0.009) (0.131,0.053)rectangle++(0.009,0.009) (0.061,0.543)rectangle++(0.009,0.009) (0.027,0.778)rectangle++(0.009,0.009) (0.531,0.371)rectangle++(0.009,0.009)
  (0.252,0.696)rectangle++(0.009,0.009) (0.119,0.332)rectangle++(0.009,0.009) (0.055,0.157)rectangle++(0.009,0.009) (0.024,0.593)rectangle++(0.009,0.009) (0.009,0.802)rectangle++(0.009,0.009)
  (0.522,0.383)rectangle++(0.009,0.009) (0.768,0.181)rectangle++(0.009,0.009) (0.366,0.605)rectangle++(0.009,0.009) (0.174,0.288)rectangle++(0.009,0.009) (0.601,0.136)rectangle++(0.009,0.009)
  (0.286,0.063)rectangle++(0.009,0.009) (0.655,0.028)rectangle++(0.009,0.009) (0.312,0.011)rectangle++(0.009,0.009) (0.147,0.523)rectangle++(0.009,0.009) (0.068,0.249)rectangle++(0.009,0.009)
  (0.031,0.637)rectangle++(0.009,0.009) (0.012,0.303)rectangle++(0.009,0.009) (0.524,0.143)rectangle++(0.009,0.009) (0.249,0.586)rectangle++(0.009,0.009) (0.637,0.279)rectangle++(0.009,0.009)
  (0.303,0.652)rectangle++(0.009,0.009) (0.663,0.31)rectangle++(0.009,0.009) (0.316,0.147)rectangle++(0.009,0.009) (0.669,0.068)rectangle++(0.009,0.009) (0.839,0.03)rectangle++(0.009,0.009);
\fill[nuage]
  (0.4,0.532)rectangle++(0.009,0.009) (0.71,0.253)rectangle++(0.009,0.009) (0.338,0.639)rectangle++(0.009,0.009) (0.16,0.304)rectangle++(0.009,0.009) (0.074,0.144)rectangle++(0.009,0.009)
  (0.033,0.587)rectangle++(0.009,0.009) (0.014,0.799)rectangle++(0.009,0.009) (0.524,0.381)rectangle++(0.009,0.009) (0.249,0.701)rectangle++(0.009,0.009) (0.637,0.334)rectangle++(0.009,0.009)
  (0.304,0.158)rectangle++(0.009,0.009) (0.143,0.593)rectangle++(0.009,0.009) (0.066,0.283)rectangle++(0.009,0.009) (0.03,0.653)rectangle++(0.009,0.009) (0.012,0.831)rectangle++(0.009,0.009)
  (0.523,0.397)rectangle++(0.009,0.009) (0.249,0.188)rectangle++(0.009,0.009) (0.117,0.608)rectangle++(0.009,0.009) (0.054,0.809)rectangle++(0.009,0.009) (0.024,0.386)rectangle++(0.009,0.009)
  (0.529,0.183)rectangle++(0.009,0.009) (0.252,0.086)rectangle++(0.009,0.009) (0.118,0.039)rectangle++(0.009,0.009) (0.054,0.536)rectangle++(0.009,0.009) (0.024,0.775)rectangle++(0.009,0.009)
  (0.529,0.37)rectangle++(0.009,0.009) (0.252,0.175)rectangle++(0.009,0.009) (0.118,0.602)rectangle++(0.009,0.009) (0.575,0.286)rectangle++(0.009,0.009) (0.273,0.655)rectangle++(0.009,0.009)
  (0.649,0.312)rectangle++(0.009,0.009) (0.309,0.147)rectangle++(0.009,0.009) (0.146,0.588)rectangle++(0.009,0.009) (0.068,0.8)rectangle++(0.009,0.009) (0.03,0.902)rectangle++(0.009,0.009)
  (0.532,0.43)rectangle++(0.009,0.009) (0.253,0.724)rectangle++(0.009,0.009) (0.639,0.345)rectangle++(0.009,0.009) (0.304,0.683)rectangle++(0.009,0.009) (0.664,0.326)rectangle++(0.009,0.009)
  (0.316,0.154)rectangle++(0.009,0.009) (0.149,0.072)rectangle++(0.009,0.009) (0.589,0.032)rectangle++(0.009,0.009) (0.281,0.533)rectangle++(0.009,0.009) (0.132,0.254)rectangle++(0.009,0.009)
  (0.581,0.119)rectangle++(0.009,0.009) (0.277,0.055)rectangle++(0.009,0.009) (0.65,0.024)rectangle++(0.009,0.009) (0.83,0.009)rectangle++(0.009,0.009) (0.396,0.522)rectangle++(0.009,0.009)
  (0.188,0.248)rectangle++(0.009,0.009) (0.088,0.637)rectangle++(0.009,0.009) (0.56,0.303)rectangle++(0.009,0.009) (0.786,0.143)rectangle++(0.009,0.009) (0.375,0.066)rectangle++(0.009,0.009)
  (0.178,0.55)rectangle++(0.009,0.009) (0.083,0.261)rectangle++(0.009,0.009) (0.037,0.123)rectangle++(0.009,0.009) (0.016,0.577)rectangle++(0.009,0.009) (0.005,0.795)rectangle++(0.009,0.009)
  (0.52,0.379)rectangle++(0.009,0.009) (0.247,0.7)rectangle++(0.009,0.009) (0.116,0.853)rectangle++(0.009,0.009) (0.054,0.407)rectangle++(0.009,0.009) (0.023,0.193)rectangle++(0.009,0.009)
  (0.529,0.09)rectangle++(0.009,0.009) (0.772,0.041)rectangle++(0.009,0.009) (0.368,0.017)rectangle++(0.009,0.009) (0.694,0.006)rectangle++(0.009,0.009) (0.851,0.001)rectangle++(0.009,0.009)
  (0.406,0.518)rectangle++(0.009,0.009) (0.713,0.246)rectangle++(0.009,0.009) (0.34,0.636)rectangle++(0.009,0.009) (0.161,0.303)rectangle++(0.009,0.009) (0.075,0.143)rectangle++(0.009,0.009)
  (0.554,0.066)rectangle++(0.009,0.009) (0.263,0.549)rectangle++(0.009,0.009) (0.124,0.781)rectangle++(0.009,0.009) (0.057,0.373)rectangle++(0.009,0.009) (0.025,0.177)rectangle++(0.009,0.009)
  (0.01,0.082)rectangle++(0.009,0.009) (0.002,0.557)rectangle++(0.009,0.009) (0.519,0.265)rectangle++(0.009,0.009) (0.767,0.125)rectangle++(0.009,0.009) (0.366,0.578)rectangle++(0.009,0.009)
  (0.173,0.795)rectangle++(0.009,0.009) (0.601,0.379)rectangle++(0.009,0.009) (0.286,0.18)rectangle++(0.009,0.009) (0.135,0.084)rectangle++(0.009,0.009) (0.062,0.038)rectangle++(0.009,0.009)
  (0.028,0.016)rectangle++(0.009,0.009) (0.531,0.005)rectangle++(0.009,0.009) (0.253,0.52)rectangle++(0.009,0.009) (0.119,0.767)rectangle++(0.009,0.009) (0.055,0.366)rectangle++(0.009,0.009)
  (0.544,0.173)rectangle++(0.009,0.009) (0.779,0.081)rectangle++(0.009,0.009) (0.891,0.036)rectangle++(0.009,0.009) (0.426,0.015)rectangle++(0.009,0.009) (0.202,0.525)rectangle++(0.009,0.009)
  (0.095,0.77)rectangle++(0.009,0.009) (0.043,0.887)rectangle++(0.009,0.009) (0.018,0.943)rectangle++(0.009,0.009) (0.526,0.451)rectangle++(0.009,0.009) (0.25,0.214)rectangle++(0.009,0.009)
  (0.118,0.1)rectangle++(0.009,0.009) (0.574,0.046)rectangle++(0.009,0.009) (0.273,0.54)rectangle++(0.009,0.009) (0.649,0.257)rectangle++(0.009,0.009) (0.829,0.121)rectangle++(0.009,0.009)
  (0.916,0.056)rectangle++(0.009,0.009) (0.437,0.544)rectangle++(0.009,0.009) (0.727,0.259)rectangle++(0.009,0.009) (0.347,0.642)rectangle++(0.009,0.009) (0.164,0.306)rectangle++(0.009,0.009)
  (0.076,0.664)rectangle++(0.009,0.009) (0.554,0.317)rectangle++(0.009,0.009) (0.264,0.67)rectangle++(0.009,0.009) (0.124,0.839)rectangle++(0.009,0.009) (0.057,0.92)rectangle++(0.009,0.009)
  (0.025,0.959)rectangle++(0.009,0.009) (0.01,0.458)rectangle++(0.009,0.009) (0.002,0.218)rectangle++(0.009,0.009) (0.519,0.102)rectangle++(0.009,0.009) (0.247,0.567)rectangle++(0.009,0.009)
  (0.116,0.27)rectangle++(0.009,0.009) (0.053,0.127)rectangle++(0.009,0.009) (0.543,0.059)rectangle++(0.009,0.009) (0.778,0.026)rectangle++(0.009,0.009) (0.371,0.01)rectangle++(0.009,0.009)
  (0.176,0.002)rectangle++(0.009,0.009) (0.082,0.519)rectangle++(0.009,0.009) (0.037,0.247)rectangle++(0.009,0.009) (0.535,0.116)rectangle++(0.009,0.009) (0.775,0.053)rectangle++(0.009,0.009)
  (0.37,0.543)rectangle++(0.009,0.009) (0.175,0.778)rectangle++(0.009,0.009) (0.602,0.371)rectangle++(0.009,0.009) (0.286,0.696)rectangle++(0.009,0.009) (0.135,0.332)rectangle++(0.009,0.009)
  (0.063,0.677)rectangle++(0.009,0.009) (0.028,0.323)rectangle++(0.009,0.009) (0.531,0.152)rectangle++(0.009,0.009) (0.253,0.071)rectangle++(0.009,0.009) (0.639,0.032)rectangle++(0.009,0.009)
  (0.304,0.533)rectangle++(0.009,0.009) (0.144,0.253)rectangle++(0.009,0.009) (0.587,0.119)rectangle++(0.009,0.009) (0.279,0.055)rectangle++(0.009,0.009) (0.132,0.544)rectangle++(0.009,0.009)
  (0.581,0.259)rectangle++(0.009,0.009) (0.276,0.642)rectangle++(0.009,0.009) (0.65,0.306)rectangle++(0.009,0.009) (0.83,0.144)rectangle++(0.009,0.009) (0.396,0.587)rectangle++(0.009,0.009)
  (0.188,0.279)rectangle++(0.009,0.009) (0.088,0.652)rectangle++(0.009,0.009) (0.04,0.311)rectangle++(0.009,0.009) (0.537,0.147)rectangle++(0.009,0.009) (0.775,0.068)rectangle++(0.009,0.009)
  (0.89,0.03)rectangle++(0.009,0.009) (0.945,0.012)rectangle++(0.009,0.009) (0.451,0.524)rectangle++(0.009,0.009) (0.214,0.769)rectangle++(0.009,0.009) (0.62,0.367)rectangle++(0.009,0.009)
  (0.295,0.174)rectangle++(0.009,0.009) (0.139,0.601)rectangle++(0.009,0.009) (0.585,0.286)rectangle++(0.009,0.009) (0.278,0.655)rectangle++(0.009,0.009) (0.131,0.312)rectangle++(0.009,0.009)
  (0.061,0.147)rectangle++(0.009,0.009) (0.027,0.588)rectangle++(0.009,0.009) (0.011,0.28)rectangle++(0.009,0.009) (0.523,0.132)rectangle++(0.009,0.009) (0.249,0.061)rectangle++(0.009,0.009)
  (0.637,0.027)rectangle++(0.009,0.009) (0.823,0.011)rectangle++(0.009,0.009) (0.913,0.003)rectangle++(0.009,0.009) (0.436,0.519)rectangle++(0.009,0.009) (0.727,0.247)rectangle++(0.009,0.009)
  (0.867,0.116)rectangle++(0.009,0.009) (0.934,0.053)rectangle++(0.009,0.009) (0.446,0.023)rectangle++(0.009,0.009) (0.212,0.009)rectangle++(0.009,0.009) (0.099,0.002)rectangle++(0.009,0.009)
  (0.565,-0.001)rectangle++(0.009,0.009) (0.269,0.517)rectangle++(0.009,0.009) (0.647,0.246)rectangle++(0.009,0.009) (0.828,0.116)rectangle++(0.009,0.009) (0.395,0.573)rectangle++(0.009,0.009)
  (0.187,0.273)rectangle++(0.009,0.009) (0.608,0.129)rectangle++(0.009,0.009) (0.809,0.059)rectangle++(0.009,0.009) (0.386,0.026)rectangle++(0.009,0.009) (0.183,0.01)rectangle++(0.009,0.009)
  (0.605,0.003)rectangle++(0.009,0.009) (0.808,-0.001)rectangle++(0.009,0.009) (0.386,-0.003)rectangle++(0.009,0.009) (0.183,-0.004)rectangle++(0.009,0.009) (0.085,0.516)rectangle++(0.009,0.009)
  (0.039,0.765)rectangle++(0.009,0.009) (0.016,0.885)rectangle++(0.009,0.009) (0.005,0.422)rectangle++(0.009,0.009) (0.52,0.2)rectangle++(0.009,0.009) (0.247,0.614)rectangle++(0.009,0.009)
  (0.636,0.292)rectangle++(0.009,0.009) (0.303,0.658)rectangle++(0.009,0.009) (0.143,0.313)rectangle++(0.009,0.009) (0.586,0.148)rectangle++(0.009,0.009) (0.279,0.069)rectangle++(0.009,0.009)
  (0.132,0.031)rectangle++(0.009,0.009) (0.061,0.532)rectangle++(0.009,0.009) (0.027,0.773)rectangle++(0.009,0.009) (0.011,0.369)rectangle++(0.009,0.009) (0.523,0.175)rectangle++(0.009,0.009)
  (0.249,0.082)rectangle++(0.009,0.009) (0.117,0.037)rectangle++(0.009,0.009) (0.054,0.535)rectangle++(0.009,0.009) (0.023,0.255)rectangle++(0.009,0.009) (0.529,0.12)rectangle++(0.009,0.009)
  (0.772,0.055)rectangle++(0.009,0.009) (0.888,0.024)rectangle++(0.009,0.009) (0.424,0.529)rectangle++(0.009,0.009) (0.721,0.252)rectangle++(0.009,0.009) (0.344,0.118)rectangle++(0.009,0.009)
  (0.683,0.055)rectangle++(0.009,0.009) (0.325,0.544)rectangle++(0.009,0.009) (0.154,0.779)rectangle++(0.009,0.009) (0.591,0.371)rectangle++(0.009,0.009) (0.282,0.696)rectangle++(0.009,0.009)
  (0.133,0.332)rectangle++(0.009,0.009) (0.061,0.157)rectangle++(0.009,0.009) (0.027,0.593)rectangle++(0.009,0.009) (0.011,0.282)rectangle++(0.009,0.009) (0.003,0.653)rectangle++(0.009,0.009)
  (-0.001,0.311)rectangle++(0.009,0.009) (0.517,0.147)rectangle++(0.009,0.009) (0.766,0.068)rectangle++(0.009,0.009) (0.365,0.55)rectangle++(0.009,0.009) (0.693,0.262)rectangle++(0.009,0.009)
  (0.33,0.643)rectangle++(0.009,0.009) (0.676,0.306)rectangle++(0.009,0.009) (0.322,0.665)rectangle++(0.009,0.009) (0.152,0.317)rectangle++(0.009,0.009) (0.071,0.67)rectangle++(0.009,0.009)
  (0.552,0.319)rectangle++(0.009,0.009) (0.782,0.151)rectangle++(0.009,0.009) (0.373,0.59)rectangle++(0.009,0.009) (0.177,0.281)rectangle++(0.009,0.009) (0.083,0.652)rectangle++(0.009,0.009)
  (0.037,0.831)rectangle++(0.009,0.009) (0.016,0.916)rectangle++(0.009,0.009) (0.525,0.438)rectangle++(0.009,0.009) (0.25,0.728)rectangle++(0.009,0.009) (0.118,0.347)rectangle++(0.009,0.009)
  (0.054,0.164)rectangle++(0.009,0.009) (0.024,0.076)rectangle++(0.009,0.009) (0.529,0.034)rectangle++(0.009,0.009) (0.252,0.534)rectangle++(0.009,0.009) (0.638,0.254)rectangle++(0.009,0.009)
  (0.304,0.64)rectangle++(0.009,0.009) (0.664,0.305)rectangle++(0.009,0.009) (0.316,0.144)rectangle++(0.009,0.009) (0.669,0.067)rectangle++(0.009,0.009) (0.319,0.55)rectangle++(0.009,0.009)
  (0.151,0.782)rectangle++(0.009,0.009) (0.07,0.893)rectangle++(0.009,0.009) (0.031,0.946)rectangle++(0.009,0.009) (0.533,0.452)rectangle++(0.009,0.009) (0.773,0.215)rectangle++(0.009,0.009)
  (0.889,0.101)rectangle++(0.009,0.009) (0.424,0.566)rectangle++(0.009,0.009) (0.201,0.789)rectangle++(0.009,0.009) (0.094,0.377)rectangle++(0.009,0.009) (0.563,0.178)rectangle++(0.009,0.009)
  (0.268,0.083)rectangle++(0.009,0.009) (0.126,0.038)rectangle++(0.009,0.009) (0.578,0.016)rectangle++(0.009,0.009) (0.275,0.525)rectangle++(0.009,0.009) (0.65,0.25)rectangle++(0.009,0.009)
  (0.31,0.638)rectangle++(0.009,0.009) (0.666,0.304)rectangle++(0.009,0.009) (0.837,0.143)rectangle++(0.009,0.009) (0.92,0.067)rectangle++(0.009,0.009) (0.439,0.55)rectangle++(0.009,0.009)
  (0.208,0.781)rectangle++(0.009,0.009) (0.098,0.893)rectangle++(0.009,0.009) (0.045,0.426)rectangle++(0.009,0.009) (0.019,0.202)rectangle++(0.009,0.009) (0.007,0.615)rectangle++(0.009,0.009)
  (0.001,0.293)rectangle++(0.009,0.009) (0.518,0.138)rectangle++(0.009,0.009) (0.246,0.584)rectangle++(0.009,0.009) (0.116,0.278)rectangle++(0.009,0.009) (0.053,0.131)rectangle++(0.009,0.009)
  (0.023,0.581)rectangle++(0.009,0.009) (0.529,0.276)rectangle++(0.009,0.009) (0.251,0.65)rectangle++(0.009,0.009) (0.118,0.31)rectangle++(0.009,0.009) (0.054,0.146)rectangle++(0.009,0.009)
  (0.544,0.068)rectangle++(0.009,0.009) (0.259,0.03)rectangle++(0.009,0.009) (0.122,0.012)rectangle++(0.009,0.009) (0.056,0.524)rectangle++(0.009,0.009) (0.545,0.249)rectangle++(0.009,0.009)
  (0.259,0.117)rectangle++(0.009,0.009) (0.122,0.054)rectangle++(0.009,0.009) (0.576,0.024)rectangle++(0.009,0.009) (0.274,0.009)rectangle++(0.009,0.009) (0.129,0.522)rectangle++(0.009,0.009)
  (0.58,0.248)rectangle++(0.009,0.009) (0.796,0.117)rectangle++(0.009,0.009) (0.38,0.574)rectangle++(0.009,0.009) (0.18,0.793)rectangle++(0.009,0.009) (0.084,0.898)rectangle++(0.009,0.009)
  (0.558,0.429)rectangle++(0.009,0.009) (0.265,0.724)rectangle++(0.009,0.009) (0.125,0.345)rectangle++(0.009,0.009) (0.578,0.163)rectangle++(0.009,0.009) (0.275,0.076)rectangle++(0.009,0.009)
  (0.13,0.554)rectangle++(0.009,0.009) (0.58,0.264)rectangle++(0.009,0.009) (0.276,0.124)rectangle++(0.009,0.009) (0.13,0.577)rectangle++(0.009,0.009) (0.06,0.795)rectangle++(0.009,0.009)
  (0.027,0.899)rectangle++(0.009,0.009) (0.53,0.429)rectangle++(0.009,0.009) (0.252,0.724)rectangle++(0.009,0.009) (0.119,0.345)rectangle++(0.009,0.009) (0.055,0.683)rectangle++(0.009,0.009)
  (0.024,0.326)rectangle++(0.009,0.009) (0.009,0.674)rectangle++(0.009,0.009) (0.002,0.321)rectangle++(0.009,0.009) (-0.001,0.672)rectangle++(0.009,0.009) (-0.003,0.84)rectangle++(0.009,0.009)
  (-0.004,0.401)rectangle++(0.009,0.009) (0.516,0.19)rectangle++(0.009,0.009) (0.245,0.089)rectangle++(0.009,0.009) (0.635,0.04)rectangle++(0.009,0.009) (0.303,0.537)rectangle++(0.009,0.009)
  (0.143,0.255)rectangle++(0.009,0.009) (0.586,0.12)rectangle++(0.009,0.009) (0.799,0.055)rectangle++(0.009,0.009) (0.901,0.024)rectangle++(0.009,0.009) (0.43,0.529)rectangle++(0.009,0.009)
  (0.724,0.252)rectangle++(0.009,0.009) (0.345,0.118)rectangle++(0.009,0.009) (0.163,0.575)rectangle++(0.009,0.009) (0.596,0.273)rectangle++(0.009,0.009) (0.284,0.129)rectangle++(0.009,0.009)
  (0.134,0.06)rectangle++(0.009,0.009) (0.062,0.026)rectangle++(0.009,0.009) (0.027,0.01)rectangle++(0.009,0.009) (0.011,0.003)rectangle++(0.009,0.009) (0.523,-0.001)rectangle++(0.009,0.009)
  (0.249,-0.003)rectangle++(0.009,0.009) (0.117,0.516)rectangle++(0.009,0.009) (0.054,0.765)rectangle++(0.009,0.009) (0.023,0.365)rectangle++(0.009,0.009) (0.009,0.173)rectangle++(0.009,0.009)
  (0.002,0.601)rectangle++(0.009,0.009) (-0.001,0.286)rectangle++(0.009,0.009) (-0.003,0.135)rectangle++(0.009,0.009) (-0.004,0.062)rectangle++(0.009,0.009) (-0.004,0.548)rectangle++(0.009,0.009)
  (0.516,0.261)rectangle++(0.009,0.009) (0.765,0.123)rectangle++(0.009,0.009) (0.885,0.057)rectangle++(0.009,0.009) (0.942,0.025)rectangle++(0.009,0.009) (0.45,0.53)rectangle++(0.009,0.009)
  (0.214,0.252)rectangle++(0.009,0.009) (0.1,0.119)rectangle++(0.009,0.009) (0.046,0.055)rectangle++(0.009,0.009) (0.02,0.544)rectangle++(0.009,0.009) (0.527,0.259)rectangle++(0.009,0.009)
  (0.251,0.122)rectangle++(0.009,0.009) (0.638,0.056)rectangle++(0.009,0.009) (0.824,0.025)rectangle++(0.009,0.009) (0.393,0.009)rectangle++(0.009,0.009) (0.186,0.002)rectangle++(0.009,0.009)
  (0.607,-0.001)rectangle++(0.009,0.009) (0.809,-0.003)rectangle++(0.009,0.009) (0.906,-0.004)rectangle++(0.009,0.009) (0.433,0.516)rectangle++(0.009,0.009) (0.205,0.765)rectangle++(0.009,0.009)
  (0.096,0.365)rectangle++(0.009,0.009) (0.044,0.173)rectangle++(0.009,0.009) (0.539,0.081)rectangle++(0.009,0.009) (0.256,0.556)rectangle++(0.009,0.009) (0.641,0.265)rectangle++(0.009,0.009)
  (0.305,0.125)rectangle++(0.009,0.009) (0.664,0.058)rectangle++(0.009,0.009) (0.836,0.025)rectangle++(0.009,0.009) (0.399,0.01)rectangle++(0.009,0.009) (0.709,0.002)rectangle++(0.009,0.009)
  (0.338,-0.001)rectangle++(0.009,0.009) (0.68,-0.003)rectangle++(0.009,0.009) (0.324,0.516)rectangle++(0.009,0.009) (0.673,0.245)rectangle++(0.009,0.009) (0.841,0.115)rectangle++(0.009,0.009)
  (0.401,0.573)rectangle++(0.009,0.009) (0.19,0.793)rectangle++(0.009,0.009) (0.089,0.378)rectangle++(0.009,0.009) (0.56,0.179)rectangle++(0.009,0.009) (0.787,0.084)rectangle++(0.009,0.009)
  (0.375,0.038)rectangle++(0.009,0.009) (0.698,0.016)rectangle++(0.009,0.009) (0.333,0.005)rectangle++(0.009,0.009) (0.677,0)rectangle++(0.009,0.009) (0.323,-0.002)rectangle++(0.009,0.009)
  (0.153,-0.003)rectangle++(0.009,0.009) (0.071,0.516)rectangle++(0.009,0.009) (0.552,0.245)rectangle++(0.009,0.009) (0.262,0.635)rectangle++(0.009,0.009) (0.644,0.303)rectangle++(0.009,0.009)
  (0.827,0.143)rectangle++(0.009,0.009) (0.394,0.066)rectangle++(0.009,0.009) (0.187,0.549)rectangle++(0.009,0.009) (0.087,0.781)rectangle++(0.009,0.009) (0.04,0.373)rectangle++(0.009,0.009)
  (0.017,0.697)rectangle++(0.009,0.009) (0.526,0.332)rectangle++(0.009,0.009) (0.25,0.157)rectangle++(0.009,0.009) (0.118,0.073)rectangle++(0.009,0.009) (0.054,0.553)rectangle++(0.009,0.009)
  (0.544,0.263)rectangle++(0.009,0.009) (0.259,0.124)rectangle++(0.009,0.009) (0.122,0.577)rectangle++(0.009,0.009) (0.056,0.795)rectangle++(0.009,0.009) (0.025,0.379)rectangle++(0.009,0.009)
  (0.529,0.18)rectangle++(0.009,0.009) (0.772,0.084)rectangle++(0.009,0.009) (0.368,0.558)rectangle++(0.009,0.009) (0.174,0.785)rectangle++(0.009,0.009) (0.081,0.895)rectangle++(0.009,0.009)
  (0.037,0.427)rectangle++(0.009,0.009) (0.015,0.723)rectangle++(0.009,0.009) (0.525,0.345)rectangle++(0.009,0.009) (0.77,0.163)rectangle++(0.009,0.009) (0.887,0.076)rectangle++(0.009,0.009)
  (0.943,0.034)rectangle++(0.009,0.009) (0.971,0.014)rectangle++(0.009,0.009) (0.464,0.004)rectangle++(0.009,0.009) (0.22,0.52)rectangle++(0.009,0.009) (0.623,0.247)rectangle++(0.009,0.009)
  (0.817,0.116)rectangle++(0.009,0.009) (0.91,0.053)rectangle++(0.009,0.009) (0.434,0.023)rectangle++(0.009,0.009) (0.206,0.529)rectangle++(0.009,0.009) (0.617,0.252)rectangle++(0.009,0.009)
  (0.814,0.118)rectangle++(0.009,0.009) (0.908,0.054)rectangle++(0.009,0.009) (0.434,0.024)rectangle++(0.009,0.009) (0.206,0.009)rectangle++(0.009,0.009) (0.096,0.522)rectangle++(0.009,0.009)
  (0.564,0.248)rectangle++(0.009,0.009) (0.268,0.637)rectangle++(0.009,0.009) (0.126,0.823)rectangle++(0.009,0.009) (0.058,0.913)rectangle++(0.009,0.009) (0.026,0.956)rectangle++(0.009,0.009)
  (0.01,0.456)rectangle++(0.009,0.009) (0.002,0.217)rectangle++(0.009,0.009) (-0.001,0.622)rectangle++(0.009,0.009) (-0.003,0.816)rectangle++(0.009,0.009) (-0.004,0.389)rectangle++(0.009,0.009)
  (-0.004,0.705)rectangle++(0.009,0.009) (-0.004,0.856)rectangle++(0.009,0.009) (-0.004,0.928)rectangle++(0.009,0.009) (0.516,0.443)rectangle++(0.009,0.009) (0.245,0.21)rectangle++(0.009,0.009)
  (0.115,0.619)rectangle++(0.009,0.009) (0.053,0.295)rectangle++(0.009,0.009) (0.023,0.139)rectangle++(0.009,0.009) (0.009,0.584)rectangle++(0.009,0.009) (0.002,0.278)rectangle++(0.009,0.009)
  (-0.001,0.131)rectangle++(0.009,0.009) (-0.003,0.581)rectangle++(0.009,0.009) (-0.004,0.796)rectangle++(0.009,0.009) (-0.004,0.38)rectangle++(0.009,0.009) (-0.004,0.18)rectangle++(0.009,0.009)
  (0.516,0.084)rectangle++(0.009,0.009) (0.245,0.558)rectangle++(0.009,0.009) (0.115,0.266)rectangle++(0.009,0.009) (0.573,0.125)rectangle++(0.009,0.009) (0.273,0.058)rectangle++(0.009,0.009)
  (0.129,0.545)rectangle++(0.009,0.009) (0.059,0.779)rectangle++(0.009,0.009) (0.026,0.892)rectangle++(0.009,0.009) (0.01,0.426)rectangle++(0.009,0.009) (0.003,0.202)rectangle++(0.009,0.009)
  (-0.001,0.615)rectangle++(0.009,0.009) (-0.003,0.813)rectangle++(0.009,0.009) (-0.004,0.908)rectangle++(0.009,0.009) (0.516,0.433)rectangle++(0.009,0.009) (0.245,0.206)rectangle++(0.009,0.009)
  (0.115,0.616)rectangle++(0.009,0.009) (0.573,0.294)rectangle++(0.009,0.009) (0.273,0.659)rectangle++(0.009,0.009) (0.129,0.314)rectangle++(0.009,0.009) (0.059,0.148)rectangle++(0.009,0.009)
  (0.026,0.069)rectangle++(0.009,0.009) (0.01,0.031)rectangle++(0.009,0.009) (0.523,0.012)rectangle++(0.009,0.009) (0.248,0.524)rectangle++(0.009,0.009) (0.117,0.769)rectangle++(0.009,0.009)
  (0.574,0.367)rectangle++(0.009,0.009) (0.273,0.174)rectangle++(0.009,0.009) (0.129,0.081)rectangle++(0.009,0.009) (0.579,0.037)rectangle++(0.009,0.009) (0.796,0.015)rectangle++(0.009,0.009)
  (0.38,0.005)rectangle++(0.009,0.009) (0.7,0)rectangle++(0.009,0.009) (0.334,-0.002)rectangle++(0.009,0.009) (0.678,-0.003)rectangle++(0.009,0.009) (0.843,-0.004)rectangle++(0.009,0.009)
  (0.402,-0.004)rectangle++(0.009,0.009) (0.191,0.516)rectangle++(0.009,0.009) (0.609,0.245)rectangle++(0.009,0.009) (0.29,0.635)rectangle++(0.009,0.009) (0.657,0.303)rectangle++(0.009,0.009)
  (0.313,0.143)rectangle++(0.009,0.009) (0.668,0.066)rectangle++(0.009,0.009) (0.318,0.549)rectangle++(0.009,0.009) (0.15,0.261)rectangle++(0.009,0.009) (0.07,0.123)rectangle++(0.009,0.009)
  (0.551,0.057)rectangle++(0.009,0.009) (0.782,0.025)rectangle++(0.009,0.009) (0.373,0.01)rectangle++(0.009,0.009) (0.177,0.522)rectangle++(0.009,0.009) (0.083,0.768)rectangle++(0.009,0.009)
  (0.037,0.886)rectangle++(0.009,0.009) (0.016,0.943)rectangle++(0.009,0.009) (0.525,0.45)rectangle++(0.009,0.009) (0.25,0.214)rectangle++(0.009,0.009) (0.118,0.1)rectangle++(0.009,0.009)
  (0.054,0.566)rectangle++(0.009,0.009) (0.024,0.789)rectangle++(0.009,0.009) (0.009,0.896)rectangle++(0.009,0.009) (0.002,0.428)rectangle++(0.009,0.009) (0.519,0.203)rectangle++(0.009,0.009)
  (0.767,0.095)rectangle++(0.009,0.009) (0.366,0.563)rectangle++(0.009,0.009) (0.173,0.268)rectangle++(0.009,0.009) (0.081,0.646)rectangle++(0.009,0.009) (0.036,0.308)rectangle++(0.009,0.009)
  (0.015,0.665)rectangle++(0.009,0.009) (0.005,0.837)rectangle++(0.009,0.009) (0.52,0.399)rectangle++(0.009,0.009) (0.247,0.709)rectangle++(0.009,0.009) (0.116,0.858)rectangle++(0.009,0.009)
  (0.574,0.41)rectangle++(0.009,0.009) (0.273,0.714)rectangle++(0.009,0.009) (0.129,0.341)rectangle++(0.009,0.009) (0.059,0.161)rectangle++(0.009,0.009) (0.026,0.075)rectangle++(0.009,0.009)
  (0.01,0.554)rectangle++(0.009,0.009) (0.003,0.263)rectangle++(0.009,0.009) (0.519,0.124)rectangle++(0.009,0.009) (0.767,0.057)rectangle++(0.009,0.009) (0.366,0.545)rectangle++(0.009,0.009)
  (0.173,0.779)rectangle++(0.009,0.009) (0.081,0.372)rectangle++(0.009,0.009) (0.556,0.176)rectangle++(0.009,0.009) (0.265,0.602)rectangle++(0.009,0.009) (0.125,0.807)rectangle++(0.009,0.009)
  (0.058,0.385)rectangle++(0.009,0.009) (0.025,0.702)rectangle++(0.009,0.009) (0.01,0.335)rectangle++(0.009,0.009) (0.002,0.158)rectangle++(0.009,0.009) (-0.001,0.594)rectangle++(0.009,0.009)
  (-0.003,0.803)rectangle++(0.009,0.009) (-0.004,0.903)rectangle++(0.009,0.009) (-0.004,0.431)rectangle++(0.009,0.009) (0.516,0.205)rectangle++(0.009,0.009) (0.765,0.096)rectangle++(0.009,0.009)
  (0.365,0.564)rectangle++(0.009,0.009) (0.173,0.268)rectangle++(0.009,0.009) (0.081,0.126)rectangle++(0.009,0.009) (0.556,0.058)rectangle++(0.009,0.009) (0.265,0.026)rectangle++(0.009,0.009)
  (0.125,0.53)rectangle++(0.009,0.009) (0.058,0.252)rectangle++(0.009,0.009) (0.025,0.119)rectangle++(0.009,0.009) (0.01,0.055)rectangle++(0.009,0.009) (0.002,0.024)rectangle++(0.009,0.009)
  (-0.001,0.529)rectangle++(0.009,0.009) (-0.003,0.252)rectangle++(0.009,0.009) (-0.004,0.118)rectangle++(0.009,0.009) (0.516,0.055)rectangle++(0.009,0.009) (0.245,0.544)rectangle++(0.009,0.009)
  (0.635,0.259)rectangle++(0.009,0.009) (0.303,0.642)rectangle++(0.009,0.009) (0.143,0.826)rectangle++(0.009,0.009) (0.066,0.394)rectangle++(0.009,0.009) (0.549,0.187)rectangle++(0.009,0.009)
  (0.261,0.607)rectangle++(0.009,0.009) (0.123,0.809)rectangle++(0.009,0.009) (0.057,0.386)rectangle++(0.009,0.009) (0.025,0.183)rectangle++(0.009,0.009) (0.53,0.085)rectangle++(0.009,0.009)
  (0.252,0.559)rectangle++(0.009,0.009) (0.119,0.266)rectangle++(0.009,0.009) (0.055,0.645)rectangle++(0.009,0.009) (0.024,0.307)rectangle++(0.009,0.009) (0.529,0.145)rectangle++(0.009,0.009)
  (0.252,0.067)rectangle++(0.009,0.009) (0.118,0.03)rectangle++(0.009,0.009) (0.055,0.532)rectangle++(0.009,0.009) (0.544,0.253)rectangle++(0.009,0.009) (0.779,0.119)rectangle++(0.009,0.009)
  (0.891,0.055)rectangle++(0.009,0.009) (0.426,0.544)rectangle++(0.009,0.009) (0.202,0.259)rectangle++(0.009,0.009) (0.615,0.122)rectangle++(0.009,0.009) (0.293,0.056)rectangle++(0.009,0.009)
  (0.658,0.025)rectangle++(0.009,0.009) (0.834,0.009)rectangle++(0.009,0.009) (0.398,0.522)rectangle++(0.009,0.009) (0.189,0.248)rectangle++(0.009,0.009) (0.608,0.117)rectangle++(0.009,0.009)
  (0.81,0.054)rectangle++(0.009,0.009) (0.386,0.543)rectangle++(0.009,0.009) (0.703,0.259)rectangle++(0.009,0.009) (0.855,0.122)rectangle++(0.009,0.009) (0.408,0.576)rectangle++(0.009,0.009)
  (0.194,0.794)rectangle++(0.009,0.009) (0.091,0.899)rectangle++(0.009,0.009) (0.041,0.949)rectangle++(0.009,0.009) (0.017,0.453)rectangle++(0.009,0.009) (0.526,0.215)rectangle++(0.009,0.009)
  (0.77,0.101)rectangle++(0.009,0.009) (0.367,0.046)rectangle++(0.009,0.009) (0.694,0.02)rectangle++(0.009,0.009) (0.331,0.007)rectangle++(0.009,0.009) (0.156,0.521)rectangle++(0.009,0.009)
  (0.073,0.248)rectangle++(0.009,0.009) (0.033,0.117)rectangle++(0.009,0.009) (0.533,0.054)rectangle++(0.009,0.009) (0.254,0.023)rectangle++(0.009,0.009) (0.119,0.529)rectangle++(0.009,0.009)
  (0.055,0.772)rectangle++(0.009,0.009) (0.024,0.888)rectangle++(0.009,0.009) (0.529,0.424)rectangle++(0.009,0.009) (0.252,0.721)rectangle++(0.009,0.009) (0.118,0.344)rectangle++(0.009,0.009)
  (0.055,0.163)rectangle++(0.009,0.009) (0.024,0.076)rectangle++(0.009,0.009) (0.009,0.554)rectangle++(0.009,0.009) (0.522,0.264)rectangle++(0.009,0.009) (0.768,0.124)rectangle++(0.009,0.009)
  (0.366,0.057)rectangle++(0.009,0.009) (0.174,0.545)rectangle++(0.009,0.009) (0.081,0.259)rectangle++(0.009,0.009) (0.557,0.122)rectangle++(0.009,0.009) (0.265,0.576)rectangle++(0.009,0.009)
  (0.645,0.274)rectangle++(0.009,0.009) (0.827,0.129)rectangle++(0.009,0.009) (0.915,0.06)rectangle++(0.009,0.009) (0.437,0.546)rectangle++(0.009,0.009) (0.727,0.26)rectangle++(0.009,0.009)
  (0.347,0.642)rectangle++(0.009,0.009) (0.164,0.306)rectangle++(0.009,0.009) (0.076,0.145)rectangle++(0.009,0.009) (0.034,0.067)rectangle++(0.009,0.009) (0.014,0.55)rectangle++(0.009,0.009)
  (0.004,0.782)rectangle++(0.009,0.009) (0,0.373)rectangle++(0.009,0.009) (0.518,0.177)rectangle++(0.009,0.009) (0.246,0.602)rectangle++(0.009,0.009) (0.116,0.807)rectangle++(0.009,0.009)
  (0.573,0.385)rectangle++(0.009,0.009) (0.793,0.182)rectangle++(0.009,0.009) (0.378,0.605)rectangle++(0.009,0.009) (0.699,0.288)rectangle++(0.009,0.009) (0.853,0.136)rectangle++(0.009,0.009)
  (0.407,0.063)rectangle++(0.009,0.009) (0.713,0.028)rectangle++(0.009,0.009) (0.34,0.531)rectangle++(0.009,0.009) (0.161,0.773)rectangle++(0.009,0.009) (0.075,0.368)rectangle++(0.009,0.009)
  (0.554,0.175)rectangle++(0.009,0.009) (0.783,0.081)rectangle++(0.009,0.009) (0.374,0.037)rectangle++(0.009,0.009) (0.177,0.015)rectangle++(0.009,0.009) (0.083,0.005)rectangle++(0.009,0.009)
  (0.037,0)rectangle++(0.009,0.009) (0.016,0.518)rectangle++(0.009,0.009) (0.005,0.246)rectangle++(0.009,0.009) (0,0.116)rectangle++(0.009,0.009) (-0.002,0.573)rectangle++(0.009,0.009)
  (-0.003,0.793)rectangle++(0.009,0.009) (0.516,0.378)rectangle++(0.009,0.009) (0.245,0.179)rectangle++(0.009,0.009) (0.115,0.084)rectangle++(0.009,0.009) (0.053,0.038)rectangle++(0.009,0.009)
  (0.023,0.536)rectangle++(0.009,0.009) (0.009,0.255)rectangle++(0.009,0.009) (0.002,0.12)rectangle++(0.009,0.009) (0.519,0.055)rectangle++(0.009,0.009) (0.247,0.544)rectangle++(0.009,0.009)
  (0.116,0.779)rectangle++(0.009,0.009) (0.053,0.372)rectangle++(0.009,0.009) (0.023,0.176)rectangle++(0.009,0.009) (0.009,0.602)rectangle++(0.009,0.009) (0.522,0.287)rectangle++(0.009,0.009)
  (0.248,0.135)rectangle++(0.009,0.009) (0.117,0.063)rectangle++(0.009,0.009) (0.054,0.028)rectangle++(0.009,0.009) (0.023,0.531)rectangle++(0.009,0.009) (0.529,0.253)rectangle++(0.009,0.009)
  (0.252,0.119)rectangle++(0.009,0.009) (0.638,0.055)rectangle++(0.009,0.009) (0.304,0.024)rectangle++(0.009,0.009) (0.144,0.529)rectangle++(0.009,0.009) (0.067,0.772)rectangle++(0.009,0.009)
  (0.03,0.888)rectangle++(0.009,0.009) (0.012,0.944)rectangle++(0.009,0.009) (0.523,0.451)rectangle++(0.009,0.009) (0.249,0.214)rectangle++(0.009,0.009) (0.117,0.1)rectangle++(0.009,0.009)
  (0.054,0.046)rectangle++(0.009,0.009) (0.024,0.02)rectangle++(0.009,0.009) (0.009,0.527)rectangle++(0.009,0.009) (0.002,0.771)rectangle++(0.009,0.009) (0.519,0.368)rectangle++(0.009,0.009)
  (0.247,0.694)rectangle++(0.009,0.009) (0.116,0.331)rectangle++(0.009,0.009) (0.053,0.676)rectangle++(0.009,0.009) (0.023,0.322)rectangle++(0.009,0.009) (0.529,0.152)rectangle++(0.009,0.009)
  (0.251,0.071)rectangle++(0.009,0.009) (0.638,0.032)rectangle++(0.009,0.009) (0.304,0.533)rectangle++(0.009,0.009) (0.144,0.773)rectangle++(0.009,0.009) (0.587,0.369)rectangle++(0.009,0.009)
  (0.279,0.175)rectangle++(0.009,0.009) (0.132,0.602)rectangle++(0.009,0.009) (0.061,0.806)rectangle++(0.009,0.009) (0.547,0.385)rectangle++(0.009,0.009) (0.26,0.182)rectangle++(0.009,0.009)
  (0.643,0.085)rectangle++(0.009,0.009) (0.306,0.039)rectangle++(0.009,0.009) (0.145,0.536)rectangle++(0.009,0.009) (0.587,0.255)rectangle++(0.009,0.009) (0.799,0.12)rectangle++(0.009,0.009)
  (0.901,0.055)rectangle++(0.009,0.009) (0.43,0.544)rectangle++(0.009,0.009) (0.204,0.779)rectangle++(0.009,0.009) (0.096,0.372)rectangle++(0.009,0.009) (0.564,0.176)rectangle++(0.009,0.009)
  (0.788,0.082)rectangle++(0.009,0.009) (0.376,0.557)rectangle++(0.009,0.009) (0.698,0.265)rectangle++(0.009,0.009) (0.853,0.125)rectangle++(0.009,0.009) (0.407,0.578)rectangle++(0.009,0.009)
  (0.713,0.275)rectangle++(0.009,0.009) (0.86,0.13)rectangle++(0.009,0.009) (0.93,0.06)rectangle++(0.009,0.009) (0.964,0.026)rectangle++(0.009,0.009) (0.461,0.53)rectangle++(0.009,0.009)
  (0.739,0.252)rectangle++(0.009,0.009) (0.352,0.119)rectangle++(0.009,0.009) (0.167,0.575)rectangle++(0.009,0.009) (0.078,0.793)rectangle++(0.009,0.009) (0.555,0.379)rectangle++(0.009,0.009)
  (0.784,0.179)rectangle++(0.009,0.009) (0.374,0.084)rectangle++(0.009,0.009) (0.697,0.038)rectangle++(0.009,0.009) (0.852,0.016)rectangle++(0.009,0.009) (0.407,0.005)rectangle++(0.009,0.009)
  (0.193,0.52)rectangle++(0.009,0.009) (0.61,0.247)rectangle++(0.009,0.009) (0.291,0.116)rectangle++(0.009,0.009) (0.657,0.054)rectangle++(0.009,0.009) (0.833,0.023)rectangle++(0.009,0.009)
  (0.398,0.529)rectangle++(0.009,0.009) (0.188,0.252)rectangle++(0.009,0.009) (0.088,0.638)rectangle++(0.009,0.009) (0.56,0.304)rectangle++(0.009,0.009) (0.266,0.664)rectangle++(0.009,0.009)
  (0.646,0.316)rectangle++(0.009,0.009) (0.828,0.149)rectangle++(0.009,0.009) (0.395,0.589)rectangle++(0.009,0.009) (0.187,0.281)rectangle++(0.009,0.009) (0.088,0.652)rectangle++(0.009,0.009)
  (0.56,0.311)rectangle++(0.009,0.009) (0.266,0.667)rectangle++(0.009,0.009) (0.125,0.318)rectangle++(0.009,0.009) (0.058,0.67)rectangle++(0.009,0.009) (0.025,0.839)rectangle++(0.009,0.009)
  (0.01,0.401)rectangle++(0.009,0.009) (0.002,0.19)rectangle++(0.009,0.009) (-0.001,0.089)rectangle++(0.009,0.009) (0.517,0.04)rectangle++(0.009,0.009) (0.246,0.017)rectangle++(0.009,0.009)
  (0.116,0.006)rectangle++(0.009,0.009) (0.573,0)rectangle++(0.009,0.009) (0.273,-0.002)rectangle++(0.009,0.009) (0.129,0.517)rectangle++(0.009,0.009) (0.059,0.246)rectangle++(0.009,0.009)
  (0.026,0.116)rectangle++(0.009,0.009) (0.01,0.053)rectangle++(0.009,0.009) (0.523,0.023)rectangle++(0.009,0.009) (0.248,0.529)rectangle++(0.009,0.009) (0.117,0.251)rectangle++(0.009,0.009)
  (0.054,0.118)rectangle++(0.009,0.009) (0.543,0.054)rectangle++(0.009,0.009) (0.259,0.024)rectangle++(0.009,0.009) (0.642,0.009)rectangle++(0.009,0.009) (0.306,0.002)rectangle++(0.009,0.009)
  (0.144,-0.001)rectangle++(0.009,0.009) (0.067,-0.003)rectangle++(0.009,0.009) (0.03,0.516)rectangle++(0.009,0.009) (0.532,0.245)rectangle++(0.009,0.009) (0.253,0.635)rectangle++(0.009,0.009)
  (0.119,0.823)rectangle++(0.009,0.009) (0.055,0.913)rectangle++(0.009,0.009) (0.024,0.956)rectangle++(0.009,0.009) (0.009,0.456)rectangle++(0.009,0.009) (0.002,0.737)rectangle++(0.009,0.009)
  (-0.001,0.351)rectangle++(0.009,0.009) (0.517,0.166)rectangle++(0.009,0.009) (0.246,0.597)rectangle++(0.009,0.009) (0.116,0.804)rectangle++(0.009,0.009) (0.053,0.904)rectangle++(0.009,0.009)
  (0.023,0.431)rectangle++(0.009,0.009) (0.009,0.205)rectangle++(0.009,0.009) (0.002,0.616)rectangle++(0.009,0.009) (0.519,0.293)rectangle++(0.009,0.009) (0.247,0.138)rectangle++(0.009,0.009)
  (0.636,0.064)rectangle++(0.009,0.009) (0.303,0.028)rectangle++(0.009,0.009) (0.663,0.011)rectangle++(0.009,0.009) (0.316,0.003)rectangle++(0.009,0.009) (0.149,-0.001)rectangle++(0.009,0.009)
  (0.069,0.517)rectangle++(0.009,0.009) (0.031,0.766)rectangle++(0.009,0.009) (0.533,0.365)rectangle++(0.009,0.009) (0.253,0.173)rectangle++(0.009,0.009) (0.119,0.601)rectangle++(0.009,0.009)
  (0.055,0.286)rectangle++(0.009,0.009) (0.024,0.655)rectangle++(0.009,0.009) (0.529,0.312)rectangle++(0.009,0.009) (0.252,0.667)rectangle++(0.009,0.009) (0.638,0.318)rectangle++(0.009,0.009)
  (0.304,0.67)rectangle++(0.009,0.009) (0.144,0.319)rectangle++(0.009,0.009) (0.067,0.151)rectangle++(0.009,0.009) (0.03,0.59)rectangle++(0.009,0.009) (0.012,0.281)rectangle++(0.009,0.009)
  (0.003,0.653)rectangle++(0.009,0.009) (0.519,0.311)rectangle++(0.009,0.009) (0.247,0.667)rectangle++(0.009,0.009) (0.116,0.318)rectangle++(0.009,0.009) (0.573,0.15)rectangle++(0.009,0.009)
  (0.273,0.07)rectangle++(0.009,0.009) (0.649,0.031)rectangle++(0.009,0.009) (0.829,0.013)rectangle++(0.009,0.009) (0.916,0.004)rectangle++(0.009,0.009) (0.957,-0.001)rectangle++(0.009,0.009)
  (0.457,-0.003)rectangle++(0.009,0.009) (0.217,0.516)rectangle++(0.009,0.009) (0.102,0.766)rectangle++(0.009,0.009) (0.567,0.365)rectangle++(0.009,0.009) (0.27,0.693)rectangle++(0.009,0.009)
  (0.647,0.33)rectangle++(0.009,0.009) (0.308,0.156)rectangle++(0.009,0.009) (0.146,0.073)rectangle++(0.009,0.009) (0.588,0.033)rectangle++(0.009,0.009) (0.28,0.013)rectangle++(0.009,0.009)
  (0.652,0.004)rectangle++(0.009,0.009) (0.311,0.52)rectangle++(0.009,0.009) (0.147,0.767)rectangle++(0.009,0.009) (0.588,0.366)rectangle++(0.009,0.009) (0.28,0.693)rectangle++(0.009,0.009)
  (0.652,0.33)rectangle++(0.009,0.009) (0.831,0.156)rectangle++(0.009,0.009) (0.916,0.073)rectangle++(0.009,0.009) (0.958,0.033)rectangle++(0.009,0.009) (0.457,0.533)rectangle++(0.009,0.009)
  (0.737,0.254)rectangle++(0.009,0.009) (0.871,0.119)rectangle++(0.009,0.009) (0.936,0.055)rectangle++(0.009,0.009) (0.967,0.024)rectangle++(0.009,0.009) (0.982,0.009)rectangle++(0.009,0.009)
  (0.989,0.002)rectangle++(0.009,0.009) (0.472,-0.001)rectangle++(0.009,0.009) (0.224,-0.003)rectangle++(0.009,0.009) (0.625,-0.004)rectangle++(0.009,0.009) (0.298,0.516)rectangle++(0.009,0.009)
  (0.141,0.245)rectangle++(0.009,0.009) (0.585,0.115)rectangle++(0.009,0.009) (0.279,0.053)rectangle++(0.009,0.009) (0.131,0.543)rectangle++(0.009,0.009) (0.061,0.258)rectangle++(0.009,0.009)
  (0.027,0.122)rectangle++(0.009,0.009) (0.531,0.056)rectangle++(0.009,0.009) (0.252,0.545)rectangle++(0.009,0.009) (0.119,0.259)rectangle++(0.009,0.009) (0.575,0.122)rectangle++(0.009,0.009)
  (0.794,0.056)rectangle++(0.009,0.009) (0.379,0.025)rectangle++(0.009,0.009) (0.699,0.009)rectangle++(0.009,0.009) (0.333,0.522)rectangle++(0.009,0.009) (0.158,0.248)rectangle++(0.009,0.009)
  (0.593,0.117)rectangle++(0.009,0.009) (0.282,0.574)rectangle++(0.009,0.009) (0.653,0.273)rectangle++(0.009,0.009) (0.831,0.129)rectangle++(0.009,0.009) (0.917,0.059)rectangle++(0.009,0.009)
  (0.438,0.026)rectangle++(0.009,0.009) (0.208,0.53)rectangle++(0.009,0.009) (0.097,0.772)rectangle++(0.009,0.009) (0.044,0.888)rectangle++(0.009,0.009) (0.019,0.424)rectangle++(0.009,0.009)
  (0.007,0.201)rectangle++(0.009,0.009) (0.001,0.614)rectangle++(0.009,0.009) (0.518,0.292)rectangle++(0.009,0.009) (0.246,0.138)rectangle++(0.009,0.009) (0.636,0.064)rectangle++(0.009,0.009)
  (0.303,0.548)rectangle++(0.009,0.009) (0.143,0.781)rectangle++(0.009,0.009) (0.586,0.372)rectangle++(0.009,0.009) (0.279,0.696)rectangle++(0.009,0.009) (0.132,0.852)rectangle++(0.009,0.009)
  (0.061,0.407)rectangle++(0.009,0.009) (0.027,0.193)rectangle++(0.009,0.009) (0.011,0.09)rectangle++(0.009,0.009) (0.003,0.041)rectangle++(0.009,0.009) (-0.001,0.537)rectangle++(0.009,0.009)
  (-0.003,0.776)rectangle++(0.009,0.009) (-0.004,0.89)rectangle++(0.009,0.009) (-0.004,0.425)rectangle++(0.009,0.009) (-0.004,0.202)rectangle++(0.009,0.009) (-0.004,0.094)rectangle++(0.009,0.009)
  (-0.004,0.043)rectangle++(0.009,0.009) (0.516,0.018)rectangle++(0.009,0.009) (0.245,0.526)rectangle++(0.009,0.009) (0.115,0.77)rectangle++(0.009,0.009) (0.573,0.367)rectangle++(0.009,0.009)
  (0.273,0.694)rectangle++(0.009,0.009) (0.129,0.851)rectangle++(0.009,0.009) (0.579,0.406)rectangle++(0.009,0.009) (0.276,0.193)rectangle++(0.009,0.009) (0.13,0.09)rectangle++(0.009,0.009)
  (0.06,0.561)rectangle++(0.009,0.009) (0.026,0.267)rectangle++(0.009,0.009) (0.01,0.646)rectangle++(0.009,0.009) (0.003,0.828)rectangle++(0.009,0.009) (-0.001,0.915)rectangle++(0.009,0.009)
  (-0.003,0.957)rectangle++(0.009,0.009) (-0.004,0.457)rectangle++(0.009,0.009) (0.516,0.217)rectangle++(0.009,0.009) (0.765,0.102)rectangle++(0.009,0.009) (0.365,0.047)rectangle++(0.009,0.009)
  (0.693,0.02)rectangle++(0.009,0.009) (0.33,0.527)rectangle++(0.009,0.009) (0.676,0.251)rectangle++(0.009,0.009) (0.842,0.118)rectangle++(0.009,0.009) (0.922,0.054)rectangle++(0.009,0.009)
  (0.44,0.024)rectangle++(0.009,0.009) (0.729,0.009)rectangle++(0.009,0.009) (0.348,0.522)rectangle++(0.009,0.009) (0.164,0.768)rectangle++(0.009,0.009) (0.077,0.886)rectangle++(0.009,0.009)
  (0.554,0.423)rectangle++(0.009,0.009) (0.264,0.721)rectangle++(0.009,0.009) (0.644,0.344)rectangle++(0.009,0.009) (0.307,0.163)rectangle++(0.009,0.009) (0.145,0.076)rectangle++(0.009,0.009)
  (0.587,0.034)rectangle++(0.009,0.009) (0.28,0.534)rectangle++(0.009,0.009) (0.132,0.254)rectangle++(0.009,0.009) (0.581,0.12)rectangle++(0.009,0.009) (0.797,0.055)rectangle++(0.009,0.009)
  (0.38,0.544)rectangle++(0.009,0.009) (0.18,0.779)rectangle++(0.009,0.009) (0.084,0.371)rectangle++(0.009,0.009) (0.038,0.696)rectangle++(0.009,0.009) (0.016,0.332)rectangle++(0.009,0.009)
  (0.005,0.677)rectangle++(0.009,0.009) (0.52,0.323)rectangle++(0.009,0.009) (0.247,0.672)rectangle++(0.009,0.009) (0.116,0.84)rectangle++(0.009,0.009) (0.054,0.401)rectangle++(0.009,0.009)
  (0.023,0.19)rectangle++(0.009,0.009) (0.529,0.089)rectangle++(0.009,0.009) (0.252,0.04)rectangle++(0.009,0.009) (0.638,0.017)rectangle++(0.009,0.009) (0.304,0.526)rectangle++(0.009,0.009)
  (0.144,0.25)rectangle++(0.009,0.009) (0.587,0.118)rectangle++(0.009,0.009) (0.279,0.054)rectangle++(0.009,0.009) (0.652,0.024)rectangle++(0.009,0.009) (0.83,0.009)rectangle++(0.009,0.009)
  (0.916,0.002)rectangle++(0.009,0.009) (0.957,-0.001)rectangle++(0.009,0.009) (0.977,-0.003)rectangle++(0.009,0.009) (0.467,-0.004)rectangle++(0.009,0.009) (0.742,-0.004)rectangle++(0.009,0.009)
  (0.874,-0.004)rectangle++(0.009,0.009) (0.417,-0.004)rectangle++(0.009,0.009) (0.718,-0.004)rectangle++(0.009,0.009) (0.342,-0.004)rectangle++(0.009,0.009) (0.162,0.516)rectangle++(0.009,0.009)
  (0.595,0.245)rectangle++(0.009,0.009) (0.803,0.115)rectangle++(0.009,0.009) (0.903,0.053)rectangle++(0.009,0.009) (0.431,0.023)rectangle++(0.009,0.009) (0.205,0.529)rectangle++(0.009,0.009)
  (0.616,0.251)rectangle++(0.009,0.009) (0.293,0.638)rectangle++(0.009,0.009) (0.138,0.824)rectangle++(0.009,0.009) (0.064,0.913)rectangle++(0.009,0.009) (0.548,0.436)rectangle++(0.009,0.009)
  (0.781,0.207)rectangle++(0.009,0.009) (0.372,0.097)rectangle++(0.009,0.009) (0.176,0.044)rectangle++(0.009,0.009) (0.602,0.019)rectangle++(0.009,0.009) (0.287,0.007)rectangle++(0.009,0.009)
  (0.135,0.521)rectangle++(0.009,0.009) (0.583,0.248)rectangle++(0.009,0.009) (0.277,0.637)rectangle++(0.009,0.009) (0.131,0.303)rectangle++(0.009,0.009) (0.06,0.663)rectangle++(0.009,0.009)
  (0.027,0.316)rectangle++(0.009,0.009) (0.01,0.149)rectangle++(0.009,0.009) (0.003,0.069)rectangle++(0.009,0.009) (-0.001,0.031)rectangle++(0.009,0.009) (0.517,0.013)rectangle++(0.009,0.009)
  (0.246,0.004)rectangle++(0.009,0.009) (0.636,-0.001)rectangle++(0.009,0.009) (0.303,-0.003)rectangle++(0.009,0.009) (0.143,-0.004)rectangle++(0.009,0.009) (0.066,0.516)rectangle++(0.009,0.009)
  (0.029,0.245)rectangle++(0.009,0.009) (0.012,0.115)rectangle++(0.009,0.009) (0.003,0.573)rectangle++(0.009,0.009) (0.519,0.273)rectangle++(0.009,0.009) (0.247,0.649)rectangle++(0.009,0.009)
  (0.116,0.829)rectangle++(0.009,0.009) (0.573,0.396)rectangle++(0.009,0.009) (0.273,0.708)rectangle++(0.009,0.009) (0.129,0.857)rectangle++(0.009,0.009) (0.579,0.409)rectangle++(0.009,0.009)
  (0.276,0.194)rectangle++(0.009,0.009) (0.13,0.091)rectangle++(0.009,0.009) (0.06,0.041)rectangle++(0.009,0.009) (0.026,0.017)rectangle++(0.009,0.009) (0.01,0.526)rectangle++(0.009,0.009)
  (0.003,0.77)rectangle++(0.009,0.009) (-0.001,0.887)rectangle++(0.009,0.009) (-0.003,0.424)rectangle++(0.009,0.009) (0.516,0.201)rectangle++(0.009,0.009) (0.765,0.094)rectangle++(0.009,0.009)
  (0.885,0.043)rectangle++(0.009,0.009) (0.943,0.018)rectangle++(0.009,0.009) (0.97,0.006)rectangle++(0.009,0.009) (0.463,0.001)rectangle++(0.009,0.009) (0.22,0.518)rectangle++(0.009,0.009)
  (0.103,0.766)rectangle++(0.009,0.009) (0.047,0.885)rectangle++(0.009,0.009) (0.02,0.943)rectangle++(0.009,0.009) (0.527,0.45)rectangle++(0.009,0.009) (0.251,0.214)rectangle++(0.009,0.009)
  (0.118,0.1)rectangle++(0.009,0.009) (0.054,0.046)rectangle++(0.009,0.009) (0.544,0.02)rectangle++(0.009,0.009) (0.259,0.527)rectangle++(0.009,0.009) (0.642,0.251)rectangle++(0.009,0.009)
  (0.826,0.118)rectangle++(0.009,0.009) (0.394,0.574)rectangle++(0.009,0.009) (0.187,0.273)rectangle++(0.009,0.009) (0.087,0.129)rectangle++(0.009,0.009) (0.04,0.06)rectangle++(0.009,0.009)
  (0.537,0.026)rectangle++(0.009,0.009) (0.775,0.01)rectangle++(0.009,0.009) (0.37,0.003)rectangle++(0.009,0.009) (0.175,0.519)rectangle++(0.009,0.009) (0.082,0.247)rectangle++(0.009,0.009)
  (0.037,0.116)rectangle++(0.009,0.009) (0.535,0.053)rectangle++(0.009,0.009) (0.255,0.543)rectangle++(0.009,0.009) (0.64,0.258)rectangle++(0.009,0.009) (0.825,0.122)rectangle++(0.009,0.009)
  (0.394,0.576)rectangle++(0.009,0.009) (0.187,0.274)rectangle++(0.009,0.009) (0.087,0.129)rectangle++(0.009,0.009) (0.56,0.06)rectangle++(0.009,0.009) (0.786,0.026)rectangle++(0.009,0.009)
  (0.375,0.53)rectangle++(0.009,0.009) (0.178,0.252)rectangle++(0.009,0.009) (0.603,0.119)rectangle++(0.009,0.009) (0.807,0.055)rectangle++(0.009,0.009) (0.905,0.024)rectangle++(0.009,0.009)
  (0.952,0.009)rectangle++(0.009,0.009) (0.455,0.002)rectangle++(0.009,0.009) (0.216,0.519)rectangle++(0.009,0.009) (0.101,0.767)rectangle++(0.009,0.009) (0.566,0.366)rectangle++(0.009,0.009)
  (0.269,0.693)rectangle++(0.009,0.009) (0.127,0.33)rectangle++(0.009,0.009) (0.579,0.156)rectangle++(0.009,0.009) (0.275,0.593)rectangle++(0.009,0.009) (0.13,0.282)rectangle++(0.009,0.009)
  (0.06,0.653)rectangle++(0.009,0.009) (0.026,0.831)rectangle++(0.009,0.009) (0.01,0.917)rectangle++(0.009,0.009) (0.003,0.438)rectangle++(0.009,0.009) (-0.001,0.728)rectangle++(0.009,0.009)
  (0.517,0.347)rectangle++(0.009,0.009) (0.246,0.164)rectangle++(0.009,0.009) (0.636,0.076)rectangle++(0.009,0.009) (0.823,0.034)rectangle++(0.009,0.009) (0.393,0.534)rectangle++(0.009,0.009)
  (0.706,0.254)rectangle++(0.009,0.009) (0.857,0.12)rectangle++(0.009,0.009) (0.409,0.055)rectangle++(0.009,0.009) (0.194,0.544)rectangle++(0.009,0.009) (0.091,0.259)rectangle++(0.009,0.009)
  (0.041,0.642)rectangle++(0.009,0.009) (0.017,0.306)rectangle++(0.009,0.009) (0.006,0.664)rectangle++(0.009,0.009) (0.521,0.317)rectangle++(0.009,0.009) (0.248,0.67)rectangle++(0.009,0.009)
  (0.636,0.319)rectangle++(0.009,0.009) (0.823,0.151)rectangle++(0.009,0.009) (0.913,0.07)rectangle++(0.009,0.009) (0.956,0.031)rectangle++(0.009,0.009) (0.456,0.533)rectangle++(0.009,0.009)
  (0.217,0.253)rectangle++(0.009,0.009) (0.622,0.119)rectangle++(0.009,0.009) (0.816,0.055)rectangle++(0.009,0.009) (0.389,0.544)rectangle++(0.009,0.009) (0.185,0.779)rectangle++(0.009,0.009)
  (0.086,0.371)rectangle++(0.009,0.009) (0.039,0.696)rectangle++(0.009,0.009) (0.016,0.332)rectangle++(0.009,0.009) (0.526,0.157)rectangle++(0.009,0.009) (0.25,0.593)rectangle++(0.009,0.009)
  (0.638,0.282)rectangle++(0.009,0.009) (0.824,0.133)rectangle++(0.009,0.009) (0.393,0.582)rectangle++(0.009,0.009) (0.186,0.797)rectangle++(0.009,0.009) (0.607,0.38)rectangle++(0.009,0.009)
  (0.809,0.18)rectangle++(0.009,0.009) (0.906,0.084)rectangle++(0.009,0.009) (0.953,0.038)rectangle++(0.009,0.009) (0.455,0.536)rectangle++(0.009,0.009) (0.216,0.775)rectangle++(0.009,0.009)
  (0.101,0.89)rectangle++(0.009,0.009) (0.566,0.425)rectangle++(0.009,0.009) (0.269,0.722)rectangle++(0.009,0.009) (0.127,0.864)rectangle++(0.009,0.009) (0.059,0.932)rectangle++(0.009,0.009)
  (0.026,0.965)rectangle++(0.009,0.009) (0.53,0.461)rectangle++(0.009,0.009) (0.772,0.219)rectangle++(0.009,0.009) (0.888,0.103)rectangle++(0.009,0.009) (0.424,0.047)rectangle++(0.009,0.009)
  (0.201,0.54)rectangle++(0.009,0.009) (0.094,0.257)rectangle++(0.009,0.009) (0.043,0.641)rectangle++(0.009,0.009) (0.538,0.305)rectangle++(0.009,0.009) (0.256,0.664)rectangle++(0.009,0.009)
  (0.641,0.316)rectangle++(0.009,0.009) (0.825,0.15)rectangle++(0.009,0.009) (0.394,0.589)rectangle++(0.009,0.009) (0.187,0.281)rectangle++(0.009,0.009) (0.607,0.132)rectangle++(0.009,0.009)
  (0.289,0.581)rectangle++(0.009,0.009) (0.136,0.797)rectangle++(0.009,0.009) (0.583,0.38)rectangle++(0.009,0.009) (0.798,0.18)rectangle++(0.009,0.009) (0.38,0.604)rectangle++(0.009,0.009)
  (0.18,0.808)rectangle++(0.009,0.009) (0.084,0.905)rectangle++(0.009,0.009) (0.558,0.432)rectangle++(0.009,0.009) (0.266,0.205)rectangle++(0.009,0.009) (0.645,0.096)rectangle++(0.009,0.009)
  (0.307,0.564)rectangle++(0.009,0.009) (0.665,0.268)rectangle++(0.009,0.009) (0.837,0.126)rectangle++(0.009,0.009) (0.919,0.058)rectangle++(0.009,0.009) (0.959,0.026)rectangle++(0.009,0.009)
  (0.458,0.01)rectangle++(0.009,0.009) (0.217,0.522)rectangle++(0.009,0.009) (0.622,0.248)rectangle++(0.009,0.009) (0.296,0.117)rectangle++(0.009,0.009) (0.66,0.054)rectangle++(0.009,0.009)
  (0.834,0.023)rectangle++(0.009,0.009) (0.918,0.009)rectangle++(0.009,0.009) (0.438,0.522)rectangle++(0.009,0.009) (0.728,0.248)rectangle++(0.009,0.009) (0.867,0.117)rectangle++(0.009,0.009)
  (0.934,0.054)rectangle++(0.009,0.009) (0.966,0.023)rectangle++(0.009,0.009) (0.461,0.529)rectangle++(0.009,0.009) (0.219,0.252)rectangle++(0.009,0.009) (0.623,0.118)rectangle++(0.009,0.009)
  (0.817,0.054)rectangle++(0.009,0.009) (0.39,0.024)rectangle++(0.009,0.009) (0.185,0.009)rectangle++(0.009,0.009) (0.086,0.522)rectangle++(0.009,0.009) (0.039,0.248)rectangle++(0.009,0.009)
  (0.016,0.117)rectangle++(0.009,0.009) (0.526,0.054)rectangle++(0.009,0.009) (0.77,0.023)rectangle++(0.009,0.009) (0.367,0.009)rectangle++(0.009,0.009) (0.174,0.002)rectangle++(0.009,0.009)
  (0.081,0.519)rectangle++(0.009,0.009) (0.557,0.247)rectangle++(0.009,0.009) (0.265,0.636)rectangle++(0.009,0.009) (0.125,0.823)rectangle++(0.009,0.009) (0.058,0.393)rectangle++(0.009,0.009)
  (0.025,0.186)rectangle++(0.009,0.009) (0.53,0.087)rectangle++(0.009,0.009) (0.252,0.559)rectangle++(0.009,0.009) (0.639,0.266)rectangle++(0.009,0.009) (0.824,0.125)rectangle++(0.009,0.009)
  (0.393,0.578)rectangle++(0.009,0.009) (0.186,0.795)rectangle++(0.009,0.009) (0.087,0.899)rectangle++(0.009,0.009) (0.039,0.949)rectangle++(0.009,0.009) (0.537,0.453)rectangle++(0.009,0.009)
  (0.255,0.735)rectangle++(0.009,0.009) (0.64,0.351)rectangle++(0.009,0.009) (0.305,0.166)rectangle++(0.009,0.009) (0.144,0.077)rectangle++(0.009,0.009) (0.067,0.035)rectangle++(0.009,0.009)
  (0.03,0.534)rectangle++(0.009,0.009) (0.532,0.254)rectangle++(0.009,0.009) (0.253,0.64)rectangle++(0.009,0.009) (0.119,0.305)rectangle++(0.009,0.009) (0.575,0.144)rectangle++(0.009,0.009)
  (0.794,0.067)rectangle++(0.009,0.009) (0.379,0.03)rectangle++(0.009,0.009) (0.179,0.012)rectangle++(0.009,0.009) (0.084,0.523)rectangle++(0.009,0.009) (0.038,0.769)rectangle++(0.009,0.009)
  (0.536,0.367)rectangle++(0.009,0.009) (0.255,0.174)rectangle++(0.009,0.009) (0.64,0.081)rectangle++(0.009,0.009) (0.825,0.037)rectangle++(0.009,0.009) (0.394,0.535)rectangle++(0.009,0.009)
  (0.187,0.775)rectangle++(0.009,0.009) (0.087,0.889)rectangle++(0.009,0.009) (0.56,0.425)rectangle++(0.009,0.009) (0.266,0.721)rectangle++(0.009,0.009) (0.125,0.344)rectangle++(0.009,0.009)
  (0.578,0.163)rectangle++(0.009,0.009) (0.275,0.596)rectangle++(0.009,0.009) (0.65,0.284)rectangle++(0.009,0.009) (0.31,0.654)rectangle++(0.009,0.009) (0.146,0.311)rectangle++(0.009,0.009)
  (0.068,0.147)rectangle++(0.009,0.009) (0.55,0.068)rectangle++(0.009,0.009) (0.262,0.03)rectangle++(0.009,0.009) (0.643,0.012)rectangle++(0.009,0.009) (0.306,0.524)rectangle++(0.009,0.009)
  (0.145,0.249)rectangle++(0.009,0.009) (0.067,0.117)rectangle++(0.009,0.009) (0.03,0.054)rectangle++(0.009,0.009) (0.532,0.024)rectangle++(0.009,0.009) (0.253,0.529)rectangle++(0.009,0.009)
  (0.119,0.252)rectangle++(0.009,0.009) (0.055,0.638)rectangle++(0.009,0.009) (0.024,0.304)rectangle++(0.009,0.009) (0.009,0.144)rectangle++(0.009,0.009) (0.522,0.067)rectangle++(0.009,0.009)
  (0.248,0.03)rectangle++(0.009,0.009) (0.117,0.012)rectangle++(0.009,0.009) (0.054,0.523)rectangle++(0.009,0.009) (0.023,0.769)rectangle++(0.009,0.009) (0.009,0.367)rectangle++(0.009,0.009)
  (0.522,0.174)rectangle++(0.009,0.009) (0.248,0.081)rectangle++(0.009,0.009) (0.117,0.037)rectangle++(0.009,0.009) (0.054,0.015)rectangle++(0.009,0.009) (0.023,0.005)rectangle++(0.009,0.009)
  (0.009,0.52)rectangle++(0.009,0.009) (0.522,0.247)rectangle++(0.009,0.009) (0.768,0.116)rectangle++(0.009,0.009) (0.886,0.054)rectangle++(0.009,0.009) (0.943,0.023)rectangle++(0.009,0.009)
  (0.45,0.009)rectangle++(0.009,0.009) (0.214,0.522)rectangle++(0.009,0.009) (0.62,0.248)rectangle++(0.009,0.009) (0.295,0.117)rectangle++(0.009,0.009) (0.139,0.054)rectangle++(0.009,0.009)
  (0.065,0.023)rectangle++(0.009,0.009) (0.029,0.009)rectangle++(0.009,0.009) (0.011,0.522)rectangle++(0.009,0.009) (0.003,0.768)rectangle++(0.009,0.009) (-0.001,0.886)rectangle++(0.009,0.009)
  (-0.003,0.423)rectangle++(0.009,0.009) (0.516,0.201)rectangle++(0.009,0.009) (0.246,0.094)rectangle++(0.009,0.009) (0.636,0.043)rectangle++(0.009,0.009) (0.303,0.018)rectangle++(0.009,0.009)
  (0.143,0.006)rectangle++(0.009,0.009) (0.586,0.001)rectangle++(0.009,0.009) (0.279,0.518)rectangle++(0.009,0.009) (0.132,0.766)rectangle++(0.009,0.009) (0.061,0.365)rectangle++(0.009,0.009)
  (0.027,0.693)rectangle++(0.009,0.009) (0.531,0.33)rectangle++(0.009,0.009) (0.252,0.156)rectangle++(0.009,0.009) (0.639,0.073)rectangle++(0.009,0.009) (0.304,0.033)rectangle++(0.009,0.009)
  (0.144,0.533)rectangle++(0.009,0.009) (0.067,0.774)rectangle++(0.009,0.009) (0.03,0.889)rectangle++(0.009,0.009) (0.532,0.424)rectangle++(0.009,0.009) (0.253,0.721)rectangle++(0.009,0.009)
  (0.119,0.864)rectangle++(0.009,0.009) (0.055,0.412)rectangle++(0.009,0.009) (0.024,0.196)rectangle++(0.009,0.009) (0.009,0.612)rectangle++(0.009,0.009) (0.002,0.811)rectangle++(0.009,0.009)
  (-0.001,0.387)rectangle++(0.009,0.009) (0.517,0.183)rectangle++(0.009,0.009) (0.766,0.086)rectangle++(0.009,0.009) (0.365,0.559)rectangle++(0.009,0.009) (0.173,0.786)rectangle++(0.009,0.009)
  (0.601,0.375)rectangle++(0.009,0.009) (0.806,0.178)rectangle++(0.009,0.009) (0.385,0.603)rectangle++(0.009,0.009) (0.182,0.287)rectangle++(0.009,0.009) (0.605,0.135)rectangle++(0.009,0.009)
  (0.808,0.063)rectangle++(0.009,0.009) (0.386,0.548)rectangle++(0.009,0.009) (0.183,0.781)rectangle++(0.009,0.009) (0.085,0.372)rectangle++(0.009,0.009) (0.039,0.176)rectangle++(0.009,0.009)
  (0.016,0.602)rectangle++(0.009,0.009) (0.005,0.287)rectangle++(0.009,0.009) (0,0.135)rectangle++(0.009,0.009) (-0.002,0.583)rectangle++(0.009,0.009) (-0.003,0.277)rectangle++(0.009,0.009)
  (-0.004,0.131)rectangle++(0.009,0.009) (-0.004,0.58)rectangle++(0.009,0.009) (-0.004,0.276)rectangle++(0.009,0.009) (-0.004,0.13)rectangle++(0.009,0.009) (0.516,0.06)rectangle++(0.009,0.009)
  (0.765,0.027)rectangle++(0.009,0.009) (0.365,0.01)rectangle++(0.009,0.009) (0.173,0.523)rectangle++(0.009,0.009) (0.081,0.249)rectangle++(0.009,0.009) (0.036,0.637)rectangle++(0.009,0.009)
  (0.015,0.823)rectangle++(0.009,0.009) (0.525,0.393)rectangle++(0.009,0.009) (0.25,0.706)rectangle++(0.009,0.009) (0.637,0.337)rectangle++(0.009,0.009) (0.304,0.159)rectangle++(0.009,0.009)
  (0.663,0.074)rectangle++(0.009,0.009) (0.316,0.553)rectangle++(0.009,0.009) (0.669,0.263)rectangle++(0.009,0.009) (0.319,0.124)rectangle++(0.009,0.009) (0.671,0.057)rectangle++(0.009,0.009)
  (0.84,0.025)rectangle++(0.009,0.009) (0.401,0.53)rectangle++(0.009,0.009) (0.19,0.252)rectangle++(0.009,0.009) (0.089,0.119)rectangle++(0.009,0.009) (0.04,0.575)rectangle++(0.009,0.009)
  (0.017,0.793)rectangle++(0.009,0.009) (0.526,0.379)rectangle++(0.009,0.009) (0.25,0.179)rectangle++(0.009,0.009) (0.118,0.084)rectangle++(0.009,0.009) (0.574,0.038)rectangle++(0.009,0.009)
  (0.273,0.016)rectangle++(0.009,0.009) (0.129,0.005)rectangle++(0.009,0.009) (0.059,0)rectangle++(0.009,0.009) (0.026,-0.002)rectangle++(0.009,0.009) (0.53,-0.003)rectangle++(0.009,0.009)
  (0.252,-0.004)rectangle++(0.009,0.009) (0.119,-0.004)rectangle++(0.009,0.009) (0.055,0.516)rectangle++(0.009,0.009) (0.024,0.245)rectangle++(0.009,0.009) (0.529,0.115)rectangle++(0.009,0.009)
  (0.252,0.573)rectangle++(0.009,0.009) (0.638,0.273)rectangle++(0.009,0.009) (0.824,0.129)rectangle++(0.009,0.009) (0.393,0.579)rectangle++(0.009,0.009) (0.706,0.276)rectangle++(0.009,0.009)
  (0.857,0.13)rectangle++(0.009,0.009) (0.409,0.58)rectangle++(0.009,0.009) (0.714,0.276)rectangle++(0.009,0.009) (0.34,0.13)rectangle++(0.009,0.009) (0.681,0.06)rectangle++(0.009,0.009)
  (0.325,0.547)rectangle++(0.009,0.009) (0.153,0.26)rectangle++(0.009,0.009) (0.071,0.642)rectangle++(0.009,0.009) (0.552,0.306)rectangle++(0.009,0.009) (0.263,0.145)rectangle++(0.009,0.009)
  (0.124,0.067)rectangle++(0.009,0.009) (0.577,0.03)rectangle++(0.009,0.009) (0.275,0.532)rectangle++(0.009,0.009) (0.129,0.773)rectangle++(0.009,0.009) (0.06,0.889)rectangle++(0.009,0.009)
  (0.026,0.424)rectangle++(0.009,0.009) (0.01,0.201)rectangle++(0.009,0.009) (0.523,0.094)rectangle++(0.009,0.009) (0.249,0.563)rectangle++(0.009,0.009) (0.117,0.788)rectangle++(0.009,0.009)
  (0.054,0.376)rectangle++(0.009,0.009) (0.543,0.178)rectangle++(0.009,0.009) (0.259,0.603)rectangle++(0.009,0.009) (0.122,0.287)rectangle++(0.009,0.009) (0.056,0.135)rectangle++(0.009,0.009)
  (0.025,0.583)rectangle++(0.009,0.009) (0.009,0.277)rectangle++(0.009,0.009) (0.522,0.131)rectangle++(0.009,0.009) (0.248,0.58)rectangle++(0.009,0.009) (0.637,0.276)rectangle++(0.009,0.009)
  (0.303,0.13)rectangle++(0.009,0.009) (0.143,0.58)rectangle++(0.009,0.009) (0.066,0.276)rectangle++(0.009,0.009) (0.55,0.13)rectangle++(0.009,0.009) (0.261,0.06)rectangle++(0.009,0.009)
  (0.643,0.027)rectangle++(0.009,0.009) (0.306,0.01)rectangle++(0.009,0.009) (0.665,0.003)rectangle++(0.009,0.009) (0.317,0.519)rectangle++(0.009,0.009) (0.15,0.767)rectangle++(0.009,0.009)
  (0.07,0.366)rectangle++(0.009,0.009) (0.031,0.173)rectangle++(0.009,0.009) (0.013,0.601)rectangle++(0.009,0.009) (0.004,0.286)rectangle++(0.009,0.009) (0.519,0.135)rectangle++(0.009,0.009)
  (0.767,0.062)rectangle++(0.009,0.009) (0.886,0.028)rectangle++(0.009,0.009) (0.423,0.011)rectangle++(0.009,0.009) (0.201,0.523)rectangle++(0.009,0.009) (0.094,0.769)rectangle++(0.009,0.009)
  (0.563,0.367)rectangle++(0.009,0.009) (0.788,0.174)rectangle++(0.009,0.009) (0.376,0.081)rectangle++(0.009,0.009) (0.698,0.037)rectangle++(0.009,0.009) (0.853,0.015)rectangle++(0.009,0.009)
  (0.407,0.005)rectangle++(0.009,0.009) (0.193,0.52)rectangle++(0.009,0.009) (0.09,0.247)rectangle++(0.009,0.009) (0.561,0.116)rectangle++(0.009,0.009) (0.267,0.574)rectangle++(0.009,0.009)
  (0.646,0.273)rectangle++(0.009,0.009) (0.828,0.129)rectangle++(0.009,0.009) (0.395,0.059)rectangle++(0.009,0.009) (0.707,0.026)rectangle++(0.009,0.009) (0.857,0.01)rectangle++(0.009,0.009)
  (0.409,0.523)rectangle++(0.009,0.009) (0.714,0.248)rectangle++(0.009,0.009) (0.34,0.117)rectangle++(0.009,0.009) (0.161,0.054)rectangle++(0.009,0.009) (0.075,0.543)rectangle++(0.009,0.009)
  (0.034,0.779)rectangle++(0.009,0.009) (0.014,0.891)rectangle++(0.009,0.009) (0.524,0.426)rectangle++(0.009,0.009) (0.249,0.202)rectangle++(0.009,0.009) (0.117,0.615)rectangle++(0.009,0.009)
  (0.574,0.293)rectangle++(0.009,0.009) (0.273,0.138)rectangle++(0.009,0.009) (0.649,0.064)rectangle++(0.009,0.009) (0.309,0.028)rectangle++(0.009,0.009) (0.146,0.531)rectangle++(0.009,0.009)
  (0.068,0.773)rectangle++(0.009,0.009) (0.03,0.889)rectangle++(0.009,0.009) (0.012,0.424)rectangle++(0.009,0.009) (0.523,0.201)rectangle++(0.009,0.009) (0.249,0.094)rectangle++(0.009,0.009)
  (0.117,0.563)rectangle++(0.009,0.009) (0.054,0.788)rectangle++(0.009,0.009) (0.544,0.376)rectangle++(0.009,0.009) (0.779,0.178)rectangle++(0.009,0.009) (0.891,0.083)rectangle++(0.009,0.009)
  (0.426,0.038)rectangle++(0.009,0.009) (0.202,0.016)rectangle++(0.009,0.009) (0.095,0.525)rectangle++(0.009,0.009) (0.043,0.25)rectangle++(0.009,0.009) (0.018,0.118)rectangle++(0.009,0.009)
  (0.006,0.574)rectangle++(0.009,0.009) (0.001,0.273)rectangle++(0.009,0.009) (-0.002,0.129)rectangle++(0.009,0.009) (-0.003,0.059)rectangle++(0.009,0.009) (-0.004,0.546)rectangle++(0.009,0.009)
  (0.516,0.26)rectangle++(0.009,0.009) (0.245,0.122)rectangle++(0.009,0.009) (0.635,0.056)rectangle++(0.009,0.009) (0.303,0.025)rectangle++(0.009,0.009) (0.143,0.01)rectangle++(0.009,0.009)
  (0.586,0.002)rectangle++(0.009,0.009) (0.799,-0.001)rectangle++(0.009,0.009) (0.901,-0.003)rectangle++(0.009,0.009) (0.95,-0.004)rectangle++(0.009,0.009) (0.974,-0.004)rectangle++(0.009,0.009)
  (0.465,0.516)rectangle++(0.009,0.009) (0.221,0.245)rectangle++(0.009,0.009) (0.624,0.115)rectangle++(0.009,0.009) (0.297,0.573)rectangle++(0.009,0.009) (0.14,0.273)rectangle++(0.009,0.009);

        \node[etiq, align=center] at (0.5,-0.2) {(a) trois copies de rapport $\num{0.48}$ :\\ disjointes};
      \end{scope}
      \begin{scope}[shift={(1.35,0)}]
        \draw[sectioncolor!60, dashed, thin] (0,0) rectangle (1,1);
        % Chevauchement : touche
% Fichier généré par generer.py : ne pas modifier à la main.
\fill[nuage]
  (0.255,0.108)rectangle++(0.009,0.009) (0.125,0.052)rectangle++(0.009,0.009) (0.06,0.024)rectangle++(0.009,0.009) (0.028,0.51)rectangle++(0.009,0.009) (0.012,0.253)rectangle++(0.009,0.009)
  (0.004,0.124)rectangle++(0.009,0.009) (0.499,0.047)rectangle++(0.009,0.009) (0.697,0.016)rectangle++(0.009,0.009) (0.346,0.506)rectangle++(0.009,0.009) (0.171,0.751)rectangle++(0.009,0.009)
  (0.083,0.373)rectangle++(0.009,0.009) (0.039,0.184)rectangle++(0.009,0.009) (0.017,0.59)rectangle++(0.009,0.009) (0.006,0.793)rectangle++(0.009,0.009) (0.001,0.394)rectangle++(0.009,0.009)
  (-0.002,0.695)rectangle++(0.009,0.009) (-0.003,0.845)rectangle++(0.009,0.009) (-0.004,0.92)rectangle++(0.009,0.009) (-0.004,0.458)rectangle++(0.009,0.009) (-0.004,0.727)rectangle++(0.009,0.009)
  (-0.004,0.861)rectangle++(0.009,0.009) (-0.004,0.428)rectangle++(0.009,0.009) (-0.004,0.712)rectangle++(0.009,0.009) (-0.004,0.354)rectangle++(0.009,0.009) (-0.004,0.175)rectangle++(0.009,0.009)
  (-0.004,0.085)rectangle++(0.009,0.009) (-0.004,0.04)rectangle++(0.009,0.009) (0.496,0.013)rectangle++(0.009,0.009) (0.246,0.504)rectangle++(0.009,0.009) (0.121,0.75)rectangle++(0.009,0.009)
  (0.058,0.873)rectangle++(0.009,0.009) (0.027,0.934)rectangle++(0.009,0.009) (0.011,0.965)rectangle++(0.009,0.009) (0.003,0.98)rectangle++(0.009,0.009) (-0.001,0.488)rectangle++(0.009,0.009)
  (-0.003,0.242)rectangle++(0.009,0.009) (-0.004,0.619)rectangle++(0.009,0.009) (0.496,0.245)rectangle++(0.009,0.009) (0.246,0.62)rectangle++(0.009,0.009) (0.121,0.308)rectangle++(0.009,0.009)
  (0.058,0.152)rectangle++(0.009,0.009) (0.027,0.574)rectangle++(0.009,0.009) (0.508,0.227)rectangle++(0.009,0.009) (0.252,0.611)rectangle++(0.009,0.009) (0.124,0.303)rectangle++(0.009,0.009)
  (0.547,0.119)rectangle++(0.009,0.009) (0.271,0.557)rectangle++(0.009,0.009) (0.133,0.776)rectangle++(0.009,0.009) (0.064,0.886)rectangle++(0.009,0.009) (0.03,0.441)rectangle++(0.009,0.009)
  (0.509,0.174)rectangle++(0.009,0.009) (0.252,0.585)rectangle++(0.009,0.009) (0.598,0.231)rectangle++(0.009,0.009) (0.297,0.613)rectangle++(0.009,0.009) (0.616,0.243)rectangle++(0.009,0.009)
  (0.306,0.119)rectangle++(0.009,0.009) (0.151,0.057)rectangle++(0.009,0.009) (0.073,0.526)rectangle++(0.009,0.009) (0.034,0.261)rectangle++(0.009,0.009) (0.015,0.128)rectangle++(0.009,0.009)
  (0.005,0.062)rectangle++(0.009,0.009) (0,0.529)rectangle++(0.009,0.009) (-0.002,0.762)rectangle++(0.009,0.009) (0.496,0.302)rectangle++(0.009,0.009) (0.696,0.118)rectangle++(0.009,0.009)
  (0.346,0.557)rectangle++(0.009,0.009) (0.171,0.276)rectangle++(0.009,0.009) (0.083,0.136)rectangle++(0.009,0.009) (0.531,0.052)rectangle++(0.009,0.009) (0.263,0.524)rectangle++(0.009,0.009)
  (0.129,0.26)rectangle++(0.009,0.009) (0.062,0.128)rectangle++(0.009,0.009) (0.029,0.062)rectangle++(0.009,0.009) (0.012,0.529)rectangle++(0.009,0.009) (0.004,0.762)rectangle++(0.009,0.009)
  (0,0.879)rectangle++(0.009,0.009) (-0.002,0.937)rectangle++(0.009,0.009) (-0.003,0.466)rectangle++(0.009,0.009) (-0.004,0.731)rectangle++(0.009,0.009) (-0.004,0.863)rectangle++(0.009,0.009)
  (-0.004,0.429)rectangle++(0.009,0.009) (-0.004,0.212)rectangle++(0.009,0.009) (-0.004,0.104)rectangle++(0.009,0.009) (-0.004,0.05)rectangle++(0.009,0.009) (-0.004,0.023)rectangle++(0.009,0.009)
  (-0.004,0.009)rectangle++(0.009,0.009) (-0.004,0.002)rectangle++(0.009,0.009) (-0.004,0.499)rectangle++(0.009,0.009) (-0.004,0.247)rectangle++(0.009,0.009) (0.496,0.096)rectangle++(0.009,0.009)
  (0.246,0.046)rectangle++(0.009,0.009) (0.121,0.021)rectangle++(0.009,0.009) (0.058,0.508)rectangle++(0.009,0.009) (0.027,0.252)rectangle++(0.009,0.009) (0.011,0.124)rectangle++(0.009,0.009)
  (0.502,0.047)rectangle++(0.009,0.009) (0.249,0.521)rectangle++(0.009,0.009) (0.597,0.206)rectangle++(0.009,0.009) (0.296,0.601)rectangle++(0.009,0.009) (0.616,0.238)rectangle++(0.009,0.009)
  (0.744,0.092)rectangle++(0.009,0.009) (0.37,0.544)rectangle++(0.009,0.009) (0.645,0.215)rectangle++(0.009,0.009) (0.32,0.605)rectangle++(0.009,0.009) (0.625,0.239)rectangle++(0.009,0.009)
  (0.31,0.117)rectangle++(0.009,0.009) (0.153,0.056)rectangle++(0.009,0.009) (0.074,0.526)rectangle++(0.009,0.009) (0.035,0.761)rectangle++(0.009,0.009) (0.015,0.878)rectangle++(0.009,0.009)
  (0.005,0.937)rectangle++(0.009,0.009) (0.499,0.372)rectangle++(0.009,0.009) (0.247,0.184)rectangle++(0.009,0.009) (0.121,0.59)rectangle++(0.009,0.009) (0.058,0.293)rectangle++(0.009,0.009)
  (0.521,0.114)rectangle++(0.009,0.009) (0.258,0.055)rectangle++(0.009,0.009) (0.601,0.019)rectangle++(0.009,0.009) (0.298,0.007)rectangle++(0.009,0.009) (0.617,0)rectangle++(0.009,0.009)
  (0.306,0.498)rectangle++(0.009,0.009) (0.62,0.196)rectangle++(0.009,0.009) (0.308,0.096)rectangle++(0.009,0.009) (0.152,0.546)rectangle++(0.009,0.009) (0.074,0.771)rectangle++(0.009,0.009)
  (0.035,0.883)rectangle++(0.009,0.009) (0.015,0.439)rectangle++(0.009,0.009) (0.503,0.173)rectangle++(0.009,0.009) (0.249,0.584)rectangle++(0.009,0.009) (0.122,0.29)rectangle++(0.009,0.009)
  (0.059,0.143)rectangle++(0.009,0.009) (0.027,0.569)rectangle++(0.009,0.009) (0.011,0.282)rectangle++(0.009,0.009) (0.502,0.11)rectangle++(0.009,0.009) (0.249,0.553)rectangle++(0.009,0.009)
  (0.597,0.218)rectangle++(0.009,0.009) (0.296,0.107)rectangle++(0.009,0.009) (0.146,0.051)rectangle++(0.009,0.009) (0.071,0.023)rectangle++(0.009,0.009) (0.033,0.509)rectangle++(0.009,0.009)
  (0.014,0.252)rectangle++(0.009,0.009) (0.005,0.624)rectangle++(0.009,0.009) (0,0.31)rectangle++(0.009,0.009) (-0.002,0.153)rectangle++(0.009,0.009) (-0.003,0.574)rectangle++(0.009,0.009)
  (0.496,0.227)rectangle++(0.009,0.009) (0.696,0.088)rectangle++(0.009,0.009) (0.346,0.542)rectangle++(0.009,0.009) (0.171,0.269)rectangle++(0.009,0.009) (0.083,0.132)rectangle++(0.009,0.009)
  (0.039,0.064)rectangle++(0.009,0.009) (0.513,0.023)rectangle++(0.009,0.009) (0.254,0.009)rectangle++(0.009,0.009) (0.599,0.001)rectangle++(0.009,0.009) (0.297,-0.002)rectangle++(0.009,0.009)
  (0.146,-0.003)rectangle++(0.009,0.009) (0.556,-0.004)rectangle++(0.009,0.009) (0.276,0.496)rectangle++(0.009,0.009) (0.136,0.246)rectangle++(0.009,0.009) (0.066,0.621)rectangle++(0.009,0.009)
  (0.524,0.246)rectangle++(0.009,0.009) (0.707,0.096)rectangle++(0.009,0.009) (0.351,0.546)rectangle++(0.009,0.009) (0.173,0.271)rectangle++(0.009,0.009) (0.567,0.106)rectangle++(0.009,0.009)
  (0.281,0.551)rectangle++(0.009,0.009) (0.61,0.218)rectangle++(0.009,0.009) (0.303,0.107)rectangle++(0.009,0.009) (0.149,0.051)rectangle++(0.009,0.009) (0.557,0.018)rectangle++(0.009,0.009)
  (0.276,0.007)rectangle++(0.009,0.009) (0.136,0.001)rectangle++(0.009,0.009) (0.552,-0.002)rectangle++(0.009,0.009) (0.274,0.497)rectangle++(0.009,0.009) (0.135,0.746)rectangle++(0.009,0.009)
  (0.551,0.296)rectangle++(0.009,0.009) (0.273,0.646)rectangle++(0.009,0.009) (0.134,0.321)rectangle++(0.009,0.009) (0.551,0.126)rectangle++(0.009,0.009) (0.273,0.561)rectangle++(0.009,0.009)
  (0.607,0.222)rectangle++(0.009,0.009) (0.301,0.109)rectangle++(0.009,0.009) (0.618,0.041)rectangle++(0.009,0.009) (0.307,0.518)rectangle++(0.009,0.009) (0.151,0.257)rectangle++(0.009,0.009)
  (0.073,0.626)rectangle++(0.009,0.009) (0.034,0.311)rectangle++(0.009,0.009) (0.015,0.653)rectangle++(0.009,0.009) (0.005,0.824)rectangle++(0.009,0.009) (0,0.41)rectangle++(0.009,0.009)
  (-0.002,0.703)rectangle++(0.009,0.009) (-0.003,0.849)rectangle++(0.009,0.009) (-0.004,0.422)rectangle++(0.009,0.009) (-0.004,0.209)rectangle++(0.009,0.009) (0.496,0.081)rectangle++(0.009,0.009)
  (0.246,0.538)rectangle++(0.009,0.009) (0.121,0.267)rectangle++(0.009,0.009) (0.058,0.131)rectangle++(0.009,0.009) (0.027,0.063)rectangle++(0.009,0.009) (0.011,0.029)rectangle++(0.009,0.009)
  (0.502,0.009)rectangle++(0.009,0.009) (0.698,0.001)rectangle++(0.009,0.009) (0.777,-0.002)rectangle++(0.009,0.009) (0.386,-0.003)rectangle++(0.009,0.009) (0.191,-0.004)rectangle++(0.009,0.009)
  (0.093,0.496)rectangle++(0.009,0.009) (0.535,0.196)rectangle++(0.009,0.009) (0.711,0.076)rectangle++(0.009,0.009) (0.782,0.028)rectangle++(0.009,0.009) (0.389,0.512)rectangle++(0.009,0.009)
  (0.192,0.754)rectangle++(0.009,0.009) (0.094,0.875)rectangle++(0.009,0.009) (0.045,0.935)rectangle++(0.009,0.009) (0.02,0.465)rectangle++(0.009,0.009) (0.008,0.23)rectangle++(0.009,0.009)
  (0.002,0.113)rectangle++(0.009,0.009) (-0.001,0.554)rectangle++(0.009,0.009) (-0.003,0.275)rectangle++(0.009,0.009) (0.496,0.107)rectangle++(0.009,0.009) (0.246,0.551)rectangle++(0.009,0.009)
  (0.121,0.273)rectangle++(0.009,0.009) (0.546,0.107)rectangle++(0.009,0.009) (0.716,0.04)rectangle++(0.009,0.009) (0.356,0.518)rectangle++(0.009,0.009) (0.176,0.757)rectangle++(0.009,0.009)
  (0.086,0.876)rectangle++(0.009,0.009) (0.041,0.436)rectangle++(0.009,0.009) (0.018,0.716)rectangle++(0.009,0.009) (0.505,0.284)rectangle++(0.009,0.009) (0.699,0.111)rectangle++(0.009,0.009)
  (0.347,0.553)rectangle++(0.009,0.009) (0.171,0.774)rectangle++(0.009,0.009) (0.083,0.885)rectangle++(0.009,0.009) (0.039,0.94)rectangle++(0.009,0.009) (0.017,0.968)rectangle++(0.009,0.009)
  (0.504,0.384)rectangle++(0.009,0.009) (0.699,0.151)rectangle++(0.009,0.009) (0.347,0.573)rectangle++(0.009,0.009) (0.636,0.227)rectangle++(0.009,0.009) (0.316,0.111)rectangle++(0.009,0.009)
  (0.624,0.042)rectangle++(0.009,0.009) (0.31,0.519)rectangle++(0.009,0.009) (0.153,0.757)rectangle++(0.009,0.009) (0.074,0.376)rectangle++(0.009,0.009) (0.527,0.148)rectangle++(0.009,0.009)
  (0.261,0.072)rectangle++(0.009,0.009) (0.128,0.034)rectangle++(0.009,0.009) (0.062,0.515)rectangle++(0.009,0.009) (0.522,0.203)rectangle++(0.009,0.009) (0.259,0.099)rectangle++(0.009,0.009)
  (0.127,0.047)rectangle++(0.009,0.009) (0.548,0.016)rectangle++(0.009,0.009) (0.272,0.506)rectangle++(0.009,0.009) (0.134,0.251)rectangle++(0.009,0.009) (0.065,0.123)rectangle++(0.009,0.009)
  (0.03,0.059)rectangle++(0.009,0.009) (0.013,0.027)rectangle++(0.009,0.009) (0.004,0.511)rectangle++(0.009,0.009) (0,0.253)rectangle++(0.009,0.009) (-0.002,0.624)rectangle++(0.009,0.009)
  (-0.003,0.31)rectangle++(0.009,0.009) (-0.004,0.653)rectangle++(0.009,0.009) (-0.004,0.324)rectangle++(0.009,0.009) (-0.004,0.66)rectangle++(0.009,0.009) (0.496,0.261)rectangle++(0.009,0.009)
  (0.246,0.128)rectangle++(0.009,0.009) (0.121,0.062)rectangle++(0.009,0.009) (0.058,0.529)rectangle++(0.009,0.009) (0.521,0.209)rectangle++(0.009,0.009) (0.706,0.081)rectangle++(0.009,0.009)
  (0.351,0.038)rectangle++(0.009,0.009) (0.173,0.517)rectangle++(0.009,0.009) (0.567,0.204)rectangle++(0.009,0.009) (0.281,0.1)rectangle++(0.009,0.009) (0.61,0.037)rectangle++(0.009,0.009)
  (0.303,0.016)rectangle++(0.009,0.009) (0.149,0.506)rectangle++(0.009,0.009) (0.557,0.2)rectangle++(0.009,0.009) (0.276,0.598)rectangle++(0.009,0.009) (0.136,0.297)rectangle++(0.009,0.009)
  (0.066,0.646)rectangle++(0.009,0.009) (0.524,0.256)rectangle++(0.009,0.009) (0.26,0.626)rectangle++(0.009,0.009) (0.128,0.811)rectangle++(0.009,0.009) (0.548,0.322)rectangle++(0.009,0.009)
  (0.717,0.126)rectangle++(0.009,0.009) (0.356,0.561)rectangle++(0.009,0.009) (0.176,0.278)rectangle++(0.009,0.009) (0.568,0.109)rectangle++(0.009,0.009) (0.282,0.552)rectangle++(0.009,0.009)
  (0.61,0.218)rectangle++(0.009,0.009) (0.741,0.085)rectangle++(0.009,0.009) (0.368,0.54)rectangle++(0.009,0.009) (0.182,0.768)rectangle++(0.009,0.009) (0.089,0.382)rectangle++(0.009,0.009)
  (0.042,0.189)rectangle++(0.009,0.009) (0.514,0.073)rectangle++(0.009,0.009) (0.255,0.034)rectangle++(0.009,0.009) (0.125,0.515)rectangle++(0.009,0.009) (0.547,0.203)rectangle++(0.009,0.009)
  (0.716,0.079)rectangle++(0.009,0.009) (0.356,0.537)rectangle++(0.009,0.009) (0.176,0.766)rectangle++(0.009,0.009) (0.086,0.381)rectangle++(0.009,0.009) (0.532,0.15)rectangle++(0.009,0.009)
  (0.264,0.073)rectangle++(0.009,0.009) (0.603,0.026)rectangle++(0.009,0.009) (0.299,0.011)rectangle++(0.009,0.009) (0.147,0.503)rectangle++(0.009,0.009) (0.071,0.249)rectangle++(0.009,0.009)
  (0.526,0.097)rectangle++(0.009,0.009) (0.261,0.546)rectangle++(0.009,0.009) (0.128,0.771)rectangle++(0.009,0.009) (0.062,0.383)rectangle++(0.009,0.009) (0.522,0.151)rectangle++(0.009,0.009)
  (0.259,0.573)rectangle++(0.009,0.009) (0.127,0.284)rectangle++(0.009,0.009) (0.061,0.14)rectangle++(0.009,0.009) (0.028,0.568)rectangle++(0.009,0.009) (0.012,0.282)rectangle++(0.009,0.009)
  (0.004,0.639)rectangle++(0.009,0.009) (0,0.817)rectangle++(0.009,0.009) (0.497,0.324)rectangle++(0.009,0.009) (0.246,0.16)rectangle++(0.009,0.009) (0.121,0.578)rectangle++(0.009,0.009)
  (0.058,0.787)rectangle++(0.009,0.009) (0.027,0.391)rectangle++(0.009,0.009) (0.011,0.693)rectangle++(0.009,0.009) (0.003,0.844)rectangle++(0.009,0.009) (-0.001,0.42)rectangle++(0.009,0.009)
  (-0.003,0.208)rectangle++(0.009,0.009) (-0.004,0.102)rectangle++(0.009,0.009) (-0.004,0.549)rectangle++(0.009,0.009) (-0.004,0.772)rectangle++(0.009,0.009) (-0.004,0.384)rectangle++(0.009,0.009)
  (-0.004,0.19)rectangle++(0.009,0.009) (-0.004,0.093)rectangle++(0.009,0.009) (0.496,0.034)rectangle++(0.009,0.009) (0.246,0.015)rectangle++(0.009,0.009) (0.121,0.005)rectangle++(0.009,0.009)
  (0.058,0)rectangle++(0.009,0.009) (0.027,0.498)rectangle++(0.009,0.009) (0.011,0.247)rectangle++(0.009,0.009) (0.502,0.096)rectangle++(0.009,0.009) (0.698,0.036)rectangle++(0.009,0.009)
  (0.347,0.016)rectangle++(0.009,0.009) (0.171,0.006)rectangle++(0.009,0.009) (0.083,0.001)rectangle++(0.009,0.009) (0.039,-0.002)rectangle++(0.009,0.009) (0.017,0.497)rectangle++(0.009,0.009)
  (0.006,0.746)rectangle++(0.009,0.009) (0.001,0.871)rectangle++(0.009,0.009) (-0.002,0.933)rectangle++(0.009,0.009) (-0.003,0.464)rectangle++(0.009,0.009) (-0.004,0.23)rectangle++(0.009,0.009)
  (-0.004,0.613)rectangle++(0.009,0.009) (0.496,0.242)rectangle++(0.009,0.009) (0.246,0.619)rectangle++(0.009,0.009) (0.121,0.307)rectangle++(0.009,0.009) (0.546,0.12)rectangle++(0.009,0.009)
  (0.271,0.558)rectangle++(0.009,0.009) (0.133,0.777)rectangle++(0.009,0.009) (0.064,0.386)rectangle++(0.009,0.009) (0.03,0.191)rectangle++(0.009,0.009) (0.013,0.093)rectangle++(0.009,0.009)
  (0.502,0.035)rectangle++(0.009,0.009) (0.249,0.515)rectangle++(0.009,0.009) (0.597,0.203)rectangle++(0.009,0.009) (0.296,0.099)rectangle++(0.009,0.009) (0.146,0.547)rectangle++(0.009,0.009)
  (0.071,0.771)rectangle++(0.009,0.009) (0.033,0.383)rectangle++(0.009,0.009) (0.014,0.689)rectangle++(0.009,0.009) (0.005,0.842)rectangle++(0.009,0.009) (0.499,0.334)rectangle++(0.009,0.009)
  (0.247,0.165)rectangle++(0.009,0.009) (0.121,0.08)rectangle++(0.009,0.009) (0.546,0.029)rectangle++(0.009,0.009) (0.271,0.012)rectangle++(0.009,0.009) (0.133,0.004)rectangle++(0.009,0.009)
  (0.064,0.5)rectangle++(0.009,0.009) (0.03,0.748)rectangle++(0.009,0.009) (0.013,0.872)rectangle++(0.009,0.009) (0.004,0.934)rectangle++(0.009,0.009) (0.499,0.371)rectangle++(0.009,0.009)
  (0.247,0.683)rectangle++(0.009,0.009) (0.121,0.339)rectangle++(0.009,0.009) (0.546,0.133)rectangle++(0.009,0.009) (0.716,0.051)rectangle++(0.009,0.009) (0.356,0.523)rectangle++(0.009,0.009)
  (0.176,0.259)rectangle++(0.009,0.009) (0.568,0.101)rectangle++(0.009,0.009) (0.282,0.048)rectangle++(0.009,0.009) (0.139,0.522)rectangle++(0.009,0.009) (0.067,0.759)rectangle++(0.009,0.009)
  (0.524,0.301)rectangle++(0.009,0.009) (0.26,0.148)rectangle++(0.009,0.009) (0.601,0.057)rectangle++(0.009,0.009) (0.738,0.02)rectangle++(0.009,0.009) (0.367,0.508)rectangle++(0.009,0.009)
  (0.181,0.752)rectangle++(0.009,0.009) (0.57,0.298)rectangle++(0.009,0.009) (0.725,0.116)rectangle++(0.009,0.009) (0.36,0.556)rectangle++(0.009,0.009) (0.178,0.276)rectangle++(0.009,0.009)
  (0.087,0.136)rectangle++(0.009,0.009) (0.041,0.566)rectangle++(0.009,0.009) (0.018,0.781)rectangle++(0.009,0.009) (0.007,0.888)rectangle++(0.009,0.009) (0.001,0.442)rectangle++(0.009,0.009)
  (-0.002,0.719)rectangle++(0.009,0.009) (-0.003,0.357)rectangle++(0.009,0.009) (0.496,0.14)rectangle++(0.009,0.009) (0.246,0.068)rectangle++(0.009,0.009) (0.596,0.024)rectangle++(0.009,0.009)
  (0.736,0.007)rectangle++(0.009,0.009) (0.366,0.501)rectangle++(0.009,0.009) (0.181,0.748)rectangle++(0.009,0.009) (0.088,0.372)rectangle++(0.009,0.009) (0.042,0.684)rectangle++(0.009,0.009)
  (0.514,0.271)rectangle++(0.009,0.009) (0.255,0.133)rectangle++(0.009,0.009) (0.125,0.064)rectangle++(0.009,0.009) (0.06,0.53)rectangle++(0.009,0.009) (0.028,0.763)rectangle++(0.009,0.009)
  (0.012,0.879)rectangle++(0.009,0.009) (0.004,0.437)rectangle++(0.009,0.009) (0,0.716)rectangle++(0.009,0.009) (0.497,0.284)rectangle++(0.009,0.009) (0.696,0.111)rectangle++(0.009,0.009)
  (0.776,0.042)rectangle++(0.009,0.009) (0.808,0.014)rectangle++(0.009,0.009) (0.402,0.505)rectangle++(0.009,0.009) (0.658,0.199)rectangle++(0.009,0.009) (0.327,0.597)rectangle++(0.009,0.009)
  (0.161,0.796)rectangle++(0.009,0.009) (0.078,0.896)rectangle++(0.009,0.009) (0.037,0.446)rectangle++(0.009,0.009) (0.016,0.221)rectangle++(0.009,0.009) (0.006,0.608)rectangle++(0.009,0.009)
  (0.5,0.241)rectangle++(0.009,0.009) (0.248,0.118)rectangle++(0.009,0.009) (0.122,0.557)rectangle++(0.009,0.009) (0.059,0.276)rectangle++(0.009,0.009) (0.027,0.136)rectangle++(0.009,0.009)
  (0.011,0.066)rectangle++(0.009,0.009) (0.003,0.531)rectangle++(0.009,0.009) (0.499,0.21)rectangle++(0.009,0.009) (0.247,0.103)rectangle++(0.009,0.009) (0.596,0.038)rectangle++(0.009,0.009)
  (0.296,0.017)rectangle++(0.009,0.009) (0.146,0.506)rectangle++(0.009,0.009) (0.071,0.751)rectangle++(0.009,0.009) (0.526,0.298)rectangle++(0.009,0.009) (0.261,0.647)rectangle++(0.009,0.009)
  (0.128,0.321)rectangle++(0.009,0.009) (0.549,0.126)rectangle++(0.009,0.009) (0.272,0.061)rectangle++(0.009,0.009) (0.134,0.028)rectangle++(0.009,0.009) (0.551,0.009)rectangle++(0.009,0.009)
  (0.273,0.002)rectangle++(0.009,0.009) (0.134,-0.001)rectangle++(0.009,0.009) (0.065,-0.003)rectangle++(0.009,0.009) (0.03,-0.004)rectangle++(0.009,0.009) (0.013,0.496)rectangle++(0.009,0.009)
  (0.004,0.246)rectangle++(0.009,0.009) (0,0.121)rectangle++(0.009,0.009) (-0.002,0.558)rectangle++(0.009,0.009) (-0.003,0.777)rectangle++(0.009,0.009) (-0.004,0.886)rectangle++(0.009,0.009)
  (-0.004,0.441)rectangle++(0.009,0.009) (-0.004,0.218)rectangle++(0.009,0.009) (0.496,0.085)rectangle++(0.009,0.009) (0.246,0.54)rectangle++(0.009,0.009) (0.596,0.213)rectangle++(0.009,0.009)
  (0.296,0.604)rectangle++(0.009,0.009) (0.146,0.3)rectangle++(0.009,0.009) (0.556,0.117)rectangle++(0.009,0.009) (0.276,0.556)rectangle++(0.009,0.009) (0.136,0.276)rectangle++(0.009,0.009)
  (0.552,0.108)rectangle++(0.009,0.009) (0.274,0.552)rectangle++(0.009,0.009) (0.135,0.774)rectangle++(0.009,0.009) (0.065,0.885)rectangle++(0.009,0.009) (0.03,0.94)rectangle++(0.009,0.009)
  (0.013,0.968)rectangle++(0.009,0.009) (0.502,0.384)rectangle++(0.009,0.009) (0.249,0.19)rectangle++(0.009,0.009) (0.122,0.093)rectangle++(0.009,0.009) (0.059,0.044)rectangle++(0.009,0.009)
  (0.027,0.52)rectangle++(0.009,0.009) (0.508,0.205)rectangle++(0.009,0.009) (0.701,0.079)rectangle++(0.009,0.009) (0.348,0.537)rectangle++(0.009,0.009) (0.172,0.766)rectangle++(0.009,0.009)
  (0.566,0.304)rectangle++(0.009,0.009) (0.724,0.119)rectangle++(0.009,0.009) (0.36,0.057)rectangle++(0.009,0.009) (0.178,0.026)rectangle++(0.009,0.009) (0.087,0.011)rectangle++(0.009,0.009)
  (0.041,0.503)rectangle++(0.009,0.009) (0.018,0.749)rectangle++(0.009,0.009) (0.505,0.297)rectangle++(0.009,0.009) (0.699,0.116)rectangle++(0.009,0.009) (0.347,0.056)rectangle++(0.009,0.009)
  (0.636,0.02)rectangle++(0.009,0.009) (0.316,0.008)rectangle++(0.009,0.009) (0.156,0.002)rectangle++(0.009,0.009) (0.076,0.499)rectangle++(0.009,0.009) (0.036,0.747)rectangle++(0.009,0.009)
  (0.016,0.371)rectangle++(0.009,0.009) (0.504,0.146)rectangle++(0.009,0.009) (0.25,0.571)rectangle++(0.009,0.009) (0.123,0.283)rectangle++(0.009,0.009) (0.059,0.139)rectangle++(0.009,0.009)
  (0.521,0.053)rectangle++(0.009,0.009) (0.258,0.024)rectangle++(0.009,0.009) (0.127,0.51)rectangle++(0.009,0.009) (0.061,0.253)rectangle++(0.009,0.009) (0.028,0.624)rectangle++(0.009,0.009)
  (0.012,0.81)rectangle++(0.009,0.009) (0.502,0.321)rectangle++(0.009,0.009) (0.249,0.158)rectangle++(0.009,0.009) (0.122,0.577)rectangle++(0.009,0.009) (0.059,0.786)rectangle++(0.009,0.009)
  (0.027,0.391)rectangle++(0.009,0.009) (0.011,0.193)rectangle++(0.009,0.009) (0.003,0.094)rectangle++(0.009,0.009) (-0.001,0.545)rectangle++(0.009,0.009) (0.497,0.215)rectangle++(0.009,0.009)
  (0.246,0.105)rectangle++(0.009,0.009) (0.121,0.55)rectangle++(0.009,0.009) (0.058,0.273)rectangle++(0.009,0.009) (0.027,0.134)rectangle++(0.009,0.009) (0.011,0.565)rectangle++(0.009,0.009)
  (0.003,0.78)rectangle++(0.009,0.009) (-0.001,0.888)rectangle++(0.009,0.009) (0.497,0.352)rectangle++(0.009,0.009) (0.246,0.674)rectangle++(0.009,0.009) (0.121,0.835)rectangle++(0.009,0.009)
  (0.058,0.415)rectangle++(0.009,0.009) (0.027,0.205)rectangle++(0.009,0.009) (0.011,0.1)rectangle++(0.009,0.009) (0.003,0.548)rectangle++(0.009,0.009) (-0.001,0.272)rectangle++(0.009,0.009)
  (-0.003,0.634)rectangle++(0.009,0.009) (-0.004,0.315)rectangle++(0.009,0.009) (-0.004,0.655)rectangle++(0.009,0.009) (0.496,0.259)rectangle++(0.009,0.009) (0.696,0.101)rectangle++(0.009,0.009)
  (0.346,0.548)rectangle++(0.009,0.009) (0.171,0.772)rectangle++(0.009,0.009) (0.566,0.306)rectangle++(0.009,0.009) (0.281,0.151)rectangle++(0.009,0.009) (0.61,0.058)rectangle++(0.009,0.009)
  (0.741,0.02)rectangle++(0.009,0.009) (0.368,0.008)rectangle++(0.009,0.009) (0.182,0.502)rectangle++(0.009,0.009) (0.089,0.249)rectangle++(0.009,0.009) (0.042,0.622)rectangle++(0.009,0.009)
  (0.019,0.309)rectangle++(0.009,0.009) (0.007,0.652)rectangle++(0.009,0.009) (0.001,0.824)rectangle++(0.009,0.009) (0.498,0.327)rectangle++(0.009,0.009) (0.247,0.661)rectangle++(0.009,0.009)
  (0.121,0.328)rectangle++(0.009,0.009) (0.058,0.162)rectangle++(0.009,0.009) (0.521,0.062)rectangle++(0.009,0.009) (0.258,0.029)rectangle++(0.009,0.009) (0.601,0.009)rectangle++(0.009,0.009)
  (0.298,0.502)rectangle++(0.009,0.009) (0.617,0.198)rectangle++(0.009,0.009) (0.744,0.077)rectangle++(0.009,0.009) (0.37,0.536)rectangle++(0.009,0.009) (0.183,0.766)rectangle++(0.009,0.009)
  (0.089,0.381)rectangle++(0.009,0.009) (0.042,0.688)rectangle++(0.009,0.009) (0.019,0.842)rectangle++(0.009,0.009) (0.007,0.419)rectangle++(0.009,0.009) (0.001,0.207)rectangle++(0.009,0.009)
  (-0.002,0.101)rectangle++(0.009,0.009) (-0.003,0.048)rectangle++(0.009,0.009) (0.496,0.017)rectangle++(0.009,0.009) (0.246,0.506)rectangle++(0.009,0.009) (0.596,0.2)rectangle++(0.009,0.009)
  (0.296,0.098)rectangle++(0.009,0.009) (0.146,0.047)rectangle++(0.009,0.009) (0.071,0.521)rectangle++(0.009,0.009) (0.526,0.206)rectangle++(0.009,0.009) (0.261,0.101)rectangle++(0.009,0.009)
  (0.128,0.548)rectangle++(0.009,0.009) (0.549,0.217)rectangle++(0.009,0.009) (0.272,0.606)rectangle++(0.009,0.009) (0.606,0.24)rectangle++(0.009,0.009) (0.301,0.118)rectangle++(0.009,0.009)
  (0.148,0.057)rectangle++(0.009,0.009) (0.072,0.026)rectangle++(0.009,0.009) (0.034,0.511)rectangle++(0.009,0.009) (0.015,0.253)rectangle++(0.009,0.009) (0.005,0.624)rectangle++(0.009,0.009)
  (0,0.81)rectangle++(0.009,0.009) (0.497,0.321)rectangle++(0.009,0.009) (0.696,0.126)rectangle++(0.009,0.009) (0.346,0.561)rectangle++(0.009,0.009) (0.171,0.778)rectangle++(0.009,0.009)
  (0.083,0.387)rectangle++(0.009,0.009) (0.039,0.191)rectangle++(0.009,0.009) (0.017,0.093)rectangle++(0.009,0.009) (0.006,0.044)rectangle++(0.009,0.009) (0.001,0.52)rectangle++(0.009,0.009)
  (-0.002,0.758)rectangle++(0.009,0.009) (-0.003,0.877)rectangle++(0.009,0.009) (0.496,0.348)rectangle++(0.009,0.009) (0.246,0.672)rectangle++(0.009,0.009) (0.596,0.266)rectangle++(0.009,0.009)
  (0.296,0.631)rectangle++(0.009,0.009) (0.616,0.25)rectangle++(0.009,0.009) (0.306,0.623)rectangle++(0.009,0.009) (0.62,0.246)rectangle++(0.009,0.009) (0.308,0.121)rectangle++(0.009,0.009)
  (0.152,0.058)rectangle++(0.009,0.009) (0.074,0.527)rectangle++(0.009,0.009) (0.035,0.761)rectangle++(0.009,0.009) (0.511,0.302)rectangle++(0.009,0.009) (0.253,0.649)rectangle++(0.009,0.009)
  (0.124,0.322)rectangle++(0.009,0.009) (0.06,0.659)rectangle++(0.009,0.009) (0.521,0.261)rectangle++(0.009,0.009) (0.258,0.128)rectangle++(0.009,0.009) (0.127,0.062)rectangle++(0.009,0.009)
  (0.061,0.029)rectangle++(0.009,0.009) (0.028,0.012)rectangle++(0.009,0.009) (0.012,0.504)rectangle++(0.009,0.009) (0.502,0.199)rectangle++(0.009,0.009) (0.249,0.097)rectangle++(0.009,0.009)
  (0.597,0.036)rectangle++(0.009,0.009) (0.296,0.516)rectangle++(0.009,0.009) (0.146,0.756)rectangle++(0.009,0.009) (0.071,0.376)rectangle++(0.009,0.009) (0.526,0.148)rectangle++(0.009,0.009)
  (0.261,0.072)rectangle++(0.009,0.009) (0.128,0.034)rectangle++(0.009,0.009) (0.062,0.515)rectangle++(0.009,0.009) (0.029,0.255)rectangle++(0.009,0.009) (0.509,0.099)rectangle++(0.009,0.009)
  (0.252,0.547)rectangle++(0.009,0.009) (0.124,0.771)rectangle++(0.009,0.009) (0.06,0.883)rectangle++(0.009,0.009) (0.028,0.939)rectangle++(0.009,0.009) (0.012,0.967)rectangle++(0.009,0.009)
  (0.004,0.481)rectangle++(0.009,0.009) (0,0.738)rectangle++(0.009,0.009) (-0.002,0.867)rectangle++(0.009,0.009) (0.496,0.344)rectangle++(0.009,0.009) (0.246,0.17)rectangle++(0.009,0.009)
  (0.596,0.065)rectangle++(0.009,0.009) (0.736,0.023)rectangle++(0.009,0.009) (0.366,0.509)rectangle++(0.009,0.009) (0.181,0.252)rectangle++(0.009,0.009) (0.088,0.624)rectangle++(0.009,0.009)
  (0.042,0.81)rectangle++(0.009,0.009) (0.514,0.321)rectangle++(0.009,0.009) (0.255,0.158)rectangle++(0.009,0.009) (0.125,0.577)rectangle++(0.009,0.009) (0.06,0.286)rectangle++(0.009,0.009)
  (0.028,0.641)rectangle++(0.009,0.009) (0.012,0.818)rectangle++(0.009,0.009) (0.502,0.325)rectangle++(0.009,0.009) (0.249,0.16)rectangle++(0.009,0.009) (0.122,0.078)rectangle++(0.009,0.009)
  (0.059,0.037)rectangle++(0.009,0.009) (0.027,0.016)rectangle++(0.009,0.009) (0.011,0.006)rectangle++(0.009,0.009) (0.502,0)rectangle++(0.009,0.009) (0.249,0.498)rectangle++(0.009,0.009)
  (0.122,0.247)rectangle++(0.009,0.009) (0.059,0.621)rectangle++(0.009,0.009) (0.027,0.308)rectangle++(0.009,0.009) (0.011,0.652)rectangle++(0.009,0.009) (0.003,0.324)rectangle++(0.009,0.009)
  (-0.001,0.16)rectangle++(0.009,0.009) (-0.003,0.578)rectangle++(0.009,0.009) (-0.004,0.787)rectangle++(0.009,0.009) (-0.004,0.391)rectangle++(0.009,0.009) (0.496,0.154)rectangle++(0.009,0.009)
  (0.246,0.575)rectangle++(0.009,0.009) (0.121,0.285)rectangle++(0.009,0.009) (0.546,0.111)rectangle++(0.009,0.009) (0.271,0.053)rectangle++(0.009,0.009) (0.133,0.024)rectangle++(0.009,0.009)
  (0.551,0.007)rectangle++(0.009,0.009) (0.273,0.001)rectangle++(0.009,0.009) (0.134,-0.002)rectangle++(0.009,0.009) (0.065,0.497)rectangle++(0.009,0.009) (0.03,0.746)rectangle++(0.009,0.009)
  (0.509,0.296)rectangle++(0.009,0.009) (0.252,0.146)rectangle++(0.009,0.009) (0.124,0.071)rectangle++(0.009,0.009) (0.06,0.533)rectangle++(0.009,0.009) (0.521,0.211)rectangle++(0.009,0.009)
  (0.258,0.103)rectangle++(0.009,0.009) (0.127,0.549)rectangle++(0.009,0.009) (0.061,0.272)rectangle++(0.009,0.009) (0.028,0.134)rectangle++(0.009,0.009) (0.012,0.565)rectangle++(0.009,0.009)
  (0.004,0.78)rectangle++(0.009,0.009) (0,0.388)rectangle++(0.009,0.009) (0.497,0.152)rectangle++(0.009,0.009) (0.246,0.074)rectangle++(0.009,0.009) (0.596,0.027)rectangle++(0.009,0.009)
  (0.296,0.511)rectangle++(0.009,0.009) (0.146,0.753)rectangle++(0.009,0.009) (0.071,0.374)rectangle++(0.009,0.009) (0.033,0.685)rectangle++(0.009,0.009) (0.014,0.84)rectangle++(0.009,0.009)
  (0.005,0.418)rectangle++(0.009,0.009) (0,0.707)rectangle++(0.009,0.009) (-0.002,0.351)rectangle++(0.009,0.009) (-0.003,0.673)rectangle++(0.009,0.009) (-0.004,0.834)rectangle++(0.009,0.009)
  (-0.004,0.915)rectangle++(0.009,0.009) (0.496,0.363)rectangle++(0.009,0.009) (0.246,0.679)rectangle++(0.009,0.009) (0.121,0.837)rectangle++(0.009,0.009) (0.546,0.332)rectangle++(0.009,0.009)
  (0.271,0.164)rectangle++(0.009,0.009) (0.133,0.58)rectangle++(0.009,0.009) (0.064,0.788)rectangle++(0.009,0.009) (0.03,0.892)rectangle++(0.009,0.009) (0.013,0.944)rectangle++(0.009,0.009)
  (0.004,0.97)rectangle++(0.009,0.009) (0,0.983)rectangle++(0.009,0.009) (-0.002,0.489)rectangle++(0.009,0.009) (0.496,0.193)rectangle++(0.009,0.009) (0.246,0.094)rectangle++(0.009,0.009)
  (0.121,0.045)rectangle++(0.009,0.009) (0.058,0.52)rectangle++(0.009,0.009) (0.521,0.205)rectangle++(0.009,0.009) (0.706,0.079)rectangle++(0.009,0.009) (0.351,0.537)rectangle++(0.009,0.009)
  (0.173,0.766)rectangle++(0.009,0.009) (0.084,0.381)rectangle++(0.009,0.009) (0.531,0.15)rectangle++(0.009,0.009) (0.71,0.057)rectangle++(0.009,0.009) (0.353,0.026)rectangle++(0.009,0.009)
  (0.174,0.511)rectangle++(0.009,0.009) (0.567,0.202)rectangle++(0.009,0.009) (0.281,0.599)rectangle++(0.009,0.009) (0.138,0.297)rectangle++(0.009,0.009) (0.067,0.146)rectangle++(0.009,0.009)
  (0.031,0.071)rectangle++(0.009,0.009) (0.51,0.026)rectangle++(0.009,0.009) (0.253,0.011)rectangle++(0.009,0.009) (0.598,0.002)rectangle++(0.009,0.009) (0.297,-0.001)rectangle++(0.009,0.009)
  (0.146,-0.003)rectangle++(0.009,0.009) (0.071,0.496)rectangle++(0.009,0.009) (0.526,0.196)rectangle++(0.009,0.009) (0.261,0.596)rectangle++(0.009,0.009) (0.602,0.236)rectangle++(0.009,0.009)
  (0.299,0.616)rectangle++(0.009,0.009) (0.617,0.244)rectangle++(0.009,0.009) (0.744,0.095)rectangle++(0.009,0.009) (0.37,0.545)rectangle++(0.009,0.009) (0.183,0.27)rectangle++(0.009,0.009)
  (0.089,0.633)rectangle++(0.009,0.009) (0.533,0.25)rectangle++(0.009,0.009) (0.264,0.623)rectangle++(0.009,0.009) (0.13,0.309)rectangle++(0.009,0.009) (0.549,0.121)rectangle++(0.009,0.009)
  (0.272,0.058)rectangle++(0.009,0.009) (0.134,0.527)rectangle++(0.009,0.009) (0.065,0.761)rectangle++(0.009,0.009) (0.03,0.878)rectangle++(0.009,0.009) (0.509,0.349)rectangle++(0.009,0.009)
  (0.252,0.172)rectangle++(0.009,0.009) (0.124,0.584)rectangle++(0.009,0.009) (0.06,0.29)rectangle++(0.009,0.009) (0.028,0.643)rectangle++(0.009,0.009) (0.508,0.254)rectangle++(0.009,0.009)
  (0.252,0.625)rectangle++(0.009,0.009) (0.124,0.81)rectangle++(0.009,0.009) (0.06,0.403)rectangle++(0.009,0.009) (0.521,0.158)rectangle++(0.009,0.009) (0.706,0.061)rectangle++(0.009,0.009)
  (0.351,0.528)rectangle++(0.009,0.009) (0.173,0.762)rectangle++(0.009,0.009) (0.084,0.879)rectangle++(0.009,0.009) (0.531,0.349)rectangle++(0.009,0.009) (0.263,0.172)rectangle++(0.009,0.009)
  (0.129,0.084)rectangle++(0.009,0.009) (0.549,0.031)rectangle++(0.009,0.009) (0.717,0.01)rectangle++(0.009,0.009) (0.356,0.503)rectangle++(0.009,0.009) (0.176,0.249)rectangle++(0.009,0.009)
  (0.086,0.122)rectangle++(0.009,0.009) (0.041,0.559)rectangle++(0.009,0.009) (0.018,0.777)rectangle++(0.009,0.009) (0.007,0.886)rectangle++(0.009,0.009) (0.001,0.441)rectangle++(0.009,0.009)
  (-0.002,0.718)rectangle++(0.009,0.009) (0.497,0.285)rectangle++(0.009,0.009) (0.696,0.111)rectangle++(0.009,0.009) (0.346,0.553)rectangle++(0.009,0.009) (0.636,0.219)rectangle++(0.009,0.009)
  (0.752,0.085)rectangle++(0.009,0.009) (0.374,0.54)rectangle++(0.009,0.009) (0.185,0.268)rectangle++(0.009,0.009) (0.571,0.104)rectangle++(0.009,0.009) (0.283,0.55)rectangle++(0.009,0.009)
  (0.139,0.773)rectangle++(0.009,0.009) (0.067,0.384)rectangle++(0.009,0.009) (0.031,0.69)rectangle++(0.009,0.009) (0.013,0.843)rectangle++(0.009,0.009) (0.503,0.334)rectangle++(0.009,0.009)
  (0.698,0.131)rectangle++(0.009,0.009) (0.347,0.063)rectangle++(0.009,0.009) (0.171,0.529)rectangle++(0.009,0.009) (0.083,0.262)rectangle++(0.009,0.009) (0.039,0.629)rectangle++(0.009,0.009)
  (0.017,0.812)rectangle++(0.009,0.009) (0.006,0.904)rectangle++(0.009,0.009) (0.001,0.95)rectangle++(0.009,0.009) (0.498,0.377)rectangle++(0.009,0.009) (0.247,0.686)rectangle++(0.009,0.009)
  (0.121,0.841)rectangle++(0.009,0.009) (0.546,0.334)rectangle++(0.009,0.009) (0.271,0.165)rectangle++(0.009,0.009) (0.133,0.08)rectangle++(0.009,0.009) (0.064,0.038)rectangle++(0.009,0.009)
  (0.523,0.012)rectangle++(0.009,0.009) (0.707,0.002)rectangle++(0.009,0.009) (0.351,0.499)rectangle++(0.009,0.009) (0.173,0.247)rectangle++(0.009,0.009) (0.567,0.096)rectangle++(0.009,0.009)
  (0.724,0.036)rectangle++(0.009,0.009) (0.36,0.516)rectangle++(0.009,0.009) (0.178,0.756)rectangle++(0.009,0.009) (0.087,0.876)rectangle++(0.009,0.009) (0.041,0.936)rectangle++(0.009,0.009)
  (0.514,0.372)rectangle++(0.009,0.009) (0.255,0.684)rectangle++(0.009,0.009) (0.125,0.34)rectangle++(0.009,0.009) (0.06,0.668)rectangle++(0.009,0.009) (0.521,0.264)rectangle++(0.009,0.009)
  (0.258,0.13)rectangle++(0.009,0.009) (0.127,0.063)rectangle++(0.009,0.009) (0.548,0.022)rectangle++(0.009,0.009) (0.272,0.009)rectangle++(0.009,0.009) (0.134,0.502)rectangle++(0.009,0.009)
  (0.551,0.198)rectangle++(0.009,0.009) (0.273,0.097)rectangle++(0.009,0.009) (0.134,0.046)rectangle++(0.009,0.009) (0.551,0.016)rectangle++(0.009,0.009) (0.273,0.506)rectangle++(0.009,0.009)
  (0.134,0.751)rectangle++(0.009,0.009) (0.551,0.298)rectangle++(0.009,0.009) (0.273,0.647)rectangle++(0.009,0.009) (0.134,0.321)rectangle++(0.009,0.009) (0.065,0.158)rectangle++(0.009,0.009)
  (0.03,0.077)rectangle++(0.009,0.009) (0.013,0.036)rectangle++(0.009,0.009) (0.004,0.016)rectangle++(0.009,0.009) (0.499,0.004)rectangle++(0.009,0.009) (0.247,0.5)rectangle++(0.009,0.009)
  (0.121,0.248)rectangle++(0.009,0.009) (0.058,0.122)rectangle++(0.009,0.009) (0.521,0.046)rectangle++(0.009,0.009) (0.258,0.521)rectangle++(0.009,0.009) (0.601,0.206)rectangle++(0.009,0.009)
  (0.298,0.601)rectangle++(0.009,0.009) (0.147,0.298)rectangle++(0.009,0.009) (0.556,0.117)rectangle++(0.009,0.009) (0.276,0.556)rectangle++(0.009,0.009) (0.136,0.776)rectangle++(0.009,0.009)
  (0.066,0.886)rectangle++(0.009,0.009) (0.031,0.441)rectangle++(0.009,0.009) (0.013,0.218)rectangle++(0.009,0.009) (0.004,0.107)rectangle++(0.009,0.009) (0,0.551)rectangle++(0.009,0.009)
  (0.497,0.218)rectangle++(0.009,0.009) (0.696,0.084)rectangle++(0.009,0.009) (0.346,0.04)rectangle++(0.009,0.009) (0.171,0.518)rectangle++(0.009,0.009) (0.083,0.757)rectangle++(0.009,0.009)
  (0.039,0.376)rectangle++(0.009,0.009) (0.017,0.686)rectangle++(0.009,0.009) (0.504,0.272)rectangle++(0.009,0.009) (0.25,0.134)rectangle++(0.009,0.009) (0.123,0.065)rectangle++(0.009,0.009)
  (0.059,0.53)rectangle++(0.009,0.009) (0.027,0.763)rectangle++(0.009,0.009) (0.508,0.302)rectangle++(0.009,0.009) (0.252,0.649)rectangle++(0.009,0.009) (0.124,0.322)rectangle++(0.009,0.009)
  (0.06,0.659)rectangle++(0.009,0.009) (0.028,0.327)rectangle++(0.009,0.009) (0.012,0.161)rectangle++(0.009,0.009) (0.004,0.078)rectangle++(0.009,0.009) (0,0.037)rectangle++(0.009,0.009)
  (-0.002,0.516)rectangle++(0.009,0.009) (0.496,0.204)rectangle++(0.009,0.009) (0.696,0.079)rectangle++(0.009,0.009) (0.346,0.537)rectangle++(0.009,0.009) (0.636,0.212)rectangle++(0.009,0.009)
  (0.316,0.104)rectangle++(0.009,0.009) (0.156,0.05)rectangle++(0.009,0.009) (0.56,0.017)rectangle++(0.009,0.009) (0.278,0.506)rectangle++(0.009,0.009) (0.137,0.751)rectangle++(0.009,0.009)
  (0.066,0.373)rectangle++(0.009,0.009) (0.031,0.184)rectangle++(0.009,0.009) (0.013,0.59)rectangle++(0.009,0.009) (0.503,0.233)rectangle++(0.009,0.009) (0.249,0.614)rectangle++(0.009,0.009)
  (0.122,0.805)rectangle++(0.009,0.009) (0.059,0.4)rectangle++(0.009,0.009) (0.027,0.698)rectangle++(0.009,0.009) (0.011,0.347)rectangle++(0.009,0.009) (0.003,0.671)rectangle++(0.009,0.009)
  (-0.001,0.333)rectangle++(0.009,0.009) (-0.003,0.164)rectangle++(0.009,0.009) (-0.004,0.58)rectangle++(0.009,0.009) (-0.004,0.788)rectangle++(0.009,0.009) (-0.004,0.892)rectangle++(0.009,0.009)
  (0.496,0.354)rectangle++(0.009,0.009) (0.246,0.175)rectangle++(0.009,0.009) (0.121,0.085)rectangle++(0.009,0.009) (0.058,0.54)rectangle++(0.009,0.009) (0.027,0.768)rectangle++(0.009,0.009)
  (0.011,0.882)rectangle++(0.009,0.009) (0.003,0.439)rectangle++(0.009,0.009) (-0.001,0.217)rectangle++(0.009,0.009) (0.497,0.084)rectangle++(0.009,0.009) (0.246,0.04)rectangle++(0.009,0.009)
  (0.596,0.013)rectangle++(0.009,0.009) (0.296,0.004)rectangle++(0.009,0.009) (0.146,0.5)rectangle++(0.009,0.009) (0.071,0.248)rectangle++(0.009,0.009) (0.526,0.096)rectangle++(0.009,0.009)
  (0.708,0.036)rectangle++(0.009,0.009) (0.352,0.016)rectangle++(0.009,0.009) (0.174,0.006)rectangle++(0.009,0.009) (0.085,0.001)rectangle++(0.009,0.009) (0.04,-0.002)rectangle++(0.009,0.009)
  (0.513,-0.003)rectangle++(0.009,0.009) (0.703,-0.004)rectangle++(0.009,0.009) (0.349,0.496)rectangle++(0.009,0.009) (0.172,0.746)rectangle++(0.009,0.009) (0.566,0.296)rectangle++(0.009,0.009)
  (0.281,0.146)rectangle++(0.009,0.009) (0.61,0.056)rectangle++(0.009,0.009) (0.303,0.026)rectangle++(0.009,0.009) (0.149,0.511)rectangle++(0.009,0.009) (0.072,0.253)rectangle++(0.009,0.009)
  (0.034,0.124)rectangle++(0.009,0.009) (0.015,0.56)rectangle++(0.009,0.009) (0.503,0.221)rectangle++(0.009,0.009) (0.249,0.108)rectangle++(0.009,0.009) (0.122,0.052)rectangle++(0.009,0.009)
  (0.059,0.024)rectangle++(0.009,0.009) (0.027,0.51)rectangle++(0.009,0.009) (0.011,0.253)rectangle++(0.009,0.009) (0.003,0.624)rectangle++(0.009,0.009) (-0.001,0.31)rectangle++(0.009,0.009)
  (-0.003,0.653)rectangle++(0.009,0.009) (-0.004,0.324)rectangle++(0.009,0.009) (-0.004,0.16)rectangle++(0.009,0.009) (-0.004,0.078)rectangle++(0.009,0.009) (-0.004,0.537)rectangle++(0.009,0.009)
  (-0.004,0.266)rectangle++(0.009,0.009) (0.496,0.104)rectangle++(0.009,0.009) (0.246,0.55)rectangle++(0.009,0.009) (0.121,0.773)rectangle++(0.009,0.009) (0.058,0.384)rectangle++(0.009,0.009)
  (0.027,0.69)rectangle++(0.009,0.009) (0.011,0.343)rectangle++(0.009,0.009) (0.003,0.169)rectangle++(0.009,0.009) (-0.001,0.082)rectangle++(0.009,0.009) (-0.003,0.039)rectangle++(0.009,0.009)
  (0.496,0.013)rectangle++(0.009,0.009) (0.246,0.504)rectangle++(0.009,0.009) (0.121,0.75)rectangle++(0.009,0.009) (0.058,0.873)rectangle++(0.009,0.009) (0.521,0.346)rectangle++(0.009,0.009)
  (0.258,0.671)rectangle++(0.009,0.009) (0.127,0.333)rectangle++(0.009,0.009) (0.061,0.664)rectangle++(0.009,0.009) (0.028,0.83)rectangle++(0.009,0.009) (0.509,0.329)rectangle++(0.009,0.009)
  (0.252,0.662)rectangle++(0.009,0.009) (0.598,0.262)rectangle++(0.009,0.009) (0.297,0.629)rectangle++(0.009,0.009) (0.146,0.812)rectangle++(0.009,0.009) (0.071,0.904)rectangle++(0.009,0.009)
  (0.033,0.95)rectangle++(0.009,0.009) (0.511,0.377)rectangle++(0.009,0.009) (0.253,0.186)rectangle++(0.009,0.009) (0.124,0.091)rectangle++(0.009,0.009) (0.06,0.043)rectangle++(0.009,0.009)
  (0.028,0.519)rectangle++(0.009,0.009) (0.012,0.257)rectangle++(0.009,0.009) (0.004,0.626)rectangle++(0.009,0.009) (0.499,0.248)rectangle++(0.009,0.009) (0.247,0.622)rectangle++(0.009,0.009)
  (0.596,0.246)rectangle++(0.009,0.009) (0.296,0.621)rectangle++(0.009,0.009) (0.146,0.308)rectangle++(0.009,0.009) (0.071,0.152)rectangle++(0.009,0.009) (0.033,0.574)rectangle++(0.009,0.009)
  (0.014,0.785)rectangle++(0.009,0.009) (0.005,0.39)rectangle++(0.009,0.009) (0.499,0.153)rectangle++(0.009,0.009) (0.247,0.574)rectangle++(0.009,0.009) (0.121,0.285)rectangle++(0.009,0.009)
  (0.058,0.64)rectangle++(0.009,0.009) (0.027,0.818)rectangle++(0.009,0.009) (0.011,0.907)rectangle++(0.009,0.009) (0.003,0.951)rectangle++(0.009,0.009) (-0.001,0.973)rectangle++(0.009,0.009)
  (-0.003,0.484)rectangle++(0.009,0.009) (-0.004,0.74)rectangle++(0.009,0.009) (-0.004,0.368)rectangle++(0.009,0.009) (-0.004,0.182)rectangle++(0.009,0.009) (-0.004,0.589)rectangle++(0.009,0.009)
  (0.496,0.233)rectangle++(0.009,0.009) (0.696,0.09)rectangle++(0.009,0.009) (0.346,0.543)rectangle++(0.009,0.009) (0.171,0.769)rectangle++(0.009,0.009) (0.083,0.382)rectangle++(0.009,0.009)
  (0.039,0.689)rectangle++(0.009,0.009) (0.017,0.842)rectangle++(0.009,0.009) (0.006,0.919)rectangle++(0.009,0.009) (0.001,0.457)rectangle++(0.009,0.009) (0.498,0.18)rectangle++(0.009,0.009)
  (0.696,0.069)rectangle++(0.009,0.009) (0.346,0.032)rectangle++(0.009,0.009) (0.171,0.514)rectangle++(0.009,0.009) (0.083,0.255)rectangle++(0.009,0.009) (0.039,0.125)rectangle++(0.009,0.009)
  (0.017,0.56)rectangle++(0.009,0.009) (0.006,0.278)rectangle++(0.009,0.009) (0.001,0.137)rectangle++(0.009,0.009) (-0.002,0.566)rectangle++(0.009,0.009) (-0.003,0.781)rectangle++(0.009,0.009)
  (0.496,0.31)rectangle++(0.009,0.009) (0.696,0.121)rectangle++(0.009,0.009) (0.346,0.058)rectangle++(0.009,0.009) (0.171,0.027)rectangle++(0.009,0.009) (0.083,0.511)rectangle++(0.009,0.009)
  (0.039,0.753)rectangle++(0.009,0.009) (0.017,0.874)rectangle++(0.009,0.009) (0.006,0.935)rectangle++(0.009,0.009) (0.5,0.371)rectangle++(0.009,0.009) (0.248,0.183)rectangle++(0.009,0.009)
  (0.122,0.589)rectangle++(0.009,0.009) (0.059,0.792)rectangle++(0.009,0.009) (0.027,0.394)rectangle++(0.009,0.009) (0.011,0.195)rectangle++(0.009,0.009) (0.003,0.095)rectangle++(0.009,0.009)
  (-0.001,0.545)rectangle++(0.009,0.009) (0.497,0.215)rectangle++(0.009,0.009) (0.246,0.105)rectangle++(0.009,0.009) (0.596,0.039)rectangle++(0.009,0.009) (0.296,0.517)rectangle++(0.009,0.009)
  (0.146,0.756)rectangle++(0.009,0.009) (0.071,0.876)rectangle++(0.009,0.009) (0.033,0.936)rectangle++(0.009,0.009) (0.014,0.466)rectangle++(0.009,0.009) (0.503,0.184)rectangle++(0.009,0.009)
  (0.699,0.071)rectangle++(0.009,0.009) (0.347,0.533)rectangle++(0.009,0.009) (0.171,0.264)rectangle++(0.009,0.009) (0.083,0.13)rectangle++(0.009,0.009) (0.039,0.063)rectangle++(0.009,0.009)
  (0.017,0.029)rectangle++(0.009,0.009) (0.006,0.012)rectangle++(0.009,0.009) (0.001,0.504)rectangle++(0.009,0.009) (-0.002,0.75)rectangle++(0.009,0.009) (-0.003,0.873)rectangle++(0.009,0.009)
  (-0.004,0.434)rectangle++(0.009,0.009) (0.496,0.171)rectangle++(0.009,0.009) (0.696,0.066)rectangle++(0.009,0.009) (0.346,0.031)rectangle++(0.009,0.009) (0.171,0.013)rectangle++(0.009,0.009)
  (0.083,0.004)rectangle++(0.009,0.009) (0.039,0)rectangle++(0.009,0.009) (0.017,-0.002)rectangle++(0.009,0.009) (0.006,0.497)rectangle++(0.009,0.009) (0.001,0.746)rectangle++(0.009,0.009)
  (-0.002,0.871)rectangle++(0.009,0.009) (-0.003,0.433)rectangle++(0.009,0.009) (-0.004,0.214)rectangle++(0.009,0.009) (-0.004,0.605)rectangle++(0.009,0.009) (-0.004,0.8)rectangle++(0.009,0.009)
  (0.496,0.317)rectangle++(0.009,0.009) (0.246,0.156)rectangle++(0.009,0.009) (0.596,0.06)rectangle++(0.009,0.009) (0.296,0.028)rectangle++(0.009,0.009) (0.146,0.012)rectangle++(0.009,0.009)
  (0.071,0.004)rectangle++(0.009,0.009) (0.033,0)rectangle++(0.009,0.009) (0.014,0.498)rectangle++(0.009,0.009) (0.005,0.747)rectangle++(0.009,0.009) (0,0.871)rectangle++(0.009,0.009)
  (0.497,0.346)rectangle++(0.009,0.009) (0.696,0.136)rectangle++(0.009,0.009) (0.346,0.066)rectangle++(0.009,0.009) (0.636,0.024)rectangle++(0.009,0.009) (0.316,0.51)rectangle++(0.009,0.009)
  (0.156,0.753)rectangle++(0.009,0.009) (0.076,0.374)rectangle++(0.009,0.009) (0.036,0.185)rectangle++(0.009,0.009) (0.016,0.59)rectangle++(0.009,0.009) (0.504,0.233)rectangle++(0.009,0.009)
  (0.699,0.091)rectangle++(0.009,0.009) (0.347,0.043)rectangle++(0.009,0.009) (0.171,0.019)rectangle++(0.009,0.009) (0.566,0.005)rectangle++(0.009,0.009) (0.281,0.5)rectangle++(0.009,0.009)
  (0.138,0.748)rectangle++(0.009,0.009) (0.067,0.372)rectangle++(0.009,0.009) (0.524,0.146)rectangle++(0.009,0.009) (0.26,0.071)rectangle++(0.009,0.009) (0.601,0.026)rectangle++(0.009,0.009)
  (0.738,0.008)rectangle++(0.009,0.009) (0.367,0.502)rectangle++(0.009,0.009) (0.644,0.198)rectangle++(0.009,0.009) (0.32,0.597)rectangle++(0.009,0.009) (0.625,0.236)rectangle++(0.009,0.009)
  (0.31,0.116)rectangle++(0.009,0.009) (0.153,0.056)rectangle++(0.009,0.009) (0.558,0.02)rectangle++(0.009,0.009) (0.277,0.008)rectangle++(0.009,0.009) (0.608,0)rectangle++(0.009,0.009)
  (0.302,0.498)rectangle++(0.009,0.009) (0.149,0.747)rectangle++(0.009,0.009) (0.072,0.871)rectangle++(0.009,0.009) (0.526,0.346)rectangle++(0.009,0.009) (0.261,0.171)rectangle++(0.009,0.009)
  (0.602,0.066)rectangle++(0.009,0.009) (0.299,0.531)rectangle++(0.009,0.009) (0.147,0.263)rectangle++(0.009,0.009) (0.071,0.629)rectangle++(0.009,0.009) (0.033,0.312)rectangle++(0.009,0.009)
  (0.014,0.654)rectangle++(0.009,0.009) (0.005,0.825)rectangle++(0.009,0.009) (0,0.91)rectangle++(0.009,0.009) (-0.002,0.953)rectangle++(0.009,0.009) (-0.003,0.474)rectangle++(0.009,0.009)
  (0.496,0.187)rectangle++(0.009,0.009) (0.246,0.591)rectangle++(0.009,0.009) (0.596,0.234)rectangle++(0.009,0.009) (0.296,0.115)rectangle++(0.009,0.009) (0.146,0.055)rectangle++(0.009,0.009)
  (0.071,0.025)rectangle++(0.009,0.009) (0.033,0.51)rectangle++(0.009,0.009) (0.014,0.253)rectangle++(0.009,0.009) (0.503,0.098)rectangle++(0.009,0.009) (0.249,0.047)rectangle++(0.009,0.009)
  (0.122,0.521)rectangle++(0.009,0.009) (0.059,0.258)rectangle++(0.009,0.009) (0.027,0.127)rectangle++(0.009,0.009) (0.011,0.561)rectangle++(0.009,0.009) (0.502,0.222)rectangle++(0.009,0.009)
  (0.698,0.086)rectangle++(0.009,0.009) (0.347,0.041)rectangle++(0.009,0.009) (0.171,0.018)rectangle++(0.009,0.009) (0.083,0.007)rectangle++(0.009,0.009) (0.531,0)rectangle++(0.009,0.009)
  (0.263,0.498)rectangle++(0.009,0.009) (0.129,0.747)rectangle++(0.009,0.009) (0.549,0.296)rectangle++(0.009,0.009) (0.717,0.116)rectangle++(0.009,0.009) (0.356,0.556)rectangle++(0.009,0.009)
  (0.176,0.776)rectangle++(0.009,0.009) (0.086,0.886)rectangle++(0.009,0.009) (0.041,0.441)rectangle++(0.009,0.009) (0.514,0.174)rectangle++(0.009,0.009) (0.255,0.085)rectangle++(0.009,0.009)
  (0.125,0.04)rectangle++(0.009,0.009) (0.06,0.518)rectangle++(0.009,0.009) (0.521,0.204)rectangle++(0.009,0.009) (0.258,0.1)rectangle++(0.009,0.009) (0.127,0.548)rectangle++(0.009,0.009)
  (0.061,0.272)rectangle++(0.009,0.009) (0.028,0.634)rectangle++(0.009,0.009) (0.012,0.815)rectangle++(0.009,0.009) (0.004,0.905)rectangle++(0.009,0.009) (0.499,0.359)rectangle++(0.009,0.009)
  (0.247,0.177)rectangle++(0.009,0.009) (0.121,0.586)rectangle++(0.009,0.009) (0.058,0.291)rectangle++(0.009,0.009) (0.521,0.114)rectangle++(0.009,0.009) (0.258,0.055)rectangle++(0.009,0.009)
  (0.127,0.025)rectangle++(0.009,0.009) (0.061,0.51)rectangle++(0.009,0.009) (0.028,0.753)rectangle++(0.009,0.009) (0.509,0.298)rectangle++(0.009,0.009) (0.252,0.147)rectangle++(0.009,0.009)
  (0.124,0.071)rectangle++(0.009,0.009) (0.06,0.033)rectangle++(0.009,0.009) (0.028,0.514)rectangle++(0.009,0.009) (0.012,0.255)rectangle++(0.009,0.009) (0.502,0.099)rectangle++(0.009,0.009)
  (0.249,0.547)rectangle++(0.009,0.009) (0.122,0.271)rectangle++(0.009,0.009) (0.546,0.106)rectangle++(0.009,0.009) (0.271,0.051)rectangle++(0.009,0.009) (0.606,0.018)rectangle++(0.009,0.009)
  (0.301,0.007)rectangle++(0.009,0.009) (0.148,0.501)rectangle++(0.009,0.009) (0.072,0.248)rectangle++(0.009,0.009) (0.034,0.622)rectangle++(0.009,0.009) (0.511,0.246)rectangle++(0.009,0.009)
  (0.702,0.096)rectangle++(0.009,0.009) (0.349,0.546)rectangle++(0.009,0.009) (0.172,0.771)rectangle++(0.009,0.009) (0.084,0.883)rectangle++(0.009,0.009) (0.04,0.439)rectangle++(0.009,0.009)
  (0.513,0.173)rectangle++(0.009,0.009) (0.254,0.084)rectangle++(0.009,0.009) (0.125,0.04)rectangle++(0.009,0.009) (0.547,0.013)rectangle++(0.009,0.009) (0.271,0.504)rectangle++(0.009,0.009)
  (0.133,0.75)rectangle++(0.009,0.009) (0.064,0.373)rectangle++(0.009,0.009) (0.03,0.684)rectangle++(0.009,0.009) (0.013,0.84)rectangle++(0.009,0.009) (0.004,0.918)rectangle++(0.009,0.009)
  (0,0.957)rectangle++(0.009,0.009) (0.497,0.38)rectangle++(0.009,0.009) (0.246,0.688)rectangle++(0.009,0.009) (0.596,0.272)rectangle++(0.009,0.009) (0.736,0.106)rectangle++(0.009,0.009)
  (0.792,0.04)rectangle++(0.009,0.009) (0.814,0.013)rectangle++(0.009,0.009) (0.823,0.003)rectangle++(0.009,0.009) (0.826,-0.002)rectangle++(0.009,0.009) (0.411,-0.003)rectangle++(0.009,0.009)
  (0.203,0.496)rectangle++(0.009,0.009) (0.099,0.246)rectangle++(0.009,0.009) (0.047,0.121)rectangle++(0.009,0.009) (0.021,0.558)rectangle++(0.009,0.009) (0.008,0.777)rectangle++(0.009,0.009)
  (0.002,0.886)rectangle++(0.009,0.009) (0.498,0.352)rectangle++(0.009,0.009) (0.247,0.174)rectangle++(0.009,0.009) (0.121,0.085)rectangle++(0.009,0.009) (0.058,0.04)rectangle++(0.009,0.009)
  (0.521,0.013)rectangle++(0.009,0.009) (0.258,0.004)rectangle++(0.009,0.009) (0.601,-0.001)rectangle++(0.009,0.009) (0.298,0.497)rectangle++(0.009,0.009) (0.147,0.746)rectangle++(0.009,0.009)
  (0.071,0.371)rectangle++(0.009,0.009) (0.526,0.146)rectangle++(0.009,0.009) (0.261,0.571)rectangle++(0.009,0.009) (0.128,0.783)rectangle++(0.009,0.009) (0.549,0.311)rectangle++(0.009,0.009)
  (0.272,0.653)rectangle++(0.009,0.009) (0.134,0.324)rectangle++(0.009,0.009) (0.065,0.16)rectangle++(0.009,0.009) (0.03,0.578)rectangle++(0.009,0.009) (0.013,0.787)rectangle++(0.009,0.009)
  (0.004,0.391)rectangle++(0.009,0.009) (0,0.693)rectangle++(0.009,0.009) (-0.002,0.844)rectangle++(0.009,0.009) (-0.003,0.42)rectangle++(0.009,0.009) (-0.004,0.708)rectangle++(0.009,0.009)
  (-0.004,0.352)rectangle++(0.009,0.009) (-0.004,0.674)rectangle++(0.009,0.009) (0.496,0.267)rectangle++(0.009,0.009) (0.246,0.631)rectangle++(0.009,0.009) (0.121,0.813)rectangle++(0.009,0.009)
  (0.058,0.404)rectangle++(0.009,0.009) (0.521,0.159)rectangle++(0.009,0.009) (0.706,0.061)rectangle++(0.009,0.009) (0.78,0.022)rectangle++(0.009,0.009) (0.388,0.509)rectangle++(0.009,0.009)
  (0.192,0.252)rectangle++(0.009,0.009) (0.094,0.624)rectangle++(0.009,0.009) (0.045,0.81)rectangle++(0.009,0.009) (0.02,0.403)rectangle++(0.009,0.009) (0.008,0.699)rectangle++(0.009,0.009)
  (0.002,0.847)rectangle++(0.009,0.009) (-0.001,0.921)rectangle++(0.009,0.009) (-0.003,0.458)rectangle++(0.009,0.009) (-0.004,0.727)rectangle++(0.009,0.009) (-0.004,0.861)rectangle++(0.009,0.009)
  (-0.004,0.928)rectangle++(0.009,0.009) (0.496,0.369)rectangle++(0.009,0.009) (0.696,0.145)rectangle++(0.009,0.009) (0.346,0.07)rectangle++(0.009,0.009) (0.171,0.033)rectangle++(0.009,0.009)
  (0.083,0.514)rectangle++(0.009,0.009) (0.039,0.755)rectangle++(0.009,0.009) (0.017,0.875)rectangle++(0.009,0.009) (0.006,0.435)rectangle++(0.009,0.009) (0.001,0.215)rectangle++(0.009,0.009)
  (0.498,0.083)rectangle++(0.009,0.009) (0.696,0.031)rectangle++(0.009,0.009) (0.346,0.513)rectangle++(0.009,0.009) (0.171,0.754)rectangle++(0.009,0.009) (0.083,0.875)rectangle++(0.009,0.009)
  (0.039,0.935)rectangle++(0.009,0.009) (0.017,0.965)rectangle++(0.009,0.009) (0.006,0.98)rectangle++(0.009,0.009) (0.001,0.988)rectangle++(0.009,0.009) (-0.002,0.492)rectangle++(0.009,0.009)
  (0.497,0.194)rectangle++(0.009,0.009) (0.246,0.595)rectangle++(0.009,0.009) (0.121,0.295)rectangle++(0.009,0.009) (0.058,0.145)rectangle++(0.009,0.009) (0.027,0.07)rectangle++(0.009,0.009)
  (0.011,0.533)rectangle++(0.009,0.009) (0.003,0.764)rectangle++(0.009,0.009) (-0.001,0.38)rectangle++(0.009,0.009) (-0.003,0.188)rectangle++(0.009,0.009) (-0.004,0.092)rectangle++(0.009,0.009)
  (-0.004,0.544)rectangle++(0.009,0.009) (-0.004,0.27)rectangle++(0.009,0.009) (-0.004,0.633)rectangle++(0.009,0.009) (-0.004,0.814)rectangle++(0.009,0.009) (-0.004,0.905)rectangle++(0.009,0.009)
  (-0.004,0.95)rectangle++(0.009,0.009) (-0.004,0.973)rectangle++(0.009,0.009) (0.496,0.386)rectangle++(0.009,0.009) (0.696,0.152)rectangle++(0.009,0.009) (0.346,0.074)rectangle++(0.009,0.009)
  (0.171,0.535)rectangle++(0.009,0.009) (0.083,0.765)rectangle++(0.009,0.009) (0.039,0.38)rectangle++(0.009,0.009) (0.017,0.188)rectangle++(0.009,0.009) (0.006,0.092)rectangle++(0.009,0.009)
  (0.001,0.544)rectangle++(0.009,0.009) (-0.002,0.27)rectangle++(0.009,0.009) (-0.003,0.133)rectangle++(0.009,0.009) (-0.004,0.564)rectangle++(0.009,0.009) (-0.004,0.28)rectangle++(0.009,0.009)
  (0.496,0.109)rectangle++(0.009,0.009) (0.246,0.552)rectangle++(0.009,0.009) (0.121,0.774)rectangle++(0.009,0.009) (0.058,0.385)rectangle++(0.009,0.009) (0.027,0.19)rectangle++(0.009,0.009)
  (0.011,0.093)rectangle++(0.009,0.009) (0.502,0.034)rectangle++(0.009,0.009) (0.249,0.515)rectangle++(0.009,0.009) (0.122,0.755)rectangle++(0.009,0.009) (0.059,0.875)rectangle++(0.009,0.009)
  (0.027,0.935)rectangle++(0.009,0.009) (0.011,0.465)rectangle++(0.009,0.009) (0.003,0.73)rectangle++(0.009,0.009) (-0.001,0.863)rectangle++(0.009,0.009) (0.497,0.342)rectangle++(0.009,0.009)
  (0.246,0.169)rectangle++(0.009,0.009) (0.121,0.082)rectangle++(0.009,0.009) (0.058,0.539)rectangle++(0.009,0.009) (0.521,0.213)rectangle++(0.009,0.009) (0.706,0.082)rectangle++(0.009,0.009)
  (0.351,0.039)rectangle++(0.009,0.009) (0.173,0.017)rectangle++(0.009,0.009) (0.084,0.006)rectangle++(0.009,0.009) (0.04,0.501)rectangle++(0.009,0.009) (0.513,0.198)rectangle++(0.009,0.009)
  (0.254,0.097)rectangle++(0.009,0.009) (0.125,0.046)rectangle++(0.009,0.009) (0.547,0.016)rectangle++(0.009,0.009) (0.271,0.006)rectangle++(0.009,0.009) (0.133,0.501)rectangle++(0.009,0.009)
  (0.064,0.748)rectangle++(0.009,0.009) (0.523,0.297)rectangle++(0.009,0.009) (0.259,0.646)rectangle++(0.009,0.009) (0.601,0.256)rectangle++(0.009,0.009) (0.298,0.626)rectangle++(0.009,0.009)
  (0.147,0.311)rectangle++(0.009,0.009) (0.071,0.653)rectangle++(0.009,0.009) (0.033,0.824)rectangle++(0.009,0.009) (0.014,0.91)rectangle++(0.009,0.009) (0.005,0.453)rectangle++(0.009,0.009)
  (0,0.224)rectangle++(0.009,0.009) (0.497,0.087)rectangle++(0.009,0.009) (0.246,0.541)rectangle++(0.009,0.009) (0.121,0.768)rectangle++(0.009,0.009) (0.058,0.382)rectangle++(0.009,0.009)
  (0.521,0.15)rectangle++(0.009,0.009) (0.258,0.573)rectangle++(0.009,0.009) (0.127,0.784)rectangle++(0.009,0.009) (0.548,0.311)rectangle++(0.009,0.009) (0.717,0.122)rectangle++(0.009,0.009)
  (0.356,0.059)rectangle++(0.009,0.009) (0.64,0.021)rectangle++(0.009,0.009) (0.318,0.008)rectangle++(0.009,0.009) (0.157,0.502)rectangle++(0.009,0.009) (0.076,0.749)rectangle++(0.009,0.009)
  (0.036,0.872)rectangle++(0.009,0.009) (0.016,0.434)rectangle++(0.009,0.009) (0.006,0.715)rectangle++(0.009,0.009) (0.001,0.355)rectangle++(0.009,0.009) (-0.002,0.675)rectangle++(0.009,0.009)
  (-0.003,0.335)rectangle++(0.009,0.009) (-0.004,0.165)rectangle++(0.009,0.009) (-0.004,0.08)rectangle++(0.009,0.009) (0.496,0.029)rectangle++(0.009,0.009) (0.246,0.512)rectangle++(0.009,0.009)
  (0.121,0.254)rectangle++(0.009,0.009) (0.058,0.125)rectangle++(0.009,0.009) (0.027,0.06)rectangle++(0.009,0.009) (0.508,0.021)rectangle++(0.009,0.009) (0.252,0.008)rectangle++(0.009,0.009)
  (0.124,0.002)rectangle++(0.009,0.009) (0.06,-0.001)rectangle++(0.009,0.009) (0.028,-0.003)rectangle++(0.009,0.009) (0.012,0.496)rectangle++(0.009,0.009) (0.502,0.196)rectangle++(0.009,0.009)
  (0.249,0.096)rectangle++(0.009,0.009) (0.122,0.546)rectangle++(0.009,0.009) (0.546,0.216)rectangle++(0.009,0.009) (0.271,0.606)rectangle++(0.009,0.009) (0.606,0.24)rectangle++(0.009,0.009)
  (0.301,0.118)rectangle++(0.009,0.009) (0.618,0.044)rectangle++(0.009,0.009) (0.307,0.02)rectangle++(0.009,0.009) (0.151,0.508)rectangle++(0.009,0.009) (0.558,0.2)rectangle++(0.009,0.009)
  (0.277,0.098)rectangle++(0.009,0.009) (0.136,0.047)rectangle++(0.009,0.009) (0.066,0.021)rectangle++(0.009,0.009) (0.031,0.008)rectangle++(0.009,0.009) (0.51,0.001)rectangle++(0.009,0.009)
  (0.253,0.498)rectangle++(0.009,0.009) (0.124,0.747)rectangle++(0.009,0.009) (0.06,0.371)rectangle++(0.009,0.009) (0.028,0.183)rectangle++(0.009,0.009) (0.012,0.089)rectangle++(0.009,0.009)
  (0.004,0.542)rectangle++(0.009,0.009) (0,0.769)rectangle++(0.009,0.009) (0.497,0.305)rectangle++(0.009,0.009) (0.696,0.119)rectangle++(0.009,0.009) (0.346,0.057)rectangle++(0.009,0.009)
  (0.636,0.02)rectangle++(0.009,0.009) (0.316,0.008)rectangle++(0.009,0.009) (0.156,0.002)rectangle++(0.009,0.009) (0.56,-0.002)rectangle++(0.009,0.009) (0.278,0.497)rectangle++(0.009,0.009)
  (0.137,0.246)rectangle++(0.009,0.009) (0.552,0.096)rectangle++(0.009,0.009) (0.718,0.036)rectangle++(0.009,0.009) (0.357,0.516)rectangle++(0.009,0.009) (0.176,0.756)rectangle++(0.009,0.009)
  (0.568,0.3)rectangle++(0.009,0.009) (0.282,0.148)rectangle++(0.009,0.009) (0.139,0.572)rectangle++(0.009,0.009) (0.067,0.284)rectangle++(0.009,0.009) (0.031,0.64)rectangle++(0.009,0.009)
  (0.013,0.818)rectangle++(0.009,0.009) (0.004,0.907)rectangle++(0.009,0.009) (0,0.951)rectangle++(0.009,0.009) (0.497,0.378)rectangle++(0.009,0.009) (0.696,0.148)rectangle++(0.009,0.009)
  (0.776,0.057)rectangle++(0.009,0.009) (0.386,0.026)rectangle++(0.009,0.009) (0.191,0.511)rectangle++(0.009,0.009) (0.093,0.253)rectangle++(0.009,0.009) (0.044,0.624)rectangle++(0.009,0.009)
  (0.02,0.81)rectangle++(0.009,0.009) (0.008,0.403)rectangle++(0.009,0.009) (0.002,0.199)rectangle++(0.009,0.009) (-0.001,0.597)rectangle++(0.009,0.009) (-0.003,0.296)rectangle++(0.009,0.009)
  (-0.004,0.646)rectangle++(0.009,0.009) (0.496,0.256)rectangle++(0.009,0.009) (0.246,0.626)rectangle++(0.009,0.009) (0.596,0.248)rectangle++(0.009,0.009) (0.296,0.622)rectangle++(0.009,0.009)
  (0.146,0.309)rectangle++(0.009,0.009) (0.071,0.652)rectangle++(0.009,0.009) (0.033,0.324)rectangle++(0.009,0.009) (0.511,0.127)rectangle++(0.009,0.009) (0.253,0.061)rectangle++(0.009,0.009)
  (0.124,0.528)rectangle++(0.009,0.009) (0.06,0.762)rectangle++(0.009,0.009) (0.028,0.379)rectangle++(0.009,0.009) (0.012,0.187)rectangle++(0.009,0.009) (0.004,0.591)rectangle++(0.009,0.009)
  (0,0.293)rectangle++(0.009,0.009) (-0.002,0.644)rectangle++(0.009,0.009) (-0.003,0.32)rectangle++(0.009,0.009) (-0.004,0.658)rectangle++(0.009,0.009) (-0.004,0.827)rectangle++(0.009,0.009)
  (-0.004,0.911)rectangle++(0.009,0.009) (-0.004,0.953)rectangle++(0.009,0.009) (0.496,0.379)rectangle++(0.009,0.009) (0.696,0.149)rectangle++(0.009,0.009) (0.346,0.072)rectangle++(0.009,0.009)
  (0.171,0.034)rectangle++(0.009,0.009) (0.566,0.011)rectangle++(0.009,0.009) (0.281,0.003)rectangle++(0.009,0.009) (0.138,0.499)rectangle++(0.009,0.009) (0.067,0.247)rectangle++(0.009,0.009);
\fill[nuage]
  (0.031,0.621)rectangle++(0.009,0.009) (0.013,0.808)rectangle++(0.009,0.009) (0.004,0.902)rectangle++(0.009,0.009) (0,0.949)rectangle++(0.009,0.009) (-0.002,0.972)rectangle++(0.009,0.009)
  (-0.003,0.984)rectangle++(0.009,0.009) (-0.004,0.49)rectangle++(0.009,0.009) (-0.004,0.243)rectangle++(0.009,0.009) (-0.004,0.119)rectangle++(0.009,0.009) (0.496,0.045)rectangle++(0.009,0.009)
  (0.246,0.52)rectangle++(0.009,0.009) (0.596,0.205)rectangle++(0.009,0.009) (0.296,0.1)rectangle++(0.009,0.009) (0.146,0.548)rectangle++(0.009,0.009) (0.071,0.272)rectangle++(0.009,0.009)
  (0.526,0.106)rectangle++(0.009,0.009) (0.261,0.051)rectangle++(0.009,0.009) (0.128,0.023)rectangle++(0.009,0.009) (0.549,0.007)rectangle++(0.009,0.009) (0.272,0.001)rectangle++(0.009,0.009)
  (0.134,-0.002)rectangle++(0.009,0.009) (0.065,0.497)rectangle++(0.009,0.009) (0.523,0.196)rectangle++(0.009,0.009) (0.259,0.096)rectangle++(0.009,0.009) (0.601,0.036)rectangle++(0.009,0.009)
  (0.298,0.516)rectangle++(0.009,0.009) (0.617,0.204)rectangle++(0.009,0.009) (0.744,0.079)rectangle++(0.009,0.009) (0.37,0.537)rectangle++(0.009,0.009) (0.183,0.266)rectangle++(0.009,0.009)
  (0.089,0.131)rectangle++(0.009,0.009) (0.042,0.063)rectangle++(0.009,0.009) (0.019,0.029)rectangle++(0.009,0.009) (0.007,0.012)rectangle++(0.009,0.009) (0.001,0.504)rectangle++(0.009,0.009)
  (0.498,0.199)rectangle++(0.009,0.009) (0.247,0.597)rectangle++(0.009,0.009) (0.121,0.296)rectangle++(0.009,0.009) (0.058,0.146)rectangle++(0.009,0.009) (0.027,0.571)rectangle++(0.009,0.009)
  (0.011,0.283)rectangle++(0.009,0.009) (0.003,0.639)rectangle++(0.009,0.009) (-0.001,0.817)rectangle++(0.009,0.009) (-0.003,0.406)rectangle++(0.009,0.009) (0.496,0.16)rectangle++(0.009,0.009)
  (0.246,0.578)rectangle++(0.009,0.009) (0.121,0.787)rectangle++(0.009,0.009) (0.058,0.391)rectangle++(0.009,0.009) (0.027,0.193)rectangle++(0.009,0.009) (0.011,0.594)rectangle++(0.009,0.009)
  (0.003,0.795)rectangle++(0.009,0.009) (-0.001,0.895)rectangle++(0.009,0.009) (0.497,0.355)rectangle++(0.009,0.009) (0.246,0.675)rectangle++(0.009,0.009) (0.121,0.335)rectangle++(0.009,0.009)
  (0.058,0.165)rectangle++(0.009,0.009) (0.027,0.08)rectangle++(0.009,0.009) (0.011,0.038)rectangle++(0.009,0.009) (0.502,0.012)rectangle++(0.009,0.009) (0.249,0.504)rectangle++(0.009,0.009)
  (0.122,0.75)rectangle++(0.009,0.009) (0.059,0.373)rectangle++(0.009,0.009) (0.027,0.184)rectangle++(0.009,0.009) (0.011,0.59)rectangle++(0.009,0.009) (0.502,0.233)rectangle++(0.009,0.009)
  (0.698,0.091)rectangle++(0.009,0.009) (0.347,0.543)rectangle++(0.009,0.009) (0.636,0.215)rectangle++(0.009,0.009) (0.316,0.605)rectangle++(0.009,0.009) (0.156,0.3)rectangle++(0.009,0.009)
  (0.56,0.117)rectangle++(0.009,0.009) (0.721,0.044)rectangle++(0.009,0.009) (0.786,0.015)rectangle++(0.009,0.009) (0.391,0.505)rectangle++(0.009,0.009) (0.193,0.75)rectangle++(0.009,0.009)
  (0.094,0.373)rectangle++(0.009,0.009) (0.045,0.184)rectangle++(0.009,0.009) (0.02,0.09)rectangle++(0.009,0.009) (0.008,0.043)rectangle++(0.009,0.009) (0.002,0.519)rectangle++(0.009,0.009)
  (-0.001,0.257)rectangle++(0.009,0.009) (-0.003,0.126)rectangle++(0.009,0.009) (0.496,0.048)rectangle++(0.009,0.009) (0.246,0.022)rectangle++(0.009,0.009) (0.121,0.509)rectangle++(0.009,0.009)
  (0.058,0.752)rectangle++(0.009,0.009) (0.027,0.374)rectangle++(0.009,0.009) (0.011,0.185)rectangle++(0.009,0.009) (0.003,0.09)rectangle++(0.009,0.009) (0.499,0.033)rectangle++(0.009,0.009)
  (0.697,0.011)rectangle++(0.009,0.009) (0.776,0.002)rectangle++(0.009,0.009) (0.386,0.499)rectangle++(0.009,0.009) (0.652,0.197)rectangle++(0.009,0.009) (0.324,0.596)rectangle++(0.009,0.009)
  (0.627,0.236)rectangle++(0.009,0.009) (0.311,0.116)rectangle++(0.009,0.009) (0.622,0.044)rectangle++(0.009,0.009) (0.746,0.015)rectangle++(0.009,0.009) (0.371,0.505)rectangle++(0.009,0.009)
  (0.183,0.25)rectangle++(0.009,0.009) (0.089,0.123)rectangle++(0.009,0.009) (0.533,0.046)rectangle++(0.009,0.009) (0.711,0.016)rectangle++(0.009,0.009) (0.782,0.004)rectangle++(0.009,0.009)
  (0.389,0.5)rectangle++(0.009,0.009) (0.192,0.248)rectangle++(0.009,0.009) (0.094,0.122)rectangle++(0.009,0.009) (0.045,0.059)rectangle++(0.009,0.009) (0.02,0.027)rectangle++(0.009,0.009)
  (0.008,0.511)rectangle++(0.009,0.009) (0.002,0.753)rectangle++(0.009,0.009) (-0.001,0.374)rectangle++(0.009,0.009) (-0.003,0.685)rectangle++(0.009,0.009) (-0.004,0.84)rectangle++(0.009,0.009)
  (0.496,0.333)rectangle++(0.009,0.009) (0.246,0.164)rectangle++(0.009,0.009) (0.121,0.58)rectangle++(0.009,0.009) (0.058,0.788)rectangle++(0.009,0.009) (0.027,0.892)rectangle++(0.009,0.009)
  (0.011,0.944)rectangle++(0.009,0.009) (0.003,0.97)rectangle++(0.009,0.009) (0.499,0.385)rectangle++(0.009,0.009) (0.247,0.69)rectangle++(0.009,0.009) (0.121,0.843)rectangle++(0.009,0.009)
  (0.546,0.334)rectangle++(0.009,0.009) (0.271,0.665)rectangle++(0.009,0.009) (0.133,0.33)rectangle++(0.009,0.009) (0.064,0.663)rectangle++(0.009,0.009) (0.03,0.829)rectangle++(0.009,0.009)
  (0.013,0.912)rectangle++(0.009,0.009) (0.004,0.454)rectangle++(0.009,0.009) (0,0.225)rectangle++(0.009,0.009) (-0.002,0.61)rectangle++(0.009,0.009) (-0.003,0.303)rectangle++(0.009,0.009)
  (-0.004,0.149)rectangle++(0.009,0.009) (-0.004,0.572)rectangle++(0.009,0.009) (-0.004,0.284)rectangle++(0.009,0.009) (-0.004,0.64)rectangle++(0.009,0.009) (-0.004,0.318)rectangle++(0.009,0.009)
  (-0.004,0.157)rectangle++(0.009,0.009) (-0.004,0.576)rectangle++(0.009,0.009) (0.496,0.228)rectangle++(0.009,0.009) (0.246,0.612)rectangle++(0.009,0.009) (0.121,0.804)rectangle++(0.009,0.009)
  (0.546,0.319)rectangle++(0.009,0.009) (0.271,0.157)rectangle++(0.009,0.009) (0.133,0.076)rectangle++(0.009,0.009) (0.064,0.536)rectangle++(0.009,0.009) (0.03,0.266)rectangle++(0.009,0.009)
  (0.013,0.631)rectangle++(0.009,0.009) (0.502,0.25)rectangle++(0.009,0.009) (0.249,0.123)rectangle++(0.009,0.009) (0.122,0.559)rectangle++(0.009,0.009) (0.059,0.777)rectangle++(0.009,0.009)
  (0.027,0.886)rectangle++(0.009,0.009) (0.011,0.441)rectangle++(0.009,0.009) (0.003,0.218)rectangle++(0.009,0.009) (0.499,0.085)rectangle++(0.009,0.009) (0.247,0.54)rectangle++(0.009,0.009)
  (0.121,0.768)rectangle++(0.009,0.009) (0.058,0.382)rectangle++(0.009,0.009) (0.027,0.689)rectangle++(0.009,0.009) (0.011,0.842)rectangle++(0.009,0.009) (0.502,0.334)rectangle++(0.009,0.009)
  (0.698,0.131)rectangle++(0.009,0.009) (0.777,0.05)rectangle++(0.009,0.009) (0.386,0.023)rectangle++(0.009,0.009) (0.191,0.009)rectangle++(0.009,0.009) (0.093,0.002)rectangle++(0.009,0.009)
  (0.044,0.499)rectangle++(0.009,0.009) (0.515,0.197)rectangle++(0.009,0.009) (0.255,0.096)rectangle++(0.009,0.009) (0.599,0.036)rectangle++(0.009,0.009) (0.737,0.012)rectangle++(0.009,0.009)
  (0.792,0.002)rectangle++(0.009,0.009) (0.394,0.499)rectangle++(0.009,0.009) (0.195,0.747)rectangle++(0.009,0.009) (0.575,0.296)rectangle++(0.009,0.009) (0.285,0.146)rectangle++(0.009,0.009)
  (0.611,0.056)rectangle++(0.009,0.009) (0.742,0.02)rectangle++(0.009,0.009) (0.794,0.005)rectangle++(0.009,0.009) (0.395,0.5)rectangle++(0.009,0.009) (0.195,0.248)rectangle++(0.009,0.009)
  (0.095,0.122)rectangle++(0.009,0.009) (0.045,0.059)rectangle++(0.009,0.009) (0.02,0.027)rectangle++(0.009,0.009) (0.008,0.011)rectangle++(0.009,0.009) (0.002,0.503)rectangle++(0.009,0.009)
  (-0.001,0.749)rectangle++(0.009,0.009) (-0.003,0.872)rectangle++(0.009,0.009) (-0.004,0.934)rectangle++(0.009,0.009) (-0.004,0.965)rectangle++(0.009,0.009) (0.496,0.383)rectangle++(0.009,0.009)
  (0.246,0.189)rectangle++(0.009,0.009) (0.121,0.592)rectangle++(0.009,0.009) (0.546,0.234)rectangle++(0.009,0.009) (0.271,0.615)rectangle++(0.009,0.009) (0.606,0.243)rectangle++(0.009,0.009)
  (0.74,0.095)rectangle++(0.009,0.009) (0.368,0.545)rectangle++(0.009,0.009) (0.182,0.27)rectangle++(0.009,0.009) (0.089,0.133)rectangle++(0.009,0.009) (0.042,0.564)rectangle++(0.009,0.009)
  (0.019,0.78)rectangle++(0.009,0.009) (0.007,0.388)rectangle++(0.009,0.009) (0.001,0.692)rectangle++(0.009,0.009) (-0.002,0.844)rectangle++(0.009,0.009) (0.497,0.335)rectangle++(0.009,0.009)
  (0.246,0.165)rectangle++(0.009,0.009) (0.121,0.08)rectangle++(0.009,0.009) (0.058,0.038)rectangle++(0.009,0.009) (0.027,0.517)rectangle++(0.009,0.009) (0.011,0.256)rectangle++(0.009,0.009)
  (0.003,0.626)rectangle++(0.009,0.009) (-0.001,0.311)rectangle++(0.009,0.009) (0.497,0.122)rectangle++(0.009,0.009) (0.696,0.046)rectangle++(0.009,0.009) (0.346,0.021)rectangle++(0.009,0.009)
  (0.636,0.006)rectangle++(0.009,0.009) (0.316,0.501)rectangle++(0.009,0.009) (0.156,0.248)rectangle++(0.009,0.009) (0.076,0.122)rectangle++(0.009,0.009) (0.528,0.046)rectangle++(0.009,0.009)
  (0.262,0.521)rectangle++(0.009,0.009) (0.129,0.758)rectangle++(0.009,0.009) (0.062,0.377)rectangle++(0.009,0.009) (0.029,0.186)rectangle++(0.009,0.009) (0.509,0.072)rectangle++(0.009,0.009)
  (0.252,0.034)rectangle++(0.009,0.009) (0.124,0.515)rectangle++(0.009,0.009) (0.06,0.755)rectangle++(0.009,0.009) (0.028,0.875)rectangle++(0.009,0.009) (0.012,0.435)rectangle++(0.009,0.009)
  (0.502,0.171)rectangle++(0.009,0.009) (0.249,0.583)rectangle++(0.009,0.009) (0.122,0.789)rectangle++(0.009,0.009) (0.546,0.313)rectangle++(0.009,0.009) (0.271,0.154)rectangle++(0.009,0.009)
  (0.133,0.575)rectangle++(0.009,0.009) (0.064,0.285)rectangle++(0.009,0.009) (0.03,0.14)rectangle++(0.009,0.009) (0.013,0.068)rectangle++(0.009,0.009) (0.004,0.032)rectangle++(0.009,0.009)
  (0,0.014)rectangle++(0.009,0.009) (-0.002,0.005)rectangle++(0.009,0.009) (-0.003,0.5)rectangle++(0.009,0.009) (-0.004,0.748)rectangle++(0.009,0.009) (-0.004,0.872)rectangle++(0.009,0.009)
  (-0.004,0.934)rectangle++(0.009,0.009) (-0.004,0.465)rectangle++(0.009,0.009) (-0.004,0.73)rectangle++(0.009,0.009) (-0.004,0.863)rectangle++(0.009,0.009) (0.496,0.342)rectangle++(0.009,0.009)
  (0.246,0.169)rectangle++(0.009,0.009) (0.596,0.065)rectangle++(0.009,0.009) (0.296,0.03)rectangle++(0.009,0.009) (0.146,0.513)rectangle++(0.009,0.009) (0.556,0.202)rectangle++(0.009,0.009)
  (0.276,0.099)rectangle++(0.009,0.009) (0.136,0.547)rectangle++(0.009,0.009) (0.066,0.771)rectangle++(0.009,0.009) (0.031,0.883)rectangle++(0.009,0.009) (0.013,0.439)rectangle++(0.009,0.009)
  (0.503,0.173)rectangle++(0.009,0.009) (0.249,0.084)rectangle++(0.009,0.009) (0.122,0.54)rectangle++(0.009,0.009) (0.546,0.213)rectangle++(0.009,0.009) (0.271,0.104)rectangle++(0.009,0.009)
  (0.133,0.55)rectangle++(0.009,0.009) (0.551,0.217)rectangle++(0.009,0.009) (0.718,0.084)rectangle++(0.009,0.009) (0.784,0.031)rectangle++(0.009,0.009) (0.811,0.01)rectangle++(0.009,0.009)
  (0.403,0.503)rectangle++(0.009,0.009) (0.659,0.198)rectangle++(0.009,0.009) (0.761,0.077)rectangle++(0.009,0.009) (0.378,0.536)rectangle++(0.009,0.009) (0.187,0.266)rectangle++(0.009,0.009)
  (0.091,0.131)rectangle++(0.009,0.009) (0.043,0.063)rectangle++(0.009,0.009) (0.019,0.529)rectangle++(0.009,0.009) (0.505,0.209)rectangle++(0.009,0.009) (0.699,0.081)rectangle++(0.009,0.009)
  (0.347,0.538)rectangle++(0.009,0.009) (0.636,0.213)rectangle++(0.009,0.009) (0.316,0.104)rectangle++(0.009,0.009) (0.624,0.039)rectangle++(0.009,0.009) (0.31,0.017)rectangle++(0.009,0.009)
  (0.153,0.006)rectangle++(0.009,0.009) (0.074,0.001)rectangle++(0.009,0.009) (0.527,-0.002)rectangle++(0.009,0.009) (0.261,0.497)rectangle++(0.009,0.009) (0.602,0.196)rectangle++(0.009,0.009)
  (0.738,0.076)rectangle++(0.009,0.009) (0.367,0.036)rectangle++(0.009,0.009) (0.181,0.516)rectangle++(0.009,0.009) (0.57,0.204)rectangle++(0.009,0.009) (0.283,0.1)rectangle++(0.009,0.009)
  (0.139,0.048)rectangle++(0.009,0.009) (0.067,0.022)rectangle++(0.009,0.009) (0.524,0.006)rectangle++(0.009,0.009) (0.26,0.001)rectangle++(0.009,0.009) (0.601,-0.002)rectangle++(0.009,0.009)
  (0.298,0.497)rectangle++(0.009,0.009) (0.147,0.246)rectangle++(0.009,0.009) (0.071,0.121)rectangle++(0.009,0.009) (0.033,0.058)rectangle++(0.009,0.009) (0.014,0.527)rectangle++(0.009,0.009)
  (0.005,0.761)rectangle++(0.009,0.009) (0,0.878)rectangle++(0.009,0.009) (-0.002,0.437)rectangle++(0.009,0.009) (0.496,0.172)rectangle++(0.009,0.009) (0.696,0.066)rectangle++(0.009,0.009)
  (0.346,0.031)rectangle++(0.009,0.009) (0.171,0.013)rectangle++(0.009,0.009) (0.566,0.003)rectangle++(0.009,0.009) (0.281,-0.001)rectangle++(0.009,0.009) (0.61,-0.003)rectangle++(0.009,0.009)
  (0.303,-0.004)rectangle++(0.009,0.009) (0.149,-0.004)rectangle++(0.009,0.009) (0.072,-0.004)rectangle++(0.009,0.009) (0.034,0.496)rectangle++(0.009,0.009) (0.015,0.246)rectangle++(0.009,0.009)
  (0.503,0.096)rectangle++(0.009,0.009) (0.249,0.546)rectangle++(0.009,0.009) (0.122,0.271)rectangle++(0.009,0.009) (0.546,0.106)rectangle++(0.009,0.009) (0.271,0.551)rectangle++(0.009,0.009)
  (0.133,0.273)rectangle++(0.009,0.009) (0.064,0.634)rectangle++(0.009,0.009) (0.523,0.251)rectangle++(0.009,0.009) (0.259,0.623)rectangle++(0.009,0.009) (0.127,0.309)rectangle++(0.009,0.009)
  (0.061,0.652)rectangle++(0.009,0.009) (0.028,0.324)rectangle++(0.009,0.009) (0.012,0.16)rectangle++(0.009,0.009) (0.004,0.078)rectangle++(0.009,0.009) (0,0.537)rectangle++(0.009,0.009)
  (0.497,0.212)rectangle++(0.009,0.009) (0.696,0.082)rectangle++(0.009,0.009) (0.776,0.03)rectangle++(0.009,0.009) (0.386,0.513)rectangle++(0.009,0.009) (0.191,0.254)rectangle++(0.009,0.009)
  (0.574,0.099)rectangle++(0.009,0.009) (0.285,0.047)rectangle++(0.009,0.009) (0.14,0.021)rectangle++(0.009,0.009) (0.553,0.006)rectangle++(0.009,0.009) (0.274,0.501)rectangle++(0.009,0.009)
  (0.135,0.748)rectangle++(0.009,0.009) (0.065,0.872)rectangle++(0.009,0.009) (0.03,0.434)rectangle++(0.009,0.009) (0.013,0.215)rectangle++(0.009,0.009) (0.502,0.083)rectangle++(0.009,0.009)
  (0.249,0.539)rectangle++(0.009,0.009) (0.122,0.267)rectangle++(0.009,0.009) (0.059,0.131)rectangle++(0.009,0.009) (0.521,0.05)rectangle++(0.009,0.009) (0.258,0.523)rectangle++(0.009,0.009)
  (0.127,0.259)rectangle++(0.009,0.009) (0.548,0.101)rectangle++(0.009,0.009) (0.272,0.548)rectangle++(0.009,0.009) (0.606,0.217)rectangle++(0.009,0.009) (0.301,0.606)rectangle++(0.009,0.009)
  (0.148,0.801)rectangle++(0.009,0.009) (0.072,0.898)rectangle++(0.009,0.009) (0.034,0.947)rectangle++(0.009,0.009) (0.015,0.971)rectangle++(0.009,0.009) (0.005,0.983)rectangle++(0.009,0.009)
  (0.499,0.391)rectangle++(0.009,0.009) (0.697,0.154)rectangle++(0.009,0.009) (0.776,0.059)rectangle++(0.009,0.009) (0.386,0.527)rectangle++(0.009,0.009) (0.652,0.208)rectangle++(0.009,0.009)
  (0.758,0.081)rectangle++(0.009,0.009) (0.377,0.538)rectangle++(0.009,0.009) (0.186,0.767)rectangle++(0.009,0.009) (0.091,0.381)rectangle++(0.009,0.009) (0.043,0.188)rectangle++(0.009,0.009)
  (0.019,0.592)rectangle++(0.009,0.009) (0.007,0.294)rectangle++(0.009,0.009) (0.5,0.115)rectangle++(0.009,0.009) (0.248,0.555)rectangle++(0.009,0.009) (0.122,0.775)rectangle++(0.009,0.009)
  (0.059,0.885)rectangle++(0.009,0.009) (0.027,0.94)rectangle++(0.009,0.009) (0.011,0.968)rectangle++(0.009,0.009) (0.502,0.384)rectangle++(0.009,0.009) (0.249,0.69)rectangle++(0.009,0.009)
  (0.122,0.343)rectangle++(0.009,0.009) (0.059,0.169)rectangle++(0.009,0.009) (0.027,0.582)rectangle++(0.009,0.009) (0.508,0.23)rectangle++(0.009,0.009) (0.252,0.113)rectangle++(0.009,0.009)
  (0.124,0.554)rectangle++(0.009,0.009) (0.06,0.275)rectangle++(0.009,0.009) (0.521,0.107)rectangle++(0.009,0.009) (0.258,0.551)rectangle++(0.009,0.009) (0.127,0.273)rectangle++(0.009,0.009)
  (0.061,0.634)rectangle++(0.009,0.009) (0.028,0.815)rectangle++(0.009,0.009) (0.509,0.323)rectangle++(0.009,0.009) (0.252,0.659)rectangle++(0.009,0.009) (0.124,0.827)rectangle++(0.009,0.009)
  (0.06,0.411)rectangle++(0.009,0.009) (0.028,0.203)rectangle++(0.009,0.009) (0.508,0.079)rectangle++(0.009,0.009) (0.701,0.029)rectangle++(0.009,0.009) (0.348,0.512)rectangle++(0.009,0.009)
  (0.172,0.754)rectangle++(0.009,0.009) (0.566,0.299)rectangle++(0.009,0.009) (0.724,0.117)rectangle++(0.009,0.009) (0.787,0.044)rectangle++(0.009,0.009) (0.391,0.02)rectangle++(0.009,0.009)
  (0.193,0.008)rectangle++(0.009,0.009) (0.575,0)rectangle++(0.009,0.009) (0.285,-0.002)rectangle++(0.009,0.009) (0.14,-0.003)rectangle++(0.009,0.009) (0.068,0.496)rectangle++(0.009,0.009)
  (0.032,0.246)rectangle++(0.009,0.009) (0.014,0.621)rectangle++(0.009,0.009) (0.503,0.246)rectangle++(0.009,0.009) (0.698,0.096)rectangle++(0.009,0.009) (0.347,0.546)rectangle++(0.009,0.009)
  (0.171,0.271)rectangle++(0.009,0.009) (0.566,0.106)rectangle++(0.009,0.009) (0.281,0.051)rectangle++(0.009,0.009) (0.138,0.023)rectangle++(0.009,0.009) (0.067,0.509)rectangle++(0.009,0.009)
  (0.031,0.252)rectangle++(0.009,0.009) (0.51,0.098)rectangle++(0.009,0.009) (0.253,0.547)rectangle++(0.009,0.009) (0.124,0.271)rectangle++(0.009,0.009) (0.06,0.633)rectangle++(0.009,0.009)
  (0.028,0.814)rectangle++(0.009,0.009) (0.012,0.905)rectangle++(0.009,0.009) (0.004,0.45)rectangle++(0.009,0.009) (0,0.223)rectangle++(0.009,0.009) (0.497,0.086)rectangle++(0.009,0.009)
  (0.246,0.041)rectangle++(0.009,0.009) (0.121,0.518)rectangle++(0.009,0.009) (0.058,0.757)rectangle++(0.009,0.009) (0.027,0.876)rectangle++(0.009,0.009) (0.508,0.348)rectangle++(0.009,0.009)
  (0.252,0.672)rectangle++(0.009,0.009) (0.598,0.266)rectangle++(0.009,0.009) (0.297,0.131)rectangle++(0.009,0.009) (0.146,0.563)rectangle++(0.009,0.009) (0.071,0.779)rectangle++(0.009,0.009)
  (0.033,0.387)rectangle++(0.009,0.009) (0.014,0.691)rectangle++(0.009,0.009) (0.005,0.343)rectangle++(0.009,0.009) (0.499,0.135)rectangle++(0.009,0.009) (0.247,0.565)rectangle++(0.009,0.009)
  (0.596,0.223)rectangle++(0.009,0.009) (0.296,0.109)rectangle++(0.009,0.009) (0.146,0.052)rectangle++(0.009,0.009) (0.071,0.524)rectangle++(0.009,0.009) (0.033,0.76)rectangle++(0.009,0.009)
  (0.511,0.301)rectangle++(0.009,0.009) (0.702,0.118)rectangle++(0.009,0.009) (0.349,0.557)rectangle++(0.009,0.009) (0.172,0.276)rectangle++(0.009,0.009) (0.084,0.636)rectangle++(0.009,0.009)
  (0.04,0.316)rectangle++(0.009,0.009) (0.018,0.656)rectangle++(0.009,0.009) (0.007,0.326)rectangle++(0.009,0.009) (0.5,0.128)rectangle++(0.009,0.009) (0.697,0.048)rectangle++(0.009,0.009)
  (0.346,0.522)rectangle++(0.009,0.009) (0.171,0.259)rectangle++(0.009,0.009) (0.566,0.101)rectangle++(0.009,0.009) (0.281,0.048)rectangle++(0.009,0.009) (0.138,0.522)rectangle++(0.009,0.009)
  (0.067,0.259)rectangle++(0.009,0.009) (0.031,0.627)rectangle++(0.009,0.009) (0.013,0.311)rectangle++(0.009,0.009) (0.503,0.122)rectangle++(0.009,0.009) (0.698,0.046)rectangle++(0.009,0.009)
  (0.347,0.521)rectangle++(0.009,0.009) (0.171,0.758)rectangle++(0.009,0.009) (0.083,0.377)rectangle++(0.009,0.009) (0.039,0.186)rectangle++(0.009,0.009) (0.017,0.591)rectangle++(0.009,0.009)
  (0.006,0.793)rectangle++(0.009,0.009) (0.001,0.394)rectangle++(0.009,0.009) (-0.002,0.195)rectangle++(0.009,0.009) (0.497,0.075)rectangle++(0.009,0.009) (0.696,0.027)rectangle++(0.009,0.009)
  (0.346,0.511)rectangle++(0.009,0.009) (0.171,0.753)rectangle++(0.009,0.009) (0.083,0.874)rectangle++(0.009,0.009) (0.039,0.935)rectangle++(0.009,0.009) (0.513,0.371)rectangle++(0.009,0.009)
  (0.254,0.683)rectangle++(0.009,0.009) (0.125,0.839)rectangle++(0.009,0.009) (0.06,0.917)rectangle++(0.009,0.009) (0.028,0.956)rectangle++(0.009,0.009) (0.508,0.38)rectangle++(0.009,0.009)
  (0.701,0.149)rectangle++(0.009,0.009) (0.348,0.572)rectangle++(0.009,0.009) (0.172,0.284)rectangle++(0.009,0.009) (0.084,0.14)rectangle++(0.009,0.009) (0.531,0.053)rectangle++(0.009,0.009)
  (0.263,0.024)rectangle++(0.009,0.009) (0.129,0.01)rectangle++(0.009,0.009) (0.062,0.503)rectangle++(0.009,0.009) (0.029,0.249)rectangle++(0.009,0.009) (0.012,0.122)rectangle++(0.009,0.009)
  (0.502,0.046)rectangle++(0.009,0.009) (0.698,0.016)rectangle++(0.009,0.009) (0.777,0.004)rectangle++(0.009,0.009) (0.386,0.5)rectangle++(0.009,0.009) (0.191,0.248)rectangle++(0.009,0.009)
  (0.093,0.122)rectangle++(0.009,0.009) (0.044,0.059)rectangle++(0.009,0.009) (0.02,0.027)rectangle++(0.009,0.009) (0.008,0.011)rectangle++(0.009,0.009) (0.5,0.002)rectangle++(0.009,0.009)
  (0.248,0.499)rectangle++(0.009,0.009) (0.122,0.747)rectangle++(0.009,0.009) (0.546,0.296)rectangle++(0.009,0.009) (0.271,0.146)rectangle++(0.009,0.009) (0.133,0.071)rectangle++(0.009,0.009)
  (0.064,0.533)rectangle++(0.009,0.009) (0.523,0.211)rectangle++(0.009,0.009) (0.259,0.103)rectangle++(0.009,0.009) (0.601,0.039)rectangle++(0.009,0.009) (0.298,0.517)rectangle++(0.009,0.009)
  (0.617,0.204)rectangle++(0.009,0.009) (0.306,0.6)rectangle++(0.009,0.009) (0.62,0.237)rectangle++(0.009,0.009) (0.308,0.116)rectangle++(0.009,0.009) (0.62,0.044)rectangle++(0.009,0.009)
  (0.308,0.52)rectangle++(0.009,0.009) (0.62,0.205)rectangle++(0.009,0.009) (0.308,0.1)rectangle++(0.009,0.009) (0.152,0.048)rectangle++(0.009,0.009) (0.074,0.522)rectangle++(0.009,0.009)
  (0.035,0.759)rectangle++(0.009,0.009) (0.015,0.877)rectangle++(0.009,0.009) (0.503,0.348)rectangle++(0.009,0.009) (0.249,0.672)rectangle++(0.009,0.009) (0.122,0.334)rectangle++(0.009,0.009)
  (0.059,0.165)rectangle++(0.009,0.009) (0.521,0.063)rectangle++(0.009,0.009) (0.258,0.529)rectangle++(0.009,0.009) (0.127,0.762)rectangle++(0.009,0.009) (0.061,0.879)rectangle++(0.009,0.009)
  (0.028,0.937)rectangle++(0.009,0.009) (0.012,0.966)rectangle++(0.009,0.009) (0.004,0.981)rectangle++(0.009,0.009) (0.499,0.39)rectangle++(0.009,0.009) (0.247,0.693)rectangle++(0.009,0.009)
  (0.596,0.274)rectangle++(0.009,0.009) (0.296,0.635)rectangle++(0.009,0.009) (0.146,0.315)rectangle++(0.009,0.009) (0.071,0.155)rectangle++(0.009,0.009) (0.033,0.075)rectangle++(0.009,0.009)
  (0.511,0.027)rectangle++(0.009,0.009) (0.253,0.011)rectangle++(0.009,0.009) (0.124,0.503)rectangle++(0.009,0.009) (0.06,0.249)rectangle++(0.009,0.009) (0.028,0.622)rectangle++(0.009,0.009)
  (0.508,0.246)rectangle++(0.009,0.009) (0.252,0.621)rectangle++(0.009,0.009) (0.124,0.808)rectangle++(0.009,0.009) (0.06,0.902)rectangle++(0.009,0.009) (0.521,0.358)rectangle++(0.009,0.009)
  (0.258,0.677)rectangle++(0.009,0.009) (0.127,0.336)rectangle++(0.009,0.009) (0.548,0.132)rectangle++(0.009,0.009) (0.272,0.564)rectangle++(0.009,0.009) (0.134,0.28)rectangle++(0.009,0.009)
  (0.065,0.638)rectangle++(0.009,0.009) (0.03,0.317)rectangle++(0.009,0.009) (0.013,0.656)rectangle++(0.009,0.009) (0.004,0.826)rectangle++(0.009,0.009) (0,0.911)rectangle++(0.009,0.009)
  (-0.002,0.453)rectangle++(0.009,0.009) (-0.003,0.224)rectangle++(0.009,0.009) (-0.004,0.11)rectangle++(0.009,0.009) (-0.004,0.553)rectangle++(0.009,0.009) (0.496,0.218)rectangle++(0.009,0.009)
  (0.246,0.107)rectangle++(0.009,0.009) (0.121,0.551)rectangle++(0.009,0.009) (0.058,0.773)rectangle++(0.009,0.009) (0.521,0.307)rectangle++(0.009,0.009) (0.258,0.651)rectangle++(0.009,0.009)
  (0.127,0.323)rectangle++(0.009,0.009) (0.061,0.659)rectangle++(0.009,0.009) (0.028,0.327)rectangle++(0.009,0.009) (0.509,0.128)rectangle++(0.009,0.009) (0.252,0.062)rectangle++(0.009,0.009)
  (0.124,0.529)rectangle++(0.009,0.009) (0.06,0.262)rectangle++(0.009,0.009) (0.521,0.102)rectangle++(0.009,0.009) (0.258,0.049)rectangle++(0.009,0.009) (0.601,0.017)rectangle++(0.009,0.009)
  (0.298,0.006)rectangle++(0.009,0.009) (0.147,0.001)rectangle++(0.009,0.009) (0.071,0.498)rectangle++(0.009,0.009) (0.526,0.197)rectangle++(0.009,0.009) (0.261,0.596)rectangle++(0.009,0.009)
  (0.128,0.796)rectangle++(0.009,0.009) (0.062,0.896)rectangle++(0.009,0.009) (0.522,0.356)rectangle++(0.009,0.009) (0.259,0.176)rectangle++(0.009,0.009) (0.127,0.586)rectangle++(0.009,0.009)
  (0.061,0.791)rectangle++(0.009,0.009) (0.522,0.314)rectangle++(0.009,0.009) (0.259,0.155)rectangle++(0.009,0.009) (0.127,0.075)rectangle++(0.009,0.009) (0.061,0.535)rectangle++(0.009,0.009)
  (0.028,0.765)rectangle++(0.009,0.009) (0.012,0.88)rectangle++(0.009,0.009) (0.004,0.438)rectangle++(0.009,0.009) (0,0.717)rectangle++(0.009,0.009) (-0.002,0.356)rectangle++(0.009,0.009)
  (-0.003,0.176)rectangle++(0.009,0.009) (-0.004,0.586)rectangle++(0.009,0.009) (-0.004,0.791)rectangle++(0.009,0.009) (-0.004,0.893)rectangle++(0.009,0.009) (-0.004,0.444)rectangle++(0.009,0.009)
  (-0.004,0.22)rectangle++(0.009,0.009) (-0.004,0.108)rectangle++(0.009,0.009) (-0.004,0.552)rectangle++(0.009,0.009) (-0.004,0.274)rectangle++(0.009,0.009) (-0.004,0.635)rectangle++(0.009,0.009)
  (-0.004,0.815)rectangle++(0.009,0.009) (-0.004,0.405)rectangle++(0.009,0.009) (0.496,0.159)rectangle++(0.009,0.009) (0.246,0.577)rectangle++(0.009,0.009) (0.121,0.786)rectangle++(0.009,0.009)
  (0.058,0.891)rectangle++(0.009,0.009) (0.027,0.443)rectangle++(0.009,0.009) (0.508,0.175)rectangle++(0.009,0.009) (0.252,0.585)rectangle++(0.009,0.009) (0.124,0.79)rectangle++(0.009,0.009)
  (0.06,0.893)rectangle++(0.009,0.009) (0.028,0.444)rectangle++(0.009,0.009) (0.012,0.22)rectangle++(0.009,0.009) (0.004,0.108)rectangle++(0.009,0.009) (0,0.552)rectangle++(0.009,0.009)
  (-0.002,0.774)rectangle++(0.009,0.009) (0.496,0.307)rectangle++(0.009,0.009) (0.246,0.151)rectangle++(0.009,0.009) (0.121,0.073)rectangle++(0.009,0.009) (0.058,0.034)rectangle++(0.009,0.009)
  (0.027,0.515)rectangle++(0.009,0.009) (0.011,0.255)rectangle++(0.009,0.009) (0.003,0.125)rectangle++(0.009,0.009) (-0.001,0.06)rectangle++(0.009,0.009) (-0.003,0.528)rectangle++(0.009,0.009)
  (0.496,0.208)rectangle++(0.009,0.009) (0.246,0.102)rectangle++(0.009,0.009) (0.121,0.549)rectangle++(0.009,0.009) (0.058,0.272)rectangle++(0.009,0.009) (0.521,0.106)rectangle++(0.009,0.009)
  (0.258,0.551)rectangle++(0.009,0.009) (0.127,0.773)rectangle++(0.009,0.009) (0.061,0.384)rectangle++(0.009,0.009) (0.522,0.151)rectangle++(0.009,0.009) (0.706,0.058)rectangle++(0.009,0.009)
  (0.351,0.027)rectangle++(0.009,0.009) (0.173,0.011)rectangle++(0.009,0.009) (0.084,0.003)rectangle++(0.009,0.009) (0.531,-0.001)rectangle++(0.009,0.009) (0.71,-0.003)rectangle++(0.009,0.009)
  (0.353,0.496)rectangle++(0.009,0.009) (0.174,0.246)rectangle++(0.009,0.009) (0.085,0.121)rectangle++(0.009,0.009) (0.531,0.046)rectangle++(0.009,0.009) (0.71,0.016)rectangle++(0.009,0.009)
  (0.353,0.006)rectangle++(0.009,0.009) (0.638,0)rectangle++(0.009,0.009) (0.317,-0.002)rectangle++(0.009,0.009) (0.156,-0.003)rectangle++(0.009,0.009) (0.56,-0.004)rectangle++(0.009,0.009)
  (0.721,-0.004)rectangle++(0.009,0.009) (0.358,-0.004)rectangle++(0.009,0.009) (0.641,-0.004)rectangle++(0.009,0.009) (0.754,-0.004)rectangle++(0.009,0.009) (0.375,-0.004)rectangle++(0.009,0.009)
  (0.185,-0.004)rectangle++(0.009,0.009) (0.09,0.496)rectangle++(0.009,0.009) (0.533,0.196)rectangle++(0.009,0.009) (0.264,0.596)rectangle++(0.009,0.009) (0.13,0.296)rectangle++(0.009,0.009)
  (0.063,0.146)rectangle++(0.009,0.009) (0.029,0.571)rectangle++(0.009,0.009) (0.012,0.283)rectangle++(0.009,0.009) (0.004,0.139)rectangle++(0.009,0.009) (0,0.567)rectangle++(0.009,0.009)
  (-0.002,0.781)rectangle++(0.009,0.009) (-0.003,0.388)rectangle++(0.009,0.009) (0.496,0.153)rectangle++(0.009,0.009) (0.246,0.574)rectangle++(0.009,0.009) (0.596,0.227)rectangle++(0.009,0.009)
  (0.296,0.611)rectangle++(0.009,0.009) (0.146,0.303)rectangle++(0.009,0.009) (0.556,0.119)rectangle++(0.009,0.009) (0.276,0.557)rectangle++(0.009,0.009) (0.136,0.276)rectangle++(0.009,0.009)
  (0.552,0.108)rectangle++(0.009,0.009) (0.274,0.052)rectangle++(0.009,0.009) (0.607,0.018)rectangle++(0.009,0.009) (0.74,0.004)rectangle++(0.009,0.009) (0.368,0)rectangle++(0.009,0.009)
  (0.644,-0.003)rectangle++(0.009,0.009) (0.32,-0.004)rectangle++(0.009,0.009) (0.158,0.496)rectangle++(0.009,0.009) (0.077,0.246)rectangle++(0.009,0.009) (0.036,0.621)rectangle++(0.009,0.009)
  (0.016,0.808)rectangle++(0.009,0.009) (0.006,0.402)rectangle++(0.009,0.009) (0.001,0.699)rectangle++(0.009,0.009) (-0.002,0.847)rectangle++(0.009,0.009) (-0.003,0.921)rectangle++(0.009,0.009)
  (-0.004,0.958)rectangle++(0.009,0.009) (0.496,0.381)rectangle++(0.009,0.009) (0.246,0.188)rectangle++(0.009,0.009) (0.121,0.092)rectangle++(0.009,0.009) (0.058,0.044)rectangle++(0.009,0.009)
  (0.027,0.52)rectangle++(0.009,0.009) (0.011,0.758)rectangle++(0.009,0.009) (0.003,0.877)rectangle++(0.009,0.009) (-0.001,0.936)rectangle++(0.009,0.009) (-0.003,0.966)rectangle++(0.009,0.009)
  (-0.004,0.481)rectangle++(0.009,0.009) (-0.004,0.238)rectangle++(0.009,0.009) (-0.004,0.117)rectangle++(0.009,0.009) (-0.004,0.556)rectangle++(0.009,0.009) (-0.004,0.776)rectangle++(0.009,0.009)
  (0.496,0.308)rectangle++(0.009,0.009) (0.246,0.652)rectangle++(0.009,0.009) (0.596,0.258)rectangle++(0.009,0.009) (0.736,0.1)rectangle++(0.009,0.009) (0.366,0.548)rectangle++(0.009,0.009)
  (0.181,0.772)rectangle++(0.009,0.009) (0.088,0.884)rectangle++(0.009,0.009) (0.533,0.351)rectangle++(0.009,0.009) (0.264,0.173)rectangle++(0.009,0.009) (0.603,0.067)rectangle++(0.009,0.009)
  (0.299,0.031)rectangle++(0.009,0.009) (0.147,0.513)rectangle++(0.009,0.009) (0.556,0.203)rectangle++(0.009,0.009) (0.276,0.599)rectangle++(0.009,0.009) (0.136,0.797)rectangle++(0.009,0.009)
  (0.066,0.396)rectangle++(0.009,0.009) (0.031,0.196)rectangle++(0.009,0.009) (0.51,0.076)rectangle++(0.009,0.009) (0.253,0.036)rectangle++(0.009,0.009) (0.124,0.016)rectangle++(0.009,0.009)
  (0.547,0.004)rectangle++(0.009,0.009) (0.271,0)rectangle++(0.009,0.009) (0.133,-0.002)rectangle++(0.009,0.009) (0.064,-0.003)rectangle++(0.009,0.009) (0.03,0.496)rectangle++(0.009,0.009)
  (0.013,0.746)rectangle++(0.009,0.009) (0.004,0.371)rectangle++(0.009,0.009) (0,0.683)rectangle++(0.009,0.009) (0.497,0.271)rectangle++(0.009,0.009) (0.246,0.133)rectangle++(0.009,0.009)
  (0.596,0.051)rectangle++(0.009,0.009) (0.736,0.018)rectangle++(0.009,0.009) (0.366,0.507)rectangle++(0.009,0.009) (0.644,0.2)rectangle++(0.009,0.009) (0.32,0.598)rectangle++(0.009,0.009)
  (0.625,0.236)rectangle++(0.009,0.009) (0.31,0.616)rectangle++(0.009,0.009) (0.153,0.306)rectangle++(0.009,0.009) (0.074,0.651)rectangle++(0.009,0.009) (0.035,0.323)rectangle++(0.009,0.009)
  (0.015,0.159)rectangle++(0.009,0.009) (0.005,0.577)rectangle++(0.009,0.009) (0,0.286)rectangle++(0.009,0.009) (0.497,0.112)rectangle++(0.009,0.009) (0.246,0.554)rectangle++(0.009,0.009)
  (0.121,0.775)rectangle++(0.009,0.009) (0.058,0.885)rectangle++(0.009,0.009) (0.027,0.44)rectangle++(0.009,0.009) (0.011,0.718)rectangle++(0.009,0.009) (0.003,0.857)rectangle++(0.009,0.009)
  (-0.001,0.926)rectangle++(0.009,0.009) (0.497,0.368)rectangle++(0.009,0.009) (0.246,0.682)rectangle++(0.009,0.009) (0.121,0.839)rectangle++(0.009,0.009) (0.058,0.417)rectangle++(0.009,0.009)
  (0.027,0.206)rectangle++(0.009,0.009) (0.011,0.101)rectangle++(0.009,0.009) (0.003,0.048)rectangle++(0.009,0.009) (-0.001,0.022)rectangle++(0.009,0.009) (-0.003,0.009)rectangle++(0.009,0.009)
  (0.496,0.001)rectangle++(0.009,0.009) (0.696,-0.002)rectangle++(0.009,0.009) (0.346,-0.003)rectangle++(0.009,0.009) (0.171,0.496)rectangle++(0.009,0.009) (0.083,0.746)rectangle++(0.009,0.009)
  (0.039,0.871)rectangle++(0.009,0.009) (0.017,0.433)rectangle++(0.009,0.009) (0.006,0.214)rectangle++(0.009,0.009) (0.001,0.105)rectangle++(0.009,0.009) (0.498,0.039)rectangle++(0.009,0.009)
  (0.247,0.517)rectangle++(0.009,0.009) (0.121,0.756)rectangle++(0.009,0.009) (0.546,0.3)rectangle++(0.009,0.009) (0.271,0.148)rectangle++(0.009,0.009) (0.133,0.572)rectangle++(0.009,0.009)
  (0.551,0.226)rectangle++(0.009,0.009) (0.273,0.611)rectangle++(0.009,0.009) (0.134,0.803)rectangle++(0.009,0.009) (0.551,0.319)rectangle++(0.009,0.009) (0.718,0.125)rectangle++(0.009,0.009)
  (0.357,0.56)rectangle++(0.009,0.009) (0.176,0.278)rectangle++(0.009,0.009) (0.086,0.637)rectangle++(0.009,0.009) (0.041,0.816)rectangle++(0.009,0.009) (0.514,0.324)rectangle++(0.009,0.009)
  (0.703,0.127)rectangle++(0.009,0.009) (0.349,0.061)rectangle++(0.009,0.009) (0.637,0.022)rectangle++(0.009,0.009) (0.316,0.509)rectangle++(0.009,0.009) (0.624,0.201)rectangle++(0.009,0.009)
  (0.31,0.598)rectangle++(0.009,0.009) (0.153,0.797)rectangle++(0.009,0.009) (0.074,0.396)rectangle++(0.009,0.009) (0.035,0.196)rectangle++(0.009,0.009) (0.511,0.076)rectangle++(0.009,0.009)
  (0.253,0.536)rectangle++(0.009,0.009) (0.599,0.212)rectangle++(0.009,0.009) (0.297,0.604)rectangle++(0.009,0.009) (0.616,0.239)rectangle++(0.009,0.009) (0.306,0.617)rectangle++(0.009,0.009)
  (0.151,0.806)rectangle++(0.009,0.009) (0.073,0.901)rectangle++(0.009,0.009) (0.034,0.448)rectangle++(0.009,0.009) (0.015,0.722)rectangle++(0.009,0.009) (0.005,0.359)rectangle++(0.009,0.009)
  (0,0.177)rectangle++(0.009,0.009) (-0.002,0.086)rectangle++(0.009,0.009) (-0.003,0.541)rectangle++(0.009,0.009) (-0.004,0.768)rectangle++(0.009,0.009) (-0.004,0.882)rectangle++(0.009,0.009)
  (0.496,0.35)rectangle++(0.009,0.009) (0.246,0.673)rectangle++(0.009,0.009) (0.121,0.334)rectangle++(0.009,0.009) (0.058,0.665)rectangle++(0.009,0.009) (0.027,0.83)rectangle++(0.009,0.009)
  (0.011,0.913)rectangle++(0.009,0.009) (0.502,0.362)rectangle++(0.009,0.009) (0.249,0.679)rectangle++(0.009,0.009) (0.122,0.837)rectangle++(0.009,0.009) (0.059,0.416)rectangle++(0.009,0.009)
  (0.027,0.706)rectangle++(0.009,0.009) (0.011,0.851)rectangle++(0.009,0.009) (0.003,0.923)rectangle++(0.009,0.009) (-0.001,0.459)rectangle++(0.009,0.009) (0.497,0.181)rectangle++(0.009,0.009)
  (0.246,0.088)rectangle++(0.009,0.009) (0.121,0.542)rectangle++(0.009,0.009) (0.058,0.269)rectangle++(0.009,0.009) (0.027,0.632)rectangle++(0.009,0.009) (0.011,0.814)rectangle++(0.009,0.009)
  (0.003,0.905)rectangle++(0.009,0.009) (-0.001,0.45)rectangle++(0.009,0.009) (-0.003,0.723)rectangle++(0.009,0.009) (-0.004,0.859)rectangle++(0.009,0.009) (0.496,0.341)rectangle++(0.009,0.009)
  (0.246,0.668)rectangle++(0.009,0.009) (0.121,0.332)rectangle++(0.009,0.009) (0.546,0.13)rectangle++(0.009,0.009) (0.271,0.563)rectangle++(0.009,0.009) (0.133,0.779)rectangle++(0.009,0.009)
  (0.551,0.309)rectangle++(0.009,0.009) (0.718,0.121)rectangle++(0.009,0.009) (0.357,0.058)rectangle++(0.009,0.009) (0.64,0.021)rectangle++(0.009,0.009) (0.318,0.008)rectangle++(0.009,0.009)
  (0.157,0.002)rectangle++(0.009,0.009) (0.076,0.499)rectangle++(0.009,0.009) (0.036,0.747)rectangle++(0.009,0.009) (0.016,0.371)rectangle++(0.009,0.009) (0.006,0.183)rectangle++(0.009,0.009)
  (0.001,0.589)rectangle++(0.009,0.009) (-0.002,0.292)rectangle++(0.009,0.009) (-0.003,0.144)rectangle++(0.009,0.009) (-0.004,0.07)rectangle++(0.009,0.009) (0.496,0.025)rectangle++(0.009,0.009)
  (0.246,0.01)rectangle++(0.009,0.009) (0.121,0.003)rectangle++(0.009,0.009) (0.058,0.499)rectangle++(0.009,0.009) (0.027,0.747)rectangle++(0.009,0.009) (0.011,0.371)rectangle++(0.009,0.009)
  (0.003,0.683)rectangle++(0.009,0.009) (-0.001,0.339)rectangle++(0.009,0.009) (-0.003,0.667)rectangle++(0.009,0.009) (0.496,0.264)rectangle++(0.009,0.009) (0.246,0.63)rectangle++(0.009,0.009)
  (0.596,0.249)rectangle++(0.009,0.009) (0.736,0.097)rectangle++(0.009,0.009) (0.366,0.546)rectangle++(0.009,0.009) (0.644,0.216)rectangle++(0.009,0.009) (0.32,0.106)rectangle++(0.009,0.009)
  (0.158,0.051)rectangle++(0.009,0.009) (0.077,0.023)rectangle++(0.009,0.009) (0.528,0.007)rectangle++(0.009,0.009) (0.262,0.001)rectangle++(0.009,0.009) (0.129,0.498)rectangle++(0.009,0.009)
  (0.062,0.747)rectangle++(0.009,0.009) (0.029,0.371)rectangle++(0.009,0.009) (0.012,0.183)rectangle++(0.009,0.009) (0.004,0.589)rectangle++(0.009,0.009) (0,0.792)rectangle++(0.009,0.009)
  (-0.002,0.894)rectangle++(0.009,0.009) (-0.003,0.945)rectangle++(0.009,0.009) (0.496,0.375)rectangle++(0.009,0.009) (0.246,0.685)rectangle++(0.009,0.009) (0.121,0.34)rectangle++(0.009,0.009)
  (0.058,0.168)rectangle++(0.009,0.009) (0.027,0.582)rectangle++(0.009,0.009) (0.011,0.789)rectangle++(0.009,0.009) (0.502,0.313)rectangle++(0.009,0.009) (0.249,0.154)rectangle++(0.009,0.009)
  (0.597,0.059)rectangle++(0.009,0.009) (0.296,0.527)rectangle++(0.009,0.009) (0.146,0.261)rectangle++(0.009,0.009) (0.071,0.628)rectangle++(0.009,0.009) (0.033,0.812)rectangle++(0.009,0.009)
  (0.014,0.904)rectangle++(0.009,0.009) (0.005,0.45)rectangle++(0.009,0.009) (0,0.723)rectangle++(0.009,0.009) (-0.002,0.859)rectangle++(0.009,0.009) (-0.003,0.927)rectangle++(0.009,0.009)
  (0.496,0.368)rectangle++(0.009,0.009) (0.246,0.682)rectangle++(0.009,0.009) (0.121,0.839)rectangle++(0.009,0.009) (0.058,0.917)rectangle++(0.009,0.009) (0.027,0.456)rectangle++(0.009,0.009)
  (0.011,0.726)rectangle++(0.009,0.009) (0.003,0.861)rectangle++(0.009,0.009) (-0.001,0.928)rectangle++(0.009,0.009) (-0.003,0.962)rectangle++(0.009,0.009) (-0.004,0.979)rectangle++(0.009,0.009)
  (0.496,0.389)rectangle++(0.009,0.009) (0.696,0.153)rectangle++(0.009,0.009) (0.776,0.058)rectangle++(0.009,0.009) (0.386,0.527)rectangle++(0.009,0.009) (0.191,0.261)rectangle++(0.009,0.009)
  (0.093,0.128)rectangle++(0.009,0.009) (0.044,0.062)rectangle++(0.009,0.009) (0.02,0.529)rectangle++(0.009,0.009) (0.008,0.262)rectangle++(0.009,0.009) (0.002,0.629)rectangle++(0.009,0.009)
  (-0.001,0.312)rectangle++(0.009,0.009) (-0.003,0.654)rectangle++(0.009,0.009) (-0.004,0.325)rectangle++(0.009,0.009) (-0.004,0.66)rectangle++(0.009,0.009) (-0.004,0.828)rectangle++(0.009,0.009)
  (0.496,0.328)rectangle++(0.009,0.009) (0.246,0.662)rectangle++(0.009,0.009) (0.121,0.329)rectangle++(0.009,0.009) (0.058,0.162)rectangle++(0.009,0.009) (0.027,0.579)rectangle++(0.009,0.009)
  (0.011,0.287)rectangle++(0.009,0.009) (0.003,0.641)rectangle++(0.009,0.009) (-0.001,0.818)rectangle++(0.009,0.009) (-0.003,0.407)rectangle++(0.009,0.009) (-0.004,0.701)rectangle++(0.009,0.009)
  (-0.004,0.348)rectangle++(0.009,0.009) (-0.004,0.672)rectangle++(0.009,0.009) (0.496,0.266)rectangle++(0.009,0.009) (0.246,0.131)rectangle++(0.009,0.009) (0.596,0.05)rectangle++(0.009,0.009)
  (0.296,0.523)rectangle++(0.009,0.009) (0.616,0.206)rectangle++(0.009,0.009) (0.744,0.08)rectangle++(0.009,0.009) (0.37,0.538)rectangle++(0.009,0.009) (0.645,0.212)rectangle++(0.009,0.009)
  (0.32,0.104)rectangle++(0.009,0.009) (0.625,0.039)rectangle++(0.009,0.009) (0.31,0.017)rectangle++(0.009,0.009) (0.153,0.506)rectangle++(0.009,0.009) (0.074,0.251)rectangle++(0.009,0.009)
  (0.527,0.098)rectangle++(0.009,0.009) (0.708,0.036)rectangle++(0.009,0.009) (0.781,0.012)rectangle++(0.009,0.009) (0.388,0.004)rectangle++(0.009,0.009) (0.192,0.5)rectangle++(0.009,0.009)
  (0.094,0.248)rectangle++(0.009,0.009) (0.045,0.122)rectangle++(0.009,0.009) (0.02,0.559)rectangle++(0.009,0.009) (0.505,0.221)rectangle++(0.009,0.009) (0.25,0.108)rectangle++(0.009,0.009)
  (0.123,0.052)rectangle++(0.009,0.009) (0.059,0.524)rectangle++(0.009,0.009) (0.521,0.207)rectangle++(0.009,0.009) (0.258,0.601)rectangle++(0.009,0.009) (0.127,0.798)rectangle++(0.009,0.009)
  (0.061,0.397)rectangle++(0.009,0.009) (0.522,0.156)rectangle++(0.009,0.009) (0.259,0.076)rectangle++(0.009,0.009) (0.601,0.028)rectangle++(0.009,0.009) (0.298,0.512)rectangle++(0.009,0.009)
  (0.147,0.254)rectangle++(0.009,0.009) (0.071,0.625)rectangle++(0.009,0.009) (0.526,0.247)rectangle++(0.009,0.009) (0.708,0.096)rectangle++(0.009,0.009) (0.352,0.046)rectangle++(0.009,0.009)
  (0.638,0.016)rectangle++(0.009,0.009) (0.752,0.004)rectangle++(0.009,0.009) (0.374,0)rectangle++(0.009,0.009) (0.185,0.498)rectangle++(0.009,0.009) (0.09,0.747)rectangle++(0.009,0.009)
  (0.043,0.371)rectangle++(0.009,0.009) (0.514,0.146)rectangle++(0.009,0.009) (0.255,0.071)rectangle++(0.009,0.009) (0.599,0.026)rectangle++(0.009,0.009) (0.297,0.511)rectangle++(0.009,0.009)
  (0.146,0.753)rectangle++(0.009,0.009) (0.071,0.374)rectangle++(0.009,0.009) (0.033,0.685)rectangle++(0.009,0.009) (0.511,0.271)rectangle++(0.009,0.009) (0.253,0.133)rectangle++(0.009,0.009)
  (0.124,0.064)rectangle++(0.009,0.009) (0.547,0.023)rectangle++(0.009,0.009) (0.271,0.009)rectangle++(0.009,0.009) (0.133,0.502)rectangle++(0.009,0.009) (0.064,0.249)rectangle++(0.009,0.009)
  (0.03,0.122)rectangle++(0.009,0.009) (0.013,0.059)rectangle++(0.009,0.009) (0.004,0.027)rectangle++(0.009,0.009) (0.499,0.008)rectangle++(0.009,0.009) (0.247,0.002)rectangle++(0.009,0.009)
  (0.121,-0.001)rectangle++(0.009,0.009) (0.058,-0.003)rectangle++(0.009,0.009) (0.027,0.496)rectangle++(0.009,0.009) (0.508,0.196)rectangle++(0.009,0.009) (0.701,0.076)rectangle++(0.009,0.009)
  (0.348,0.536)rectangle++(0.009,0.009) (0.172,0.266)rectangle++(0.009,0.009) (0.084,0.131)rectangle++(0.009,0.009) (0.04,0.563)rectangle++(0.009,0.009) (0.018,0.279)rectangle++(0.009,0.009)
  (0.007,0.637)rectangle++(0.009,0.009) (0.5,0.252)rectangle++(0.009,0.009) (0.697,0.098)rectangle++(0.009,0.009) (0.346,0.047)rectangle++(0.009,0.009) (0.636,0.016)rectangle++(0.009,0.009)
  (0.752,0.004)rectangle++(0.009,0.009) (0.374,0.5)rectangle++(0.009,0.009) (0.185,0.248)rectangle++(0.009,0.009) (0.571,0.096)rectangle++(0.009,0.009) (0.283,0.046)rectangle++(0.009,0.009)
  (0.611,0.016)rectangle++(0.009,0.009) (0.303,0.506)rectangle++(0.009,0.009) (0.149,0.751)rectangle++(0.009,0.009) (0.072,0.373)rectangle++(0.009,0.009) (0.526,0.147)rectangle++(0.009,0.009)
  (0.708,0.056)rectangle++(0.009,0.009) (0.352,0.526)rectangle++(0.009,0.009) (0.174,0.261)rectangle++(0.009,0.009) (0.085,0.128)rectangle++(0.009,0.009) (0.04,0.062)rectangle++(0.009,0.009)
  (0.018,0.029)rectangle++(0.009,0.009) (0.007,0.512)rectangle++(0.009,0.009) (0.001,0.254)rectangle++(0.009,0.009) (-0.002,0.625)rectangle++(0.009,0.009) (-0.003,0.81)rectangle++(0.009,0.009)
  (-0.004,0.903)rectangle++(0.009,0.009) (0.496,0.358)rectangle++(0.009,0.009) (0.246,0.677)rectangle++(0.009,0.009) (0.121,0.836)rectangle++(0.009,0.009) (0.058,0.416)rectangle++(0.009,0.009)
  (0.027,0.706)rectangle++(0.009,0.009) (0.011,0.351)rectangle++(0.009,0.009) (0.502,0.138)rectangle++(0.009,0.009) (0.698,0.052)rectangle++(0.009,0.009) (0.347,0.524)rectangle++(0.009,0.009)
  (0.171,0.76)rectangle++(0.009,0.009) (0.083,0.878)rectangle++(0.009,0.009) (0.039,0.437)rectangle++(0.009,0.009) (0.017,0.216)rectangle++(0.009,0.009) (0.006,0.606)rectangle++(0.009,0.009)
  (0.5,0.24)rectangle++(0.009,0.009) (0.248,0.618)rectangle++(0.009,0.009) (0.122,0.307)rectangle++(0.009,0.009) (0.546,0.12)rectangle++(0.009,0.009) (0.271,0.058)rectangle++(0.009,0.009)
  (0.606,0.02)rectangle++(0.009,0.009) (0.301,0.508)rectangle++(0.009,0.009) (0.148,0.752)rectangle++(0.009,0.009) (0.072,0.874)rectangle++(0.009,0.009) (0.034,0.435)rectangle++(0.009,0.009)
  (0.015,0.715)rectangle++(0.009,0.009) (0.005,0.855)rectangle++(0.009,0.009) (0,0.425)rectangle++(0.009,0.009) (-0.002,0.71)rectangle++(0.009,0.009) (0.496,0.281)rectangle++(0.009,0.009)
  (0.246,0.138)rectangle++(0.009,0.009) (0.121,0.067)rectangle++(0.009,0.009) (0.058,0.531)rectangle++(0.009,0.009) (0.521,0.21)rectangle++(0.009,0.009) (0.258,0.103)rectangle++(0.009,0.009)
  (0.127,0.549)rectangle++(0.009,0.009) (0.061,0.272)rectangle++(0.009,0.009) (0.028,0.134)rectangle++(0.009,0.009) (0.509,0.051)rectangle++(0.009,0.009) (0.252,0.023)rectangle++(0.009,0.009)
  (0.124,0.009)rectangle++(0.009,0.009) (0.06,0.002)rectangle++(0.009,0.009) (0.028,0.499)rectangle++(0.009,0.009) (0.012,0.747)rectangle++(0.009,0.009) (0.004,0.871)rectangle++(0.009,0.009)
  (0,0.433)rectangle++(0.009,0.009) (0.497,0.171)rectangle++(0.009,0.009) (0.696,0.066)rectangle++(0.009,0.009) (0.776,0.024)rectangle++(0.009,0.009) (0.386,0.51)rectangle++(0.009,0.009)
  (0.652,0.201)rectangle++(0.009,0.009) (0.324,0.598)rectangle++(0.009,0.009) (0.16,0.297)rectangle++(0.009,0.009) (0.078,0.646)rectangle++(0.009,0.009) (0.037,0.821)rectangle++(0.009,0.009)
  (0.016,0.408)rectangle++(0.009,0.009) (0.006,0.202)rectangle++(0.009,0.009) (0.001,0.599)rectangle++(0.009,0.009) (0.498,0.237)rectangle++(0.009,0.009) (0.247,0.616)rectangle++(0.009,0.009)
  (0.121,0.306)rectangle++(0.009,0.009) (0.058,0.151)rectangle++(0.009,0.009) (0.027,0.073)rectangle++(0.009,0.009) (0.011,0.034)rectangle++(0.009,0.009) (0.502,0.011)rectangle++(0.009,0.009)
  (0.249,0.503)rectangle++(0.009,0.009) (0.597,0.199)rectangle++(0.009,0.009) (0.296,0.097)rectangle++(0.009,0.009) (0.146,0.546)rectangle++(0.009,0.009) (0.556,0.216)rectangle++(0.009,0.009)
  (0.276,0.106)rectangle++(0.009,0.009) (0.136,0.051)rectangle++(0.009,0.009) (0.552,0.018)rectangle++(0.009,0.009) (0.718,0.004)rectangle++(0.009,0.009) (0.357,0)rectangle++(0.009,0.009)
  (0.176,0.498)rectangle++(0.009,0.009) (0.568,0.196)rectangle++(0.009,0.009) (0.282,0.096)rectangle++(0.009,0.009) (0.139,0.546)rectangle++(0.009,0.009) (0.067,0.271)rectangle++(0.009,0.009)
  (0.031,0.633)rectangle++(0.009,0.009) (0.013,0.814)rectangle++(0.009,0.009) (0.004,0.405)rectangle++(0.009,0.009) (0,0.2)rectangle++(0.009,0.009) (-0.002,0.098)rectangle++(0.009,0.009)
  (-0.003,0.047)rectangle++(0.009,0.009) (0.496,0.016)rectangle++(0.009,0.009) (0.246,0.506)rectangle++(0.009,0.009) (0.121,0.751)rectangle++(0.009,0.009) (0.058,0.873)rectangle++(0.009,0.009)
  (0.027,0.934)rectangle++(0.009,0.009) (0.011,0.465)rectangle++(0.009,0.009) (0.502,0.183)rectangle++(0.009,0.009) (0.249,0.589)rectangle++(0.009,0.009) (0.122,0.292)rectangle++(0.009,0.009)
  (0.059,0.644)rectangle++(0.009,0.009) (0.521,0.255)rectangle++(0.009,0.009) (0.258,0.625)rectangle++(0.009,0.009) (0.127,0.31)rectangle++(0.009,0.009) (0.548,0.121)rectangle++(0.009,0.009)
  (0.717,0.046)rectangle++(0.009,0.009) (0.356,0.021)rectangle++(0.009,0.009) (0.64,0.006)rectangle++(0.009,0.009) (0.318,0.501)rectangle++(0.009,0.009) (0.157,0.748)rectangle++(0.009,0.009)
  (0.076,0.872)rectangle++(0.009,0.009) (0.036,0.934)rectangle++(0.009,0.009) (0.016,0.965)rectangle++(0.009,0.009) (0.006,0.98)rectangle++(0.009,0.009) (0.001,0.988)rectangle++(0.009,0.009)
  (-0.002,0.992)rectangle++(0.009,0.009) (-0.003,0.994)rectangle++(0.009,0.009) (-0.004,0.995)rectangle++(0.009,0.009) (-0.004,0.495)rectangle++(0.009,0.009) (-0.004,0.745)rectangle++(0.009,0.009)
  (-0.004,0.87)rectangle++(0.009,0.009) (-0.004,0.933)rectangle++(0.009,0.009) (-0.004,0.464)rectangle++(0.009,0.009) (-0.004,0.73)rectangle++(0.009,0.009) (-0.004,0.863)rectangle++(0.009,0.009)
  (-0.004,0.929)rectangle++(0.009,0.009) (-0.004,0.962)rectangle++(0.009,0.009) (-0.004,0.979)rectangle++(0.009,0.009) (0.496,0.389)rectangle++(0.009,0.009) (0.246,0.192)rectangle++(0.009,0.009)
  (0.596,0.074)rectangle++(0.009,0.009) (0.296,0.035)rectangle++(0.009,0.009) (0.146,0.515)rectangle++(0.009,0.009) (0.071,0.755)rectangle++(0.009,0.009) (0.526,0.299)rectangle++(0.009,0.009)
  (0.261,0.647)rectangle++(0.009,0.009) (0.128,0.821)rectangle++(0.009,0.009) (0.062,0.908)rectangle++(0.009,0.009) (0.029,0.452)rectangle++(0.009,0.009) (0.012,0.724)rectangle++(0.009,0.009)
  (0.004,0.86)rectangle++(0.009,0.009) (0,0.428)rectangle++(0.009,0.009) (-0.002,0.712)rectangle++(0.009,0.009) (-0.003,0.854)rectangle++(0.009,0.009) (0.496,0.339)rectangle++(0.009,0.009)
  (0.696,0.133)rectangle++(0.009,0.009) (0.776,0.05)rectangle++(0.009,0.009) (0.808,0.017)rectangle++(0.009,0.009) (0.402,0.506)rectangle++(0.009,0.009) (0.199,0.751)rectangle++(0.009,0.009)
  (0.097,0.373)rectangle++(0.009,0.009) (0.046,0.684)rectangle++(0.009,0.009) (0.021,0.34)rectangle++(0.009,0.009) (0.008,0.668)rectangle++(0.009,0.009) (0.002,0.332)rectangle++(0.009,0.009)
  (0.498,0.13)rectangle++(0.009,0.009) (0.247,0.063)rectangle++(0.009,0.009) (0.121,0.529)rectangle++(0.009,0.009) (0.546,0.209)rectangle++(0.009,0.009) (0.271,0.602)rectangle++(0.009,0.009)
  (0.133,0.799)rectangle++(0.009,0.009) (0.064,0.897)rectangle++(0.009,0.009) (0.03,0.946)rectangle++(0.009,0.009) (0.013,0.471)rectangle++(0.009,0.009) (0.004,0.233)rectangle++(0.009,0.009)
  (0,0.114)rectangle++(0.009,0.009) (-0.002,0.555)rectangle++(0.009,0.009) (0.496,0.219)rectangle++(0.009,0.009) (0.696,0.085)rectangle++(0.009,0.009) (0.346,0.04)rectangle++(0.009,0.009)
  (0.171,0.018)rectangle++(0.009,0.009) (0.083,0.507)rectangle++(0.009,0.009) (0.039,0.751)rectangle++(0.009,0.009) (0.017,0.873)rectangle++(0.009,0.009) (0.504,0.347)rectangle++(0.009,0.009)
  (0.25,0.171)rectangle++(0.009,0.009) (0.597,0.066)rectangle++(0.009,0.009) (0.296,0.031)rectangle++(0.009,0.009) (0.146,0.513)rectangle++(0.009,0.009) (0.071,0.254)rectangle++(0.009,0.009)
  (0.033,0.625)rectangle++(0.009,0.009) (0.014,0.31)rectangle++(0.009,0.009) (0.005,0.653)rectangle++(0.009,0.009) (0,0.324)rectangle++(0.009,0.009) (0.497,0.127)rectangle++(0.009,0.009)
  (0.246,0.561)rectangle++(0.009,0.009) (0.121,0.778)rectangle++(0.009,0.009) (0.058,0.387)rectangle++(0.009,0.009) (0.027,0.691)rectangle++(0.009,0.009) (0.011,0.843)rectangle++(0.009,0.009)
  (0.003,0.419)rectangle++(0.009,0.009) (-0.001,0.707)rectangle++(0.009,0.009) (-0.003,0.851)rectangle++(0.009,0.009) (0.496,0.338)rectangle++(0.009,0.009) (0.246,0.167)rectangle++(0.009,0.009)
  (0.596,0.064)rectangle++(0.009,0.009) (0.296,0.53)rectangle++(0.009,0.009) (0.146,0.763)rectangle++(0.009,0.009) (0.071,0.379)rectangle++(0.009,0.009) (0.526,0.149)rectangle++(0.009,0.009)
  (0.261,0.072)rectangle++(0.009,0.009) (0.128,0.534)rectangle++(0.009,0.009) (0.062,0.765)rectangle++(0.009,0.009) (0.522,0.303)rectangle++(0.009,0.009) (0.259,0.649)rectangle++(0.009,0.009)
  (0.127,0.322)rectangle++(0.009,0.009) (0.061,0.659)rectangle++(0.009,0.009) (0.028,0.827)rectangle++(0.009,0.009) (0.012,0.411)rectangle++(0.009,0.009) (0.004,0.703)rectangle++(0.009,0.009)
  (0,0.349)rectangle++(0.009,0.009) (-0.002,0.172)rectangle++(0.009,0.009) (-0.003,0.584)rectangle++(0.009,0.009) (-0.004,0.29)rectangle++(0.009,0.009) (-0.004,0.143)rectangle++(0.009,0.009)
  (-0.004,0.569)rectangle++(0.009,0.009) (-0.004,0.782)rectangle++(0.009,0.009) (0.496,0.31)rectangle++(0.009,0.009) (0.246,0.653)rectangle++(0.009,0.009) (0.121,0.324)rectangle++(0.009,0.009)
  (0.058,0.66)rectangle++(0.009,0.009) (0.027,0.328)rectangle++(0.009,0.009) (0.011,0.162)rectangle++(0.009,0.009) (0.502,0.062)rectangle++(0.009,0.009) (0.249,0.029)rectangle++(0.009,0.009)
  (0.122,0.012)rectangle++(0.009,0.009) (0.059,0.504)rectangle++(0.009,0.009) (0.027,0.75)rectangle++(0.009,0.009) (0.011,0.873)rectangle++(0.009,0.009) (0.003,0.434)rectangle++(0.009,0.009)
  (-0.001,0.215)rectangle++(0.009,0.009) (-0.003,0.105)rectangle++(0.009,0.009) (-0.004,0.05)rectangle++(0.009,0.009) (-0.004,0.023)rectangle++(0.009,0.009) (-0.004,0.509)rectangle++(0.009,0.009)
  (-0.004,0.752)rectangle++(0.009,0.009) (-0.004,0.874)rectangle++(0.009,0.009) (-0.004,0.435)rectangle++(0.009,0.009) (-0.004,0.715)rectangle++(0.009,0.009) (-0.004,0.355)rectangle++(0.009,0.009)
  (0.496,0.139)rectangle++(0.009,0.009) (0.696,0.053)rectangle++(0.009,0.009) (0.346,0.524)rectangle++(0.009,0.009) (0.171,0.26)rectangle++(0.009,0.009) (0.083,0.628)rectangle++(0.009,0.009)
  (0.531,0.248)rectangle++(0.009,0.009) (0.263,0.122)rectangle++(0.009,0.009) (0.129,0.059)rectangle++(0.009,0.009) (0.062,0.527)rectangle++(0.009,0.009) (0.029,0.261)rectangle++(0.009,0.009)
  (0.012,0.128)rectangle++(0.009,0.009) (0.004,0.562)rectangle++(0.009,0.009) (0,0.779)rectangle++(0.009,0.009) (-0.002,0.887)rectangle++(0.009,0.009) (-0.003,0.441)rectangle++(0.009,0.009)
  (-0.004,0.718)rectangle++(0.009,0.009) (-0.004,0.357)rectangle++(0.009,0.009) (-0.004,0.676)rectangle++(0.009,0.009) (-0.004,0.836)rectangle++(0.009,0.009) (-0.004,0.416)rectangle++(0.009,0.009)
  (-0.004,0.706)rectangle++(0.009,0.009) (0.496,0.28)rectangle++(0.009,0.009) (0.246,0.638)rectangle++(0.009,0.009) (0.121,0.817)rectangle++(0.009,0.009) (0.546,0.324)rectangle++(0.009,0.009)
  (0.716,0.127)rectangle++(0.009,0.009) (0.356,0.061)rectangle++(0.009,0.009) (0.64,0.022)rectangle++(0.009,0.009) (0.318,0.509)rectangle++(0.009,0.009) (0.157,0.752)rectangle++(0.009,0.009)
  (0.076,0.374)rectangle++(0.009,0.009) (0.036,0.185)rectangle++(0.009,0.009) (0.512,0.071)rectangle++(0.009,0.009) (0.254,0.533)rectangle++(0.009,0.009) (0.125,0.264)rectangle++(0.009,0.009)
  (0.547,0.103)rectangle++(0.009,0.009) (0.271,0.049)rectangle++(0.009,0.009) (0.133,0.522)rectangle++(0.009,0.009) (0.064,0.259)rectangle++(0.009,0.009) (0.03,0.127)rectangle++(0.009,0.009)
  (0.013,0.561)rectangle++(0.009,0.009) (0.004,0.778)rectangle++(0.009,0.009) (0,0.387)rectangle++(0.009,0.009) (-0.002,0.691)rectangle++(0.009,0.009) (-0.003,0.843)rectangle++(0.009,0.009)
  (0.496,0.335)rectangle++(0.009,0.009) (0.246,0.165)rectangle++(0.009,0.009) (0.121,0.08)rectangle++(0.009,0.009) (0.058,0.038)rectangle++(0.009,0.009) (0.521,0.012)rectangle++(0.009,0.009)
  (0.258,0.504)rectangle++(0.009,0.009) (0.127,0.25)rectangle++(0.009,0.009) (0.061,0.623)rectangle++(0.009,0.009) (0.028,0.809)rectangle++(0.009,0.009) (0.012,0.402)rectangle++(0.009,0.009)
  (0.004,0.199)rectangle++(0.009,0.009) (0,0.597)rectangle++(0.009,0.009) (-0.002,0.796)rectangle++(0.009,0.009) (-0.003,0.396)rectangle++(0.009,0.009) (-0.004,0.196)rectangle++(0.009,0.009)
  (-0.004,0.596)rectangle++(0.009,0.009) (-0.004,0.796)rectangle++(0.009,0.009) (-0.004,0.396)rectangle++(0.009,0.009) (-0.004,0.696)rectangle++(0.009,0.009) (-0.004,0.846)rectangle++(0.009,0.009)
  (-0.004,0.921)rectangle++(0.009,0.009) (0.496,0.366)rectangle++(0.009,0.009) (0.246,0.181)rectangle++(0.009,0.009) (0.121,0.088)rectangle++(0.009,0.009) (0.058,0.042)rectangle++(0.009,0.009)
  (0.027,0.519)rectangle++(0.009,0.009) (0.011,0.257)rectangle++(0.009,0.009) (0.003,0.626)rectangle++(0.009,0.009) (-0.001,0.811)rectangle++(0.009,0.009) (0.497,0.322)rectangle++(0.009,0.009)
  (0.246,0.659)rectangle++(0.009,0.009) (0.121,0.327)rectangle++(0.009,0.009) (0.058,0.661)rectangle++(0.009,0.009) (0.027,0.328)rectangle++(0.009,0.009) (0.508,0.129)rectangle++(0.009,0.009)
  (0.252,0.562)rectangle++(0.009,0.009) (0.598,0.222)rectangle++(0.009,0.009) (0.297,0.609)rectangle++(0.009,0.009) (0.146,0.802)rectangle++(0.009,0.009) (0.071,0.899)rectangle++(0.009,0.009)
  (0.033,0.947)rectangle++(0.009,0.009) (0.014,0.471)rectangle++(0.009,0.009) (0.503,0.186)rectangle++(0.009,0.009) (0.249,0.091)rectangle++(0.009,0.009) (0.122,0.543)rectangle++(0.009,0.009)
  (0.546,0.215)rectangle++(0.009,0.009) (0.716,0.083)rectangle++(0.009,0.009) (0.356,0.039)rectangle++(0.009,0.009) (0.64,0.013)rectangle++(0.009,0.009) (0.753,0.003)rectangle++(0.009,0.009)
  (0.374,-0.001)rectangle++(0.009,0.009) (0.185,-0.003)rectangle++(0.009,0.009) (0.571,-0.004)rectangle++(0.009,0.009) (0.283,0.496)rectangle++(0.009,0.009) (0.139,0.246)rectangle++(0.009,0.009)
  (0.067,0.621)rectangle++(0.009,0.009) (0.031,0.308)rectangle++(0.009,0.009) (0.013,0.652)rectangle++(0.009,0.009) (0.503,0.258)rectangle++(0.009,0.009) (0.249,0.627)rectangle++(0.009,0.009)
  (0.122,0.811)rectangle++(0.009,0.009) (0.059,0.403)rectangle++(0.009,0.009) (0.521,0.159)rectangle++(0.009,0.009) (0.258,0.577)rectangle++(0.009,0.009) (0.601,0.228)rectangle++(0.009,0.009)
  (0.738,0.089)rectangle++(0.009,0.009) (0.367,0.542)rectangle++(0.009,0.009) (0.181,0.769)rectangle++(0.009,0.009) (0.088,0.382)rectangle++(0.009,0.009) (0.042,0.689)rectangle++(0.009,0.009)
  (0.019,0.342)rectangle++(0.009,0.009) (0.007,0.169)rectangle++(0.009,0.009) (0.5,0.065)rectangle++(0.009,0.009) (0.248,0.03)rectangle++(0.009,0.009) (0.122,0.013)rectangle++(0.009,0.009)
  (0.059,0.504)rectangle++(0.009,0.009) (0.027,0.25)rectangle++(0.009,0.009) (0.011,0.123)rectangle++(0.009,0.009) (0.003,0.559)rectangle++(0.009,0.009) (-0.001,0.277)rectangle++(0.009,0.009)
  (-0.003,0.636)rectangle++(0.009,0.009) (0.496,0.252)rectangle++(0.009,0.009) (0.246,0.624)rectangle++(0.009,0.009) (0.121,0.31)rectangle++(0.009,0.009) (0.058,0.653)rectangle++(0.009,0.009)
  (0.027,0.824)rectangle++(0.009,0.009) (0.011,0.91)rectangle++(0.009,0.009) (0.003,0.953)rectangle++(0.009,0.009) (-0.001,0.474)rectangle++(0.009,0.009) (-0.003,0.235)rectangle++(0.009,0.009)
  (-0.004,0.615)rectangle++(0.009,0.009) (-0.004,0.305)rectangle++(0.009,0.009) (-0.004,0.65)rectangle++(0.009,0.009) (0.496,0.257)rectangle++(0.009,0.009) (0.246,0.626)rectangle++(0.009,0.009)
  (0.121,0.311)rectangle++(0.009,0.009) (0.058,0.153)rectangle++(0.009,0.009) (0.027,0.574)rectangle++(0.009,0.009) (0.011,0.785)rectangle++(0.009,0.009) (0.003,0.39)rectangle++(0.009,0.009)
  (0.499,0.153)rectangle++(0.009,0.009) (0.697,0.059)rectangle++(0.009,0.009) (0.346,0.027)rectangle++(0.009,0.009) (0.636,0.008)rectangle++(0.009,0.009) (0.752,0.001)rectangle++(0.009,0.009)
  (0.798,-0.002)rectangle++(0.009,0.009) (0.816,-0.004)rectangle++(0.009,0.009) (0.406,0.496)rectangle++(0.009,0.009) (0.201,0.246)rectangle++(0.009,0.009) (0.098,0.621)rectangle++(0.009,0.009)
  (0.537,0.246)rectangle++(0.009,0.009) (0.712,0.096)rectangle++(0.009,0.009) (0.354,0.046)rectangle++(0.009,0.009) (0.175,0.021)rectangle++(0.009,0.009) (0.085,0.508)rectangle++(0.009,0.009)
  (0.04,0.252)rectangle++(0.009,0.009) (0.513,0.098)rectangle++(0.009,0.009) (0.703,0.037)rectangle++(0.009,0.009) (0.349,0.516)rectangle++(0.009,0.009) (0.172,0.756)rectangle++(0.009,0.009)
  (0.566,0.3)rectangle++(0.009,0.009) (0.724,0.117)rectangle++(0.009,0.009) (0.36,0.556)rectangle++(0.009,0.009) (0.641,0.22)rectangle++(0.009,0.009) (0.318,0.108)rectangle++(0.009,0.009)
  (0.157,0.052)rectangle++(0.009,0.009) (0.076,0.524)rectangle++(0.009,0.009) (0.036,0.76)rectangle++(0.009,0.009) (0.512,0.301)rectangle++(0.009,0.009) (0.254,0.148)rectangle++(0.009,0.009)
  (0.599,0.057)rectangle++(0.009,0.009) (0.737,0.02)rectangle++(0.009,0.009) (0.366,0.008)rectangle++(0.009,0.009) (0.181,0.502)rectangle++(0.009,0.009) (0.088,0.249)rectangle++(0.009,0.009)
  (0.533,0.097)rectangle++(0.009,0.009) (0.264,0.546)rectangle++(0.009,0.009) (0.13,0.271)rectangle++(0.009,0.009) (0.549,0.106)rectangle++(0.009,0.009) (0.272,0.551)rectangle++(0.009,0.009)
  (0.134,0.273)rectangle++(0.009,0.009) (0.065,0.634)rectangle++(0.009,0.009) (0.03,0.815)rectangle++(0.009,0.009) (0.013,0.905)rectangle++(0.009,0.009) (0.004,0.45)rectangle++(0.009,0.009)
  (0,0.223)rectangle++(0.009,0.009) (-0.002,0.109)rectangle++(0.009,0.009) (-0.003,0.052)rectangle++(0.009,0.009) (0.496,0.018)rectangle++(0.009,0.009) (0.246,0.507)rectangle++(0.009,0.009);
\fill[nuage]
  (0.596,0.2)rectangle++(0.009,0.009) (0.296,0.098)rectangle++(0.009,0.009) (0.146,0.547)rectangle++(0.009,0.009) (0.071,0.771)rectangle++(0.009,0.009) (0.033,0.883)rectangle++(0.009,0.009)
  (0.511,0.351)rectangle++(0.009,0.009) (0.253,0.673)rectangle++(0.009,0.009) (0.124,0.334)rectangle++(0.009,0.009) (0.06,0.165)rectangle++(0.009,0.009) (0.028,0.58)rectangle++(0.009,0.009)
  (0.508,0.229)rectangle++(0.009,0.009) (0.252,0.612)rectangle++(0.009,0.009) (0.124,0.304)rectangle++(0.009,0.009) (0.547,0.119)rectangle++(0.009,0.009) (0.271,0.557)rectangle++(0.009,0.009)
  (0.133,0.276)rectangle++(0.009,0.009) (0.064,0.136)rectangle++(0.009,0.009) (0.03,0.566)rectangle++(0.009,0.009) (0.509,0.224)rectangle++(0.009,0.009) (0.252,0.61)rectangle++(0.009,0.009)
  (0.124,0.303)rectangle++(0.009,0.009) (0.06,0.149)rectangle++(0.009,0.009) (0.028,0.072)rectangle++(0.009,0.009) (0.012,0.034)rectangle++(0.009,0.009) (0.004,0.015)rectangle++(0.009,0.009)
  (0,0.005)rectangle++(0.009,0.009) (-0.002,0.5)rectangle++(0.009,0.009) (-0.003,0.248)rectangle++(0.009,0.009) (-0.004,0.622)rectangle++(0.009,0.009) (-0.004,0.309)rectangle++(0.009,0.009)
  (0.496,0.121)rectangle++(0.009,0.009) (0.696,0.046)rectangle++(0.009,0.009) (0.346,0.521)rectangle++(0.009,0.009) (0.171,0.258)rectangle++(0.009,0.009) (0.083,0.127)rectangle++(0.009,0.009)
  (0.039,0.061)rectangle++(0.009,0.009) (0.017,0.528)rectangle++(0.009,0.009) (0.504,0.209)rectangle++(0.009,0.009) (0.25,0.602)rectangle++(0.009,0.009) (0.123,0.799)rectangle++(0.009,0.009)
  (0.546,0.317)rectangle++(0.009,0.009) (0.716,0.124)rectangle++(0.009,0.009) (0.356,0.56)rectangle++(0.009,0.009) (0.176,0.778)rectangle++(0.009,0.009) (0.086,0.887)rectangle++(0.009,0.009)
  (0.041,0.441)rectangle++(0.009,0.009) (0.514,0.174)rectangle++(0.009,0.009) (0.703,0.067)rectangle++(0.009,0.009) (0.349,0.531)rectangle++(0.009,0.009) (0.637,0.21)rectangle++(0.009,0.009)
  (0.316,0.603)rectangle++(0.009,0.009) (0.156,0.299)rectangle++(0.009,0.009) (0.076,0.147)rectangle++(0.009,0.009) (0.036,0.071)rectangle++(0.009,0.009) (0.512,0.026)rectangle++(0.009,0.009)
  (0.254,0.011)rectangle++(0.009,0.009) (0.125,0.503)rectangle++(0.009,0.009) (0.06,0.749)rectangle++(0.009,0.009) (0.028,0.872)rectangle++(0.009,0.009) (0.012,0.434)rectangle++(0.009,0.009)
  (0.004,0.215)rectangle++(0.009,0.009) (0.499,0.083)rectangle++(0.009,0.009) (0.247,0.039)rectangle++(0.009,0.009) (0.596,0.013)rectangle++(0.009,0.009) (0.296,0.004)rectangle++(0.009,0.009)
  (0.146,0)rectangle++(0.009,0.009) (0.071,-0.002)rectangle++(0.009,0.009) (0.033,-0.003)rectangle++(0.009,0.009) (0.014,0.496)rectangle++(0.009,0.009) (0.005,0.246)rectangle++(0.009,0.009)
  (0,0.121)rectangle++(0.009,0.009) (-0.002,0.058)rectangle++(0.009,0.009) (-0.003,0.527)rectangle++(0.009,0.009) (-0.004,0.761)rectangle++(0.009,0.009) (-0.004,0.378)rectangle++(0.009,0.009)
  (-0.004,0.687)rectangle++(0.009,0.009) (-0.004,0.341)rectangle++(0.009,0.009) (0.496,0.134)rectangle++(0.009,0.009) (0.246,0.565)rectangle++(0.009,0.009) (0.121,0.28)rectangle++(0.009,0.009)
  (0.058,0.638)rectangle++(0.009,0.009) (0.027,0.817)rectangle++(0.009,0.009) (0.011,0.906)rectangle++(0.009,0.009) (0.502,0.36)rectangle++(0.009,0.009) (0.698,0.141)rectangle++(0.009,0.009)
  (0.347,0.568)rectangle++(0.009,0.009) (0.171,0.782)rectangle++(0.009,0.009) (0.083,0.389)rectangle++(0.009,0.009) (0.531,0.153)rectangle++(0.009,0.009) (0.71,0.058)rectangle++(0.009,0.009)
  (0.353,0.027)rectangle++(0.009,0.009) (0.174,0.511)rectangle++(0.009,0.009) (0.085,0.753)rectangle++(0.009,0.009) (0.04,0.874)rectangle++(0.009,0.009) (0.018,0.935)rectangle++(0.009,0.009)
  (0.007,0.965)rectangle++(0.009,0.009) (0.001,0.48)rectangle++(0.009,0.009) (-0.002,0.738)rectangle++(0.009,0.009) (-0.003,0.867)rectangle++(0.009,0.009) (0.496,0.344)rectangle++(0.009,0.009)
  (0.696,0.135)rectangle++(0.009,0.009) (0.346,0.065)rectangle++(0.009,0.009) (0.171,0.03)rectangle++(0.009,0.009) (0.083,0.013)rectangle++(0.009,0.009) (0.039,0.004)rectangle++(0.009,0.009)
  (0.017,0)rectangle++(0.009,0.009) (0.006,-0.002)rectangle++(0.009,0.009) (0.001,-0.003)rectangle++(0.009,0.009) (-0.002,0.496)rectangle++(0.009,0.009) (-0.003,0.246)rectangle++(0.009,0.009)
  (-0.004,0.121)rectangle++(0.009,0.009) (-0.004,0.558)rectangle++(0.009,0.009) (-0.004,0.277)rectangle++(0.009,0.009) (-0.004,0.636)rectangle++(0.009,0.009) (-0.004,0.316)rectangle++(0.009,0.009)
  (-0.004,0.656)rectangle++(0.009,0.009) (0.496,0.26)rectangle++(0.009,0.009) (0.246,0.128)rectangle++(0.009,0.009) (0.121,0.062)rectangle++(0.009,0.009) (0.546,0.022)rectangle++(0.009,0.009)
  (0.271,0.509)rectangle++(0.009,0.009) (0.133,0.752)rectangle++(0.009,0.009) (0.064,0.374)rectangle++(0.009,0.009) (0.03,0.685)rectangle++(0.009,0.009) (0.013,0.34)rectangle++(0.009,0.009)
  (0.004,0.168)rectangle++(0.009,0.009) (0,0.082)rectangle++(0.009,0.009) (-0.002,0.539)rectangle++(0.009,0.009) (-0.003,0.767)rectangle++(0.009,0.009) (0.496,0.304)rectangle++(0.009,0.009)
  (0.246,0.15)rectangle++(0.009,0.009) (0.596,0.057)rectangle++(0.009,0.009) (0.296,0.526)rectangle++(0.009,0.009) (0.146,0.761)rectangle++(0.009,0.009) (0.556,0.302)rectangle++(0.009,0.009)
  (0.72,0.118)rectangle++(0.009,0.009) (0.358,0.057)rectangle++(0.009,0.009) (0.177,0.526)rectangle++(0.009,0.009) (0.568,0.208)rectangle++(0.009,0.009) (0.282,0.602)rectangle++(0.009,0.009)
  (0.139,0.299)rectangle++(0.009,0.009) (0.067,0.147)rectangle++(0.009,0.009) (0.031,0.571)rectangle++(0.009,0.009) (0.013,0.283)rectangle++(0.009,0.009) (0.004,0.639)rectangle++(0.009,0.009)
  (0,0.317)rectangle++(0.009,0.009) (-0.002,0.156)rectangle++(0.009,0.009) (-0.003,0.076)rectangle++(0.009,0.009) (-0.004,0.036)rectangle++(0.009,0.009) (-0.004,0.016)rectangle++(0.009,0.009)
  (-0.004,0.506)rectangle++(0.009,0.009) (-0.004,0.751)rectangle++(0.009,0.009) (-0.004,0.873)rectangle++(0.009,0.009) (-0.004,0.934)rectangle++(0.009,0.009) (0.496,0.371)rectangle++(0.009,0.009)
  (0.246,0.683)rectangle++(0.009,0.009) (0.121,0.339)rectangle++(0.009,0.009) (0.058,0.667)rectangle++(0.009,0.009) (0.027,0.831)rectangle++(0.009,0.009) (0.011,0.413)rectangle++(0.009,0.009)
  (0.003,0.704)rectangle++(0.009,0.009) (-0.001,0.35)rectangle++(0.009,0.009) (-0.003,0.673)rectangle++(0.009,0.009) (-0.004,0.334)rectangle++(0.009,0.009) (0.496,0.131)rectangle++(0.009,0.009)
  (0.246,0.563)rectangle++(0.009,0.009) (0.596,0.223)rectangle++(0.009,0.009) (0.296,0.109)rectangle++(0.009,0.009) (0.146,0.552)rectangle++(0.009,0.009) (0.071,0.774)rectangle++(0.009,0.009)
  (0.033,0.885)rectangle++(0.009,0.009) (0.014,0.44)rectangle++(0.009,0.009) (0.005,0.218)rectangle++(0.009,0.009) (0.499,0.084)rectangle++(0.009,0.009) (0.247,0.54)rectangle++(0.009,0.009)
  (0.121,0.268)rectangle++(0.009,0.009) (0.058,0.132)rectangle++(0.009,0.009) (0.027,0.064)rectangle++(0.009,0.009) (0.011,0.53)rectangle++(0.009,0.009) (0.502,0.209)rectangle++(0.009,0.009)
  (0.698,0.081)rectangle++(0.009,0.009) (0.347,0.538)rectangle++(0.009,0.009) (0.171,0.267)rectangle++(0.009,0.009) (0.566,0.104)rectangle++(0.009,0.009) (0.281,0.55)rectangle++(0.009,0.009)
  (0.138,0.773)rectangle++(0.009,0.009) (0.067,0.884)rectangle++(0.009,0.009) (0.031,0.94)rectangle++(0.009,0.009) (0.013,0.468)rectangle++(0.009,0.009) (0.004,0.732)rectangle++(0.009,0.009)
  (0,0.364)rectangle++(0.009,0.009) (-0.002,0.68)rectangle++(0.009,0.009) (-0.003,0.338)rectangle++(0.009,0.009) (-0.004,0.167)rectangle++(0.009,0.009) (0.496,0.064)rectangle++(0.009,0.009)
  (0.246,0.03)rectangle++(0.009,0.009) (0.121,0.013)rectangle++(0.009,0.009) (0.546,0.002)rectangle++(0.009,0.009) (0.716,-0.002)rectangle++(0.009,0.009) (0.356,0.497)rectangle++(0.009,0.009)
  (0.64,0.196)rectangle++(0.009,0.009) (0.318,0.596)rectangle++(0.009,0.009) (0.157,0.296)rectangle++(0.009,0.009) (0.076,0.146)rectangle++(0.009,0.009) (0.528,0.056)rectangle++(0.009,0.009)
  (0.262,0.026)rectangle++(0.009,0.009) (0.129,0.511)rectangle++(0.009,0.009) (0.062,0.753)rectangle++(0.009,0.009) (0.029,0.374)rectangle++(0.009,0.009) (0.012,0.185)rectangle++(0.009,0.009)
  (0.502,0.071)rectangle++(0.009,0.009) (0.249,0.033)rectangle++(0.009,0.009) (0.597,0.011)rectangle++(0.009,0.009) (0.296,0.003)rectangle++(0.009,0.009) (0.146,0.499)rectangle++(0.009,0.009)
  (0.556,0.197)rectangle++(0.009,0.009) (0.276,0.096)rectangle++(0.009,0.009) (0.136,0.546)rectangle++(0.009,0.009) (0.066,0.271)rectangle++(0.009,0.009) (0.031,0.633)rectangle++(0.009,0.009)
  (0.013,0.814)rectangle++(0.009,0.009) (0.503,0.323)rectangle++(0.009,0.009) (0.249,0.659)rectangle++(0.009,0.009) (0.122,0.327)rectangle++(0.009,0.009) (0.059,0.161)rectangle++(0.009,0.009)
  (0.027,0.578)rectangle++(0.009,0.009) (0.508,0.229)rectangle++(0.009,0.009) (0.252,0.112)rectangle++(0.009,0.009) (0.124,0.054)rectangle++(0.009,0.009) (0.06,0.025)rectangle++(0.009,0.009)
  (0.028,0.51)rectangle++(0.009,0.009) (0.508,0.201)rectangle++(0.009,0.009) (0.252,0.598)rectangle++(0.009,0.009) (0.124,0.297)rectangle++(0.009,0.009) (0.547,0.116)rectangle++(0.009,0.009)
  (0.271,0.056)rectangle++(0.009,0.009) (0.133,0.026)rectangle++(0.009,0.009) (0.064,0.511)rectangle++(0.009,0.009) (0.523,0.202)rectangle++(0.009,0.009) (0.259,0.099)rectangle++(0.009,0.009)
  (0.127,0.047)rectangle++(0.009,0.009) (0.061,0.521)rectangle++(0.009,0.009) (0.028,0.758)rectangle++(0.009,0.009) (0.012,0.377)rectangle++(0.009,0.009) (0.004,0.186)rectangle++(0.009,0.009)
  (0,0.591)rectangle++(0.009,0.009) (0.497,0.234)rectangle++(0.009,0.009) (0.246,0.115)rectangle++(0.009,0.009) (0.121,0.055)rectangle++(0.009,0.009) (0.058,0.525)rectangle++(0.009,0.009)
  (0.027,0.76)rectangle++(0.009,0.009) (0.011,0.878)rectangle++(0.009,0.009) (0.003,0.437)rectangle++(0.009,0.009) (-0.001,0.716)rectangle++(0.009,0.009) (-0.003,0.356)rectangle++(0.009,0.009)
  (-0.004,0.676)rectangle++(0.009,0.009) (-0.004,0.836)rectangle++(0.009,0.009) (-0.004,0.916)rectangle++(0.009,0.009) (-0.004,0.956)rectangle++(0.009,0.009) (-0.004,0.476)rectangle++(0.009,0.009)
  (-0.004,0.736)rectangle++(0.009,0.009) (-0.004,0.866)rectangle++(0.009,0.009) (-0.004,0.931)rectangle++(0.009,0.009) (-0.004,0.963)rectangle++(0.009,0.009) (-0.004,0.979)rectangle++(0.009,0.009)
  (0.496,0.389)rectangle++(0.009,0.009) (0.246,0.192)rectangle++(0.009,0.009) (0.121,0.094)rectangle++(0.009,0.009) (0.058,0.545)rectangle++(0.009,0.009) (0.027,0.27)rectangle++(0.009,0.009)
  (0.508,0.105)rectangle++(0.009,0.009) (0.252,0.05)rectangle++(0.009,0.009) (0.124,0.023)rectangle++(0.009,0.009) (0.06,0.509)rectangle++(0.009,0.009) (0.028,0.252)rectangle++(0.009,0.009)
  (0.508,0.098)rectangle++(0.009,0.009) (0.252,0.047)rectangle++(0.009,0.009) (0.124,0.521)rectangle++(0.009,0.009) (0.06,0.758)rectangle++(0.009,0.009) (0.521,0.301)rectangle++(0.009,0.009)
  (0.258,0.148)rectangle++(0.009,0.009) (0.127,0.572)rectangle++(0.009,0.009) (0.061,0.784)rectangle++(0.009,0.009) (0.522,0.311)rectangle++(0.009,0.009) (0.259,0.653)rectangle++(0.009,0.009)
  (0.127,0.324)rectangle++(0.009,0.009) (0.061,0.16)rectangle++(0.009,0.009) (0.028,0.078)rectangle++(0.009,0.009) (0.012,0.537)rectangle++(0.009,0.009) (0.004,0.766)rectangle++(0.009,0.009)
  (0,0.381)rectangle++(0.009,0.009) (-0.002,0.188)rectangle++(0.009,0.009) (-0.003,0.592)rectangle++(0.009,0.009) (-0.004,0.794)rectangle++(0.009,0.009) (-0.004,0.895)rectangle++(0.009,0.009)
  (-0.004,0.945)rectangle++(0.009,0.009) (0.496,0.375)rectangle++(0.009,0.009) (0.696,0.147)rectangle++(0.009,0.009) (0.776,0.056)rectangle++(0.009,0.009) (0.808,0.02)rectangle++(0.009,0.009)
  (0.402,0.008)rectangle++(0.009,0.009) (0.199,0.002)rectangle++(0.009,0.009) (0.097,0.499)rectangle++(0.009,0.009) (0.536,0.197)rectangle++(0.009,0.009) (0.712,0.076)rectangle++(0.009,0.009)
  (0.354,0.536)rectangle++(0.009,0.009) (0.175,0.766)rectangle++(0.009,0.009) (0.567,0.304)rectangle++(0.009,0.009) (0.281,0.65)rectangle++(0.009,0.009) (0.138,0.323)rectangle++(0.009,0.009)
  (0.067,0.159)rectangle++(0.009,0.009) (0.031,0.077)rectangle++(0.009,0.009) (0.51,0.028)rectangle++(0.009,0.009) (0.253,0.512)rectangle++(0.009,0.009) (0.124,0.754)rectangle++(0.009,0.009)
  (0.06,0.375)rectangle++(0.009,0.009) (0.521,0.147)rectangle++(0.009,0.009) (0.258,0.071)rectangle++(0.009,0.009) (0.127,0.033)rectangle++(0.009,0.009) (0.548,0.011)rectangle++(0.009,0.009)
  (0.717,0.002)rectangle++(0.009,0.009) (0.784,-0.002)rectangle++(0.009,0.009) (0.811,-0.004)rectangle++(0.009,0.009) (0.403,0.496)rectangle++(0.009,0.009) (0.199,0.746)rectangle++(0.009,0.009)
  (0.577,0.296)rectangle++(0.009,0.009) (0.286,0.646)rectangle++(0.009,0.009) (0.141,0.821)rectangle++(0.009,0.009) (0.068,0.908)rectangle++(0.009,0.009) (0.525,0.361)rectangle++(0.009,0.009)
  (0.26,0.678)rectangle++(0.009,0.009) (0.128,0.337)rectangle++(0.009,0.009) (0.062,0.666)rectangle++(0.009,0.009) (0.029,0.331)rectangle++(0.009,0.009) (0.509,0.13)rectangle++(0.009,0.009)
  (0.252,0.063)rectangle++(0.009,0.009) (0.124,0.529)rectangle++(0.009,0.009) (0.547,0.209)rectangle++(0.009,0.009) (0.716,0.081)rectangle++(0.009,0.009) (0.356,0.038)rectangle++(0.009,0.009)
  (0.64,0.013)rectangle++(0.009,0.009) (0.318,0.004)rectangle++(0.009,0.009) (0.157,0)rectangle++(0.009,0.009) (0.56,-0.003)rectangle++(0.009,0.009) (0.278,0.496)rectangle++(0.009,0.009)
  (0.608,0.196)rectangle++(0.009,0.009) (0.741,0.076)rectangle++(0.009,0.009) (0.368,0.036)rectangle++(0.009,0.009) (0.182,0.516)rectangle++(0.009,0.009) (0.089,0.756)rectangle++(0.009,0.009)
  (0.042,0.876)rectangle++(0.009,0.009) (0.514,0.348)rectangle++(0.009,0.009) (0.255,0.172)rectangle++(0.009,0.009) (0.599,0.066)rectangle++(0.009,0.009) (0.297,0.531)rectangle++(0.009,0.009)
  (0.146,0.763)rectangle++(0.009,0.009) (0.556,0.303)rectangle++(0.009,0.009) (0.276,0.149)rectangle++(0.009,0.009) (0.136,0.572)rectangle++(0.009,0.009) (0.066,0.784)rectangle++(0.009,0.009)
  (0.031,0.39)rectangle++(0.009,0.009) (0.013,0.193)rectangle++(0.009,0.009) (0.004,0.094)rectangle++(0.009,0.009) (0,0.045)rectangle++(0.009,0.009) (-0.002,0.52)rectangle++(0.009,0.009)
  (-0.003,0.258)rectangle++(0.009,0.009) (-0.004,0.127)rectangle++(0.009,0.009) (-0.004,0.561)rectangle++(0.009,0.009) (0.496,0.222)rectangle++(0.009,0.009) (0.696,0.086)rectangle++(0.009,0.009)
  (0.776,0.032)rectangle++(0.009,0.009) (0.386,0.514)rectangle++(0.009,0.009) (0.191,0.255)rectangle++(0.009,0.009) (0.093,0.625)rectangle++(0.009,0.009) (0.044,0.81)rectangle++(0.009,0.009)
  (0.02,0.403)rectangle++(0.009,0.009) (0.008,0.199)rectangle++(0.009,0.009) (0.002,0.597)rectangle++(0.009,0.009) (-0.001,0.796)rectangle++(0.009,0.009) (-0.003,0.396)rectangle++(0.009,0.009)
  (-0.004,0.196)rectangle++(0.009,0.009) (-0.004,0.096)rectangle++(0.009,0.009) (-0.004,0.046)rectangle++(0.009,0.009) (-0.004,0.521)rectangle++(0.009,0.009) (-0.004,0.258)rectangle++(0.009,0.009)
  (0.496,0.101)rectangle++(0.009,0.009) (0.246,0.048)rectangle++(0.009,0.009) (0.121,0.022)rectangle++(0.009,0.009) (0.546,0.006)rectangle++(0.009,0.009) (0.271,0.001)rectangle++(0.009,0.009)
  (0.133,0.498)rectangle++(0.009,0.009) (0.551,0.197)rectangle++(0.009,0.009) (0.273,0.596)rectangle++(0.009,0.009) (0.134,0.796)rectangle++(0.009,0.009) (0.065,0.396)rectangle++(0.009,0.009)
  (0.523,0.156)rectangle++(0.009,0.009) (0.707,0.06)rectangle++(0.009,0.009) (0.351,0.028)rectangle++(0.009,0.009) (0.173,0.512)rectangle++(0.009,0.009) (0.084,0.754)rectangle++(0.009,0.009)
  (0.04,0.875)rectangle++(0.009,0.009) (0.513,0.347)rectangle++(0.009,0.009) (0.254,0.671)rectangle++(0.009,0.009) (0.125,0.833)rectangle++(0.009,0.009) (0.06,0.914)rectangle++(0.009,0.009)
  (0.521,0.363)rectangle++(0.009,0.009) (0.706,0.143)rectangle++(0.009,0.009) (0.351,0.569)rectangle++(0.009,0.009) (0.173,0.282)rectangle++(0.009,0.009) (0.084,0.639)rectangle++(0.009,0.009)
  (0.531,0.253)rectangle++(0.009,0.009) (0.71,0.098)rectangle++(0.009,0.009) (0.353,0.547)rectangle++(0.009,0.009) (0.174,0.771)rectangle++(0.009,0.009) (0.567,0.306)rectangle++(0.009,0.009)
  (0.281,0.651)rectangle++(0.009,0.009) (0.138,0.323)rectangle++(0.009,0.009) (0.067,0.659)rectangle++(0.009,0.009) (0.031,0.327)rectangle++(0.009,0.009) (0.013,0.161)rectangle++(0.009,0.009)
  (0.004,0.078)rectangle++(0.009,0.009) (0,0.037)rectangle++(0.009,0.009) (0.497,0.012)rectangle++(0.009,0.009) (0.246,0.504)rectangle++(0.009,0.009) (0.121,0.25)rectangle++(0.009,0.009)
  (0.058,0.623)rectangle++(0.009,0.009) (0.521,0.246)rectangle++(0.009,0.009) (0.706,0.096)rectangle++(0.009,0.009) (0.351,0.046)rectangle++(0.009,0.009) (0.173,0.521)rectangle++(0.009,0.009)
  (0.084,0.758)rectangle++(0.009,0.009) (0.531,0.301)rectangle++(0.009,0.009) (0.263,0.148)rectangle++(0.009,0.009) (0.129,0.572)rectangle++(0.009,0.009) (0.062,0.284)rectangle++(0.009,0.009)
  (0.029,0.64)rectangle++(0.009,0.009) (0.012,0.818)rectangle++(0.009,0.009) (0.004,0.407)rectangle++(0.009,0.009) (0,0.701)rectangle++(0.009,0.009) (0.497,0.278)rectangle++(0.009,0.009)
  (0.696,0.108)rectangle++(0.009,0.009) (0.346,0.552)rectangle++(0.009,0.009) (0.171,0.774)rectangle++(0.009,0.009) (0.083,0.385)rectangle++(0.009,0.009) (0.039,0.19)rectangle++(0.009,0.009)
  (0.017,0.093)rectangle++(0.009,0.009) (0.006,0.044)rectangle++(0.009,0.009) (0.001,0.52)rectangle++(0.009,0.009) (-0.002,0.258)rectangle++(0.009,0.009) (0.497,0.1)rectangle++(0.009,0.009)
  (0.246,0.048)rectangle++(0.009,0.009) (0.121,0.022)rectangle++(0.009,0.009) (0.058,0.009)rectangle++(0.009,0.009) (0.027,0.002)rectangle++(0.009,0.009) (0.011,0.499)rectangle++(0.009,0.009)
  (0.003,0.747)rectangle++(0.009,0.009) (-0.001,0.371)rectangle++(0.009,0.009) (0.497,0.146)rectangle++(0.009,0.009) (0.246,0.571)rectangle++(0.009,0.009) (0.121,0.283)rectangle++(0.009,0.009)
  (0.058,0.639)rectangle++(0.009,0.009) (0.027,0.317)rectangle++(0.009,0.009) (0.011,0.156)rectangle++(0.009,0.009) (0.003,0.076)rectangle++(0.009,0.009) (-0.001,0.536)rectangle++(0.009,0.009)
  (-0.003,0.766)rectangle++(0.009,0.009) (-0.004,0.381)rectangle++(0.009,0.009) (-0.004,0.188)rectangle++(0.009,0.009) (-0.004,0.592)rectangle++(0.009,0.009) (0.496,0.234)rectangle++(0.009,0.009)
  (0.696,0.091)rectangle++(0.009,0.009) (0.776,0.034)rectangle++(0.009,0.009) (0.808,0.011)rectangle++(0.009,0.009) (0.402,0.003)rectangle++(0.009,0.009) (0.199,-0.001)rectangle++(0.009,0.009)
  (0.097,-0.003)rectangle++(0.009,0.009) (0.046,0.496)rectangle++(0.009,0.009) (0.021,0.246)rectangle++(0.009,0.009) (0.008,0.121)rectangle++(0.009,0.009) (0.002,0.558)rectangle++(0.009,0.009)
  (-0.001,0.777)rectangle++(0.009,0.009) (-0.003,0.886)rectangle++(0.009,0.009) (-0.004,0.441)rectangle++(0.009,0.009) (0.496,0.174)rectangle++(0.009,0.009) (0.246,0.585)rectangle++(0.009,0.009)
  (0.121,0.79)rectangle++(0.009,0.009) (0.058,0.893)rectangle++(0.009,0.009) (0.521,0.354)rectangle++(0.009,0.009) (0.706,0.139)rectangle++(0.009,0.009) (0.78,0.053)rectangle++(0.009,0.009)
  (0.388,0.024)rectangle++(0.009,0.009) (0.192,0.51)rectangle++(0.009,0.009) (0.574,0.201)rectangle++(0.009,0.009) (0.285,0.598)rectangle++(0.009,0.009) (0.14,0.797)rectangle++(0.009,0.009)
  (0.068,0.896)rectangle++(0.009,0.009) (0.032,0.446)rectangle++(0.009,0.009) (0.014,0.721)rectangle++(0.009,0.009) (0.005,0.358)rectangle++(0.009,0.009) (0.499,0.141)rectangle++(0.009,0.009)
  (0.697,0.054)rectangle++(0.009,0.009) (0.776,0.019)rectangle++(0.009,0.009) (0.386,0.507)rectangle++(0.009,0.009) (0.191,0.251)rectangle++(0.009,0.009) (0.093,0.123)rectangle++(0.009,0.009)
  (0.044,0.559)rectangle++(0.009,0.009) (0.02,0.777)rectangle++(0.009,0.009) (0.008,0.386)rectangle++(0.009,0.009) (0.5,0.152)rectangle++(0.009,0.009) (0.248,0.574)rectangle++(0.009,0.009)
  (0.596,0.227)rectangle++(0.009,0.009) (0.296,0.611)rectangle++(0.009,0.009) (0.616,0.242)rectangle++(0.009,0.009) (0.306,0.119)rectangle++(0.009,0.009) (0.151,0.557)rectangle++(0.009,0.009)
  (0.073,0.276)rectangle++(0.009,0.009) (0.527,0.108)rectangle++(0.009,0.009) (0.261,0.552)rectangle++(0.009,0.009) (0.128,0.774)rectangle++(0.009,0.009) (0.062,0.885)rectangle++(0.009,0.009)
  (0.029,0.94)rectangle++(0.009,0.009) (0.012,0.468)rectangle++(0.009,0.009) (0.004,0.232)rectangle++(0.009,0.009) (0,0.114)rectangle++(0.009,0.009) (0.497,0.043)rectangle++(0.009,0.009)
  (0.246,0.519)rectangle++(0.009,0.009) (0.121,0.757)rectangle++(0.009,0.009) (0.058,0.876)rectangle++(0.009,0.009) (0.521,0.348)rectangle++(0.009,0.009) (0.258,0.172)rectangle++(0.009,0.009)
  (0.601,0.066)rectangle++(0.009,0.009) (0.738,0.024)rectangle++(0.009,0.009) (0.367,0.01)rectangle++(0.009,0.009) (0.644,0.001)rectangle++(0.009,0.009) (0.755,-0.002)rectangle++(0.009,0.009)
  (0.375,0.497)rectangle++(0.009,0.009) (0.647,0.196)rectangle++(0.009,0.009) (0.321,0.596)rectangle++(0.009,0.009) (0.158,0.296)rectangle++(0.009,0.009) (0.077,0.146)rectangle++(0.009,0.009)
  (0.036,0.071)rectangle++(0.009,0.009) (0.016,0.033)rectangle++(0.009,0.009) (0.006,0.014)rectangle++(0.009,0.009) (0.001,0.005)rectangle++(0.009,0.009) (0.498,-0.001)rectangle++(0.009,0.009)
  (0.247,0.497)rectangle++(0.009,0.009) (0.596,0.196)rectangle++(0.009,0.009) (0.296,0.596)rectangle++(0.009,0.009) (0.146,0.296)rectangle++(0.009,0.009) (0.071,0.646)rectangle++(0.009,0.009)
  (0.033,0.321)rectangle++(0.009,0.009) (0.014,0.658)rectangle++(0.009,0.009) (0.503,0.261)rectangle++(0.009,0.009) (0.249,0.128)rectangle++(0.009,0.009) (0.122,0.062)rectangle++(0.009,0.009)
  (0.546,0.022)rectangle++(0.009,0.009) (0.716,0.006)rectangle++(0.009,0.009) (0.356,0.001)rectangle++(0.009,0.009) (0.176,0.498)rectangle++(0.009,0.009) (0.568,0.197)rectangle++(0.009,0.009)
  (0.282,0.096)rectangle++(0.009,0.009) (0.139,0.046)rectangle++(0.009,0.009) (0.067,0.021)rectangle++(0.009,0.009) (0.524,0.006)rectangle++(0.009,0.009) (0.26,0.001)rectangle++(0.009,0.009)
  (0.128,-0.002)rectangle++(0.009,0.009) (0.062,0.497)rectangle++(0.009,0.009) (0.029,0.246)rectangle++(0.009,0.009) (0.012,0.121)rectangle++(0.009,0.009) (0.004,0.058)rectangle++(0.009,0.009)
  (0,0.527)rectangle++(0.009,0.009) (0.497,0.208)rectangle++(0.009,0.009) (0.246,0.102)rectangle++(0.009,0.009) (0.121,0.049)rectangle++(0.009,0.009) (0.058,0.522)rectangle++(0.009,0.009)
  (0.027,0.259)rectangle++(0.009,0.009) (0.011,0.127)rectangle++(0.009,0.009) (0.502,0.048)rectangle++(0.009,0.009) (0.249,0.022)rectangle++(0.009,0.009) (0.597,0.006)rectangle++(0.009,0.009)
  (0.296,0.001)rectangle++(0.009,0.009) (0.146,0.498)rectangle++(0.009,0.009) (0.071,0.747)rectangle++(0.009,0.009) (0.526,0.296)rectangle++(0.009,0.009) (0.261,0.646)rectangle++(0.009,0.009)
  (0.128,0.321)rectangle++(0.009,0.009) (0.062,0.158)rectangle++(0.009,0.009) (0.029,0.577)rectangle++(0.009,0.009) (0.012,0.786)rectangle++(0.009,0.009) (0.004,0.391)rectangle++(0.009,0.009)
  (0,0.193)rectangle++(0.009,0.009) (-0.002,0.094)rectangle++(0.009,0.009) (-0.003,0.545)rectangle++(0.009,0.009) (-0.004,0.27)rectangle++(0.009,0.009) (0.496,0.105)rectangle++(0.009,0.009)
  (0.246,0.05)rectangle++(0.009,0.009) (0.596,0.017)rectangle++(0.009,0.009) (0.736,0.004)rectangle++(0.009,0.009) (0.366,0)rectangle++(0.009,0.009) (0.181,0.498)rectangle++(0.009,0.009)
  (0.088,0.747)rectangle++(0.009,0.009) (0.042,0.871)rectangle++(0.009,0.009) (0.514,0.346)rectangle++(0.009,0.009) (0.703,0.136)rectangle++(0.009,0.009) (0.778,0.052)rectangle++(0.009,0.009)
  (0.387,0.024)rectangle++(0.009,0.009) (0.191,0.01)rectangle++(0.009,0.009) (0.093,0.503)rectangle++(0.009,0.009) (0.044,0.749)rectangle++(0.009,0.009) (0.02,0.372)rectangle++(0.009,0.009)
  (0.008,0.184)rectangle++(0.009,0.009) (0.002,0.59)rectangle++(0.009,0.009) (0.498,0.233)rectangle++(0.009,0.009) (0.247,0.114)rectangle++(0.009,0.009) (0.121,0.055)rectangle++(0.009,0.009)
  (0.058,0.525)rectangle++(0.009,0.009) (0.027,0.76)rectangle++(0.009,0.009) (0.011,0.378)rectangle++(0.009,0.009) (0.003,0.187)rectangle++(0.009,0.009) (0.499,0.072)rectangle++(0.009,0.009)
  (0.247,0.534)rectangle++(0.009,0.009) (0.121,0.265)rectangle++(0.009,0.009) (0.058,0.13)rectangle++(0.009,0.009) (0.027,0.063)rectangle++(0.009,0.009) (0.011,0.029)rectangle++(0.009,0.009)
  (0.003,0.012)rectangle++(0.009,0.009) (-0.001,0.004)rectangle++(0.009,0.009) (-0.003,0.5)rectangle++(0.009,0.009) (0.496,0.197)rectangle++(0.009,0.009) (0.246,0.596)rectangle++(0.009,0.009)
  (0.121,0.296)rectangle++(0.009,0.009) (0.058,0.646)rectangle++(0.009,0.009) (0.521,0.256)rectangle++(0.009,0.009) (0.258,0.126)rectangle++(0.009,0.009) (0.127,0.561)rectangle++(0.009,0.009)
  (0.061,0.778)rectangle++(0.009,0.009) (0.028,0.887)rectangle++(0.009,0.009) (0.012,0.441)rectangle++(0.009,0.009) (0.004,0.218)rectangle++(0.009,0.009) (0.499,0.085)rectangle++(0.009,0.009)
  (0.247,0.54)rectangle++(0.009,0.009) (0.121,0.768)rectangle++(0.009,0.009) (0.546,0.304)rectangle++(0.009,0.009) (0.271,0.15)rectangle++(0.009,0.009) (0.606,0.057)rectangle++(0.009,0.009)
  (0.301,0.026)rectangle++(0.009,0.009) (0.148,0.011)rectangle++(0.009,0.009) (0.072,0.003)rectangle++(0.009,0.009) (0.034,0.499)rectangle++(0.009,0.009) (0.015,0.247)rectangle++(0.009,0.009)
  (0.005,0.121)rectangle++(0.009,0.009) (0.499,0.046)rectangle++(0.009,0.009) (0.697,0.016)rectangle++(0.009,0.009) (0.346,0.506)rectangle++(0.009,0.009) (0.636,0.2)rectangle++(0.009,0.009)
  (0.752,0.077)rectangle++(0.009,0.009) (0.374,0.536)rectangle++(0.009,0.009) (0.185,0.766)rectangle++(0.009,0.009) (0.571,0.304)rectangle++(0.009,0.009) (0.283,0.15)rectangle++(0.009,0.009)
  (0.139,0.073)rectangle++(0.009,0.009) (0.067,0.534)rectangle++(0.009,0.009) (0.031,0.265)rectangle++(0.009,0.009) (0.51,0.103)rectangle++(0.009,0.009) (0.253,0.549)rectangle++(0.009,0.009)
  (0.124,0.272)rectangle++(0.009,0.009) (0.547,0.106)rectangle++(0.009,0.009) (0.271,0.551)rectangle++(0.009,0.009) (0.606,0.218)rectangle++(0.009,0.009) (0.301,0.607)rectangle++(0.009,0.009)
  (0.148,0.801)rectangle++(0.009,0.009) (0.072,0.898)rectangle++(0.009,0.009) (0.034,0.447)rectangle++(0.009,0.009) (0.511,0.176)rectangle++(0.009,0.009) (0.253,0.086)rectangle++(0.009,0.009)
  (0.124,0.541)rectangle++(0.009,0.009) (0.06,0.768)rectangle++(0.009,0.009) (0.521,0.305)rectangle++(0.009,0.009) (0.258,0.15)rectangle++(0.009,0.009) (0.127,0.573)rectangle++(0.009,0.009)
  (0.061,0.284)rectangle++(0.009,0.009) (0.028,0.14)rectangle++(0.009,0.009) (0.509,0.053)rectangle++(0.009,0.009) (0.252,0.024)rectangle++(0.009,0.009) (0.124,0.01)rectangle++(0.009,0.009)
  (0.06,0.003)rectangle++(0.009,0.009) (0.028,-0.001)rectangle++(0.009,0.009) (0.508,-0.003)rectangle++(0.009,0.009) (0.252,0.496)rectangle++(0.009,0.009) (0.124,0.746)rectangle++(0.009,0.009)
  (0.547,0.296)rectangle++(0.009,0.009) (0.716,0.116)rectangle++(0.009,0.009) (0.784,0.044)rectangle++(0.009,0.009) (0.39,0.02)rectangle++(0.009,0.009) (0.653,0.005)rectangle++(0.009,0.009)
  (0.324,0)rectangle++(0.009,0.009) (0.627,-0.003)rectangle++(0.009,0.009) (0.311,0.496)rectangle++(0.009,0.009) (0.153,0.246)rectangle++(0.009,0.009) (0.074,0.121)rectangle++(0.009,0.009)
  (0.035,0.558)rectangle++(0.009,0.009) (0.015,0.777)rectangle++(0.009,0.009) (0.005,0.386)rectangle++(0.009,0.009) (0,0.191)rectangle++(0.009,0.009) (-0.002,0.093)rectangle++(0.009,0.009)
  (-0.003,0.544)rectangle++(0.009,0.009) (-0.004,0.77)rectangle++(0.009,0.009) (-0.004,0.383)rectangle++(0.009,0.009) (-0.004,0.689)rectangle++(0.009,0.009) (-0.004,0.842)rectangle++(0.009,0.009)
  (-0.004,0.419)rectangle++(0.009,0.009) (-0.004,0.207)rectangle++(0.009,0.009) (-0.004,0.101)rectangle++(0.009,0.009) (0.496,0.038)rectangle++(0.009,0.009) (0.246,0.517)rectangle++(0.009,0.009)
  (0.121,0.256)rectangle++(0.009,0.009) (0.058,0.126)rectangle++(0.009,0.009) (0.521,0.048)rectangle++(0.009,0.009) (0.258,0.022)rectangle++(0.009,0.009) (0.127,0.509)rectangle++(0.009,0.009)
  (0.061,0.752)rectangle++(0.009,0.009) (0.028,0.874)rectangle++(0.009,0.009) (0.012,0.435)rectangle++(0.009,0.009) (0.004,0.215)rectangle++(0.009,0.009) (0.499,0.083)rectangle++(0.009,0.009)
  (0.247,0.039)rectangle++(0.009,0.009) (0.121,0.517)rectangle++(0.009,0.009) (0.058,0.256)rectangle++(0.009,0.009) (0.027,0.126)rectangle++(0.009,0.009) (0.011,0.561)rectangle++(0.009,0.009)
  (0.003,0.278)rectangle++(0.009,0.009) (-0.001,0.637)rectangle++(0.009,0.009) (-0.003,0.316)rectangle++(0.009,0.009) (0.496,0.124)rectangle++(0.009,0.009) (0.696,0.047)rectangle++(0.009,0.009)
  (0.346,0.021)rectangle++(0.009,0.009) (0.171,0.008)rectangle++(0.009,0.009) (0.083,0.502)rectangle++(0.009,0.009) (0.531,0.198)rectangle++(0.009,0.009) (0.263,0.097)rectangle++(0.009,0.009)
  (0.603,0.036)rectangle++(0.009,0.009) (0.738,0.012)rectangle++(0.009,0.009) (0.367,0.004)rectangle++(0.009,0.009) (0.181,0)rectangle++(0.009,0.009) (0.088,-0.002)rectangle++(0.009,0.009)
  (0.042,-0.003)rectangle++(0.009,0.009) (0.019,0.496)rectangle++(0.009,0.009) (0.007,0.746)rectangle++(0.009,0.009) (0.5,0.296)rectangle++(0.009,0.009) (0.248,0.646)rectangle++(0.009,0.009)
  (0.122,0.321)rectangle++(0.009,0.009) (0.059,0.158)rectangle++(0.009,0.009) (0.027,0.077)rectangle++(0.009,0.009) (0.011,0.536)rectangle++(0.009,0.009) (0.003,0.766)rectangle++(0.009,0.009)
  (-0.001,0.381)rectangle++(0.009,0.009) (-0.003,0.188)rectangle++(0.009,0.009) (-0.004,0.592)rectangle++(0.009,0.009) (0.496,0.234)rectangle++(0.009,0.009) (0.246,0.115)rectangle++(0.009,0.009)
  (0.121,0.055)rectangle++(0.009,0.009) (0.546,0.019)rectangle++(0.009,0.009) (0.271,0.507)rectangle++(0.009,0.009) (0.606,0.2)rectangle++(0.009,0.009) (0.74,0.077)rectangle++(0.009,0.009)
  (0.368,0.536)rectangle++(0.009,0.009) (0.182,0.266)rectangle++(0.009,0.009) (0.089,0.131)rectangle++(0.009,0.009) (0.533,0.05)rectangle++(0.009,0.009) (0.264,0.523)rectangle++(0.009,0.009)
  (0.13,0.759)rectangle++(0.009,0.009) (0.549,0.301)rectangle++(0.009,0.009) (0.272,0.148)rectangle++(0.009,0.009) (0.134,0.072)rectangle++(0.009,0.009) (0.065,0.034)rectangle++(0.009,0.009)
  (0.523,0.011)rectangle++(0.009,0.009) (0.259,0.003)rectangle++(0.009,0.009) (0.127,0.499)rectangle++(0.009,0.009) (0.548,0.197)rectangle++(0.009,0.009) (0.272,0.596)rectangle++(0.009,0.009)
  (0.134,0.296)rectangle++(0.009,0.009) (0.065,0.646)rectangle++(0.009,0.009) (0.03,0.821)rectangle++(0.009,0.009) (0.509,0.326)rectangle++(0.009,0.009) (0.252,0.661)rectangle++(0.009,0.009)
  (0.124,0.328)rectangle++(0.009,0.009) (0.06,0.162)rectangle++(0.009,0.009) (0.521,0.062)rectangle++(0.009,0.009) (0.258,0.029)rectangle++(0.009,0.009) (0.127,0.512)rectangle++(0.009,0.009)
  (0.548,0.202)rectangle++(0.009,0.009) (0.272,0.099)rectangle++(0.009,0.009) (0.134,0.547)rectangle++(0.009,0.009) (0.065,0.771)rectangle++(0.009,0.009) (0.03,0.883)rectangle++(0.009,0.009)
  (0.013,0.939)rectangle++(0.009,0.009) (0.004,0.467)rectangle++(0.009,0.009) (0,0.731)rectangle++(0.009,0.009) (0.497,0.29)rectangle++(0.009,0.009) (0.246,0.643)rectangle++(0.009,0.009)
  (0.121,0.319)rectangle++(0.009,0.009) (0.546,0.125)rectangle++(0.009,0.009) (0.271,0.56)rectangle++(0.009,0.009) (0.133,0.278)rectangle++(0.009,0.009) (0.064,0.137)rectangle++(0.009,0.009)
  (0.03,0.066)rectangle++(0.009,0.009) (0.013,0.531)rectangle++(0.009,0.009) (0.004,0.263)rectangle++(0.009,0.009) (0,0.129)rectangle++(0.009,0.009) (0.497,0.049)rectangle++(0.009,0.009)
  (0.246,0.522)rectangle++(0.009,0.009) (0.121,0.259)rectangle++(0.009,0.009) (0.058,0.127)rectangle++(0.009,0.009) (0.027,0.561)rectangle++(0.009,0.009) (0.508,0.222)rectangle++(0.009,0.009)
  (0.252,0.609)rectangle++(0.009,0.009) (0.598,0.241)rectangle++(0.009,0.009) (0.297,0.118)rectangle++(0.009,0.009) (0.616,0.045)rectangle++(0.009,0.009) (0.306,0.02)rectangle++(0.009,0.009)
  (0.62,0.005)rectangle++(0.009,0.009) (0.308,0.5)rectangle++(0.009,0.009) (0.62,0.197)rectangle++(0.009,0.009) (0.745,0.076)rectangle++(0.009,0.009) (0.37,0.036)rectangle++(0.009,0.009)
  (0.183,0.016)rectangle++(0.009,0.009) (0.089,0.506)rectangle++(0.009,0.009) (0.042,0.251)rectangle++(0.009,0.009) (0.019,0.123)rectangle++(0.009,0.009) (0.007,0.059)rectangle++(0.009,0.009)
  (0.001,0.027)rectangle++(0.009,0.009) (0.498,0.008)rectangle++(0.009,0.009) (0.247,0.002)rectangle++(0.009,0.009) (0.121,-0.001)rectangle++(0.009,0.009) (0.058,0.497)rectangle++(0.009,0.009)
  (0.027,0.246)rectangle++(0.009,0.009) (0.508,0.096)rectangle++(0.009,0.009) (0.252,0.546)rectangle++(0.009,0.009) (0.124,0.271)rectangle++(0.009,0.009) (0.06,0.633)rectangle++(0.009,0.009)
  (0.521,0.251)rectangle++(0.009,0.009) (0.706,0.098)rectangle++(0.009,0.009) (0.78,0.036)rectangle++(0.009,0.009) (0.388,0.016)rectangle++(0.009,0.009) (0.192,0.006)rectangle++(0.009,0.009)
  (0.094,0.001)rectangle++(0.009,0.009) (0.535,-0.002)rectangle++(0.009,0.009) (0.265,-0.003)rectangle++(0.009,0.009) (0.13,0.496)rectangle++(0.009,0.009) (0.063,0.246)rectangle++(0.009,0.009)
  (0.029,0.621)rectangle++(0.009,0.009) (0.509,0.246)rectangle++(0.009,0.009) (0.701,0.096)rectangle++(0.009,0.009) (0.778,0.036)rectangle++(0.009,0.009) (0.808,0.012)rectangle++(0.009,0.009)
  (0.821,0.002)rectangle++(0.009,0.009) (0.408,0.499)rectangle++(0.009,0.009) (0.661,0.197)rectangle++(0.009,0.009) (0.762,0.076)rectangle++(0.009,0.009) (0.379,0.536)rectangle++(0.009,0.009)
  (0.187,0.266)rectangle++(0.009,0.009) (0.091,0.131)rectangle++(0.009,0.009) (0.043,0.063)rectangle++(0.009,0.009) (0.019,0.029)rectangle++(0.009,0.009) (0.505,0.009)rectangle++(0.009,0.009)
  (0.699,0.001)rectangle++(0.009,0.009) (0.777,-0.002)rectangle++(0.009,0.009) (0.386,-0.003)rectangle++(0.009,0.009) (0.191,0.496)rectangle++(0.009,0.009) (0.574,0.196)rectangle++(0.009,0.009)
  (0.285,0.596)rectangle++(0.009,0.009) (0.14,0.796)rectangle++(0.009,0.009) (0.068,0.396)rectangle++(0.009,0.009) (0.032,0.696)rectangle++(0.009,0.009) (0.014,0.846)rectangle++(0.009,0.009)
  (0.005,0.921)rectangle++(0.009,0.009) (0,0.458)rectangle++(0.009,0.009) (-0.002,0.727)rectangle++(0.009,0.009) (-0.003,0.361)rectangle++(0.009,0.009) (0.496,0.142)rectangle++(0.009,0.009)
  (0.246,0.569)rectangle++(0.009,0.009) (0.121,0.782)rectangle++(0.009,0.009) (0.058,0.889)rectangle++(0.009,0.009) (0.027,0.942)rectangle++(0.009,0.009) (0.011,0.969)rectangle++(0.009,0.009)
  (0.502,0.385)rectangle++(0.009,0.009) (0.249,0.19)rectangle++(0.009,0.009) (0.122,0.093)rectangle++(0.009,0.009) (0.059,0.544)rectangle++(0.009,0.009) (0.521,0.215)rectangle++(0.009,0.009)
  (0.706,0.083)rectangle++(0.009,0.009) (0.351,0.539)rectangle++(0.009,0.009) (0.638,0.213)rectangle++(0.009,0.009) (0.317,0.604)rectangle++(0.009,0.009) (0.156,0.3)rectangle++(0.009,0.009)
  (0.56,0.117)rectangle++(0.009,0.009) (0.721,0.044)rectangle++(0.009,0.009) (0.358,0.52)rectangle++(0.009,0.009) (0.177,0.258)rectangle++(0.009,0.009) (0.086,0.127)rectangle++(0.009,0.009)
  (0.041,0.061)rectangle++(0.009,0.009) (0.018,0.528)rectangle++(0.009,0.009) (0.007,0.262)rectangle++(0.009,0.009) (0.001,0.129)rectangle++(0.009,0.009) (-0.002,0.562)rectangle++(0.009,0.009)
  (-0.003,0.279)rectangle++(0.009,0.009) (0.496,0.109)rectangle++(0.009,0.009) (0.246,0.552)rectangle++(0.009,0.009) (0.596,0.218)rectangle++(0.009,0.009) (0.296,0.607)rectangle++(0.009,0.009)
  (0.146,0.301)rectangle++(0.009,0.009) (0.071,0.648)rectangle++(0.009,0.009) (0.033,0.822)rectangle++(0.009,0.009) (0.014,0.409)rectangle++(0.009,0.009) (0.005,0.202)rectangle++(0.009,0.009)
  (0,0.099)rectangle++(0.009,0.009) (-0.002,0.547)rectangle++(0.009,0.009) (-0.003,0.271)rectangle++(0.009,0.009) (0.496,0.106)rectangle++(0.009,0.009) (0.696,0.04)rectangle++(0.009,0.009)
  (0.346,0.018)rectangle++(0.009,0.009) (0.636,0.004)rectangle++(0.009,0.009) (0.316,0.5)rectangle++(0.009,0.009) (0.156,0.248)rectangle++(0.009,0.009) (0.56,0.096)rectangle++(0.009,0.009)
  (0.278,0.546)rectangle++(0.009,0.009) (0.608,0.216)rectangle++(0.009,0.009) (0.302,0.606)rectangle++(0.009,0.009) (0.149,0.301)rectangle++(0.009,0.009) (0.072,0.148)rectangle++(0.009,0.009)
  (0.034,0.572)rectangle++(0.009,0.009) (0.015,0.284)rectangle++(0.009,0.009) (0.503,0.111)rectangle++(0.009,0.009) (0.249,0.053)rectangle++(0.009,0.009) (0.122,0.024)rectangle++(0.009,0.009)
  (0.059,0.01)rectangle++(0.009,0.009) (0.027,0.003)rectangle++(0.009,0.009) (0.508,-0.002)rectangle++(0.009,0.009) (0.701,-0.003)rectangle++(0.009,0.009) (0.348,-0.004)rectangle++(0.009,0.009)
  (0.172,-0.004)rectangle++(0.009,0.009) (0.084,0.496)rectangle++(0.009,0.009) (0.04,0.746)rectangle++(0.009,0.009) (0.513,0.296)rectangle++(0.009,0.009) (0.254,0.646)rectangle++(0.009,0.009)
  (0.125,0.821)rectangle++(0.009,0.009) (0.547,0.326)rectangle++(0.009,0.009) (0.716,0.128)rectangle++(0.009,0.009) (0.356,0.062)rectangle++(0.009,0.009) (0.176,0.029)rectangle++(0.009,0.009)
  (0.086,0.012)rectangle++(0.009,0.009) (0.532,0.002)rectangle++(0.009,0.009) (0.71,-0.002)rectangle++(0.009,0.009) (0.353,-0.003)rectangle++(0.009,0.009) (0.174,0.496)rectangle++(0.009,0.009)
  (0.567,0.196)rectangle++(0.009,0.009) (0.724,0.076)rectangle++(0.009,0.009) (0.36,0.536)rectangle++(0.009,0.009) (0.178,0.266)rectangle++(0.009,0.009) (0.087,0.631)rectangle++(0.009,0.009)
  (0.041,0.813)rectangle++(0.009,0.009) (0.514,0.323)rectangle++(0.009,0.009) (0.255,0.159)rectangle++(0.009,0.009) (0.125,0.577)rectangle++(0.009,0.009) (0.06,0.286)rectangle++(0.009,0.009)
  (0.028,0.141)rectangle++(0.009,0.009) (0.012,0.068)rectangle++(0.009,0.009) (0.004,0.032)rectangle++(0.009,0.009) (0,0.014)rectangle++(0.009,0.009) (-0.002,0.505)rectangle++(0.009,0.009)
  (-0.003,0.25)rectangle++(0.009,0.009) (-0.004,0.623)rectangle++(0.009,0.009) (-0.004,0.309)rectangle++(0.009,0.009) (-0.004,0.152)rectangle++(0.009,0.009) (-0.004,0.074)rectangle++(0.009,0.009)
  (-0.004,0.535)rectangle++(0.009,0.009) (-0.004,0.765)rectangle++(0.009,0.009) (-0.004,0.38)rectangle++(0.009,0.009) (0.496,0.149)rectangle++(0.009,0.009) (0.246,0.072)rectangle++(0.009,0.009)
  (0.121,0.534)rectangle++(0.009,0.009) (0.058,0.765)rectangle++(0.009,0.009) (0.521,0.303)rectangle++(0.009,0.009) (0.258,0.649)rectangle++(0.009,0.009) (0.127,0.322)rectangle++(0.009,0.009)
  (0.061,0.159)rectangle++(0.009,0.009) (0.028,0.577)rectangle++(0.009,0.009) (0.012,0.786)rectangle++(0.009,0.009) (0.502,0.312)rectangle++(0.009,0.009) (0.249,0.154)rectangle++(0.009,0.009)
  (0.597,0.059)rectangle++(0.009,0.009) (0.296,0.527)rectangle++(0.009,0.009) (0.616,0.208)rectangle++(0.009,0.009) (0.744,0.081)rectangle++(0.009,0.009) (0.37,0.038)rectangle++(0.009,0.009)
  (0.183,0.517)rectangle++(0.009,0.009) (0.089,0.256)rectangle++(0.009,0.009) (0.042,0.626)rectangle++(0.009,0.009) (0.514,0.248)rectangle++(0.009,0.009) (0.255,0.122)rectangle++(0.009,0.009)
  (0.125,0.559)rectangle++(0.009,0.009) (0.06,0.777)rectangle++(0.009,0.009) (0.521,0.308)rectangle++(0.009,0.009) (0.706,0.121)rectangle++(0.009,0.009) (0.351,0.558)rectangle++(0.009,0.009)
  (0.173,0.777)rectangle++(0.009,0.009) (0.084,0.386)rectangle++(0.009,0.009) (0.04,0.191)rectangle++(0.009,0.009) (0.018,0.093)rectangle++(0.009,0.009) (0.007,0.544)rectangle++(0.009,0.009)
  (0.5,0.215)rectangle++(0.009,0.009) (0.248,0.605)rectangle++(0.009,0.009) (0.122,0.8)rectangle++(0.009,0.009) (0.059,0.898)rectangle++(0.009,0.009) (0.027,0.447)rectangle++(0.009,0.009)
  (0.011,0.721)rectangle++(0.009,0.009) (0.003,0.858)rectangle++(0.009,0.009) (0.499,0.341)rectangle++(0.009,0.009) (0.247,0.168)rectangle++(0.009,0.009) (0.596,0.065)rectangle++(0.009,0.009)
  (0.296,0.53)rectangle++(0.009,0.009) (0.616,0.209)rectangle++(0.009,0.009) (0.744,0.081)rectangle++(0.009,0.009) (0.37,0.038)rectangle++(0.009,0.009) (0.645,0.013)rectangle++(0.009,0.009)
  (0.32,0.004)rectangle++(0.009,0.009) (0.625,-0.001)rectangle++(0.009,0.009) (0.31,0.497)rectangle++(0.009,0.009) (0.621,0.196)rectangle++(0.009,0.009) (0.746,0.076)rectangle++(0.009,0.009)
  (0.371,0.036)rectangle++(0.009,0.009) (0.183,0.516)rectangle++(0.009,0.009) (0.571,0.204)rectangle++(0.009,0.009) (0.283,0.1)rectangle++(0.009,0.009) (0.139,0.048)rectangle++(0.009,0.009)
  (0.553,0.016)rectangle++(0.009,0.009) (0.274,0.006)rectangle++(0.009,0.009) (0.135,0.001)rectangle++(0.009,0.009) (0.065,-0.002)rectangle++(0.009,0.009) (0.03,0.497)rectangle++(0.009,0.009)
  (0.509,0.196)rectangle++(0.009,0.009) (0.252,0.596)rectangle++(0.009,0.009) (0.124,0.796)rectangle++(0.009,0.009) (0.547,0.316)rectangle++(0.009,0.009) (0.716,0.124)rectangle++(0.009,0.009)
  (0.356,0.56)rectangle++(0.009,0.009) (0.176,0.278)rectangle++(0.009,0.009) (0.086,0.637)rectangle++(0.009,0.009) (0.041,0.316)rectangle++(0.009,0.009) (0.018,0.156)rectangle++(0.009,0.009)
  (0.007,0.576)rectangle++(0.009,0.009) (0.001,0.786)rectangle++(0.009,0.009) (-0.002,0.891)rectangle++(0.009,0.009) (-0.003,0.943)rectangle++(0.009,0.009) (0.496,0.375)rectangle++(0.009,0.009)
  (0.246,0.685)rectangle++(0.009,0.009) (0.596,0.271)rectangle++(0.009,0.009) (0.296,0.133)rectangle++(0.009,0.009) (0.146,0.064)rectangle++(0.009,0.009) (0.071,0.53)rectangle++(0.009,0.009)
  (0.033,0.763)rectangle++(0.009,0.009) (0.014,0.379)rectangle++(0.009,0.009) (0.005,0.187)rectangle++(0.009,0.009) (0.499,0.072)rectangle++(0.009,0.009) (0.247,0.034)rectangle++(0.009,0.009)
  (0.121,0.015)rectangle++(0.009,0.009) (0.058,0.005)rectangle++(0.009,0.009) (0.521,-0.001)rectangle++(0.009,0.009) (0.706,-0.003)rectangle++(0.009,0.009) (0.351,0.496)rectangle++(0.009,0.009)
  (0.638,0.196)rectangle++(0.009,0.009) (0.317,0.596)rectangle++(0.009,0.009) (0.156,0.296)rectangle++(0.009,0.009) (0.56,0.116)rectangle++(0.009,0.009) (0.278,0.056)rectangle++(0.009,0.009)
  (0.137,0.026)rectangle++(0.009,0.009) (0.066,0.511)rectangle++(0.009,0.009) (0.031,0.753)rectangle++(0.009,0.009) (0.013,0.874)rectangle++(0.009,0.009) (0.004,0.935)rectangle++(0.009,0.009)
  (0,0.965)rectangle++(0.009,0.009) (-0.002,0.48)rectangle++(0.009,0.009) (-0.003,0.738)rectangle++(0.009,0.009) (-0.004,0.367)rectangle++(0.009,0.009) (-0.004,0.681)rectangle++(0.009,0.009)
  (0.496,0.27)rectangle++(0.009,0.009) (0.696,0.105)rectangle++(0.009,0.009) (0.346,0.55)rectangle++(0.009,0.009) (0.636,0.217)rectangle++(0.009,0.009) (0.316,0.606)rectangle++(0.009,0.009)
  (0.624,0.24)rectangle++(0.009,0.009) (0.31,0.118)rectangle++(0.009,0.009) (0.153,0.557)rectangle++(0.009,0.009) (0.558,0.22)rectangle++(0.009,0.009) (0.277,0.108)rectangle++(0.009,0.009)
  (0.136,0.552)rectangle++(0.009,0.009) (0.066,0.774)rectangle++(0.009,0.009) (0.031,0.885)rectangle++(0.009,0.009) (0.013,0.94)rectangle++(0.009,0.009) (0.004,0.968)rectangle++(0.009,0.009)
  (0.499,0.384)rectangle++(0.009,0.009) (0.247,0.69)rectangle++(0.009,0.009) (0.121,0.343)rectangle++(0.009,0.009) (0.058,0.669)rectangle++(0.009,0.009) (0.027,0.832)rectangle++(0.009,0.009)
  (0.011,0.914)rectangle++(0.009,0.009) (0.502,0.363)rectangle++(0.009,0.009) (0.249,0.179)rectangle++(0.009,0.009) (0.122,0.087)rectangle++(0.009,0.009) (0.546,0.032)rectangle++(0.009,0.009)
  (0.271,0.514)rectangle++(0.009,0.009) (0.133,0.255)rectangle++(0.009,0.009) (0.064,0.625)rectangle++(0.009,0.009) (0.523,0.247)rectangle++(0.009,0.009) (0.707,0.096)rectangle++(0.009,0.009)
  (0.351,0.546)rectangle++(0.009,0.009) (0.173,0.271)rectangle++(0.009,0.009) (0.084,0.133)rectangle++(0.009,0.009) (0.531,0.051)rectangle++(0.009,0.009) (0.263,0.023)rectangle++(0.009,0.009)
  (0.603,0.007)rectangle++(0.009,0.009) (0.299,0.001)rectangle++(0.009,0.009) (0.147,-0.002)rectangle++(0.009,0.009) (0.071,0.497)rectangle++(0.009,0.009) (0.526,0.196)rectangle++(0.009,0.009)
  (0.261,0.596)rectangle++(0.009,0.009) (0.128,0.296)rectangle++(0.009,0.009) (0.549,0.116)rectangle++(0.009,0.009) (0.272,0.556)rectangle++(0.009,0.009) (0.134,0.276)rectangle++(0.009,0.009)
  (0.551,0.108)rectangle++(0.009,0.009) (0.273,0.552)rectangle++(0.009,0.009) (0.134,0.774)rectangle++(0.009,0.009) (0.065,0.385)rectangle++(0.009,0.009) (0.03,0.19)rectangle++(0.009,0.009)
  (0.013,0.093)rectangle++(0.009,0.009) (0.004,0.044)rectangle++(0.009,0.009) (0,0.02)rectangle++(0.009,0.009) (-0.002,0.508)rectangle++(0.009,0.009) (-0.003,0.252)rectangle++(0.009,0.009)
  (-0.004,0.624)rectangle++(0.009,0.009) (-0.004,0.31)rectangle++(0.009,0.009) (-0.004,0.653)rectangle++(0.009,0.009) (-0.004,0.824)rectangle++(0.009,0.009) (0.496,0.327)rectangle++(0.009,0.009)
  (0.696,0.128)rectangle++(0.009,0.009) (0.346,0.562)rectangle++(0.009,0.009) (0.171,0.779)rectangle++(0.009,0.009) (0.083,0.387)rectangle++(0.009,0.009) (0.039,0.691)rectangle++(0.009,0.009)
  (0.017,0.843)rectangle++(0.009,0.009) (0.006,0.919)rectangle++(0.009,0.009) (0.001,0.957)rectangle++(0.009,0.009) (-0.002,0.476)rectangle++(0.009,0.009) (0.497,0.188)rectangle++(0.009,0.009)
  (0.246,0.092)rectangle++(0.009,0.009) (0.121,0.044)rectangle++(0.009,0.009) (0.058,0.02)rectangle++(0.009,0.009) (0.027,0.008)rectangle++(0.009,0.009) (0.011,0.502)rectangle++(0.009,0.009)
  (0.502,0.198)rectangle++(0.009,0.009) (0.249,0.097)rectangle++(0.009,0.009) (0.122,0.546)rectangle++(0.009,0.009) (0.059,0.771)rectangle++(0.009,0.009) (0.521,0.306)rectangle++(0.009,0.009)
  (0.706,0.12)rectangle++(0.009,0.009) (0.351,0.558)rectangle++(0.009,0.009) (0.638,0.22)rectangle++(0.009,0.009) (0.317,0.608)rectangle++(0.009,0.009) (0.156,0.802)rectangle++(0.009,0.009)
  (0.076,0.399)rectangle++(0.009,0.009) (0.036,0.197)rectangle++(0.009,0.009) (0.016,0.096)rectangle++(0.009,0.009) (0.006,0.546)rectangle++(0.009,0.009) (0.001,0.771)rectangle++(0.009,0.009)
  (-0.002,0.383)rectangle++(0.009,0.009) (-0.003,0.189)rectangle++(0.009,0.009) (-0.004,0.592)rectangle++(0.009,0.009) (-0.004,0.794)rectangle++(0.009,0.009) (0.496,0.315)rectangle++(0.009,0.009)
  (0.246,0.155)rectangle++(0.009,0.009) (0.596,0.059)rectangle++(0.009,0.009) (0.296,0.527)rectangle++(0.009,0.009) (0.146,0.761)rectangle++(0.009,0.009) (0.556,0.302)rectangle++(0.009,0.009)
  (0.276,0.649)rectangle++(0.009,0.009) (0.136,0.822)rectangle++(0.009,0.009) (0.066,0.409)rectangle++(0.009,0.009) (0.031,0.702)rectangle++(0.009,0.009) (0.013,0.849)rectangle++(0.009,0.009)
  (0.004,0.922)rectangle++(0.009,0.009) (0,0.959)rectangle++(0.009,0.009) (0.497,0.381)rectangle++(0.009,0.009) (0.246,0.188)rectangle++(0.009,0.009) (0.596,0.073)rectangle++(0.009,0.009)
  (0.296,0.034)rectangle++(0.009,0.009) (0.146,0.515)rectangle++(0.009,0.009) (0.071,0.255)rectangle++(0.009,0.009) (0.526,0.099)rectangle++(0.009,0.009) (0.708,0.037)rectangle++(0.009,0.009)
  (0.78,0.012)rectangle++(0.009,0.009) (0.388,0.004)rectangle++(0.009,0.009) (0.192,0)rectangle++(0.009,0.009) (0.574,-0.003)rectangle++(0.009,0.009) (0.285,-0.004)rectangle++(0.009,0.009)
  (0.611,-0.004)rectangle++(0.009,0.009) (0.303,0.496)rectangle++(0.009,0.009) (0.619,0.196)rectangle++(0.009,0.009) (0.307,0.096)rectangle++(0.009,0.009) (0.151,0.046)rectangle++(0.009,0.009)
  (0.558,0.016)rectangle++(0.009,0.009) (0.277,0.506)rectangle++(0.009,0.009) (0.608,0.2)rectangle++(0.009,0.009) (0.302,0.598)rectangle++(0.009,0.009) (0.149,0.297)rectangle++(0.009,0.009)
  (0.072,0.646)rectangle++(0.009,0.009) (0.526,0.256)rectangle++(0.009,0.009) (0.261,0.126)rectangle++(0.009,0.009) (0.128,0.061)rectangle++(0.009,0.009) (0.062,0.528)rectangle++(0.009,0.009)
  (0.029,0.262)rectangle++(0.009,0.009) (0.012,0.629)rectangle++(0.009,0.009) (0.004,0.312)rectangle++(0.009,0.009) (0.499,0.122)rectangle++(0.009,0.009) (0.247,0.559)rectangle++(0.009,0.009)
  (0.121,0.277)rectangle++(0.009,0.009) (0.058,0.136)rectangle++(0.009,0.009) (0.521,0.052)rectangle++(0.009,0.009) (0.258,0.524)rectangle++(0.009,0.009) (0.127,0.26)rectangle++(0.009,0.009)
  (0.061,0.128)rectangle++(0.009,0.009) (0.028,0.062)rectangle++(0.009,0.009) (0.012,0.529)rectangle++(0.009,0.009) (0.004,0.262)rectangle++(0.009,0.009) (0,0.129)rectangle++(0.009,0.009)
  (-0.002,0.062)rectangle++(0.009,0.009) (-0.003,0.529)rectangle++(0.009,0.009) (-0.004,0.262)rectangle++(0.009,0.009) (0.496,0.102)rectangle++(0.009,0.009) (0.246,0.549)rectangle++(0.009,0.009)
  (0.121,0.272)rectangle++(0.009,0.009) (0.058,0.634)rectangle++(0.009,0.009) (0.027,0.815)rectangle++(0.009,0.009) (0.011,0.905)rectangle++(0.009,0.009) (0.003,0.45)rectangle++(0.009,0.009)
  (0.499,0.177)rectangle++(0.009,0.009) (0.247,0.086)rectangle++(0.009,0.009) (0.121,0.541)rectangle++(0.009,0.009) (0.058,0.768)rectangle++(0.009,0.009) (0.027,0.382)rectangle++(0.009,0.009)
  (0.508,0.15)rectangle++(0.009,0.009) (0.252,0.573)rectangle++(0.009,0.009) (0.124,0.784)rectangle++(0.009,0.009) (0.06,0.39)rectangle++(0.009,0.009) (0.028,0.693)rectangle++(0.009,0.009)
  (0.012,0.344)rectangle++(0.009,0.009) (0.004,0.17)rectangle++(0.009,0.009) (0,0.083)rectangle++(0.009,0.009) (-0.002,0.539)rectangle++(0.009,0.009) (-0.003,0.267)rectangle++(0.009,0.009)
  (-0.004,0.631)rectangle++(0.009,0.009) (-0.004,0.813)rectangle++(0.009,0.009) (-0.004,0.904)rectangle++(0.009,0.009) (0.496,0.359)rectangle++(0.009,0.009) (0.696,0.141)rectangle++(0.009,0.009)
  (0.346,0.068)rectangle++(0.009,0.009) (0.171,0.032)rectangle++(0.009,0.009) (0.083,0.014)rectangle++(0.009,0.009) (0.039,0.505)rectangle++(0.009,0.009) (0.513,0.199)rectangle++(0.009,0.009)
  (0.254,0.097)rectangle++(0.009,0.009) (0.125,0.046)rectangle++(0.009,0.009) (0.06,0.521)rectangle++(0.009,0.009) (0.028,0.758)rectangle++(0.009,0.009) (0.012,0.377)rectangle++(0.009,0.009)
  (0.004,0.686)rectangle++(0.009,0.009) (0.499,0.272)rectangle++(0.009,0.009) (0.697,0.106)rectangle++(0.009,0.009) (0.346,0.551)rectangle++(0.009,0.009) (0.171,0.273)rectangle++(0.009,0.009)
  (0.083,0.134)rectangle++(0.009,0.009) (0.039,0.065)rectangle++(0.009,0.009) (0.017,0.53)rectangle++(0.009,0.009) (0.006,0.763)rectangle++(0.009,0.009) (0.001,0.379)rectangle++(0.009,0.009)
  (-0.002,0.187)rectangle++(0.009,0.009) (-0.003,0.091)rectangle++(0.009,0.009) (-0.004,0.543)rectangle++(0.009,0.009) (0.496,0.215)rectangle++(0.009,0.009) (0.696,0.083)rectangle++(0.009,0.009)
  (0.346,0.039)rectangle++(0.009,0.009) (0.171,0.517)rectangle++(0.009,0.009) (0.083,0.756)rectangle++(0.009,0.009) (0.039,0.376)rectangle++(0.009,0.009) (0.017,0.186)rectangle++(0.009,0.009)
  (0.006,0.091)rectangle++(0.009,0.009) (0.5,0.034)rectangle++(0.009,0.009) (0.248,0.515)rectangle++(0.009,0.009) (0.596,0.203)rectangle++(0.009,0.009) (0.296,0.599)rectangle++(0.009,0.009)
  (0.146,0.297)rectangle++(0.009,0.009) (0.071,0.146)rectangle++(0.009,0.009) (0.033,0.071)rectangle++(0.009,0.009) (0.014,0.533)rectangle++(0.009,0.009) (0.005,0.764)rectangle++(0.009,0.009)
  (0,0.38)rectangle++(0.009,0.009) (-0.002,0.188)rectangle++(0.009,0.009) (-0.003,0.592)rectangle++(0.009,0.009) (-0.004,0.794)rectangle++(0.009,0.009) (-0.004,0.895)rectangle++(0.009,0.009)
  (-0.004,0.445)rectangle++(0.009,0.009) (-0.004,0.22)rectangle++(0.009,0.009) (-0.004,0.608)rectangle++(0.009,0.009) (0.496,0.24)rectangle++(0.009,0.009) (0.246,0.618)rectangle++(0.009,0.009)
  (0.596,0.244)rectangle++(0.009,0.009) (0.296,0.12)rectangle++(0.009,0.009) (0.616,0.045)rectangle++(0.009,0.009) (0.306,0.52)rectangle++(0.009,0.009) (0.151,0.758)rectangle++(0.009,0.009)
  (0.558,0.3)rectangle++(0.009,0.009) (0.72,0.117)rectangle++(0.009,0.009) (0.785,0.044)rectangle++(0.009,0.009) (0.39,0.52)rectangle++(0.009,0.009) (0.193,0.258)rectangle++(0.009,0.009)
  (0.094,0.627)rectangle++(0.009,0.009) (0.045,0.811)rectangle++(0.009,0.009) (0.02,0.903)rectangle++(0.009,0.009) (0.008,0.949)rectangle++(0.009,0.009) (0.5,0.377)rectangle++(0.009,0.009)
  (0.697,0.148)rectangle++(0.009,0.009) (0.346,0.072)rectangle++(0.009,0.009) (0.171,0.034)rectangle++(0.009,0.009) (0.083,0.015)rectangle++(0.009,0.009) (0.039,0.505)rectangle++(0.009,0.009)
  (0.017,0.75)rectangle++(0.009,0.009) (0.006,0.873)rectangle++(0.009,0.009) (0.001,0.934)rectangle++(0.009,0.009) (-0.002,0.465)rectangle++(0.009,0.009) (0.497,0.183)rectangle++(0.009,0.009)
  (0.696,0.071)rectangle++(0.009,0.009) (0.346,0.033)rectangle++(0.009,0.009) (0.171,0.514)rectangle++(0.009,0.009) (0.083,0.755)rectangle++(0.009,0.009) (0.039,0.875)rectangle++(0.009,0.009)
  (0.017,0.435)rectangle++(0.009,0.009) (0.006,0.715)rectangle++(0.009,0.009) (0.5,0.283)rectangle++(0.009,0.009) (0.248,0.139)rectangle++(0.009,0.009) (0.596,0.053)rectangle++(0.009,0.009)
  (0.296,0.524)rectangle++(0.009,0.009) (0.146,0.76)rectangle++(0.009,0.009) (0.556,0.301)rectangle++(0.009,0.009) (0.276,0.148)rectangle++(0.009,0.009) (0.136,0.072)rectangle++(0.009,0.009)
  (0.552,0.026)rectangle++(0.009,0.009) (0.274,0.011)rectangle++(0.009,0.009) (0.607,0.002)rectangle++(0.009,0.009) (0.74,-0.002)rectangle++(0.009,0.009) (0.368,0.497)rectangle++(0.009,0.009)
  (0.182,0.246)rectangle++(0.009,0.009) (0.089,0.121)rectangle++(0.009,0.009) (0.042,0.058)rectangle++(0.009,0.009) (0.514,0.021)rectangle++(0.009,0.009) (0.255,0.508)rectangle++(0.009,0.009)
  (0.125,0.752)rectangle++(0.009,0.009) (0.06,0.374)rectangle++(0.009,0.009) (0.028,0.185)rectangle++(0.009,0.009) (0.012,0.09)rectangle++(0.009,0.009) (0.004,0.043)rectangle++(0.009,0.009)
  (0,0.019)rectangle++(0.009,0.009) (-0.002,0.007)rectangle++(0.009,0.009) (0.496,0)rectangle++(0.009,0.009) (0.246,0.498)rectangle++(0.009,0.009) (0.121,0.747)rectangle++(0.009,0.009)
  (0.546,0.296)rectangle++(0.009,0.009) (0.271,0.646)rectangle++(0.009,0.009) (0.133,0.821)rectangle++(0.009,0.009) (0.064,0.408)rectangle++(0.009,0.009) (0.03,0.702)rectangle++(0.009,0.009)
  (0.509,0.278)rectangle++(0.009,0.009) (0.252,0.637)rectangle++(0.009,0.009) (0.124,0.316)rectangle++(0.009,0.009) (0.547,0.124)rectangle++(0.009,0.009) (0.716,0.047)rectangle++(0.009,0.009)
  (0.356,0.021)rectangle++(0.009,0.009) (0.176,0.008)rectangle++(0.009,0.009) (0.086,0.002)rectangle++(0.009,0.009) (0.532,-0.002)rectangle++(0.009,0.009) (0.264,-0.003)rectangle++(0.009,0.009)
  (0.13,-0.004)rectangle++(0.009,0.009) (0.063,0.496)rectangle++(0.009,0.009) (0.522,0.196)rectangle++(0.009,0.009) (0.259,0.096)rectangle++(0.009,0.009) (0.127,0.546)rectangle++(0.009,0.009)
  (0.061,0.771)rectangle++(0.009,0.009) (0.028,0.383)rectangle++(0.009,0.009) (0.509,0.151)rectangle++(0.009,0.009) (0.252,0.073)rectangle++(0.009,0.009) (0.124,0.034)rectangle++(0.009,0.009)
  (0.06,0.015)rectangle++(0.009,0.009) (0.521,0.003)rectangle++(0.009,0.009) (0.258,-0.001)rectangle++(0.009,0.009) (0.601,-0.003)rectangle++(0.009,0.009) (0.738,-0.004)rectangle++(0.009,0.009)
  (0.792,-0.004)rectangle++(0.009,0.009) (0.394,0.496)rectangle++(0.009,0.009) (0.195,0.746)rectangle++(0.009,0.009) (0.095,0.371)rectangle++(0.009,0.009) (0.045,0.183)rectangle++(0.009,0.009)
  (0.515,0.071)rectangle++(0.009,0.009) (0.255,0.033)rectangle++(0.009,0.009) (0.125,0.514)rectangle++(0.009,0.009) (0.06,0.755)rectangle++(0.009,0.009) (0.028,0.375)rectangle++(0.009,0.009)
  (0.012,0.685)rectangle++(0.009,0.009) (0.004,0.84)rectangle++(0.009,0.009) (0.499,0.333)rectangle++(0.009,0.009) (0.247,0.664)rectangle++(0.009,0.009) (0.121,0.83)rectangle++(0.009,0.009)
  (0.058,0.913)rectangle++(0.009,0.009) (0.027,0.954)rectangle++(0.009,0.009) (0.011,0.975)rectangle++(0.009,0.009) (0.003,0.985)rectangle++(0.009,0.009) (-0.001,0.49)rectangle++(0.009,0.009)
  (-0.003,0.243)rectangle++(0.009,0.009) (-0.004,0.119)rectangle++(0.009,0.009) (0.496,0.045)rectangle++(0.009,0.009) (0.246,0.02)rectangle++(0.009,0.009) (0.596,0.005)rectangle++(0.009,0.009)
  (0.736,-0.001)rectangle++(0.009,0.009) (0.366,0.497)rectangle++(0.009,0.009) (0.181,0.246)rectangle++(0.009,0.009) (0.088,0.121)rectangle++(0.009,0.009) (0.533,0.046)rectangle++(0.009,0.009)
  (0.264,0.521)rectangle++(0.009,0.009) (0.13,0.758)rectangle++(0.009,0.009) (0.549,0.301)rectangle++(0.009,0.009) (0.272,0.148)rectangle++(0.009,0.009) (0.134,0.072)rectangle++(0.009,0.009)
  (0.551,0.026)rectangle++(0.009,0.009) (0.273,0.011)rectangle++(0.009,0.009) (0.607,0.002)rectangle++(0.009,0.009) (0.74,-0.002)rectangle++(0.009,0.009) (0.368,-0.003)rectangle++(0.009,0.009)
  (0.182,0.496)rectangle++(0.009,0.009) (0.089,0.246)rectangle++(0.009,0.009) (0.533,0.096)rectangle++(0.009,0.009) (0.71,0.036)rectangle++(0.009,0.009) (0.353,0.016)rectangle++(0.009,0.009)
  (0.638,0.004)rectangle++(0.009,0.009) (0.753,-0.001)rectangle++(0.009,0.009) (0.374,-0.003)rectangle++(0.009,0.009) (0.185,-0.004)rectangle++(0.009,0.009) (0.09,-0.004)rectangle++(0.009,0.009)
  (0.043,-0.004)rectangle++(0.009,0.009) (0.019,-0.004)rectangle++(0.009,0.009) (0.505,-0.004)rectangle++(0.009,0.009) (0.25,-0.004)rectangle++(0.009,0.009) (0.597,-0.004)rectangle++(0.009,0.009)
  (0.736,-0.004)rectangle++(0.009,0.009) (0.792,-0.004)rectangle++(0.009,0.009) (0.394,0.496)rectangle++(0.009,0.009) (0.655,0.196)rectangle++(0.009,0.009) (0.325,0.096)rectangle++(0.009,0.009)
  (0.16,0.046)rectangle++(0.009,0.009) (0.561,0.016)rectangle++(0.009,0.009) (0.278,0.006)rectangle++(0.009,0.009) (0.137,0.501)rectangle++(0.009,0.009) (0.552,0.198)rectangle++(0.009,0.009)
  (0.274,0.097)rectangle++(0.009,0.009) (0.135,0.046)rectangle++(0.009,0.009) (0.065,0.521)rectangle++(0.009,0.009) (0.03,0.758)rectangle++(0.009,0.009) (0.509,0.301)rectangle++(0.009,0.009)
  (0.252,0.648)rectangle++(0.009,0.009) (0.124,0.322)rectangle++(0.009,0.009) (0.06,0.159)rectangle++(0.009,0.009) (0.028,0.577)rectangle++(0.009,0.009) (0.012,0.786)rectangle++(0.009,0.009)
  (0.502,0.312)rectangle++(0.009,0.009) (0.249,0.654)rectangle++(0.009,0.009) (0.122,0.825)rectangle++(0.009,0.009) (0.059,0.41)rectangle++(0.009,0.009) (0.027,0.203)rectangle++(0.009,0.009)
  (0.011,0.099)rectangle++(0.009,0.009) (0.502,0.037)rectangle++(0.009,0.009) (0.249,0.016)rectangle++(0.009,0.009) (0.122,0.506)rectangle++(0.009,0.009) (0.059,0.251)rectangle++(0.009,0.009)
  (0.027,0.623)rectangle++(0.009,0.009) (0.011,0.309)rectangle++(0.009,0.009) (0.003,0.652)rectangle++(0.009,0.009) (-0.001,0.824)rectangle++(0.009,0.009) (0.497,0.327)rectangle++(0.009,0.009)
  (0.246,0.661)rectangle++(0.009,0.009) (0.121,0.828)rectangle++(0.009,0.009) (0.058,0.412)rectangle++(0.009,0.009) (0.521,0.162)rectangle++(0.009,0.009) (0.706,0.062)rectangle++(0.009,0.009);
\fill[nuage]
  (0.351,0.529)rectangle++(0.009,0.009) (0.173,0.262)rectangle++(0.009,0.009) (0.084,0.629)rectangle++(0.009,0.009) (0.04,0.312)rectangle++(0.009,0.009) (0.018,0.154)rectangle++(0.009,0.009)
  (0.007,0.575)rectangle++(0.009,0.009) (0.001,0.785)rectangle++(0.009,0.009) (-0.002,0.89)rectangle++(0.009,0.009) (-0.003,0.943)rectangle++(0.009,0.009) (0.496,0.374)rectangle++(0.009,0.009)
  (0.246,0.185)rectangle++(0.009,0.009) (0.121,0.59)rectangle++(0.009,0.009) (0.058,0.293)rectangle++(0.009,0.009) (0.027,0.644)rectangle++(0.009,0.009) (0.011,0.82)rectangle++(0.009,0.009)
  (0.502,0.325)rectangle++(0.009,0.009) (0.249,0.16)rectangle++(0.009,0.009) (0.122,0.578)rectangle++(0.009,0.009) (0.059,0.787)rectangle++(0.009,0.009) (0.027,0.391)rectangle++(0.009,0.009)
  (0.011,0.193)rectangle++(0.009,0.009) (0.003,0.094)rectangle++(0.009,0.009) (-0.001,0.045)rectangle++(0.009,0.009) (-0.003,0.52)rectangle++(0.009,0.009) (-0.004,0.758)rectangle++(0.009,0.009)
  (0.496,0.3)rectangle++(0.009,0.009) (0.246,0.148)rectangle++(0.009,0.009) (0.121,0.572)rectangle++(0.009,0.009) (0.058,0.284)rectangle++(0.009,0.009) (0.027,0.64)rectangle++(0.009,0.009)
  (0.508,0.253)rectangle++(0.009,0.009) (0.252,0.124)rectangle++(0.009,0.009) (0.124,0.56)rectangle++(0.009,0.009) (0.06,0.778)rectangle++(0.009,0.009) (0.028,0.887)rectangle++(0.009,0.009)
  (0.012,0.941)rectangle++(0.009,0.009) (0.004,0.968)rectangle++(0.009,0.009) (0.499,0.385)rectangle++(0.009,0.009) (0.247,0.69)rectangle++(0.009,0.009) (0.121,0.843)rectangle++(0.009,0.009)
  (0.058,0.419)rectangle++(0.009,0.009) (0.027,0.207)rectangle++(0.009,0.009) (0.011,0.601)rectangle++(0.009,0.009) (0.003,0.798)rectangle++(0.009,0.009) (-0.001,0.397)rectangle++(0.009,0.009)
  (-0.003,0.696)rectangle++(0.009,0.009) (-0.004,0.346)rectangle++(0.009,0.009) (-0.004,0.671)rectangle++(0.009,0.009) (0.496,0.266)rectangle++(0.009,0.009) (0.246,0.631)rectangle++(0.009,0.009)
  (0.121,0.313)rectangle++(0.009,0.009) (0.058,0.654)rectangle++(0.009,0.009) (0.521,0.259)rectangle++(0.009,0.009) (0.706,0.101)rectangle++(0.009,0.009) (0.351,0.048)rectangle++(0.009,0.009)
  (0.173,0.522)rectangle++(0.009,0.009) (0.084,0.259)rectangle++(0.009,0.009) (0.04,0.127)rectangle++(0.009,0.009) (0.018,0.561)rectangle++(0.009,0.009) (0.007,0.778)rectangle++(0.009,0.009)
  (0.5,0.309)rectangle++(0.009,0.009) (0.248,0.652)rectangle++(0.009,0.009) (0.122,0.824)rectangle++(0.009,0.009) (0.059,0.41)rectangle++(0.009,0.009) (0.027,0.203)rectangle++(0.009,0.009)
  (0.011,0.599)rectangle++(0.009,0.009) (0.502,0.237)rectangle++(0.009,0.009) (0.249,0.116)rectangle++(0.009,0.009) (0.122,0.556)rectangle++(0.009,0.009) (0.546,0.22)rectangle++(0.009,0.009)
  (0.271,0.608)rectangle++(0.009,0.009) (0.606,0.24)rectangle++(0.009,0.009) (0.301,0.618)rectangle++(0.009,0.009) (0.148,0.307)rectangle++(0.009,0.009) (0.072,0.151)rectangle++(0.009,0.009)
  (0.034,0.073)rectangle++(0.009,0.009) (0.015,0.534)rectangle++(0.009,0.009) (0.005,0.765)rectangle++(0.009,0.009) (0,0.38)rectangle++(0.009,0.009) (-0.002,0.188)rectangle++(0.009,0.009)
  (-0.003,0.092)rectangle++(0.009,0.009) (-0.004,0.544)rectangle++(0.009,0.009) (-0.004,0.27)rectangle++(0.009,0.009) (-0.004,0.133)rectangle++(0.009,0.009) (-0.004,0.564)rectangle++(0.009,0.009)
  (-0.004,0.78)rectangle++(0.009,0.009) (0.496,0.309)rectangle++(0.009,0.009) (0.246,0.152)rectangle++(0.009,0.009) (0.121,0.074)rectangle++(0.009,0.009) (0.058,0.035)rectangle++(0.009,0.009)
  (0.027,0.015)rectangle++(0.009,0.009) (0.508,0.003)rectangle++(0.009,0.009) (0.252,0.499)rectangle++(0.009,0.009) (0.124,0.747)rectangle++(0.009,0.009) (0.06,0.371)rectangle++(0.009,0.009)
  (0.521,0.146)rectangle++(0.009,0.009) (0.706,0.056)rectangle++(0.009,0.009) (0.78,0.02)rectangle++(0.009,0.009) (0.388,0.008)rectangle++(0.009,0.009) (0.192,0.502)rectangle++(0.009,0.009)
  (0.094,0.749)rectangle++(0.009,0.009) (0.045,0.872)rectangle++(0.009,0.009) (0.02,0.934)rectangle++(0.009,0.009) (0.505,0.371)rectangle++(0.009,0.009) (0.25,0.183)rectangle++(0.009,0.009)
  (0.123,0.089)rectangle++(0.009,0.009) (0.546,0.033)rectangle++(0.009,0.009) (0.271,0.514)rectangle++(0.009,0.009) (0.606,0.203)rectangle++(0.009,0.009) (0.301,0.599)rectangle++(0.009,0.009)
  (0.148,0.297)rectangle++(0.009,0.009) (0.072,0.646)rectangle++(0.009,0.009) (0.526,0.256)rectangle++(0.009,0.009) (0.261,0.626)rectangle++(0.009,0.009) (0.128,0.311)rectangle++(0.009,0.009)
  (0.062,0.653)rectangle++(0.009,0.009) (0.522,0.259)rectangle++(0.009,0.009) (0.259,0.627)rectangle++(0.009,0.009) (0.127,0.811)rectangle++(0.009,0.009) (0.061,0.903)rectangle++(0.009,0.009)
  (0.028,0.949)rectangle++(0.009,0.009) (0.012,0.472)rectangle++(0.009,0.009) (0.004,0.234)rectangle++(0.009,0.009) (0.499,0.091)rectangle++(0.009,0.009) (0.247,0.543)rectangle++(0.009,0.009)
  (0.121,0.269)rectangle++(0.009,0.009) (0.058,0.132)rectangle++(0.009,0.009) (0.521,0.05)rectangle++(0.009,0.009) (0.706,0.017)rectangle++(0.009,0.009) (0.351,0.006)rectangle++(0.009,0.009)
  (0.173,0.001)rectangle++(0.009,0.009) (0.084,0.498)rectangle++(0.009,0.009) (0.04,0.247)rectangle++(0.009,0.009) (0.513,0.096)rectangle++(0.009,0.009) (0.703,0.036)rectangle++(0.009,0.009)
  (0.349,0.516)rectangle++(0.009,0.009) (0.172,0.756)rectangle++(0.009,0.009) (0.084,0.876)rectangle++(0.009,0.009) (0.04,0.936)rectangle++(0.009,0.009) (0.018,0.466)rectangle++(0.009,0.009)
  (0.007,0.731)rectangle++(0.009,0.009) (0.001,0.363)rectangle++(0.009,0.009) (0.498,0.143)rectangle++(0.009,0.009) (0.247,0.069)rectangle++(0.009,0.009) (0.121,0.532)rectangle++(0.009,0.009)
  (0.058,0.764)rectangle++(0.009,0.009) (0.027,0.38)rectangle++(0.009,0.009) (0.508,0.149)rectangle++(0.009,0.009) (0.252,0.072)rectangle++(0.009,0.009) (0.124,0.534)rectangle++(0.009,0.009)
  (0.06,0.765)rectangle++(0.009,0.009) (0.028,0.88)rectangle++(0.009,0.009) (0.012,0.438)rectangle++(0.009,0.009) (0.502,0.172)rectangle++(0.009,0.009) (0.249,0.584)rectangle++(0.009,0.009)
  (0.122,0.29)rectangle++(0.009,0.009) (0.059,0.643)rectangle++(0.009,0.009) (0.027,0.319)rectangle++(0.009,0.009) (0.508,0.125)rectangle++(0.009,0.009) (0.701,0.047)rectangle++(0.009,0.009)
  (0.778,0.016)rectangle++(0.009,0.009) (0.808,0.004)rectangle++(0.009,0.009) (0.402,0.5)rectangle++(0.009,0.009) (0.199,0.748)rectangle++(0.009,0.009) (0.577,0.296)rectangle++(0.009,0.009)
  (0.286,0.146)rectangle++(0.009,0.009) (0.141,0.571)rectangle++(0.009,0.009) (0.554,0.226)rectangle++(0.009,0.009) (0.275,0.611)rectangle++(0.009,0.009) (0.135,0.303)rectangle++(0.009,0.009)
  (0.065,0.149)rectangle++(0.009,0.009) (0.03,0.572)rectangle++(0.009,0.009) (0.013,0.284)rectangle++(0.009,0.009) (0.004,0.14)rectangle++(0.009,0.009) (0,0.068)rectangle++(0.009,0.009)
  (0.497,0.024)rectangle++(0.009,0.009) (0.696,0.007)rectangle++(0.009,0.009) (0.346,0.501)rectangle++(0.009,0.009) (0.171,0.748)rectangle++(0.009,0.009) (0.566,0.297)rectangle++(0.009,0.009)
  (0.724,0.116)rectangle++(0.009,0.009) (0.787,0.044)rectangle++(0.009,0.009) (0.391,0.02)rectangle++(0.009,0.009) (0.193,0.008)rectangle++(0.009,0.009) (0.094,0.002)rectangle++(0.009,0.009)
  (0.535,-0.002)rectangle++(0.009,0.009) (0.265,0.497)rectangle++(0.009,0.009) (0.603,0.196)rectangle++(0.009,0.009) (0.739,0.076)rectangle++(0.009,0.009) (0.367,0.536)rectangle++(0.009,0.009)
  (0.181,0.266)rectangle++(0.009,0.009) (0.088,0.631)rectangle++(0.009,0.009) (0.533,0.25)rectangle++(0.009,0.009) (0.264,0.123)rectangle++(0.009,0.009) (0.13,0.059)rectangle++(0.009,0.009)
  (0.549,0.021)rectangle++(0.009,0.009) (0.717,0.006)rectangle++(0.009,0.009) (0.356,0.001)rectangle++(0.009,0.009) (0.176,-0.002)rectangle++(0.009,0.009) (0.086,0.497)rectangle++(0.009,0.009)
  (0.041,0.746)rectangle++(0.009,0.009) (0.018,0.871)rectangle++(0.009,0.009) (0.007,0.433)rectangle++(0.009,0.009) (0.5,0.171)rectangle++(0.009,0.009) (0.248,0.583)rectangle++(0.009,0.009)
  (0.596,0.231)rectangle++(0.009,0.009) (0.296,0.613)rectangle++(0.009,0.009) (0.146,0.304)rectangle++(0.009,0.009) (0.556,0.119)rectangle++(0.009,0.009) (0.276,0.057)rectangle++(0.009,0.009)
  (0.136,0.026)rectangle++(0.009,0.009) (0.066,0.511)rectangle++(0.009,0.009) (0.031,0.753)rectangle++(0.009,0.009) (0.013,0.374)rectangle++(0.009,0.009) (0.503,0.147)rectangle++(0.009,0.009)
  (0.249,0.071)rectangle++(0.009,0.009) (0.122,0.033)rectangle++(0.009,0.009) (0.059,0.514)rectangle++(0.009,0.009) (0.027,0.255)rectangle++(0.009,0.009) (0.011,0.625)rectangle++(0.009,0.009)
  (0.502,0.247)rectangle++(0.009,0.009) (0.249,0.621)rectangle++(0.009,0.009) (0.122,0.808)rectangle++(0.009,0.009) (0.546,0.321)rectangle++(0.009,0.009) (0.271,0.158)rectangle++(0.009,0.009)
  (0.606,0.061)rectangle++(0.009,0.009) (0.301,0.528)rectangle++(0.009,0.009) (0.148,0.762)rectangle++(0.009,0.009) (0.072,0.379)rectangle++(0.009,0.009) (0.034,0.687)rectangle++(0.009,0.009)
  (0.015,0.341)rectangle++(0.009,0.009) (0.005,0.168)rectangle++(0.009,0.009) (0,0.582)rectangle++(0.009,0.009) (-0.002,0.289)rectangle++(0.009,0.009) (-0.003,0.642)rectangle++(0.009,0.009)
  (-0.004,0.319)rectangle++(0.009,0.009) (-0.004,0.657)rectangle++(0.009,0.009) (-0.004,0.826)rectangle++(0.009,0.009) (-0.004,0.911)rectangle++(0.009,0.009) (-0.004,0.953)rectangle++(0.009,0.009)
  (-0.004,0.974)rectangle++(0.009,0.009) (0.496,0.387)rectangle++(0.009,0.009) (0.246,0.691)rectangle++(0.009,0.009) (0.121,0.343)rectangle++(0.009,0.009) (0.058,0.669)rectangle++(0.009,0.009)
  (0.521,0.265)rectangle++(0.009,0.009) (0.706,0.103)rectangle++(0.009,0.009) (0.351,0.549)rectangle++(0.009,0.009) (0.173,0.272)rectangle++(0.009,0.009) (0.084,0.634)rectangle++(0.009,0.009)
  (0.04,0.815)rectangle++(0.009,0.009) (0.018,0.905)rectangle++(0.009,0.009) (0.007,0.95)rectangle++(0.009,0.009) (0.001,0.973)rectangle++(0.009,0.009) (-0.002,0.484)rectangle++(0.009,0.009)
  (-0.003,0.24)rectangle++(0.009,0.009) (-0.004,0.118)rectangle++(0.009,0.009) (0.496,0.044)rectangle++(0.009,0.009) (0.246,0.52)rectangle++(0.009,0.009) (0.121,0.758)rectangle++(0.009,0.009)
  (0.058,0.877)rectangle++(0.009,0.009) (0.521,0.348)rectangle++(0.009,0.009) (0.258,0.172)rectangle++(0.009,0.009) (0.601,0.066)rectangle++(0.009,0.009) (0.298,0.531)rectangle++(0.009,0.009)
  (0.147,0.763)rectangle++(0.009,0.009) (0.071,0.879)rectangle++(0.009,0.009) (0.033,0.937)rectangle++(0.009,0.009) (0.511,0.372)rectangle++(0.009,0.009) (0.702,0.146)rectangle++(0.009,0.009)
  (0.349,0.071)rectangle++(0.009,0.009) (0.172,0.533)rectangle++(0.009,0.009) (0.084,0.764)rectangle++(0.009,0.009) (0.04,0.38)rectangle++(0.009,0.009) (0.018,0.688)rectangle++(0.009,0.009)
  (0.007,0.342)rectangle++(0.009,0.009) (0.001,0.169)rectangle++(0.009,0.009) (0.498,0.065)rectangle++(0.009,0.009) (0.247,0.53)rectangle++(0.009,0.009) (0.121,0.263)rectangle++(0.009,0.009)
  (0.058,0.629)rectangle++(0.009,0.009) (0.027,0.812)rectangle++(0.009,0.009) (0.508,0.322)rectangle++(0.009,0.009) (0.701,0.126)rectangle++(0.009,0.009) (0.348,0.561)rectangle++(0.009,0.009)
  (0.172,0.778)rectangle++(0.009,0.009) (0.084,0.887)rectangle++(0.009,0.009) (0.04,0.441)rectangle++(0.009,0.009) (0.018,0.218)rectangle++(0.009,0.009) (0.007,0.607)rectangle++(0.009,0.009)
  (0.001,0.301)rectangle++(0.009,0.009) (0.498,0.118)rectangle++(0.009,0.009) (0.247,0.557)rectangle++(0.009,0.009) (0.121,0.276)rectangle++(0.009,0.009) (0.058,0.136)rectangle++(0.009,0.009)
  (0.027,0.566)rectangle++(0.009,0.009) (0.508,0.224)rectangle++(0.009,0.009) (0.252,0.61)rectangle++(0.009,0.009) (0.124,0.303)rectangle++(0.009,0.009) (0.06,0.149)rectangle++(0.009,0.009)
  (0.028,0.572)rectangle++(0.009,0.009) (0.012,0.284)rectangle++(0.009,0.009) (0.004,0.14)rectangle++(0.009,0.009) (0,0.568)rectangle++(0.009,0.009) (-0.002,0.282)rectangle++(0.009,0.009)
  (-0.003,0.139)rectangle++(0.009,0.009) (-0.004,0.067)rectangle++(0.009,0.009) (-0.004,0.031)rectangle++(0.009,0.009) (-0.004,0.013)rectangle++(0.009,0.009) (-0.004,0.504)rectangle++(0.009,0.009)
  (-0.004,0.75)rectangle++(0.009,0.009) (-0.004,0.373)rectangle++(0.009,0.009) (-0.004,0.684)rectangle++(0.009,0.009) (-0.004,0.84)rectangle++(0.009,0.009) (-0.004,0.918)rectangle++(0.009,0.009)
  (-0.004,0.457)rectangle++(0.009,0.009) (-0.004,0.726)rectangle++(0.009,0.009) (-0.004,0.361)rectangle++(0.009,0.009) (0.496,0.142)rectangle++(0.009,0.009) (0.246,0.069)rectangle++(0.009,0.009)
  (0.121,0.532)rectangle++(0.009,0.009) (0.546,0.21)rectangle++(0.009,0.009) (0.271,0.103)rectangle++(0.009,0.009) (0.133,0.549)rectangle++(0.009,0.009) (0.064,0.772)rectangle++(0.009,0.009)
  (0.03,0.884)rectangle++(0.009,0.009) (0.509,0.351)rectangle++(0.009,0.009) (0.252,0.673)rectangle++(0.009,0.009) (0.124,0.334)rectangle++(0.009,0.009) (0.06,0.665)rectangle++(0.009,0.009)
  (0.028,0.33)rectangle++(0.009,0.009) (0.012,0.663)rectangle++(0.009,0.009) (0.004,0.329)rectangle++(0.009,0.009) (0,0.662)rectangle++(0.009,0.009) (-0.002,0.829)rectangle++(0.009,0.009)
  (-0.003,0.412)rectangle++(0.009,0.009) (0.496,0.162)rectangle++(0.009,0.009) (0.246,0.079)rectangle++(0.009,0.009) (0.596,0.029)rectangle++(0.009,0.009) (0.296,0.512)rectangle++(0.009,0.009)
  (0.146,0.254)rectangle++(0.009,0.009) (0.556,0.099)rectangle++(0.009,0.009) (0.276,0.547)rectangle++(0.009,0.009) (0.136,0.771)rectangle++(0.009,0.009) (0.066,0.883)rectangle++(0.009,0.009)
  (0.031,0.439)rectangle++(0.009,0.009) (0.013,0.217)rectangle++(0.009,0.009) (0.004,0.606)rectangle++(0.009,0.009) (0.499,0.24)rectangle++(0.009,0.009) (0.247,0.118)rectangle++(0.009,0.009)
  (0.121,0.057)rectangle++(0.009,0.009) (0.058,0.026)rectangle++(0.009,0.009) (0.027,0.011)rectangle++(0.009,0.009) (0.011,0.003)rectangle++(0.009,0.009) (0.003,-0.001)rectangle++(0.009,0.009)
  (-0.001,-0.003)rectangle++(0.009,0.009) (-0.003,0.496)rectangle++(0.009,0.009) (-0.004,0.746)rectangle++(0.009,0.009) (-0.004,0.371)rectangle++(0.009,0.009) (-0.004,0.183)rectangle++(0.009,0.009)
  (-0.004,0.589)rectangle++(0.009,0.009) (-0.004,0.292)rectangle++(0.009,0.009) (-0.004,0.144)rectangle++(0.009,0.009) (-0.004,0.07)rectangle++(0.009,0.009) (-0.004,0.533)rectangle++(0.009,0.009)
  (0.496,0.21)rectangle++(0.009,0.009) (0.246,0.103)rectangle++(0.009,0.009) (0.596,0.038)rectangle++(0.009,0.009) (0.296,0.017)rectangle++(0.009,0.009) (0.146,0.506)rectangle++(0.009,0.009)
  (0.071,0.251)rectangle++(0.009,0.009) (0.033,0.123)rectangle++(0.009,0.009) (0.014,0.059)rectangle++(0.009,0.009) (0.005,0.527)rectangle++(0.009,0.009) (0.499,0.208)rectangle++(0.009,0.009)
  (0.247,0.102)rectangle++(0.009,0.009) (0.596,0.038)rectangle++(0.009,0.009) (0.736,0.013)rectangle++(0.009,0.009) (0.366,0.004)rectangle++(0.009,0.009) (0.181,0)rectangle++(0.009,0.009)
  (0.57,-0.003)rectangle++(0.009,0.009) (0.725,-0.004)rectangle++(0.009,0.009) (0.787,-0.004)rectangle++(0.009,0.009) (0.391,0.496)rectangle++(0.009,0.009) (0.193,0.746)rectangle++(0.009,0.009)
  (0.094,0.371)rectangle++(0.009,0.009) (0.045,0.183)rectangle++(0.009,0.009) (0.515,0.071)rectangle++(0.009,0.009) (0.255,0.533)rectangle++(0.009,0.009) (0.599,0.211)rectangle++(0.009,0.009)
  (0.297,0.103)rectangle++(0.009,0.009) (0.616,0.039)rectangle++(0.009,0.009) (0.744,0.013)rectangle++(0.009,0.009) (0.37,0.004)rectangle++(0.009,0.009) (0.645,-0.001)rectangle++(0.009,0.009)
  (0.32,-0.003)rectangle++(0.009,0.009) (0.158,-0.004)rectangle++(0.009,0.009) (0.077,0.496)rectangle++(0.009,0.009) (0.528,0.196)rectangle++(0.009,0.009) (0.708,0.076)rectangle++(0.009,0.009)
  (0.352,0.536)rectangle++(0.009,0.009) (0.174,0.766)rectangle++(0.009,0.009) (0.085,0.381)rectangle++(0.009,0.009) (0.04,0.688)rectangle++(0.009,0.009) (0.513,0.273)rectangle++(0.009,0.009)
  (0.254,0.134)rectangle++(0.009,0.009) (0.599,0.051)rectangle++(0.009,0.009) (0.297,0.023)rectangle++(0.009,0.009) (0.616,0.007)rectangle++(0.009,0.009) (0.306,0.001)rectangle++(0.009,0.009)
  (0.151,-0.002)rectangle++(0.009,0.009) (0.073,0.497)rectangle++(0.009,0.009) (0.527,0.196)rectangle++(0.009,0.009) (0.261,0.596)rectangle++(0.009,0.009) (0.602,0.236)rectangle++(0.009,0.009)
  (0.738,0.092)rectangle++(0.009,0.009) (0.367,0.044)rectangle++(0.009,0.009) (0.181,0.52)rectangle++(0.009,0.009) (0.088,0.758)rectangle++(0.009,0.009) (0.042,0.377)rectangle++(0.009,0.009)
  (0.019,0.686)rectangle++(0.009,0.009) (0.007,0.841)rectangle++(0.009,0.009) (0.001,0.418)rectangle++(0.009,0.009) (-0.002,0.207)rectangle++(0.009,0.009) (-0.003,0.601)rectangle++(0.009,0.009)
  (0.496,0.238)rectangle++(0.009,0.009) (0.246,0.117)rectangle++(0.009,0.009) (0.121,0.556)rectangle++(0.009,0.009) (0.058,0.776)rectangle++(0.009,0.009) (0.027,0.386)rectangle++(0.009,0.009)
  (0.508,0.152)rectangle++(0.009,0.009) (0.701,0.058)rectangle++(0.009,0.009) (0.348,0.527)rectangle++(0.009,0.009) (0.172,0.761)rectangle++(0.009,0.009) (0.084,0.878)rectangle++(0.009,0.009)
  (0.04,0.437)rectangle++(0.009,0.009) (0.018,0.716)rectangle++(0.009,0.009) (0.504,0.284)rectangle++(0.009,0.009) (0.699,0.111)rectangle++(0.009,0.009) (0.777,0.042)rectangle++(0.009,0.009)
  (0.808,0.014)rectangle++(0.009,0.009) (0.821,0.003)rectangle++(0.009,0.009) (0.408,-0.001)rectangle++(0.009,0.009) (0.202,0.497)rectangle++(0.009,0.009) (0.578,0.196)rectangle++(0.009,0.009)
  (0.287,0.596)rectangle++(0.009,0.009) (0.612,0.236)rectangle++(0.009,0.009) (0.304,0.116)rectangle++(0.009,0.009) (0.15,0.556)rectangle++(0.009,0.009) (0.557,0.22)rectangle++(0.009,0.009)
  (0.72,0.085)rectangle++(0.009,0.009) (0.785,0.031)rectangle++(0.009,0.009) (0.39,0.013)rectangle++(0.009,0.009) (0.193,0.004)rectangle++(0.009,0.009) (0.094,0.5)rectangle++(0.009,0.009)
  (0.535,0.197)rectangle++(0.009,0.009) (0.265,0.596)rectangle++(0.009,0.009) (0.13,0.796)rectangle++(0.009,0.009) (0.063,0.896)rectangle++(0.009,0.009) (0.029,0.946)rectangle++(0.009,0.009)
  (0.012,0.471)rectangle++(0.009,0.009) (0.004,0.233)rectangle++(0.009,0.009) (0,0.614)rectangle++(0.009,0.009) (-0.002,0.805)rectangle++(0.009,0.009) (-0.003,0.4)rectangle++(0.009,0.009)
  (-0.004,0.698)rectangle++(0.009,0.009) (-0.004,0.847)rectangle++(0.009,0.009) (-0.004,0.921)rectangle++(0.009,0.009) (0.496,0.366)rectangle++(0.009,0.009) (0.246,0.181)rectangle++(0.009,0.009)
  (0.121,0.588)rectangle++(0.009,0.009) (0.058,0.292)rectangle++(0.009,0.009) (0.027,0.144)rectangle++(0.009,0.009) (0.011,0.57)rectangle++(0.009,0.009) (0.003,0.283)rectangle++(0.009,0.009)
  (-0.001,0.139)rectangle++(0.009,0.009) (-0.003,0.567)rectangle++(0.009,0.009) (-0.004,0.781)rectangle++(0.009,0.009) (-0.004,0.388)rectangle++(0.009,0.009) (-0.004,0.192)rectangle++(0.009,0.009)
  (-0.004,0.594)rectangle++(0.009,0.009) (-0.004,0.795)rectangle++(0.009,0.009) (-0.004,0.395)rectangle++(0.009,0.009) (0.496,0.155)rectangle++(0.009,0.009) (0.246,0.075)rectangle++(0.009,0.009)
  (0.121,0.535)rectangle++(0.009,0.009) (0.058,0.765)rectangle++(0.009,0.009) (0.027,0.88)rectangle++(0.009,0.009) (0.011,0.438)rectangle++(0.009,0.009) (0.003,0.217)rectangle++(0.009,0.009)
  (-0.001,0.606)rectangle++(0.009,0.009) (-0.003,0.801)rectangle++(0.009,0.009) (-0.004,0.898)rectangle++(0.009,0.009) (-0.004,0.947)rectangle++(0.009,0.009) (-0.004,0.471)rectangle++(0.009,0.009)
  (-0.004,0.733)rectangle++(0.009,0.009) (0.496,0.291)rectangle++(0.009,0.009) (0.246,0.643)rectangle++(0.009,0.009) (0.121,0.319)rectangle++(0.009,0.009) (0.058,0.157)rectangle++(0.009,0.009)
  (0.027,0.076)rectangle++(0.009,0.009) (0.011,0.036)rectangle++(0.009,0.009) (0.502,0.012)rectangle++(0.009,0.009) (0.249,0.504)rectangle++(0.009,0.009) (0.122,0.75)rectangle++(0.009,0.009)
  (0.546,0.297)rectangle++(0.009,0.009) (0.271,0.146)rectangle++(0.009,0.009) (0.133,0.071)rectangle++(0.009,0.009) (0.064,0.033)rectangle++(0.009,0.009) (0.03,0.514)rectangle++(0.009,0.009)
  (0.013,0.255)rectangle++(0.009,0.009) (0.004,0.625)rectangle++(0.009,0.009) (0,0.31)rectangle++(0.009,0.009) (-0.002,0.653)rectangle++(0.009,0.009) (0.496,0.258)rectangle++(0.009,0.009)
  (0.246,0.127)rectangle++(0.009,0.009) (0.121,0.561)rectangle++(0.009,0.009) (0.058,0.278)rectangle++(0.009,0.009) (0.027,0.637)rectangle++(0.009,0.009) (0.508,0.252)rectangle++(0.009,0.009)
  (0.252,0.124)rectangle++(0.009,0.009) (0.598,0.047)rectangle++(0.009,0.009) (0.297,0.521)rectangle++(0.009,0.009) (0.146,0.258)rectangle++(0.009,0.009) (0.071,0.127)rectangle++(0.009,0.009)
  (0.526,0.048)rectangle++(0.009,0.009) (0.708,0.017)rectangle++(0.009,0.009) (0.352,0.006)rectangle++(0.009,0.009) (0.174,0.501)rectangle++(0.009,0.009) (0.085,0.748)rectangle++(0.009,0.009)
  (0.04,0.872)rectangle++(0.009,0.009) (0.018,0.934)rectangle++(0.009,0.009) (0.007,0.965)rectangle++(0.009,0.009) (0.001,0.48)rectangle++(0.009,0.009) (-0.002,0.238)rectangle++(0.009,0.009)
  (-0.003,0.617)rectangle++(0.009,0.009) (-0.004,0.806)rectangle++(0.009,0.009) (-0.004,0.901)rectangle++(0.009,0.009) (-0.004,0.448)rectangle++(0.009,0.009) (0.496,0.177)rectangle++(0.009,0.009)
  (0.246,0.086)rectangle++(0.009,0.009) (0.121,0.541)rectangle++(0.009,0.009) (0.058,0.268)rectangle++(0.009,0.009) (0.027,0.632)rectangle++(0.009,0.009) (0.011,0.314)rectangle++(0.009,0.009)
  (0.003,0.655)rectangle++(0.009,0.009) (-0.001,0.825)rectangle++(0.009,0.009) (-0.003,0.91)rectangle++(0.009,0.009) (-0.004,0.953)rectangle++(0.009,0.009) (-0.004,0.974)rectangle++(0.009,0.009)
  (-0.004,0.485)rectangle++(0.009,0.009) (-0.004,0.74)rectangle++(0.009,0.009) (-0.004,0.368)rectangle++(0.009,0.009) (-0.004,0.182)rectangle++(0.009,0.009) (-0.004,0.089)rectangle++(0.009,0.009)
  (-0.004,0.542)rectangle++(0.009,0.009) (-0.004,0.269)rectangle++(0.009,0.009) (-0.004,0.632)rectangle++(0.009,0.009) (0.496,0.25)rectangle++(0.009,0.009) (0.246,0.623)rectangle++(0.009,0.009)
  (0.121,0.809)rectangle++(0.009,0.009) (0.058,0.402)rectangle++(0.009,0.009) (0.521,0.158)rectangle++(0.009,0.009) (0.258,0.577)rectangle++(0.009,0.009) (0.127,0.786)rectangle++(0.009,0.009)
  (0.061,0.391)rectangle++(0.009,0.009) (0.028,0.693)rectangle++(0.009,0.009) (0.012,0.344)rectangle++(0.009,0.009) (0.004,0.17)rectangle++(0.009,0.009) (0,0.583)rectangle++(0.009,0.009)
  (-0.002,0.789)rectangle++(0.009,0.009) (-0.003,0.892)rectangle++(0.009,0.009) (-0.004,0.444)rectangle++(0.009,0.009) (-0.004,0.72)rectangle++(0.009,0.009) (-0.004,0.858)rectangle++(0.009,0.009)
  (-0.004,0.927)rectangle++(0.009,0.009) (-0.004,0.461)rectangle++(0.009,0.009) (-0.004,0.228)rectangle++(0.009,0.009) (-0.004,0.112)rectangle++(0.009,0.009) (-0.004,0.054)rectangle++(0.009,0.009)
  (-0.004,0.525)rectangle++(0.009,0.009) (-0.004,0.26)rectangle++(0.009,0.009) (-0.004,0.128)rectangle++(0.009,0.009) (-0.004,0.062)rectangle++(0.009,0.009) (-0.004,0.029)rectangle++(0.009,0.009)
  (-0.004,0.512)rectangle++(0.009,0.009) (-0.004,0.254)rectangle++(0.009,0.009) (-0.004,0.125)rectangle++(0.009,0.009) (0.496,0.047)rectangle++(0.009,0.009) (0.246,0.521)rectangle++(0.009,0.009)
  (0.596,0.206)rectangle++(0.009,0.009) (0.296,0.601)rectangle++(0.009,0.009) (0.146,0.798)rectangle++(0.009,0.009) (0.071,0.397)rectangle++(0.009,0.009) (0.526,0.156)rectangle++(0.009,0.009)
  (0.261,0.576)rectangle++(0.009,0.009) (0.128,0.786)rectangle++(0.009,0.009) (0.062,0.391)rectangle++(0.009,0.009) (0.029,0.193)rectangle++(0.009,0.009) (0.012,0.094)rectangle++(0.009,0.009)
  (0.004,0.545)rectangle++(0.009,0.009) (0,0.27)rectangle++(0.009,0.009) (-0.002,0.633)rectangle++(0.009,0.009) (-0.003,0.314)rectangle++(0.009,0.009) (0.496,0.123)rectangle++(0.009,0.009)
  (0.246,0.059)rectangle++(0.009,0.009) (0.121,0.027)rectangle++(0.009,0.009) (0.058,0.511)rectangle++(0.009,0.009) (0.521,0.202)rectangle++(0.009,0.009) (0.706,0.078)rectangle++(0.009,0.009)
  (0.78,0.029)rectangle++(0.009,0.009) (0.388,0.512)rectangle++(0.009,0.009) (0.192,0.254)rectangle++(0.009,0.009) (0.574,0.099)rectangle++(0.009,0.009) (0.285,0.047)rectangle++(0.009,0.009)
  (0.14,0.021)rectangle++(0.009,0.009) (0.553,0.006)rectangle++(0.009,0.009) (0.274,0.501)rectangle++(0.009,0.009) (0.135,0.248)rectangle++(0.009,0.009) (0.065,0.122)rectangle++(0.009,0.009)
  (0.523,0.046)rectangle++(0.009,0.009) (0.259,0.521)rectangle++(0.009,0.009) (0.601,0.206)rectangle++(0.009,0.009) (0.298,0.601)rectangle++(0.009,0.009) (0.147,0.798)rectangle++(0.009,0.009)
  (0.071,0.897)rectangle++(0.009,0.009) (0.033,0.946)rectangle++(0.009,0.009) (0.014,0.971)rectangle++(0.009,0.009) (0.005,0.483)rectangle++(0.009,0.009) (0,0.239)rectangle++(0.009,0.009)
  (0.497,0.093)rectangle++(0.009,0.009) (0.246,0.044)rectangle++(0.009,0.009) (0.596,0.015)rectangle++(0.009,0.009) (0.296,0.005)rectangle++(0.009,0.009) (0.146,0.5)rectangle++(0.009,0.009)
  (0.071,0.248)rectangle++(0.009,0.009) (0.033,0.122)rectangle++(0.009,0.009) (0.511,0.046)rectangle++(0.009,0.009) (0.253,0.021)rectangle++(0.009,0.009) (0.124,0.508)rectangle++(0.009,0.009)
  (0.06,0.752)rectangle++(0.009,0.009) (0.028,0.874)rectangle++(0.009,0.009) (0.012,0.935)rectangle++(0.009,0.009) (0.004,0.965)rectangle++(0.009,0.009) (0,0.48)rectangle++(0.009,0.009)
  (-0.002,0.238)rectangle++(0.009,0.009) (-0.003,0.117)rectangle++(0.009,0.009) (-0.004,0.556)rectangle++(0.009,0.009) (-0.004,0.276)rectangle++(0.009,0.009) (0.496,0.108)rectangle++(0.009,0.009)
  (0.246,0.052)rectangle++(0.009,0.009) (0.121,0.524)rectangle++(0.009,0.009) (0.058,0.26)rectangle++(0.009,0.009) (0.027,0.128)rectangle++(0.009,0.009) (0.011,0.562)rectangle++(0.009,0.009)
  (0.003,0.779)rectangle++(0.009,0.009) (0.499,0.309)rectangle++(0.009,0.009) (0.697,0.121)rectangle++(0.009,0.009) (0.346,0.558)rectangle++(0.009,0.009) (0.636,0.221)rectangle++(0.009,0.009)
  (0.316,0.608)rectangle++(0.009,0.009) (0.156,0.302)rectangle++(0.009,0.009) (0.076,0.149)rectangle++(0.009,0.009) (0.036,0.072)rectangle++(0.009,0.009) (0.016,0.534)rectangle++(0.009,0.009)
  (0.006,0.765)rectangle++(0.009,0.009) (0.001,0.38)rectangle++(0.009,0.009) (0.498,0.149)rectangle++(0.009,0.009) (0.247,0.572)rectangle++(0.009,0.009) (0.121,0.784)rectangle++(0.009,0.009)
  (0.058,0.39)rectangle++(0.009,0.009) (0.521,0.153)rectangle++(0.009,0.009) (0.258,0.574)rectangle++(0.009,0.009) (0.601,0.227)rectangle++(0.009,0.009) (0.738,0.088)rectangle++(0.009,0.009)
  (0.367,0.042)rectangle++(0.009,0.009) (0.644,0.014)rectangle++(0.009,0.009) (0.32,0.505)rectangle++(0.009,0.009) (0.158,0.75)rectangle++(0.009,0.009) (0.077,0.373)rectangle++(0.009,0.009)
  (0.528,0.146)rectangle++(0.009,0.009) (0.708,0.056)rectangle++(0.009,0.009) (0.352,0.026)rectangle++(0.009,0.009) (0.174,0.011)rectangle++(0.009,0.009) (0.085,0.003)rectangle++(0.009,0.009)
  (0.04,-0.001)rectangle++(0.009,0.009) (0.018,0.497)rectangle++(0.009,0.009) (0.007,0.246)rectangle++(0.009,0.009) (0.001,0.121)rectangle++(0.009,0.009) (-0.002,0.558)rectangle++(0.009,0.009)
  (-0.003,0.777)rectangle++(0.009,0.009) (0.496,0.308)rectangle++(0.009,0.009) (0.246,0.152)rectangle++(0.009,0.009) (0.121,0.074)rectangle++(0.009,0.009) (0.058,0.035)rectangle++(0.009,0.009)
  (0.027,0.515)rectangle++(0.009,0.009) (0.011,0.255)rectangle++(0.009,0.009) (0.003,0.125)rectangle++(0.009,0.009) (0.499,0.047)rectangle++(0.009,0.009) (0.247,0.521)rectangle++(0.009,0.009)
  (0.121,0.758)rectangle++(0.009,0.009) (0.058,0.377)rectangle++(0.009,0.009) (0.027,0.186)rectangle++(0.009,0.009) (0.011,0.591)rectangle++(0.009,0.009) (0.502,0.234)rectangle++(0.009,0.009)
  (0.249,0.115)rectangle++(0.009,0.009) (0.122,0.055)rectangle++(0.009,0.009) (0.059,0.025)rectangle++(0.009,0.009) (0.027,0.51)rectangle++(0.009,0.009) (0.508,0.201)rectangle++(0.009,0.009)
  (0.252,0.098)rectangle++(0.009,0.009) (0.598,0.037)rectangle++(0.009,0.009) (0.297,0.016)rectangle++(0.009,0.009) (0.146,0.506)rectangle++(0.009,0.009) (0.071,0.751)rectangle++(0.009,0.009)
  (0.033,0.873)rectangle++(0.009,0.009) (0.014,0.934)rectangle++(0.009,0.009) (0.005,0.465)rectangle++(0.009,0.009) (0,0.23)rectangle++(0.009,0.009) (-0.002,0.113)rectangle++(0.009,0.009)
  (-0.003,0.054)rectangle++(0.009,0.009) (-0.004,0.025)rectangle++(0.009,0.009) (-0.004,0.51)rectangle++(0.009,0.009) (-0.004,0.753)rectangle++(0.009,0.009) (0.496,0.298)rectangle++(0.009,0.009)
  (0.246,0.647)rectangle++(0.009,0.009) (0.121,0.321)rectangle++(0.009,0.009) (0.058,0.658)rectangle++(0.009,0.009) (0.027,0.327)rectangle++(0.009,0.009) (0.508,0.128)rectangle++(0.009,0.009)
  (0.252,0.062)rectangle++(0.009,0.009) (0.598,0.022)rectangle++(0.009,0.009) (0.297,0.509)rectangle++(0.009,0.009) (0.146,0.752)rectangle++(0.009,0.009) (0.071,0.874)rectangle++(0.009,0.009)
  (0.033,0.435)rectangle++(0.009,0.009) (0.014,0.715)rectangle++(0.009,0.009) (0.005,0.855)rectangle++(0.009,0.009) (0.499,0.339)rectangle++(0.009,0.009) (0.247,0.167)rectangle++(0.009,0.009)
  (0.121,0.581)rectangle++(0.009,0.009) (0.058,0.288)rectangle++(0.009,0.009) (0.027,0.642)rectangle++(0.009,0.009) (0.508,0.254)rectangle++(0.009,0.009) (0.701,0.099)rectangle++(0.009,0.009)
  (0.778,0.037)rectangle++(0.009,0.009) (0.387,0.516)rectangle++(0.009,0.009) (0.191,0.756)rectangle++(0.009,0.009) (0.093,0.376)rectangle++(0.009,0.009) (0.535,0.148)rectangle++(0.009,0.009)
  (0.711,0.056)rectangle++(0.009,0.009) (0.353,0.526)rectangle++(0.009,0.009) (0.174,0.261)rectangle++(0.009,0.009) (0.567,0.102)rectangle++(0.009,0.009) (0.281,0.549)rectangle++(0.009,0.009)
  (0.61,0.217)rectangle++(0.009,0.009) (0.741,0.084)rectangle++(0.009,0.009) (0.794,0.031)rectangle++(0.009,0.009) (0.815,0.01)rectangle++(0.009,0.009) (0.405,0.503)rectangle++(0.009,0.009)
  (0.659,0.198)rectangle++(0.009,0.009) (0.327,0.097)rectangle++(0.009,0.009) (0.161,0.546)rectangle++(0.009,0.009) (0.078,0.771)rectangle++(0.009,0.009) (0.529,0.306)rectangle++(0.009,0.009)
  (0.709,0.12)rectangle++(0.009,0.009) (0.352,0.058)rectangle++(0.009,0.009) (0.638,0.02)rectangle++(0.009,0.009) (0.753,0.005)rectangle++(0.009,0.009) (0.374,0)rectangle++(0.009,0.009)
  (0.185,0.498)rectangle++(0.009,0.009) (0.571,0.196)rectangle++(0.009,0.009) (0.283,0.096)rectangle++(0.009,0.009) (0.611,0.036)rectangle++(0.009,0.009) (0.742,0.012)rectangle++(0.009,0.009)
  (0.369,0.504)rectangle++(0.009,0.009) (0.182,0.25)rectangle++(0.009,0.009) (0.089,0.623)rectangle++(0.009,0.009) (0.533,0.246)rectangle++(0.009,0.009) (0.264,0.621)rectangle++(0.009,0.009)
  (0.603,0.246)rectangle++(0.009,0.009) (0.738,0.096)rectangle++(0.009,0.009) (0.367,0.546)rectangle++(0.009,0.009) (0.181,0.271)rectangle++(0.009,0.009) (0.088,0.633)rectangle++(0.009,0.009)
  (0.533,0.251)rectangle++(0.009,0.009) (0.264,0.623)rectangle++(0.009,0.009) (0.13,0.309)rectangle++(0.009,0.009) (0.063,0.652)rectangle++(0.009,0.009) (0.029,0.824)rectangle++(0.009,0.009)
  (0.012,0.41)rectangle++(0.009,0.009) (0.004,0.203)rectangle++(0.009,0.009) (0,0.099)rectangle++(0.009,0.009) (0.497,0.037)rectangle++(0.009,0.009) (0.246,0.016)rectangle++(0.009,0.009)
  (0.121,0.006)rectangle++(0.009,0.009) (0.058,0.001)rectangle++(0.009,0.009) (0.027,-0.002)rectangle++(0.009,0.009) (0.011,0.497)rectangle++(0.009,0.009) (0.003,0.246)rectangle++(0.009,0.009)
  (-0.001,0.121)rectangle++(0.009,0.009) (-0.003,0.058)rectangle++(0.009,0.009) (-0.004,0.027)rectangle++(0.009,0.009) (-0.004,0.511)rectangle++(0.009,0.009) (-0.004,0.253)rectangle++(0.009,0.009)
  (-0.004,0.124)rectangle++(0.009,0.009) (-0.004,0.06)rectangle++(0.009,0.009) (-0.004,0.028)rectangle++(0.009,0.009) (-0.004,0.012)rectangle++(0.009,0.009) (-0.004,0.004)rectangle++(0.009,0.009)
  (-0.004,0)rectangle++(0.009,0.009) (-0.004,-0.002)rectangle++(0.009,0.009) (-0.004,0.497)rectangle++(0.009,0.009) (0.496,0.196)rectangle++(0.009,0.009) (0.246,0.596)rectangle++(0.009,0.009)
  (0.121,0.796)rectangle++(0.009,0.009) (0.058,0.896)rectangle++(0.009,0.009) (0.027,0.946)rectangle++(0.009,0.009) (0.011,0.471)rectangle++(0.009,0.009) (0.003,0.733)rectangle++(0.009,0.009)
  (-0.001,0.364)rectangle++(0.009,0.009) (0.497,0.143)rectangle++(0.009,0.009) (0.246,0.569)rectangle++(0.009,0.009) (0.121,0.782)rectangle++(0.009,0.009) (0.058,0.889)rectangle++(0.009,0.009)
  (0.027,0.442)rectangle++(0.009,0.009) (0.011,0.219)rectangle++(0.009,0.009) (0.003,0.607)rectangle++(0.009,0.009) (0.499,0.24)rectangle++(0.009,0.009) (0.247,0.118)rectangle++(0.009,0.009)
  (0.596,0.044)rectangle++(0.009,0.009) (0.296,0.02)rectangle++(0.009,0.009) (0.146,0.008)rectangle++(0.009,0.009) (0.071,0.002)rectangle++(0.009,0.009) (0.033,-0.001)rectangle++(0.009,0.009)
  (0.014,0.497)rectangle++(0.009,0.009) (0.005,0.746)rectangle++(0.009,0.009) (0.499,0.296)rectangle++(0.009,0.009) (0.247,0.146)rectangle++(0.009,0.009) (0.121,0.571)rectangle++(0.009,0.009)
  (0.058,0.283)rectangle++(0.009,0.009) (0.027,0.639)rectangle++(0.009,0.009) (0.508,0.253)rectangle++(0.009,0.009) (0.252,0.624)rectangle++(0.009,0.009) (0.598,0.247)rectangle++(0.009,0.009)
  (0.297,0.621)rectangle++(0.009,0.009) (0.146,0.308)rectangle++(0.009,0.009) (0.071,0.152)rectangle++(0.009,0.009) (0.033,0.574)rectangle++(0.009,0.009) (0.014,0.285)rectangle++(0.009,0.009)
  (0.005,0.64)rectangle++(0.009,0.009) (0.499,0.253)rectangle++(0.009,0.009) (0.247,0.624)rectangle++(0.009,0.009) (0.121,0.31)rectangle++(0.009,0.009) (0.546,0.121)rectangle++(0.009,0.009)
  (0.271,0.058)rectangle++(0.009,0.009) (0.606,0.021)rectangle++(0.009,0.009) (0.74,0.006)rectangle++(0.009,0.009) (0.793,0)rectangle++(0.009,0.009) (0.394,-0.002)rectangle++(0.009,0.009)
  (0.195,-0.003)rectangle++(0.009,0.009) (0.095,0.496)rectangle++(0.009,0.009) (0.045,0.746)rectangle++(0.009,0.009) (0.515,0.296)rectangle++(0.009,0.009) (0.255,0.646)rectangle++(0.009,0.009)
  (0.599,0.256)rectangle++(0.009,0.009) (0.297,0.126)rectangle++(0.009,0.009) (0.146,0.061)rectangle++(0.009,0.009) (0.556,0.022)rectangle++(0.009,0.009) (0.276,0.009)rectangle++(0.009,0.009)
  (0.608,0.001)rectangle++(0.009,0.009) (0.302,0.498)rectangle++(0.009,0.009) (0.149,0.747)rectangle++(0.009,0.009) (0.072,0.871)rectangle++(0.009,0.009) (0.034,0.933)rectangle++(0.009,0.009)
  (0.511,0.371)rectangle++(0.009,0.009) (0.702,0.146)rectangle++(0.009,0.009) (0.349,0.571)rectangle++(0.009,0.009) (0.637,0.226)rectangle++(0.009,0.009) (0.316,0.611)rectangle++(0.009,0.009)
  (0.624,0.242)rectangle++(0.009,0.009) (0.747,0.094)rectangle++(0.009,0.009) (0.371,0.545)rectangle++(0.009,0.009) (0.646,0.215)rectangle++(0.009,0.009) (0.756,0.083)rectangle++(0.009,0.009)
  (0.376,0.039)rectangle++(0.009,0.009) (0.186,0.017)rectangle++(0.009,0.009) (0.091,0.006)rectangle++(0.009,0.009) (0.534,0)rectangle++(0.009,0.009) (0.265,0.498)rectangle++(0.009,0.009)
  (0.13,0.247)rectangle++(0.009,0.009) (0.549,0.096)rectangle++(0.009,0.009) (0.272,0.046)rectangle++(0.009,0.009) (0.134,0.521)rectangle++(0.009,0.009) (0.065,0.258)rectangle++(0.009,0.009)
  (0.03,0.127)rectangle++(0.009,0.009) (0.509,0.048)rectangle++(0.009,0.009) (0.252,0.522)rectangle++(0.009,0.009) (0.124,0.259)rectangle++(0.009,0.009) (0.547,0.101)rectangle++(0.009,0.009)
  (0.271,0.048)rectangle++(0.009,0.009) (0.133,0.022)rectangle++(0.009,0.009) (0.064,0.509)rectangle++(0.009,0.009) (0.03,0.752)rectangle++(0.009,0.009) (0.013,0.374)rectangle++(0.009,0.009)
  (0.502,0.147)rectangle++(0.009,0.009) (0.249,0.571)rectangle++(0.009,0.009) (0.122,0.783)rectangle++(0.009,0.009) (0.546,0.311)rectangle++(0.009,0.009) (0.716,0.122)rectangle++(0.009,0.009)
  (0.356,0.059)rectangle++(0.009,0.009) (0.176,0.527)rectangle++(0.009,0.009) (0.086,0.761)rectangle++(0.009,0.009) (0.041,0.878)rectangle++(0.009,0.009) (0.018,0.437)rectangle++(0.009,0.009)
  (0.007,0.216)rectangle++(0.009,0.009) (0.001,0.606)rectangle++(0.009,0.009) (-0.002,0.801)rectangle++(0.009,0.009) (-0.003,0.398)rectangle++(0.009,0.009) (0.496,0.157)rectangle++(0.009,0.009)
  (0.246,0.576)rectangle++(0.009,0.009) (0.121,0.786)rectangle++(0.009,0.009) (0.546,0.312)rectangle++(0.009,0.009) (0.271,0.654)rectangle++(0.009,0.009) (0.133,0.825)rectangle++(0.009,0.009)
  (0.064,0.41)rectangle++(0.009,0.009) (0.03,0.203)rectangle++(0.009,0.009) (0.013,0.099)rectangle++(0.009,0.009) (0.004,0.047)rectangle++(0.009,0.009) (0,0.521)rectangle++(0.009,0.009)
  (-0.002,0.758)rectangle++(0.009,0.009) (-0.003,0.877)rectangle++(0.009,0.009) (-0.004,0.436)rectangle++(0.009,0.009) (-0.004,0.216)rectangle++(0.009,0.009) (-0.004,0.106)rectangle++(0.009,0.009)
  (-0.004,0.051)rectangle++(0.009,0.009) (0.496,0.018)rectangle++(0.009,0.009) (0.246,0.507)rectangle++(0.009,0.009) (0.121,0.751)rectangle++(0.009,0.009) (0.546,0.298)rectangle++(0.009,0.009)
  (0.271,0.647)rectangle++(0.009,0.009) (0.133,0.821)rectangle++(0.009,0.009) (0.551,0.326)rectangle++(0.009,0.009) (0.273,0.161)rectangle++(0.009,0.009) (0.134,0.078)rectangle++(0.009,0.009)
  (0.065,0.537)rectangle++(0.009,0.009) (0.03,0.266)rectangle++(0.009,0.009) (0.013,0.631)rectangle++(0.009,0.009) (0.004,0.813)rectangle++(0.009,0.009) (0,0.904)rectangle++(0.009,0.009)
  (-0.002,0.95)rectangle++(0.009,0.009) (-0.003,0.473)rectangle++(0.009,0.009) (0.496,0.186)rectangle++(0.009,0.009) (0.696,0.072)rectangle++(0.009,0.009) (0.346,0.034)rectangle++(0.009,0.009)
  (0.171,0.515)rectangle++(0.009,0.009) (0.083,0.755)rectangle++(0.009,0.009) (0.531,0.299)rectangle++(0.009,0.009) (0.71,0.117)rectangle++(0.009,0.009) (0.781,0.044)rectangle++(0.009,0.009)
  (0.388,0.02)rectangle++(0.009,0.009) (0.653,0.005)rectangle++(0.009,0.009) (0.324,0.5)rectangle++(0.009,0.009) (0.16,0.748)rectangle++(0.009,0.009) (0.078,0.872)rectangle++(0.009,0.009)
  (0.037,0.434)rectangle++(0.009,0.009) (0.016,0.715)rectangle++(0.009,0.009) (0.504,0.283)rectangle++(0.009,0.009) (0.25,0.139)rectangle++(0.009,0.009) (0.123,0.067)rectangle++(0.009,0.009)
  (0.059,0.031)rectangle++(0.009,0.009) (0.027,0.513)rectangle++(0.009,0.009) (0.011,0.254)rectangle++(0.009,0.009) (0.502,0.099)rectangle++(0.009,0.009) (0.698,0.037)rectangle++(0.009,0.009)
  (0.347,0.516)rectangle++(0.009,0.009) (0.171,0.756)rectangle++(0.009,0.009) (0.083,0.376)rectangle++(0.009,0.009) (0.039,0.686)rectangle++(0.009,0.009) (0.017,0.341)rectangle++(0.009,0.009)
  (0.006,0.668)rectangle++(0.009,0.009) (0.001,0.832)rectangle++(0.009,0.009) (-0.002,0.914)rectangle++(0.009,0.009) (-0.003,0.955)rectangle++(0.009,0.009) (-0.004,0.475)rectangle++(0.009,0.009)
  (-0.004,0.235)rectangle++(0.009,0.009) (0.496,0.091)rectangle++(0.009,0.009) (0.246,0.043)rectangle++(0.009,0.009) (0.596,0.015)rectangle++(0.009,0.009) (0.296,0.505)rectangle++(0.009,0.009)
  (0.146,0.25)rectangle++(0.009,0.009) (0.556,0.097)rectangle++(0.009,0.009) (0.276,0.046)rectangle++(0.009,0.009) (0.608,0.016)rectangle++(0.009,0.009) (0.302,0.006)rectangle++(0.009,0.009)
  (0.149,0.001)rectangle++(0.009,0.009) (0.072,0.498)rectangle++(0.009,0.009) (0.526,0.197)rectangle++(0.009,0.009) (0.261,0.096)rectangle++(0.009,0.009) (0.128,0.046)rectangle++(0.009,0.009)
  (0.549,0.016)rectangle++(0.009,0.009) (0.272,0.006)rectangle++(0.009,0.009) (0.606,0)rectangle++(0.009,0.009) (0.301,-0.002)rectangle++(0.009,0.009) (0.148,0.497)rectangle++(0.009,0.009)
  (0.072,0.746)rectangle++(0.009,0.009) (0.526,0.296)rectangle++(0.009,0.009) (0.708,0.116)rectangle++(0.009,0.009) (0.352,0.056)rectangle++(0.009,0.009) (0.174,0.526)rectangle++(0.009,0.009)
  (0.567,0.208)rectangle++(0.009,0.009) (0.281,0.602)rectangle++(0.009,0.009) (0.138,0.799)rectangle++(0.009,0.009) (0.067,0.897)rectangle++(0.009,0.009) (0.524,0.356)rectangle++(0.009,0.009)
  (0.26,0.676)rectangle++(0.009,0.009) (0.128,0.336)rectangle++(0.009,0.009) (0.062,0.166)rectangle++(0.009,0.009) (0.522,0.064)rectangle++(0.009,0.009) (0.259,0.03)rectangle++(0.009,0.009)
  (0.127,0.513)rectangle++(0.009,0.009) (0.548,0.202)rectangle++(0.009,0.009) (0.272,0.599)rectangle++(0.009,0.009) (0.134,0.297)rectangle++(0.009,0.009) (0.065,0.646)rectangle++(0.009,0.009)
  (0.03,0.321)rectangle++(0.009,0.009) (0.013,0.158)rectangle++(0.009,0.009) (0.004,0.077)rectangle++(0.009,0.009) (0,0.036)rectangle++(0.009,0.009) (0.497,0.012)rectangle++(0.009,0.009)
  (0.246,0.004)rectangle++(0.009,0.009) (0.121,0.5)rectangle++(0.009,0.009) (0.058,0.248)rectangle++(0.009,0.009) (0.027,0.122)rectangle++(0.009,0.009) (0.011,0.559)rectangle++(0.009,0.009)
  (0.003,0.277)rectangle++(0.009,0.009) (-0.001,0.136)rectangle++(0.009,0.009) (-0.003,0.566)rectangle++(0.009,0.009) (0.496,0.224)rectangle++(0.009,0.009) (0.246,0.61)rectangle++(0.009,0.009)
  (0.121,0.803)rectangle++(0.009,0.009) (0.546,0.318)rectangle++(0.009,0.009) (0.271,0.657)rectangle++(0.009,0.009) (0.133,0.826)rectangle++(0.009,0.009) (0.064,0.911)rectangle++(0.009,0.009)
  (0.03,0.453)rectangle++(0.009,0.009) (0.013,0.224)rectangle++(0.009,0.009) (0.004,0.11)rectangle++(0.009,0.009) (0,0.053)rectangle++(0.009,0.009) (-0.002,0.524)rectangle++(0.009,0.009)
  (-0.003,0.76)rectangle++(0.009,0.009) (-0.004,0.878)rectangle++(0.009,0.009) (-0.004,0.437)rectangle++(0.009,0.009) (0.496,0.172)rectangle++(0.009,0.009) (0.696,0.066)rectangle++(0.009,0.009)
  (0.776,0.024)rectangle++(0.009,0.009) (0.808,0.007)rectangle++(0.009,0.009) (0.82,0)rectangle++(0.009,0.009) (0.408,-0.002)rectangle++(0.009,0.009) (0.202,0.497)rectangle++(0.009,0.009)
  (0.099,0.746)rectangle++(0.009,0.009) (0.047,0.871)rectangle++(0.009,0.009) (0.021,0.933)rectangle++(0.009,0.009) (0.506,0.371)rectangle++(0.009,0.009) (0.251,0.183)rectangle++(0.009,0.009)
  (0.123,0.089)rectangle++(0.009,0.009) (0.059,0.042)rectangle++(0.009,0.009) (0.027,0.519)rectangle++(0.009,0.009) (0.011,0.757)rectangle++(0.009,0.009) (0.502,0.3)rectangle++(0.009,0.009)
  (0.249,0.648)rectangle++(0.009,0.009) (0.122,0.822)rectangle++(0.009,0.009) (0.059,0.409)rectangle++(0.009,0.009) (0.027,0.202)rectangle++(0.009,0.009) (0.011,0.099)rectangle++(0.009,0.009)
  (0.003,0.547)rectangle++(0.009,0.009) (-0.001,0.771)rectangle++(0.009,0.009) (-0.003,0.383)rectangle++(0.009,0.009) (-0.004,0.689)rectangle++(0.009,0.009) (-0.004,0.342)rectangle++(0.009,0.009)
  (-0.004,0.169)rectangle++(0.009,0.009) (-0.004,0.582)rectangle++(0.009,0.009) (-0.004,0.789)rectangle++(0.009,0.009) (0.496,0.313)rectangle++(0.009,0.009) (0.696,0.122)rectangle++(0.009,0.009)
  (0.346,0.559)rectangle++(0.009,0.009) (0.171,0.277)rectangle++(0.009,0.009) (0.083,0.136)rectangle++(0.009,0.009) (0.531,0.052)rectangle++(0.009,0.009) (0.263,0.024)rectangle++(0.009,0.009)
  (0.129,0.51)rectangle++(0.009,0.009) (0.062,0.253)rectangle++(0.009,0.009) (0.522,0.098)rectangle++(0.009,0.009) (0.706,0.037)rectangle++(0.009,0.009) (0.78,0.012)rectangle++(0.009,0.009)
  (0.809,0.002)rectangle++(0.009,0.009) (0.402,-0.001)rectangle++(0.009,0.009) (0.199,0.497)rectangle++(0.009,0.009) (0.097,0.746)rectangle++(0.009,0.009) (0.536,0.296)rectangle++(0.009,0.009)
  (0.266,0.646)rectangle++(0.009,0.009) (0.131,0.321)rectangle++(0.009,0.009) (0.55,0.126)rectangle++(0.009,0.009) (0.273,0.561)rectangle++(0.009,0.009) (0.134,0.278)rectangle++(0.009,0.009)
  (0.065,0.637)rectangle++(0.009,0.009) (0.03,0.816)rectangle++(0.009,0.009) (0.013,0.906)rectangle++(0.009,0.009) (0.004,0.451)rectangle++(0.009,0.009) (0,0.723)rectangle++(0.009,0.009)
  (0.497,0.287)rectangle++(0.009,0.009) (0.246,0.141)rectangle++(0.009,0.009) (0.596,0.054)rectangle++(0.009,0.009) (0.736,0.019)rectangle++(0.009,0.009) (0.366,0.507)rectangle++(0.009,0.009)
  (0.644,0.2)rectangle++(0.009,0.009) (0.755,0.077)rectangle++(0.009,0.009) (0.375,0.036)rectangle++(0.009,0.009) (0.185,0.516)rectangle++(0.009,0.009) (0.09,0.256)rectangle++(0.009,0.009)
  (0.043,0.626)rectangle++(0.009,0.009) (0.019,0.311)rectangle++(0.009,0.009) (0.007,0.653)rectangle++(0.009,0.009) (0.001,0.324)rectangle++(0.009,0.009) (-0.002,0.66)rectangle++(0.009,0.009)
  (0.497,0.261)rectangle++(0.009,0.009) (0.696,0.102)rectangle++(0.009,0.009) (0.776,0.038)rectangle++(0.009,0.009) (0.386,0.017)rectangle++(0.009,0.009) (0.191,0.506)rectangle++(0.009,0.009)
  (0.093,0.251)rectangle++(0.009,0.009) (0.535,0.098)rectangle++(0.009,0.009) (0.711,0.036)rectangle++(0.009,0.009) (0.353,0.516)rectangle++(0.009,0.009) (0.174,0.756)rectangle++(0.009,0.009)
  (0.085,0.376)rectangle++(0.009,0.009) (0.04,0.686)rectangle++(0.009,0.009) (0.018,0.341)rectangle++(0.009,0.009) (0.504,0.134)rectangle++(0.009,0.009) (0.25,0.565)rectangle++(0.009,0.009)
  (0.597,0.223)rectangle++(0.009,0.009) (0.736,0.087)rectangle++(0.009,0.009) (0.366,0.541)rectangle++(0.009,0.009) (0.181,0.768)rectangle++(0.009,0.009) (0.088,0.382)rectangle++(0.009,0.009)
  (0.533,0.15)rectangle++(0.009,0.009) (0.71,0.057)rectangle++(0.009,0.009) (0.781,0.02)rectangle++(0.009,0.009) (0.388,0.508)rectangle++(0.009,0.009) (0.192,0.752)rectangle++(0.009,0.009)
  (0.094,0.874)rectangle++(0.009,0.009) (0.045,0.435)rectangle++(0.009,0.009) (0.02,0.715)rectangle++(0.009,0.009) (0.008,0.855)rectangle++(0.009,0.009) (0.002,0.925)rectangle++(0.009,0.009)
  (-0.001,0.96)rectangle++(0.009,0.009) (0.497,0.381)rectangle++(0.009,0.009) (0.696,0.15)rectangle++(0.009,0.009) (0.776,0.057)rectangle++(0.009,0.009) (0.386,0.026)rectangle++(0.009,0.009)
  (0.191,0.511)rectangle++(0.009,0.009) (0.093,0.253)rectangle++(0.009,0.009) (0.044,0.624)rectangle++(0.009,0.009) (0.515,0.247)rectangle++(0.009,0.009) (0.255,0.621)rectangle++(0.009,0.009)
  (0.599,0.246)rectangle++(0.009,0.009) (0.297,0.121)rectangle++(0.009,0.009) (0.146,0.558)rectangle++(0.009,0.009) (0.071,0.277)rectangle++(0.009,0.009) (0.526,0.108)rectangle++(0.009,0.009)
  (0.261,0.552)rectangle++(0.009,0.009) (0.128,0.774)rectangle++(0.009,0.009) (0.549,0.307)rectangle++(0.009,0.009) (0.272,0.151)rectangle++(0.009,0.009) (0.134,0.573)rectangle++(0.009,0.009)
  (0.065,0.784)rectangle++(0.009,0.009) (0.03,0.89)rectangle++(0.009,0.009) (0.509,0.353)rectangle++(0.009,0.009) (0.252,0.174)rectangle++(0.009,0.009) (0.598,0.067)rectangle++(0.009,0.009)
  (0.297,0.531)rectangle++(0.009,0.009) (0.616,0.21)rectangle++(0.009,0.009) (0.306,0.603)rectangle++(0.009,0.009) (0.62,0.238)rectangle++(0.009,0.009) (0.745,0.093)rectangle++(0.009,0.009)
  (0.37,0.044)rectangle++(0.009,0.009) (0.183,0.52)rectangle++(0.009,0.009) (0.57,0.205)rectangle++(0.009,0.009) (0.283,0.1)rectangle++(0.009,0.009) (0.139,0.048)rectangle++(0.009,0.009)
  (0.553,0.016)rectangle++(0.009,0.009) (0.718,0.004)rectangle++(0.009,0.009) (0.357,0.5)rectangle++(0.009,0.009) (0.64,0.197)rectangle++(0.009,0.009) (0.753,0.076)rectangle++(0.009,0.009)
  (0.374,0.536)rectangle++(0.009,0.009) (0.185,0.266)rectangle++(0.009,0.009) (0.09,0.631)rectangle++(0.009,0.009) (0.043,0.313)rectangle++(0.009,0.009) (0.514,0.123)rectangle++(0.009,0.009)
  (0.703,0.046)rectangle++(0.009,0.009) (0.349,0.021)rectangle++(0.009,0.009) (0.172,0.008)rectangle++(0.009,0.009) (0.084,0.502)rectangle++(0.009,0.009) (0.04,0.249)rectangle++(0.009,0.009)
  (0.018,0.122)rectangle++(0.009,0.009) (0.504,0.046)rectangle++(0.009,0.009) (0.25,0.521)rectangle++(0.009,0.009) (0.123,0.258)rectangle++(0.009,0.009) (0.059,0.127)rectangle++(0.009,0.009)
  (0.027,0.561)rectangle++(0.009,0.009) (0.011,0.278)rectangle++(0.009,0.009) (0.003,0.637)rectangle++(0.009,0.009) (-0.001,0.816)rectangle++(0.009,0.009) (-0.003,0.406)rectangle++(0.009,0.009)
  (-0.004,0.201)rectangle++(0.009,0.009) (-0.004,0.598)rectangle++(0.009,0.009) (-0.004,0.797)rectangle++(0.009,0.009) (0.496,0.316)rectangle++(0.009,0.009) (0.696,0.124)rectangle++(0.009,0.009)
  (0.346,0.56)rectangle++(0.009,0.009) (0.171,0.778)rectangle++(0.009,0.009) (0.083,0.887)rectangle++(0.009,0.009) (0.039,0.941)rectangle++(0.009,0.009) (0.513,0.374)rectangle++(0.009,0.009)
  (0.254,0.685)rectangle++(0.009,0.009) (0.599,0.271)rectangle++(0.009,0.009) (0.297,0.133)rectangle++(0.009,0.009) (0.146,0.064)rectangle++(0.009,0.009) (0.071,0.03)rectangle++(0.009,0.009)
  (0.033,0.513)rectangle++(0.009,0.009) (0.511,0.202)rectangle++(0.009,0.009) (0.253,0.599)rectangle++(0.009,0.009) (0.124,0.297)rectangle++(0.009,0.009) (0.547,0.116)rectangle++(0.009,0.009)
  (0.716,0.044)rectangle++(0.009,0.009) (0.356,0.02)rectangle++(0.009,0.009) (0.176,0.008)rectangle++(0.009,0.009) (0.086,0.502)rectangle++(0.009,0.009) (0.041,0.749)rectangle++(0.009,0.009)
  (0.514,0.297)rectangle++(0.009,0.009) (0.255,0.146)rectangle++(0.009,0.009) (0.599,0.056)rectangle++(0.009,0.009) (0.297,0.526)rectangle++(0.009,0.009) (0.146,0.761)rectangle++(0.009,0.009)
  (0.071,0.878)rectangle++(0.009,0.009) (0.033,0.937)rectangle++(0.009,0.009) (0.511,0.372)rectangle++(0.009,0.009) (0.253,0.684)rectangle++(0.009,0.009) (0.124,0.34)rectangle++(0.009,0.009)
  (0.06,0.168)rectangle++(0.009,0.009) (0.028,0.582)rectangle++(0.009,0.009) (0.012,0.789)rectangle++(0.009,0.009) (0.004,0.892)rectangle++(0.009,0.009) (0,0.444)rectangle++(0.009,0.009)
  (-0.002,0.22)rectangle++(0.009,0.009) (-0.003,0.608)rectangle++(0.009,0.009) (-0.004,0.302)rectangle++(0.009,0.009) (0.496,0.118)rectangle++(0.009,0.009) (0.246,0.557)rectangle++(0.009,0.009)
  (0.121,0.276)rectangle++(0.009,0.009) (0.058,0.136)rectangle++(0.009,0.009) (0.027,0.066)rectangle++(0.009,0.009) (0.508,0.024)rectangle++(0.009,0.009) (0.252,0.51)rectangle++(0.009,0.009)
  (0.124,0.253)rectangle++(0.009,0.009) (0.06,0.624)rectangle++(0.009,0.009) (0.028,0.31)rectangle++(0.009,0.009) (0.012,0.153)rectangle++(0.009,0.009) (0.502,0.058)rectangle++(0.009,0.009)
  (0.249,0.027)rectangle++(0.009,0.009) (0.122,0.011)rectangle++(0.009,0.009) (0.059,0.503)rectangle++(0.009,0.009) (0.027,0.749)rectangle++(0.009,0.009) (0.011,0.372)rectangle++(0.009,0.009)
  (0.502,0.146)rectangle++(0.009,0.009) (0.249,0.071)rectangle++(0.009,0.009) (0.122,0.033)rectangle++(0.009,0.009) (0.059,0.014)rectangle++(0.009,0.009) (0.027,0.005)rectangle++(0.009,0.009)
  (0.011,0.5)rectangle++(0.009,0.009) (0.003,0.748)rectangle++(0.009,0.009) (0.499,0.296)rectangle++(0.009,0.009) (0.697,0.116)rectangle++(0.009,0.009) (0.346,0.556)rectangle++(0.009,0.009)
  (0.171,0.276)rectangle++(0.009,0.009) (0.083,0.636)rectangle++(0.009,0.009) (0.531,0.252)rectangle++(0.009,0.009) (0.263,0.124)rectangle++(0.009,0.009) (0.129,0.06)rectangle++(0.009,0.009)
  (0.062,0.028)rectangle++(0.009,0.009) (0.029,0.012)rectangle++(0.009,0.009) (0.012,0.504)rectangle++(0.009,0.009) (0.004,0.75)rectangle++(0.009,0.009) (0,0.873)rectangle++(0.009,0.009)
  (-0.002,0.434)rectangle++(0.009,0.009) (0.496,0.171)rectangle++(0.009,0.009) (0.246,0.083)rectangle++(0.009,0.009) (0.596,0.031)rectangle++(0.009,0.009) (0.296,0.013)rectangle++(0.009,0.009)
  (0.146,0.004)rectangle++(0.009,0.009) (0.071,0)rectangle++(0.009,0.009) (0.033,0.498)rectangle++(0.009,0.009) (0.014,0.747)rectangle++(0.009,0.009) (0.005,0.371)rectangle++(0.009,0.009)
  (0,0.683)rectangle++(0.009,0.009) (-0.002,0.839)rectangle++(0.009,0.009) (-0.003,0.417)rectangle++(0.009,0.009) (0.496,0.164)rectangle++(0.009,0.009) (0.246,0.08)rectangle++(0.009,0.009)
  (0.121,0.538)rectangle++(0.009,0.009) (0.058,0.767)rectangle++(0.009,0.009) (0.027,0.881)rectangle++(0.009,0.009) (0.508,0.35)rectangle++(0.009,0.009) (0.252,0.673)rectangle++(0.009,0.009)
  (0.124,0.834)rectangle++(0.009,0.009) (0.06,0.415)rectangle++(0.009,0.009) (0.028,0.205)rectangle++(0.009,0.009) (0.012,0.6)rectangle++(0.009,0.009) (0.004,0.798)rectangle++(0.009,0.009)
  (0,0.397)rectangle++(0.009,0.009) (-0.002,0.196)rectangle++(0.009,0.009) (-0.003,0.596)rectangle++(0.009,0.009) (-0.004,0.796)rectangle++(0.009,0.009) (-0.004,0.896)rectangle++(0.009,0.009)
  (0.496,0.356)rectangle++(0.009,0.009) (0.696,0.14)rectangle++(0.009,0.009) (0.346,0.568)rectangle++(0.009,0.009) (0.171,0.282)rectangle++(0.009,0.009) (0.083,0.139)rectangle++(0.009,0.009)
  (0.039,0.567)rectangle++(0.009,0.009) (0.017,0.781)rectangle++(0.009,0.009) (0.006,0.888)rectangle++(0.009,0.009) (0.001,0.442)rectangle++(0.009,0.009) (-0.002,0.219)rectangle++(0.009,0.009)
  (-0.003,0.607)rectangle++(0.009,0.009) (-0.004,0.301)rectangle++(0.009,0.009) (-0.004,0.148)rectangle++(0.009,0.009) (-0.004,0.572)rectangle++(0.009,0.009) (-0.004,0.284)rectangle++(0.009,0.009)
  (-0.004,0.14)rectangle++(0.009,0.009) (-0.004,0.568)rectangle++(0.009,0.009) (-0.004,0.282)rectangle++(0.009,0.009) (-0.004,0.139)rectangle++(0.009,0.009) (0.496,0.053)rectangle++(0.009,0.009)
  (0.696,0.018)rectangle++(0.009,0.009) (0.346,0.007)rectangle++(0.009,0.009) (0.171,0.501)rectangle++(0.009,0.009) (0.083,0.248)rectangle++(0.009,0.009) (0.039,0.622)rectangle++(0.009,0.009)
  (0.017,0.809)rectangle++(0.009,0.009) (0.504,0.321)rectangle++(0.009,0.009) (0.25,0.658)rectangle++(0.009,0.009) (0.597,0.261)rectangle++(0.009,0.009) (0.296,0.128)rectangle++(0.009,0.009)
  (0.146,0.562)rectangle++(0.009,0.009) (0.071,0.779)rectangle++(0.009,0.009) (0.526,0.309)rectangle++(0.009,0.009) (0.261,0.152)rectangle++(0.009,0.009) (0.602,0.058)rectangle++(0.009,0.009)
  (0.738,0.021)rectangle++(0.009,0.009) (0.367,0.508)rectangle++(0.009,0.009) (0.181,0.252)rectangle++(0.009,0.009) (0.088,0.124)rectangle++(0.009,0.009) (0.042,0.56)rectangle++(0.009,0.009)
  (0.019,0.778)rectangle++(0.009,0.009) (0.505,0.308)rectangle++(0.009,0.009) (0.25,0.152)rectangle++(0.009,0.009) (0.123,0.074)rectangle++(0.009,0.009) (0.546,0.027)rectangle++(0.009,0.009)
  (0.271,0.011)rectangle++(0.009,0.009) (0.133,0.003)rectangle++(0.009,0.009) (0.064,-0.001)rectangle++(0.009,0.009) (0.03,-0.003)rectangle++(0.009,0.009) (0.509,-0.004)rectangle++(0.009,0.009)
  (0.252,-0.004)rectangle++(0.009,0.009) (0.124,-0.004)rectangle++(0.009,0.009) (0.06,0.496)rectangle++(0.009,0.009) (0.028,0.246)rectangle++(0.009,0.009) (0.508,0.096)rectangle++(0.009,0.009)
  (0.252,0.546)rectangle++(0.009,0.009) (0.598,0.216)rectangle++(0.009,0.009) (0.737,0.084)rectangle++(0.009,0.009) (0.366,0.54)rectangle++(0.009,0.009) (0.644,0.213)rectangle++(0.009,0.009)
  (0.755,0.083)rectangle++(0.009,0.009) (0.375,0.539)rectangle++(0.009,0.009) (0.647,0.213)rectangle++(0.009,0.009) (0.321,0.104)rectangle++(0.009,0.009) (0.626,0.039)rectangle++(0.009,0.009)
  (0.311,0.517)rectangle++(0.009,0.009) (0.153,0.256)rectangle++(0.009,0.009) (0.074,0.626)rectangle++(0.009,0.009) (0.527,0.248)rectangle++(0.009,0.009) (0.261,0.122)rectangle++(0.009,0.009)
  (0.128,0.059)rectangle++(0.009,0.009) (0.062,0.027)rectangle++(0.009,0.009) (0.029,0.511)rectangle++(0.009,0.009) (0.012,0.753)rectangle++(0.009,0.009) (0.004,0.874)rectangle++(0.009,0.009)
  (0,0.435)rectangle++(0.009,0.009) (-0.002,0.215)rectangle++(0.009,0.009) (-0.003,0.605)rectangle++(0.009,0.009) (-0.004,0.8)rectangle++(0.009,0.009) (-0.004,0.898)rectangle++(0.009,0.009)
  (-0.004,0.447)rectangle++(0.009,0.009) (0.496,0.176)rectangle++(0.009,0.009) (0.246,0.586)rectangle++(0.009,0.009) (0.121,0.291)rectangle++(0.009,0.009) (0.058,0.143)rectangle++(0.009,0.009)
  (0.027,0.569)rectangle++(0.009,0.009) (0.011,0.282)rectangle++(0.009,0.009) (0.502,0.11)rectangle++(0.009,0.009) (0.249,0.553)rectangle++(0.009,0.009) (0.597,0.218)rectangle++(0.009,0.009)
  (0.296,0.107)rectangle++(0.009,0.009) (0.146,0.551)rectangle++(0.009,0.009) (0.071,0.273)rectangle++(0.009,0.009) (0.033,0.634)rectangle++(0.009,0.009) (0.014,0.315)rectangle++(0.009,0.009)
  (0.503,0.123)rectangle++(0.009,0.009) (0.249,0.059)rectangle++(0.009,0.009) (0.122,0.527)rectangle++(0.009,0.009) (0.059,0.761)rectangle++(0.009,0.009) (0.027,0.878)rectangle++(0.009,0.009)
  (0.011,0.437)rectangle++(0.009,0.009) (0.003,0.216)rectangle++(0.009,0.009) (-0.001,0.606)rectangle++(0.009,0.009) (-0.003,0.301)rectangle++(0.009,0.009) (0.496,0.118)rectangle++(0.009,0.009)
  (0.696,0.044)rectangle++(0.009,0.009) (0.346,0.02)rectangle++(0.009,0.009) (0.171,0.008)rectangle++(0.009,0.009) (0.083,0.502)rectangle++(0.009,0.009) (0.039,0.749)rectangle++(0.009,0.009)
  (0.017,0.372)rectangle++(0.009,0.009) (0.504,0.146)rectangle++(0.009,0.009) (0.25,0.071)rectangle++(0.009,0.009) (0.123,0.033)rectangle++(0.009,0.009) (0.546,0.011)rectangle++(0.009,0.009)
  (0.271,0.003)rectangle++(0.009,0.009) (0.133,0.499)rectangle++(0.009,0.009) (0.064,0.247)rectangle++(0.009,0.009) (0.523,0.096)rectangle++(0.009,0.009) (0.259,0.546)rectangle++(0.009,0.009)
  (0.601,0.216)rectangle++(0.009,0.009) (0.738,0.084)rectangle++(0.009,0.009) (0.367,0.04)rectangle++(0.009,0.009) (0.644,0.013)rectangle++(0.009,0.009) (0.32,0.004)rectangle++(0.009,0.009)
  (0.158,0.5)rectangle++(0.009,0.009) (0.077,0.248)rectangle++(0.009,0.009) (0.036,0.122)rectangle++(0.009,0.009) (0.016,0.059)rectangle++(0.009,0.009) (0.006,0.527)rectangle++(0.009,0.009)
  (0.001,0.761)rectangle++(0.009,0.009) (-0.002,0.878)rectangle++(0.009,0.009) (0.497,0.349)rectangle++(0.009,0.009) (0.246,0.172)rectangle++(0.009,0.009) (0.121,0.584)rectangle++(0.009,0.009)
  (0.546,0.231)rectangle++(0.009,0.009) (0.271,0.113)rectangle++(0.009,0.009) (0.606,0.043)rectangle++(0.009,0.009) (0.301,0.019)rectangle++(0.009,0.009) (0.148,0.507)rectangle++(0.009,0.009)
  (0.072,0.751)rectangle++(0.009,0.009) (0.034,0.873)rectangle++(0.009,0.009) (0.015,0.434)rectangle++(0.009,0.009) (0.503,0.171)rectangle++(0.009,0.009) (0.249,0.083)rectangle++(0.009,0.009)
  (0.122,0.539)rectangle++(0.009,0.009) (0.059,0.767)rectangle++(0.009,0.009) (0.521,0.304)rectangle++(0.009,0.009) (0.706,0.119)rectangle++(0.009,0.009) (0.78,0.045)rectangle++(0.009,0.009)
  (0.388,0.02)rectangle++(0.009,0.009) (0.192,0.008)rectangle++(0.009,0.009) (0.094,0.502)rectangle++(0.009,0.009) (0.045,0.249)rectangle++(0.009,0.009) (0.02,0.122)rectangle++(0.009,0.009)
  (0.008,0.559)rectangle++(0.009,0.009) (0.002,0.277)rectangle++(0.009,0.009) (-0.001,0.136)rectangle++(0.009,0.009) (-0.003,0.066)rectangle++(0.009,0.009) (-0.004,0.531)rectangle++(0.009,0.009)
  (0.496,0.21)rectangle++(0.009,0.009) (0.246,0.103)rectangle++(0.009,0.009) (0.596,0.038)rectangle++(0.009,0.009) (0.296,0.017)rectangle++(0.009,0.009) (0.146,0.006)rectangle++(0.009,0.009)
  (0.556,0)rectangle++(0.009,0.009) (0.276,-0.002)rectangle++(0.009,0.009) (0.608,-0.004)rectangle++(0.009,0.009) (0.74,-0.004)rectangle++(0.009,0.009) (0.793,-0.004)rectangle++(0.009,0.009)
  (0.394,0.496)rectangle++(0.009,0.009) (0.195,0.246)rectangle++(0.009,0.009) (0.575,0.096)rectangle++(0.009,0.009) (0.285,0.546)rectangle++(0.009,0.009) (0.14,0.271)rectangle++(0.009,0.009);

        \node[etiq, align=center] at (0.5,-0.2) {(b) deux copies de rapport $\tfrac12$,\\ une de rapport $\num{0.4}$ : elles se touchent};
      \end{scope}
      \begin{scope}[shift={(2.7,0)}]
        \foreach \x/\y in {0/0, 0.5/0, 0/0.5, 0.5/0.5}
          \fill[steelblue4, fill opacity=0.35] (\x,\y) rectangle ++(0.5,0.5);
        \fill[firebrick3, fill opacity=0.45] (0.25,0.25) rectangle ++(0.5,0.5);
        \draw[sectioncolor!60, dashed, thin] (0,0) rectangle (1,1);
        \node[etiq, align=center] at (0.5,-0.2) {(c) cinq copies de rapport $\tfrac12$ :\\ elles se chevauchent};
      \end{scope}
    \end{tikzpicture}}
  \small Le carré unité ouvert $V$ en pointillé. En (a) et (b), la
  condition de l'ouvert est satisfaite ; en (c), la cinquième image
  chevauche les quatre autres.
\end{frame}

\begin{frame}{Exemple 3.9 --- Pourquoi la condition de l'ouvert est nécessaire}
  \small Cas (c) : cinq homothéties de rapport $\tfrac12$, qui envoient
  le carré unité $C$ sur ses quatre quarts et sur le carré de côté
  $\tfrac12$ centré en $\big(\tfrac12,\tfrac12\big)$.
  \begin{itemize}
    \item Les quatre coins reconstituent déjà $C$, et la cinquième copie
      est contenue dans $C$ : $W(C)=C$, l'attracteur est le carré plein,
      de dimension $2$.
    \item Moran donnerait $5\big(\tfrac12\big)^d=1$, soit
      $d=\ln5/\ln2\approx\num{2.32}$ --- impossible pour un ensemble plan.
    \item Le chevauchement fait compter plusieurs fois certaines parties
      du carré.
  \end{itemize}
\end{frame}

\begin{frame}{Exemple 3.10 --- Le calcul numérique d'une dimension}
  \small Cas (b) : on résout
  $f(d) = 2\big(\tfrac12\big)^d + (\num{0.4})^d - 1 = 0$ par la méthode
  de Newton (chapitre 2), avec
  $f'(d) = 2\big(\tfrac12\big)^d\ln\tfrac12 + (\num{0.4})^d\ln(\num{0.4})$
  et $d_0=1$ :
  \begin{center}
  \small
  \begin{tabular}{ccc}
    \toprule
    $n$ & $d_n$ & $f(d_n)$ \\
    \midrule
    $0$ & $\num{1.000000000}$ & $\num{4.0e-1}$ \\
    $1$ & $\num{1.377478332}$ & $\num{5.3e-2}$ \\
    $2$ & $\num{1.444094153}$ & $\num{1.3e-3}$ \\
    $3$ & $\num{1.445852445}$ & $\num{8.9e-7}$ \\
    $4$ & $\num{1.445853630}$ & $\num{4.0e-13}$ \\
    \bottomrule
  \end{tabular}
  \end{center}
  \small $d\approx\num{1.44585}$ : entre $\ln3/\ln(\num{2.5})\approx\num{1.199}$
  (trois copies de rapport $\num{0.4}$) et $\ln3/\ln2\approx\num{1.585}$
  (trois copies de rapport $\tfrac12$).
\end{frame}

\begin{frame}{Remarque 3.8 --- Et pour les SIF qui ne sont pas des similitudes ?}
  \begin{Rem}{Remarque 3.8}
    La définition 3.5 ne vaut que pour des similitudes. L'arbre, la
    fougère, tout attracteur d'applications affines quelconques : une
    application qui écrase davantage dans une direction ne multiplie pas
    une « taille $d$-dimensionnelle » par un facteur bien défini.
    \vspace{4pt}

    La dimension de ces ensembles auto-affines demande d'autres outils,
    et fait encore l'objet de recherches actives.
  \end{Rem}
\end{frame}
